\documentclass[final,12pt,a4paper,twoside,fleqn]{extbook}

\usepackage[spanish]{babel}
\decimalpoint
\def\@makechapterhead#1{%
  \vspace*{50\p@}%
  {\parindent \z@ \raggedright \normalfont
    \ifnum \c@secnumdepth >\m@ne
      \if@mainmatter
        %\huge\bfseries \@chapapp\space \thechapter
        \huge\bfseries \thechapter.\space%
        %\par\nobreak
        %\vskip 20\p@
      \fi
    \fi
    \interlinepenalty\@M
    \huge \bfseries #1\par\nobreak
    \vskip 40\p@
  }
}
\makeatother
\makeatletter
\addto\captionsspanish{
    \renewcommand{\partname}{}
    \renewcommand{\chaptername}{}
}
\makeatother
\makeatletter
\usepackage[backend=biber]{biblatex}
\addbibresource{bibliografia.bib}
\usepackage{hyperref}
\usepackage{listings}
\usepackage[T1,LGR]{fontenc}
\usepackage[utf8]{inputenc}
%\usepackage[noenc]{tipa}
\usepackage[fleqn]{mathtools}
\usepackage{lplfitch}
\usepackage{12many}
\usepackage{abraces}
%\usepackage{abstract}
\usepackage{yfonts}
\usepackage{color}
\usepackage[usenames,dvipsnames,svgnames,table]{xcolor}
\usepackage{xxcolor}
\usepackage{verbatim}
\usepackage{anyfontsize}
\usepackage{anysize}
\usepackage{helvet}
\usepackage{textcomp}
\usepackage{dsfont}
\usepackage{fancybox}
\usepackage{pgf}
\usepackage{tikz}
\usetikzlibrary{arrows,decorations.shapes,automata,backgrounds,petri,cd}
\usetikzlibrary{positioning}
\usetikzlibrary{shadows}
\usetikzlibrary{chains}
\usepackage{bbm}
\usepackage[inline]{enumitem}
\usepackage{mathdesign}
\usepackage{calrsfs}
\usepackage{mathrsfs}
\usepackage{mathalfa}
\usepackage{amssymb}
%\usepackage{amsrefs}
\usepackage{amsmath}
\usepackage{amsthm}
%\usepackage{amsaddr}
\usepackage{amsfonts}
\usepackage{amsxtra}
\usepackage{amscd}
\usepackage{bm}
%\usepackage{wasysym}
\usepackage{bold-extra}
\usepackage{bbm}
\usepackage[mathscr]{euscript}
\usepackage{color}
\usepackage{units}
\usepackage{nomencl}
\providecommand{\printnomenclature}{\printglossary}
\providecommand{\makenomenclature}{\makeglossary}
\makenomenclature
\usepackage{microtype}
\usepackage{graphicx}
\makeatletter
%% EPISTEMOLOGIA DE LAS MATEM\'ATICAS
\newcommand{\sforall}{\ensuremath{\forall}}
\newcommand{\sexists}{\ensuremath{         \exists         }}
\newcommand{\snexists}{\ensuremath{         \nexists         }}
\newcommand{\sand}{\ensuremath{         \wedge         }}
\newcommand{\sor}{\ensuremath{ ∨}}
\newcommand{\sthen}{\ensuremath{         \Rightarrow         }}
\newcommand{\sif}{\ensuremath{         \Leftarrow         }}
\newcommand{\ssi}{\ensuremath{\Leftrightarrow}}
\newcommand{\setomega}{\ensuremath{\omega }}
\newcommand{\setmu}{\ensuremath{\mu }}
\newcommand{\gensig}{\ensuremath{\sigma }}
\newcommand{\bneg}{\ensuremath{         \sim }}
\newcommand{\ssimeq}{\ensuremath{         \simeq         }}
\newcommand{\sunion}{\ensuremath{         \cup         }}
\newcommand{\sintersection}{\ensuremath{ ∩}}
\newcommand{\sproduct}{\ensuremath{         \times         }}
\newcommand{\sminus}{\ensuremath{\setminus}}
\newcommand{\ssymdiff}{\ensuremath{\triangle}}
\newcommand{\sbin}{\ensuremath{ ∈}}
\newcommand{\snin}{\ensuremath{         \notin         }}
\newcommand{\seq}{\ensuremath{=}}
\newcommand{\sneq}{\ensuremath{         \neq         }}
\newcommand{\eps}{\ensuremath{\varepsilon }}
\newcommand{\beps}{\ensuremath{\backepsilon}}
\newcommand{\toward}{\ensuremath{         \longrightarrow         }}
\newcommand{\elmapstoel}{\ensuremath{         \mapsto }}
\newcommand{\absoluto}{\ensuremath{\boldsymbol{\mathbf{V}}}}
\newcommand{\nada}{\ensuremath{\boldsymbol{\mathbf{\Lambda }}}}
\newcommand{\nta}{\ensuremath{\varphi }}
\newcommand{\ntb}{\ensuremath{\psi }}
\newcommand{\ntc}{\ensuremath{\vartheta }}
\newcommand{\vnta}{\ensuremath{\Phi }}
\newcommand{\vntb}{\ensuremath{\Psi }}
\newcommand{\vntc}{\ensuremath{\Theta }}
\newcommand{\chtitletext}[1]{\textbf{\emph{#1}}}
\newcommand{\memph}[1]{\textbf{\emph{#1}}}
\newcommand{\mmemph}[1]{\large{\memph{#1}}}
\newcommand{\bfs}[1]{\ensuremath{\boldsymbol{#1}}}
\newcommand{\can}[1]{\ensuremath{\mathfrak{#1}}}
\newcommand{\ntn}[1]{\ensuremath{\bfs{#1}}}
\newcommand{\exn}[1]{\ensuremath{\bfs{\mathscr{#1}}}}
\newcommand{\notas}[1]{\ensuremath{{
{}^{\displaystyle{\bfs{nt}}}{\!}{\bfs{#1}}
}}}
\newcommand{\casos}[1]{\ensuremath{{
    {}^{\displaystyle{\bfs{cs}}}{\!}{\bfs{#1}}
}}}
\newcommand{\notasn}[1]{\ensuremath{
    {\notas{\mathscr{#1}}}
}}
\newcommand{\casosn}[1]{\ensuremath{
    {\casos{\mathscr{#1}}}
}}
\newcommand{\Gen}[1]{\ensuremath{\bfs{
    {\mathcal{G}\mathbf{e}\mathbf{n}}\set{\exn{#1}}
}}}
\newcommand{\GenTo}[2]{\ensuremath{
    \set{#1}
    \bfs{\mathcal{G}\mathbf{e}\mathbf{n}{\_}\mathbf{a}}
    \set{#2}
}}
\newcommand{\ExtGenTo}[2]{\ensuremath{
    \GenTo{\exn{#1}}{\exn{#2}}
}}
\newcommand{\Par}[1]{\ensuremath{
        \bfs{
            {\mathcal{P}\mathbf{a}\mathbf{r}}\set{\mathscr{#1}}
        }
}}
\newcommand{\ParTo}[2]{\ensuremath{
    \set{#1}
    \bfs{\mathcal{P}\mathbf{a}\mathbf{r}\bfs{\_}\mathbf{a}}
    \set{#2}
}}
\newcommand{\ExtParTo}[2]{\ensuremath{
    \ParTo{\exn{#1}}{\exn{#2}}
}}
\newcommand{\tty}{\ensuremath{\mathtt{y}}}
\newcommand{\ttY}{\ensuremath{\mathtt{Y}}}
\newcommand{\ttx}{\ensuremath{\mathtt{x}}}
\newcommand{\entreparentesis}[1]{\ensuremath{
        \left({#1}\right)
}}
\newcommand{\entrecorchetes}[1]{\ensuremath{
        \left\lbrack{#1}\right\rbrack
}}
\newcommand{\entredoblescorchetes}[1]{\ensuremath{
        \left[\!{\left[{#1}\right]}\!\right]
}}
\newcommand{\entreangulos}[1]{\ensuremath{\left\langle{#1}\right\rangle}}
\newcommand{\doblelanglegrande}{\ensuremath{{\big{
                \left\langle\!\!\left\langle
}}}}
\newcommand{\dobleranglegrande}{\ensuremath{{\big{
                \right\rangle\!\!\right\rangle
}}}}
\newcommand{\entregrandesdoblesangulos}[1]{\ensuremath{
    \left\doblelanglegrande
    {#1}
    \right\dobleranglegrande
}}
\newcommand{\entrellaves}[1]{\ensuremath{\left\lbrace{#1}\right\rbrace}}
\newcommand{\entredoblesllaves}[1]{\ensuremath{
        \left\lbrace\!\left\lbrace{#1}\right\rbrace\!\right\rbrace
}}
\newcommand{\desederivaque}[2]{\ensuremath{{#1}\bfs{\;\vdash\;}{#2}}}
\newcommand{\denosederivaque}[2]{\ensuremath{{#1}\bfs{\;\nvdash\;}{#2}}}
\newcommand{\parts}{\ensuremath{\bfs{\mathcal{P}}}}
\newcommand{\partsof}[1]{\ensuremath{\parts{\entreparentesis{\exn{#1}}}}}
\newcommand{\universe}{\ensuremath{\bfs{\mathfrak{U}}}}
\newcommand{\semptyset}{\ensuremath{\emptyset}}
\newcommand{\cna}{\ensuremath{\can{a}}}
\newcommand{\cnb}{\ensuremath{\can{b}}}
\newcommand{\cnc}{\ensuremath{\can{c}}}
\newcommand{\cnd}{\ensuremath{\can{d}}}
\newcommand{\cne}{\ensuremath{\can{e}}}
\newcommand{\cnf}{\ensuremath{\can{f}}}
\newcommand{\cng}{\ensuremath{\can{g}}}
\newcommand{\cnh}{\ensuremath{\can{h}}}
\newcommand{\cni}{\ensuremath{\can{i}}}
\newcommand{\cnj}{\ensuremath{\can{j}}}
\newcommand{\cnk}{\ensuremath{\can{k}}}
\newcommand{\cnl}{\ensuremath{\can{l}}}
\newcommand{\cnm}{\ensuremath{\can{m}}}
\newcommand{\cnn}{\ensuremath{\can{n}}}
\newcommand{\cno}{\ensuremath{\can{o}}}
\newcommand{\cnp}{\ensuremath{\can{p}}}
\newcommand{\cnq}{\ensuremath{\can{q}}}
\newcommand{\cnr}{\ensuremath{\can{r}}}
\newcommand{\cns}{\ensuremath{\can{s}}}
\newcommand{\cnt}{\ensuremath{\can{t}}}
\newcommand{\cnu}{\ensuremath{\can{u}}}
\newcommand{\cnw}{\ensuremath{\can{w}}}
\newcommand{\cnx}{\ensuremath{\can{x}}}
\newcommand{\cny}{\ensuremath{\can{y}}}
\newcommand{\cnz}{\ensuremath{\can{z}}}
\newcommand{\enA}{\ensuremath{\exn{A}}}
\newcommand{\enB}{\ensuremath{\exn{B}}}
\newcommand{\enC}{\ensuremath{\exn{C}}}
\newcommand{\enD}{\ensuremath{\exn{D}}}
\newcommand{\enE}{\ensuremath{\exn{E}}}
\newcommand{\enF}{\ensuremath{\exn{F}}}
\newcommand{\enG}{\ensuremath{\exn{G}}}
\newcommand{\enH}{\ensuremath{\exn{H}}}
\newcommand{\enI}{\ensuremath{\exn{I}}}
\newcommand{\enJ}{\ensuremath{\exn{A}}}
\newcommand{\enK}{\ensuremath{\exn{B}}}
\newcommand{\enL}{\ensuremath{\exn{C}}}
\newcommand{\enM}{\ensuremath{\exn{A}}}
\newcommand{\enN}{\ensuremath{\exn{B}}}
\newcommand{\enO}{\ensuremath{\exn{C}}}
\newcommand{\enP}{\ensuremath{\exn{A}}}
\newcommand{\enQ}{\ensuremath{\exn{B}}}
\newcommand{\enR}{\ensuremath{\exn{C}}}
\newcommand{\enS}{\ensuremath{\exn{A}}}
\newcommand{\enT}{\ensuremath{\exn{B}}}
\newcommand{\enU}{\ensuremath{\exn{C}}}
\newcommand{\enV}{\ensuremath{\exn{A}}}
\newcommand{\enW}{\ensuremath{\exn{B}}}
\newcommand{\enX}{\ensuremath{\exn{C}}}
\newcommand{\enY}{\ensuremath{\exn{A}}}
\newcommand{\enZ}{\ensuremath{\exn{B}}}
\newcommand{\predSet}{\ensuremath{\can{S}\!\can{e}\!\can{t}}}
\newcommand{\predClass}{\ensuremath{
        \can{C}\!\can{l}\!\can{a}\!\can{s}\!\can{s}
}}
\newcommand{\mathdef}{\ensuremath{\stackrel{\mathsf{def}}{{}={}}}}
\newcommand{\pfmathdef}[2]{\ensuremath{{#1}{\mathdef}{#2}}}
\newcommand{\logicdef}{\ensuremath{
        \stackrel
            {\mathsf{def}}
            {\Longleftrightarrow}
}}
\newcommand{\pflogicdef}[2]{\ensuremath{{#1}{\logicdef}{#2}}}
\newcommand{\Real}{\ensuremath{\mathds{R}}}
\newcommand{\logicand}{\ensuremath{\sand }}
\newcommand{\logicor}{\ensuremath{\sor}}
\newcommand{\then}{\ensuremath{\sthen}}
\newcommand{\entonces}{\ensuremath{\sthen}}
\newcommand{\pflogicand}[2]{\ensuremath{{#1}\logicand{#2}}}
\newcommand{\pflogicor}[2]{\ensuremath{{#1}\logicor{#2}}}
\newcommand{\pfthen}[2]{\ensuremath{{#1}\then{#2}}}
\newcommand{\pfentonces}[2]{\ensuremath{{#1}\entonces{#2}}}
\newcommand{\pfsetunion}[2]{\ensuremath{{#1}\sunion{#2}}}
\newcommand{\pfexnunionexn}[2]{\ensuremath{
        \pfsetunion{\exn{#1}}{\exn{#2}}
}}
\newcommand{\pfcsnexnunionexn}[2]{\ensuremath{
        \pfsetunion{\casosn{#1}}{\casosn{#2}}
}}
\newcommand{\pfsetintersection}[2]{\ensuremath{
        {#1}\sintersection{#2}
}}
\newcommand{\pfexnintersectionexn}[2]{\ensuremath{
        \pfsetintersection{\exn{#1}}{\exn{#2}}
}}
\newcommand{\pfcsnexnintersectionexn}[2]{\ensuremath{
        \pfsetintersection{\casosn{#1}}{\casosn{#2}}
}}
\newcommand{\pfsetproduct}[2]{\ensuremath{{#1}{\sproduct}{#2}}}
\newcommand{\pfexnproductexn}[2]{\ensuremath{
        \pfsetproduct{\exn{#1}}{\exn{#2}}
}}
\newcommand{\pfcsnexnproductexn}[2]{\ensuremath{
        \pfsetproduct{\casosn{#1}}{\casosn{#2}}
}}
\newcommand{\pfsetminus}[2]{\ensuremath{{#1}{\sminus}{#2}}}
\newcommand{\pfexnminusexn}[2]{\ensuremath{
        \pfsetminus{\exn{#1}}{\exn{#2}}
}}
\newcommand{\pfcsnexnminusexn}[2]{\ensuremath{
        \pfsetminus{\casosn{#1}}{\casosn{#2}}
}}
\newcommand{\pfsetsymdiff}[2]{\ensuremath{{#1}{\ssymdiff}{#2}}}
\newcommand{\pfexnsymdiffexn}[2]{\ensuremath{
        \pfsetsymdiff{\exn{#1}}{\exn{#2}}
}}
\newcommand{\pfcsnexnsymdiffexn}[2]{\ensuremath{
        \pfsetsymdiff{\casosn{#1}}{\casosn{#2}}
}}
\newcommand{\pfin}[2]{\ensuremath{{#1}\sbin{#2}}}
\newcommand{\pfcsnin}[2]{\ensuremath{{#1}\sbin\casosn{#2}}}
\newcommand{\pfnotin}[2]{\ensuremath{{#1}\snin{#2}}}
\newcommand{\pfcsnnotin}[2]{\ensuremath{{#1}\snin\casosn{#2}}}
\newcommand{\pfeq}[2]{\ensuremath{{#1}\seq{#2}}}
\newcommand{\pfneq}[2]{\ensuremath{{#1}\sneq{#2}}}
\newcommand{\pfexninexn}[2]{\ensuremath{\pfin{\exn{#1}}{\exn{#2}}}}
\newcommand{\pfcaninexn}[2]{\ensuremath{\pfin{\can{#1}}{\exn{#2}}}}
\newcommand{\pfcanincan}[2]{\ensuremath{\pfin{\can{#1}}{\can{#2}}}}
\newcommand{\pfexnincan}[2]{\ensuremath{\pfin{\exn{#1}}{\can{#2}}}}
\newcommand{\pfexneqexn}[2]{\ensuremath{\pfeq{\exn{#1}}{\exn{#2}}}}
\newcommand{\pfcaneqexn}[2]{\ensuremath{\pfeq{\can{#1}}{\exn{#2}}}}
\newcommand{\pfcaneqcan}[2]{\ensuremath{\pfeq{\can{#1}}{\can{#2}}}}
\newcommand{\pfexneqcan}[2]{\ensuremath{\pfeq{\exn{#1}}{\can{#2}}}}
\newcommand{\pfexnnotinexn}[2]{\ensuremath{\pfnotin{\exn{#1}}{\exn{#2}}}}
\newcommand{\pfcannotinexn}[2]{\ensuremath{\pfnotin{\can{#1}}{\exn{#2}}}}
\newcommand{\pfcannotincan}[2]{\ensuremath{\pfnotin{\can{#1}}{\can{#2}}}}
\newcommand{\pfexnnotincan}[2]{\ensuremath{\pfnotin{\exn{#1}}{\can{#2}}}}
\newcommand{\pfexnneqexn}[2]{\ensuremath{\pfneq{\exn{#1}}{\exn{#2}}}}
\newcommand{\pfcanneqexn}[2]{\ensuremath{\pfneq{\can{#1}}{\exn{#2}}}}
\newcommand{\pfcanneqcan}[2]{\ensuremath{\pfneq{\can{#1}}{\can{#2}}}}
\newcommand{\pfexnneqcan}[2]{\ensuremath{\pfneq{\exn{#1}}{\can{#2}}}}
\newcommand{\pfsubseteq}[2]{\ensuremath{{#1}         \subseteq {#2}}}
\newcommand{\pfnsubseteq}[2]{\ensuremath{{#1}         \nsubseteq {#2}}}
\newcommand{\pfsupseteq}[2]{\ensuremath{{#1}         \supseteq {#2}}}
\newcommand{\pfnsupseteq}[2]{\ensuremath{{#1}\nsupseteq{#2}}}
\newcommand{\pfexnsubexneq}[2]{\ensuremath{\pfsubseteq{\exn{#1}}{\exn{#2}}}}
\newcommand{\pfcansubexneq}[2]{\ensuremath{\pfsubseteq{\can{#1}}{\exn{#2}}}}
\newcommand{\pfcansubcaneq}[2]{\ensuremath{\pfsubseteq{\can{#1}}{\can{#2}}}}
\newcommand{\pfexnsubcaneq}[2]{\ensuremath{\pfsubseteq{\exn{#1}}{\can{#2}}}}
\newcommand{\pfexnsupexneq}[2]{\ensuremath{\pfsupseteq{\exn{#1}}{\exn{#2}}}}
\newcommand{\pfcansupexneq}[2]{\ensuremath{\pfsupseteq{\can{#1}}{\exn{#2}}}}
\newcommand{\pfcansupcaneq}[2]{\ensuremath{\pfsupseteq{\can{#1}}{\can{#2}}}}
\newcommand{\pfexnsupcaneq}[2]{\ensuremath{\pfsupseteq{\exn{#1}}{\can{#2}}}}
\newcommand{\pfcansubcan}[2]{\ensuremath{\pfsubset{\can{#1}}{\can{#2}}}}
\newcommand{\pfcansubexn}[2]{\ensuremath{\pfsubset{\can{#1}}{\exn{#2}}}}
\newcommand{\pfexnsubcan}[2]{\ensuremath{\pfsubset{\exn{#1}}{\can{#2}}}}
\newcommand{\pfexnsubexn}[2]{\ensuremath{\pfsubset{\exn{#1}}{\exn{#2}}}}
\newcommand{\pfcansupcan}[2]{\ensuremath{\pfsupset{\can{#1}}{\can{#2}}}}
\newcommand{\pfcansupexn}[2]{\ensuremath{\pfsupset{\can{#1}}{\exn{#2}}}}
\newcommand{\pfexnsupcan}[2]{\ensuremath{\pfsupset{\exn{#1}}{\can{#2}}}}
\newcommand{\pfexnsupexn}[2]{\ensuremath{\pfsupset{\exn{#1}}{\exn{#2}}}}
\newcommand{\existecan}{\ensuremath{\sexists}}
\newcommand{\existecande}[1]{\ensuremath{\existecan{\can{#1}}}}
\newcommand{\existecanen}[2]{\ensuremath{
    \entreparentesis{\existecande{#1}}
    \entreparentesis{\pfin{\can{#1}}{#2}}
}}
\newcommand{\existeexnen}[2]{\ensuremath{
    \entreparentesis{\existecande{#1}}
    \entreparentesis{\pfin{\exn{#1}}{#2}}
}}
\newcommand{\existeexn}{\ensuremath{\sexists}}
\newcommand{\existeexnde}[1]{\ensuremath{\existeexn{\exn{#1}}}}
\newcommand{\existeexnencan}[2]{\ensuremath{
    \left(\existeexnde{#1}\right)
    \left(\pfexnincan{#1}{#2}\right)
}}
\newcommand{\existeexnencanque}[3]{\ensuremath{
\left(\existeexnde{#1}\right)
\entreparentesis{
        \pflogicand{\left(\pfexnincan{#1}{#2}\right)}
        {\entreparentesis{#3}}
    }
}}
\newcommand{\existecanenexn}[2]{\ensuremath{
    \left(\existecande{#1}\right)
    \left(\pfcaninexn{#1}{#2}\right)
}}
\newcommand{\existecanenexnque}[3]{\ensuremath{
\left(\existecande{#1}\right)
\entreparentesis{
        \pflogicand{\left(\pfcaninexn{#1}{#2}\right)}
        {\entreparentesis{#3}}
    }
}}
\newcommand{\existecanencan}[2]{\ensuremath{
    \left(\existecande{#1}\right)
    \left(\pfcanincan{#1}{#2}\right)
}}
\newcommand{\existecanencanque}[3]{\ensuremath{
    \left(\existecande{#1}\right)
    \entreparentesis{
        \pflogicand{\left(\pfcanincan{#1}{#2}\right)}
        {\entreparentesis{#3}}
    }
}}
\newcommand{\existeexnenexn}[2]{\ensuremath{
    \left(\existeexnde{#1}\right)
    \left(\pfexninexn{#1}{#2}\right)
}}
\newcommand{\existeexnenexnque}[3]{\ensuremath{
    \left(\existeexnde{#1}\right)
    \entreparentesis{
          \pflogicand{\left(\pfexninexn{#1}{#2}\right)}
          {\entreparentesis{#3}}
    }
}}
\newcommand{\noexistecanen}[2]{\ensuremath{\noexistecan\pfin{#1}{#2}}}
\newcommand{\noexisteexten}[2]{\ensuremath{\noexistecasoen{\exn{#1}}{\exn{#2}}}}
\newcommand{\noexistecaso}{\ensuremath{\snexists}}
\newcommand{\noexistecasode}[1]{\ensuremath{\noexistecaso{#1}}}
\newcommand{\noexisteextde}[1]{\ensuremath{\noexistecasode{\exn{#1}}}}
\newcommand{\paratodocaso}{\ensuremath{\sforall}}
\newcommand{\paratodocan}{\ensuremath{\sforall}}
\newcommand{\paratodocasode}[1]{\ensuremath{\paratodocaso{#1}}}
\newcommand{\paratodocande}[1]{\ensuremath{\paratodocan{\can{#1}}}}
\newcommand{\paratodocasonombde}[1]{\ensuremath{\paratodocasode{\can{#1}}}}
\newcommand{\paratodocasoextde}[1]{\ensuremath{\paratodocasode{\exn{#1}}}}
\newcommand{\paraninguncasode}[1]{\ensuremath{\paraninguncaso{#1}}}
\newcommand{\paraningunaextde}[1]{\ensuremath{\paraninguncasode{\exn{#1}}}}
\newcommand{\paratodocasoen}[2]{\ensuremath{\paratodocasode{#1}\pfin{#1}{#2}}}
\newcommand{\paratodocanen}[2]{\ensuremath{
        \entreparentesis{
            \paratodocande{#1}
        }
        \entreparentesis{
            \pfin{\can{#1}}{#2}
        }
}}
\newcommand{\esoytalque}[2]{\ensuremath{
        \entreparentesis{
            \pflogicand{#1}{#2}
        }
}}
\newcommand{\esoimplicaque}[2]{\ensuremath{
    \entreparentesis{
        \pfentonces{#1}{#2}
    }
}}
\newcommand{\paratodocanenexnque}[3]{\ensuremath{
    \entreparentesis{\paratodocasode{\can{#1}}}
    \esoytalque{\entreparentesis{\pfin{\can{#1}}{\exn{#2}}}}{#3}
}}
\newcommand{\paratodocanenexnimplicaque}[3]{\ensuremath{
    \entreparentesis{\paratodocasode{\can{#1}}}
    \esoimplicaque{\entreparentesis{\pfin{\can{#1}}{\exn{#2}}}}{#3}
}}
\newcommand{\paratodocanencanque}[3]{\ensuremath{
    \entreparentesis{\paratodocasode{\can{#1}}}
    \esoytalque{\entreparentesis{\pfin{\can{#1}}{\can{#2}}}}{#3}
}}
\newcommand{\paratodoexnencanque}[3]{\ensuremath{
    \entreparentesis{\paratodocasode{\exn{#1}}}
    \esoytalque{\entreparentesis{\pfin{\exn{#1}}{\can{#2}}}}{#3}
}}
\newcommand{\paratodoexnenexnque}[3]{\ensuremath{
\entreparentesis{\paratodocasode{\exn{#1}}}
\esoytalque{\entreparentesis{\pfin{\exn{#1}}{\exn{#2}}}}{#3}
}}
\newcommand{\paratodocanenexn}[2]{\ensuremath{
    \entreparentesis{\paratodocasode{\can{#1}}}
    \entreparentesis{\pfin{\can{#1}}{\exn{#2}}}
}}
\newcommand{\paratodocanencan}[2]{\ensuremath{
    \entreparentesis{\paratodocasode{\can{#1}}}
    \entreparentesis{\pfin{\can{#1}}{\can{#2}}}
}}
\newcommand{\paratodoexnencan}[2]{\ensuremath{
    \entreparentesis{\paratodocasode{\exn{#1}}}
    \entreparentesis{\pfin{\exn{#1}}{\can{#2}}}
}}
\newcommand{\paratodoexnenexn}[2]{\ensuremath{
    \entreparentesis{
        \paratodocasode{\exn{#1}}
    }
    \entreparentesis{
        \pfin{\exn{#1}}{\exn{#2}}
    }
}}
\newcommand{\paraninguncasoen}[2]{\ensuremath{
        \paratodocaso\pfnotin{#1}{#2}
}}
\newcommand{\pfgensig}[1]{\ensuremath{\gensig{\left({#1}\right)}}}
\newcommand{\set}[1]{\ensuremath{\left\lbrace{#1}\right\rbrace}}
\newcommand{\extset}[2]{\ensuremath{
        \left\lbrace{#1}\,\vert\,{#2}\right\rbrace
}}
\newcommand{\ext}[1]{\ensuremath{\left\lbrace{#1}\right\rbrace}}
\newcommand{\extdext}[1]{\ensuremath{\left\lbrace{\exn{#1}}\right\rbrace}}
\newcommand{\predext}[2]{\ensuremath{\left\lbrace{#1}\,∣\,{#2}\right\rbrace}}
%\newcommand{\rfrac}[2]{\ensuremath{{{}^{#1}}{\!}{{/{\!}}_{#2}}}}
\newcommand{\norm}[1]{\ensuremath{\left\Vert#1\right\Vert}}
\newcommand{\abs}[1]{\ensuremath{\left\vert#1\right\vert}}
\newcommand{\BX}{\ensuremath{\bfs{B}(X)}}
\newcommand{\A}{\ensuremath{\mathcal{A}}}
\newcommand{\setN}{\ensuremath{\mathds{N}}}
\newcommand{\setR}{\ensuremath{\mathds{R}}}
\newcommand{\setQ}{\ensuremath{\mathds{Q}}}
\newcommand{\setZ}{\ensuremath{\mathbb{Z}}}
\newcommand{\setC}{\ensuremath{\mathbb{C}}}
\newcommand{\setH}{\ensuremath{\mathbb{H}}}
\newcommand{\setO}{\ensuremath{\mathbb{O}}}
\newcommand{\setB}{\ensuremath{\mathbb{B}}}
\newcommand{\relR}{\ensuremath{\mathscr{R}}}
\newcommand{\ntna}{\ensuremath{\bfs{\nta}}}
\newcommand{\ntnb}{\ensuremath{\bfs{\ntb}}}
\newcommand{\ntnc}{\ensuremath{\bfs{\ntc}}}
\newcommand{\vntna}{\ensuremath{\bfs{\vnta}}}
\newcommand{\vntnb}{\ensuremath{\bfs{\vntb}}}
\newcommand{\vntnc}{\ensuremath{\bfs{\vntc}}}
\makeatother
\makeatletter
% ANILLOS Y \'ALGEBRAS
\renewcommand{\norm}[1]{\ensuremath{\left\Vert#1\right\Vert}}
\renewcommand{\abs}[1]{\ensuremath{\left\vert#1\right\vert}}
\newcommand{\card}[1]{\ensuremath{\abs{#1}}}
\newcommand{\conj}[1]{\ensuremath{\left\lbrace{#1}\right\rbrace}}
\newcommand{\extconj}[2] {\ensuremath{
        \left\lbrace{#1}\,∣\,{#2}\right\rbrace
}}
\newcommand{\seps}{\ensuremath{\varepsilon }}
\newcommand{\To}{\ensuremath{         \longrightarrow         }}
\newcommand{\setfunct}[2]{\ensuremath{\mathbf{#1}{\left({#2}\right)}}}
\newcommand{\bcomplement}[1]{\ensuremath{\complement\left[{#1}\right]}}
\newcommand{\NullaryFunct}[1]{\ensuremath{#1\left(\right)\tighht}}
\newcommand{\UnaryFunct}[2]{\ensuremath{#1\left(#2\right)}}
\newcommand{\Idealz}{\ensuremath{\left(\mathfrak{0}\right)}}
\newcommand{\BinaryFunct}[3]{\ensuremath{#1\left(#2,#3\right)}}
\newcommand{\ThreearyFunct}[4]{\ensuremath{#1\left(#2,#3,#4\right)}}
\newcommand{\FouraryFunct}[5]{\ensuremath{#1\left(#2,#3,#4,#5\right)}}
\newcommand{\RingLocalizedOn}[2]{\ensuremath{
        \NullaryFunct{{#1}_{\Ideal{#2}}}
}}
\newcommand{\Ideal}[1]{\ensuremath{\mathfrak{#1}}}
\newcommand{\Idealx}[2]{\ensuremath{{{\Ideal{#1}}_{#2}}}}
\newcommand{\IdealGeneratorWithSet}[1] {\ensuremath{
        \left\langle\left\lbrace{#1}\right\rbrace\right\rangle
}}
\newcommand{\IdealGeneratorWithExtSet}[2]
{\ensuremath{
        \left\langle{\left\lbrace{#1}\,|\,{#2}\right\rbrace}\right\rangle
}}
\newcommand{\IdealGenerator}[1] {\ensuremath{
        \left\langle{#1}\right\rangle
}}
\newcommand{\IdealGeneratorExt}[2] {\ensuremath{
        \left\langle{{#1}\,|\,{#2}}\right\rangle
}}
\newcommand{\idealof}[2]{\ensuremath{\Ideal{#1}{\unlhd}\,{#2}}}
\newcommand{\primeidealof}[2]{\ensuremath{\Ideal{#1}{\lhd}_{prime}}{#2}}
\newcommand{\maximalidealof}[2]{\ensuremath{\Ideal{#1}{\lhd}_{max}}{#2}}
\newcommand{\idealofx}[3]{\ensuremath{\Idealx{#1}{#2}{\unlhd}\,{#3}}}
\newcommand{\primeidealofx}[3]{\ensuremath{\Idealx{#1}{#2}{\lhd}_{prime}}{#3}}
\newcommand{\maximalidealofx}[3]{\ensuremath{\Idealx{#1}{#2}{\lhd}_{max}}{#3}}
\newcommand{\Krullht}{\ensuremath{\textbf{ht}}}
\newcommand{\KrullHeight}[1]{\ensuremath{\Krullht\left({#1}\right)}}
\newcommand{\Kht}[1]{\ensuremath{\KrullHeight{\Ideal{#1}}}}
\newcommand{\Khtx}[2]{\ensuremath{\Krullht\left(\Idealx{#1}{#2}\right)}}
\newcommand{\ChainOfIdealsOnV}[2]{\ensuremath{\mathcal{#1}\left({#2}\right)}}
\newcommand{\ChainOfIdealsOn}[3] {\ensuremath{\mathcal{#1}\left({#2},\Ideal{#3}\right)}}
\newcommand{\ChainOfPrimeIdealsOn}[3] {\ensuremath{\mathcal{#1}_{prime}\left({#2},\Ideal{#3}\right)}}
\newcommand{\IsAChain}[3] {\ensuremath{\textsc{Chain}\left\lbrace \ChainOfIdealsOn{#1}{#2}{#3}\right\rbrace}}
\newcommand{\AscendentChain} {\ensuremath{\textsc{AscChain}}}
\newcommand{\AscendentChainV} {\ensuremath{\AscendentChain}}
\newcommand{\IsAnAscendentChainV}[2] {\ensuremath{\AscendentChainV\left\lbrace \ChainOfIdealsOnV{#1}{#2}\right\rbrace}}
\newcommand{\IsAnAscendentChain}[3] {\ensuremath{\AscendentChain\left\lbrace \ChainOfIdealsOn{#1}{#2}{\Ideal{#3}}\right\rbrace}}
\newcommand{\IsAnAscendentChainx}[4] {\ensuremath{\AscendentChain\left\lbrace \ChainOfIdealsOn{#1}{#2}{\Idealx{#3}{#4}}\right\rbrace}}
\newcommand{\IsAPrimeChain}[3] {\ensuremath{\textsc{Chain}_{prime} \left\lbrace\mathcal{#1}\left({#2},\Ideal{#3}\right)\right\rbrace}}
\newcommand{\IsAPrimeChainV}[2] {\ensuremath{\textsc{Chain}_{prime} \left\lbrace\mathcal{#1}\left({#2}\right)\right\rbrace}}
\newcommand{\FamilyOfChainsOf}[2] {\ensuremath{\mathcal{F}\mathcal{C}_{\left({#1},{\Ideal{#2}}\right)}}}
\newcommand{\FamilyOfChainsOfV}[1] {\ensuremath{\mathcal{F}\mathcal{C}_{\left({#1}\right)}}}
\newcommand{\FamilyOfPrimeChainsOf}[2] {\ensuremath{
        \mathcal{F}
        \mathcal{C}_{
            \left({#1},{\Ideal{#2}}\right)}^{prime}
}}
\newcommand{\FamilyOfPrimeChainsOfV}[1]{\ensuremath{
        \mathcal{F}
        \mathcal{C}_{\left({#1}\right)}^{prime}
}}
\newcommand{\PowerSetOf}[1]{\ensuremath{
        \boldmath{\mathfrak{P}}
        \left({#1}\right)
}}
\newcommand{\IdealsOf}[1]{\ensuremath{\mathcal{I}{deals}\left({#1}\right)}}
\newcommand{\SpecOf}[1]{\ensuremath{\mathcal{S}{pec}\left({#1}\right)}}
\newcommand{\SpecOnOf}[2]{\ensuremath{\mathcal{V}_{#2}\left({#1}\right)}}
\newcommand{\MaxSpecOf}[1]{\ensuremath{\mathcal{M}{ax}\left({#1}\right)}}
\newcommand{\MaxSpecOnOf}[2]{\ensuremath{\mathcal{M}_{#2}\left({#1}\right)}}
\newcommand{\Then}{\ensuremath{          \Rightarrow          }}
\newcommand{\If}{\ensuremath{          \Leftarrow          }}
\newcommand{\Iff}{\ensuremath{\Leftrightarrow }}
\newcommand{\sbneg}{\ensuremath{∼}}
\newcommand{\IsAnUCRing}[1] {\ensuremath{
        \textsc{UCRng}
        \left\lbrace
            {#1}
        \right\rbrace
}}
\newcommand{\IsACRing}[1]{\ensuremath{
        \textsc{CRng}
        \left\lbrace{#1}\right\rbrace
}}
\newcommand{\IsARing}[1] {\ensuremath{
        \textsc{Rng}
        \left\lbrace{#1}\right\rbrace
}}
\newcommand{\IsAField}[1]{\ensuremath{
        \textsc{Field}\left\lbrace{#1}\right\rbrace
}}
\newcommand{\IsAGroup}[1] {\ensuremath{
        \textsc{Grp}
        \left\lbrace{#1}\right\rbrace
}}
\newcommand{\IsACGroup}[1] {\ensuremath{
        \textsc{AbGrp}
        \left\lbrace{#1}\right\rbrace
}}
\newcommand{\IsAModuleOver}[2] {\ensuremath{
        {#2}-\textsc{Mod}
        \left\lbrace{#1}\right\rbrace
}}
\newcommand{\IsALinearVectorialSpaceOver}[2] {\ensuremath{
        \textsc{LinVect}_{#2}
        \left\lbrace{#1}\right\rbrace
}}
\newcommand{\IsAnAlgebraOver}[3] {\ensuremath{
        \textsc{Alg}_{\IsAModuleOver{#2}{#3}}
        \left\lbrace{#1}\right\rbrace
}}
\newcommand{\FamilyOfAllChainsOf}[1] {\ensuremath{
        \mathcal{F}\mathcal{F}\mathcal{C}\left({#1}\right)
}}
\newcommand{\FamilyOfAllPrimeChainsOf}[1] {\ensuremath{
        {\mathcal{F}\mathcal{F}\mathcal{C}}^{prime}
        \left({#1}\right)
}}
\newcommand{\IsNoether}[1] {\ensuremath{
        \textsc{Noether}
        \left\lbrace{#1}\right\rbrace
}}
\newcommand{\IsNoetherian}[1] {\ensuremath{
        \textsc{Noether}
        \left\lbrace{#1}\right\rbrace
}}
\newcommand{\dimK}[1] {\ensuremath{
        {\textbf{dim}}_{\textsf{Krull}}
        \left(R\right)
}}
\newcommand{\lo}{\ensuremath{ ∨}}
\newcommand{\ly}{\ensuremath{         \wedge         }}
\newcommand{\Union}{\ensuremath{\bigcup}}
\newcommand{\DisjointUnion}{\ensuremath{\biguplus}}
\newcommand{\Intersection}{\ensuremath{\bigcap}}
\newcommand{\DirectSum}{\ensuremath{\bigoplus}}
\newcommand{\TensorProduct}{\ensuremath{\bigotimes}}
\newcommand{\FractionsOfRngOn}[2]{\ensuremath{\RingLocalizedOn{#1}{#2}}}
\newcommand{\NaturalInclusion}[1]{\ensuremath{{\imath}_{\Ideal{#1}}}}
\newcommand{\NaturalInclusionMapsTo}[2] {\ensuremath{
        \UnaryFunct{{\imath}_{\Ideal{#1}}}
        {#2}
}}
\newcommand{\KernelOfNaturalInclusion}[1] {\ensuremath{
        \UnaryFunct{\Ideal{K}}
        {\Ideal{#1}}
}}
\newcommand{\KernelOf}[1]{\ensuremath{
        \UnaryFunct{\textbf{\ker}}
        {#1}
}}
\newcommand{\length}[1]{\ensuremath{\mathbf{length}{\left({#1}\right)}}}
\newcommand{\IsADomain}[1]{\ensuremath{\textsc{Dom}{\lbrace{#1}\rbrace}}}
\newcommand{\IsSubNoetherian}[1]{\ensuremath{
        \textsc{SubNoether}
        {\lbrace{#1}\rbrace}
}}
\newcommand{\IsFG}[2]{\ensuremath{{
            {\boldsymbol{\textsc{fg}}}_{{#2}
            \textsf{-mod}}}{\lbrace{#1}\rbrace}
}}
\newcommand{\IsASet}[1]{\ensuremath{
        {{\textsc{Set}}
        {\lbrace{#1}\rbrace}}
}}
\newcommand{\GPolynomialRingOfOnOver}[3]{\ensuremath{
        {#1}
        \left[
            {
            \lbrace{
                {\mathtt{#2}}_{\imath}
                }
            \rbrace
            }_{
                \imath          \in {\mathcal{#3}}
                }
        \right]
}}
\newcommand{\quotientofideals}[2]{\ensuremath{
        \rfrac{\Ideal{#2}}{\Ideal{#1}}
}}
\newcommand{\quotientringbyideal}[2]{\ensuremath{
        \rfrac{{#2}}{\Ideal{#1}}
}}
\newcommand{\TuplesOfNearbyPrimes}[2]{\ensuremath{
        \mathfrak{#2}{-}\mathcal{N}{ear}\mathcal{P}{rimes}
        {\left({#1}\right)}
}}
\newcommand{\SubChain}[3]{\ensuremath{
        \overline{{\mathcal{#1}}_{\Ideal{#2}}^{\Ideal{#3}}}
}}
\newcommand{\myfrac}[2]{\ensuremath{
        \genfrac{}{}{0pt}{0}{#1}{#2}
}}
\newcommand{\MakeFiniteChains}[3]
{	
    \ensuremath{
		\textsc{FiniteChains}
			{\lbrace{
				{#1}_{\Ideal{#2}}^{\Ideal{#3}}
			}\rbrace}
	}
}
\newcommand{\AnilloGenerado}{\ensuremath{\boldsymbol{\mathcal{A}}}}
\newcommand{\AnilloGeneradoPor}[2]{\ensuremath{
        {\boldsymbol{\mathcal{A}}}_{\left(#1,\mathcal{#2}\right)}
}}
\newcommand{\SetOfSetsOfIdealsOn}[2]{
	\ensuremath{
		\boldsymbol{\mathcal{P}}
		\left(
			\IdealsOf{\left({\AnilloGeneradoPor{#1}{#2}}\right)}
		\right)
	}
}
\newcommand{\bmid}{\ensuremath{\boldsymbol{\;∣\;}}}
\newcommand{\ascchain}[1]{\ensuremath{{\boldsymbol{\mathcal{#1}}}}}
\newcommand{\TOrderWMin}{\ensuremath{
    \underset{{with\,\mathcal{M}{in}}}
    {\mathcal{L}{in}\mathcal{O}{rder}}
}}
\newcommand{\SetOfAscChainsOn}[1]{\ensuremath{
        \mathcal{A}\mathcal{C}\left({#1}\right)
}}
\newcommand{\OrderOn}[3]
{
	\ensuremath
	{
		\left(
			{
				{#1},{#2},{#3}
			}
		\right)
	}
}
\newcommand{\OrderOnIndexes}[2]{\ensuremath{
        \left({\mathcal{#1},{#2}_{0},\leq}\right)
}}
\newcommand{\QueCumplaCon}{\ensuremath{
        \,\boldsymbol{:}\,
}}
\newcommand{\OrderOnIdealsChain}[2]{\ensuremath{
        \left(\ascchain{#1},{\mathfrak{#1}}_{{#2}_{0}},         \subseteq \right)
}}
\newcommand{\OrderIsomorphism}{\ensuremath{
        {         \cong         }_{\TOrderWMin}
}}
\newcommand{\sumab}[3]{\ensuremath{
        \overset{#3}{\underset{{#1}={#2}}{\sum}}
}}
\newcommand{\diff}[2]{\ensuremath{\mathbf{d}{{#1}^{#2}}}}
\newcommand{\tangentfibered}[1]{\ensuremath{
        \mathbf{T}{\left({#1}\right)}
}}
\newcommand{\tangentspace}[2]{\ensuremath{
        {{\mathbf{T}}_{\mathbf{#1}}}{\left({#2}\right)}
}}
\newcommand{\dualtangentspace}[2]{\ensuremath{
        {{\mathbf{T}}_{\mathbf{#1}}}^{\boldsymbol{*}}
        {\left({#2}\right)}
}}
\newcommand{\diffof}[3]{\ensuremath{
        \frac{\boldsymbol{\partial}{#1}}
        {\boldsymbol{\partial}{{#2}^{#3}}}
}}
\newcommand{\esfera}[1]{\ensuremath{\mathbb{S}^{#1}}}
\newcommand{\bola}[1]{\ensuremath{\mathbb{B}^{#1}}}
\newcommand{\espacioproyectivoreal}[1]{\ensuremath{{
            \mathbb{R}\mathbb{P}}^{#1}
}}
\newcommand{\espacioproyectivocomplejo}[1]{\ensuremath{
        {\mathbb{C}\mathbb{P}}^{#1}
}}
% TEOR\'IA DE GRUPOS FINITOS
\newcommand{\MW}[1]{\ensuremath{{\mathsf{MW}}_{#1}}}
\newcommand{\nil}[1]{\ensuremath{{\mathtt{nil}}-{#1}}}
\newcommand{\setP}{\ensuremath{}\mathbb{P}}
\newcommand{\setPgt}[1]{\ensuremath{{\setP}_{>{#1}}}}
\newcommand{\setNgt}[1]{\ensuremath{{\setN}_{>{#1}}}}
\newcommand{\generadoporZitems}{\ensuremath{
        \entreangulos{\pbm{1}}
}}
\newcommand{\generadoporUitem}[1]{\ensuremath{
        \entreangulos{#1}
}}
\newcommand{\generadopor}[1]{\ensuremath{
        \generadoporUitem{#1}
}}
\newcommand{\generadoporDitems}[2]{\ensuremath{
        \entreangulos{#1,#2}
}}
\newcommand{\generadoporTitems}[3]{\ensuremath{
        \entreangulos{ #1,#2,#3}
}}
\newcommand{\generadoporNitems}[3]{\ensuremath{
        \entreangulos{#1,#2,\ldots,#3}
}}
\newcommand{\grA}{\ensuremath{\mathbf{A}}}
\newcommand{\grB}{\ensuremath{\mathbf{B}}}
\newcommand{\grC}{\ensuremath{\mathbf{C}}}
\newcommand{\grD}{\ensuremath{\mathbf{D}}}
\newcommand{\grE}{\ensuremath{\mathbf{E}}}
\newcommand{\grF}{\ensuremath{\mathbf{F}}}
\newcommand{\grG}{\ensuremath{\mathbf{G}}}
\newcommand{\grH}{\ensuremath{\mathbf{H}}}
\newcommand{\grI}{\ensuremath{\mathbf{I}}}
\newcommand{\grJ}{\ensuremath{\mathbf{J}}}
\newcommand{\grK}{\ensuremath{\mathbf{K}}}
\newcommand{\grN}{\ensuremath{\mathbf{N}}}
\newcommand{\grM}{\ensuremath{\mathbf{M}}}
\newcommand{\grL}{\ensuremath{\mathbf{L}}}
\newcommand{\grP}{\ensuremath{\mathbf{P}}}
\newcommand{\grQ}{\ensuremath{\mathbf{Q}}}
\newcommand{\grR}{\ensuremath{\mathbf{R}}}
\newcommand{\grS}{\ensuremath{\mathbf{S}}}
\newcommand{\grT}{\ensuremath{\mathbf{T}}}
\newcommand{\grU}{\ensuremath{\mathbf{U}}}
\newcommand{\subgr}{\ensuremath{<}}
\newcommand{\subgreq}{\ensuremath{\leqslant}}
\newcommand{\charsubgr}{\ensuremath{\stackrel{\mathtt{char}}{\_}}}
\newcommand{\charsubgreq}{\ensuremath{\stackrel{\mathtt{char}}{{}_{=}}}}
\newcommand{\normsubgr}{\ensuremath{         \vartriangleleft         }}
\newcommand{\normsubgreq}{\ensuremath{\trianglelefteq}}
\newcommand{\setofsets}[1]{\ensuremath{\mathcal{#1}}}
\newcommand{\biy}[1]{\ensuremath{{\mathtt{Biy}}{\left({#1}\right)}}}
\newcommand{\biyuno}[1]{\ensuremath{
        {{\mathtt{Biy}}_{1\mapsto1}}
        {\left({#1}\right)}
}}
\newcommand{\graut}[1]{\ensuremath{{\mathtt{Aut}}_{\mathtt{Grp}}{\left({#1}\right)}}}
\newcommand{\grinn}[1]{\ensuremath{{\mathtt{Inn}}_{\mathtt{Grp}}{\left({#1}\right)}}}
\newcommand{\grout}[1]{\ensuremath{{\mathtt{Out}}_{\mathtt{Grp}}{\left({#1}\right)}}}
\newcommand{\grendo}[1]{\ensuremath{{\mathtt{End}}_{\mathtt{Grp}}{\left({#1}\right)}}}
\newcommand{\grmonoendo}[1]{\ensuremath{
        {{\mathtt{End}}_{\mathtt{Mono}}}^{\mathtt{Grp}}{\left({#1}\right)}
}}
\newcommand{\grhom}[2]{\ensuremath{
        {{\mathtt{Hom}}_{Grp}}
        {\left({#1,#2}\right)}
}}
\newcommand{\griso}[2]{\ensuremath{
        {{\mathtt{Iso}}_{Grp}}
        {\left({#1,#2}\right)}
}}
\newcommand{\grbiy}[2]{\ensuremath{
        {\mathtt{Biy}}
        {\left({#1,#2}\right)}
}}
\newcommand{\grlattice}[1]{\ensuremath{
        {\mathcal{L}}_{#1}
}}
\newcommand{\pgrlattice}[2]{\ensuremath{
        {\mathcal{P}}_{#2}
        \entreparentesis{#1}
}}
\newcommand{\pnormgrlattice}[2]{\ensuremath{
        {\mathcal{N}}_{#2}
        \entreparentesis{#1}
}}
\newcommand{\psylowlattice}[2]{\ensuremath{
        {{\mathcal{S}\mathsf{yl}}_{#2}}
        \entreparentesis{#1}
}}
\newcommand{\grcoc}[2]{\ensuremath{
        \nicefrac{#1}{#2}
}}
\newcommand{\normgrlattice}[1]{\ensuremath{{\mathcal{N}}_{#1}}}
\newcommand{\chargrlattice}[1]{\ensuremath{{\mathcal{C}har}_{#1}}}
\newcommand{\maxgrlattice}[1]{\ensuremath{{\mathcal{M}ax}_{#1}}}
\newcommand{\mingrlattice}[1]{\ensuremath{{\mathcal{M}in}_{#1}}}
\newcommand{\supcoregrlattice}[2]{\ensuremath{
        \overline{\mathcal{K}_{#1}}
        {\left({#2}\right)}
}}
\newcommand{\subcoregrlattice}[2]{\ensuremath{
        \underline{\mathcal{K}_{#1}}
        {\left({#2}\right)}
}}
\newcommand{\pSyllattice}[2]{\ensuremath{
        {{\mathcal{S}yl}_{#2}}
        {\left({#1}\right)}
}}
\newcommand{\grclass}[1]{\ensuremath{\mathsf{#1}}}
\newcommand{\cyclicgr}[1]{\ensuremath{{\grclass{C}}_{#1}}}
\newcommand{\trivialgr}{\ensuremath{{\grclass{C}}_{\mathbf{1}}}}
\newcommand{\pd}[2]{\ensuremath{{#1}\!          \times          \!{#2}}}
\newcommand{\pdpot}[2]{\ensuremath{{\left({#1}\right)}^{#2}}}
\newcommand{\psd}[2]{\ensuremath{{#1}\!\rtimes\!{#2}}}
\newcommand{\psdef}[3]{\ensuremath{{#1}\!{\rtimes}_{#3}\!{#2}}}
\newcommand{\pactsd}[3]{\ensuremath{{#1}\!{\rtimes}_{#3}\!{#2}}}
\newcommand{\pxd}[2]{\ensuremath{{#1}\!\bowtie\!{#2}}}
\newcommand{\pactxd}[4]{\ensuremath{{#1}\! {_{#3\!}}{\bowtie}_{#4}\!{#2}}}
\newcommand{\pwr}[2]{\ensuremath{{#1}\wr{#2}}}
\newcommand{\pw}[2]{\ensuremath{{#1}\wr{#2}}}
\newcommand{\pdcyclics}[2]{\ensuremath{\pd{\cyclicgr{#1}}{\cyclicgr{#2}}}}
\newcommand{\pdtrescyclics}[3]{\ensuremath{{\cyclicgr{#1}}\!         \times         \!{\cyclicgr{#2}}\!         \times         \!{\cyclicgr{#3}}}}
\newcommand{\pdquatrocyclics}[4]{\ensuremath{{\cyclicgr{#1}}\!         \times         \!{\cyclicgr{#2}}\!         \times         \!{\cyclicgr{#3}}\!         \times         \!{\cyclicgr{#4}}}}
\newcommand{\psdcyclics}[2]{\ensuremath{\psd{\cyclicgr{#1}}{\cyclicgr{#2}}}}
\newcommand{\psdefcyclics}[3]{\ensuremath{\psdef{\cyclicgr{#1}}{\cyclicgr{#2}}{#3}}}
\newcommand{\pxdcyclics}[2]{\ensuremath{\pxd{\cyclicgr{#1}}{\cyclicgr{#2}}}}
\newcommand{\pwcyclics}[2]{\ensuremath{\pw{\cyclicgr{#1}}{\cyclicgr{#2}}}}
\newcommand{\dihgr}[1]{\ensuremath{{\grclass{Dih}}_{#1}}}
\newcommand{\dicgr}[1]{\ensuremath{{\grclass{Dic}}_{#1}}}
\newcommand{\gdihgr}[2]{\ensuremath{{\grclass{GDih}}_{#1,#2}}}
\newcommand{\sdihgr}[2]{\ensuremath{{\grclass{SDih}}_{#1,#2}}}
\newcommand{\gdicgr}[2]{\ensuremath{{\grclass{GDic}}_{#1,#2}}}
\newcommand{\simgr}[2]{\ensuremath{{\grclass{Sym}}_{#1,#2}}}
\newcommand{\altgr}[2]{\ensuremath{{\grclass{Alt}}_{#1,#2}}}
\newcommand{\gKleingr}[3]{\ensuremath{{\grclass{gK}}_{#1,#2}^{#3}}}
\newcommand{\Egr}[2]{\ensuremath{{\grclass{E}}_{#1}^{#2}}}
\newcommand{\quaterniongr}[1]{\ensuremath{{\grclass{Q}}_{#1}}}
\newcommand{\squaterniongr}[1]{\ensuremath{{\grclass{SQ}}_{#1}}}
\newcommand{\gquaterniongr}[1]{\ensuremath{{\grclass{GQ}}_{#1}}}
\newcommand{\normalizador}[2]{\ensuremath{
        {{\mathsf{N}}_{#1}}
        {\left({#2}\right)}
}}
\newcommand{\centralizador}[2]{\ensuremath{
        {{\mathsf{C}}_{#1}}
        {\left({#2}\right)}
}}
\newcommand{\centro}[1]{\ensuremath{
        {\mathsf{Z}}
        {\left({#1}\right)}
}}
\newcommand{\centronum}[2]{\ensuremath{
        {{\mathsf{Z}}_{#2}}
        {\left({#1}\right)}
}}
\newcommand{\frattini}[1]{\ensuremath{{\mathbf{\Phi }}_{#1}}}
\newcommand{\fitting}[1]{\ensuremath{{\bm{\mathsf{F}}}_{#1}}}
\newcommand{\gfitting}[1]{\ensuremath{{\bm{\mathsf{F}^{*}}}_{#1}}}
\newcommand{\indice}[2]{\ensuremath{\entrecorchetes{{#1}\;:\;{#2}}}}
\newcommand{\grderivado}[2]{\ensuremath{
        {#1}^{\entreparentesis{#2}}
}}
\newcommand{\grLderivado}[2]{\ensuremath{
        {{\mathsf{L}}_{\entreparentesis{#2}}}
        {\entreparentesis{#1}}
}}
\newcommand{\grabelianizado}[1]{\ensuremath{
        {#1}^{\mathsf{ab}}
}}
\newcommand{\conmutador}[2]{\ensuremath{
        \entrecorchetes{#1,#2}
}}
\newcommand{\grconmutador}[2]{\ensuremath{
        \generadopor{
            \entrecorchetes{#1,#2}
        }
}}
\newcommand{\supcore}[2]{\ensuremath{
        {\overline{{\bm{\mathsf{Core}}}_{#1}}}
        {\left({#2}\right)}
}}
\newcommand{\subcore}[2]{\ensuremath{
        {\underline{{\bm{\mathsf{Core}}}_{#1}}}
        {\left({#2}\right)}
}}
\newcommand{\neutromult}{\ensuremath{\mathbf{u}}}
\newcommand{\unity}{\ensuremath{\neutromult}}
\newcommand{\neutroadd}{\ensuremath{\mathbf{e}}}
\newcommand{\zero}{\ensuremath{\neutroadd}}
\newcommand {\gfrac}[2]{\ensuremath{\genfrac{}{}{0pt}{0}{#1}{#2}}}
\setcounter{secnumdepth}{4}
\setcounter{tocdepth}{4}
\synctex=1
\newtheorem{thm}{Teorema}[chapter]
\newtheorem{example}[thm]{Ejemplo}
\newtheorem{exercise}[thm]{Ejercicio}
\newtheorem{axiom}[thm]{Axioma}
\newtheorem{cor}[thm]{Corolario}
\newtheorem{lem}[thm]{Lema}
\newtheorem{prop}[thm]{Proposici\'on}
\theoremstyle{definition}
\newtheorem{defn}[thm]{Definici\'on}
\theoremstyle{remark}
\newtheorem{rem}[thm]{Nota}
\newtheorem{fixed}[thm]{Fijaci\'on de Notaci\'on}
\newtheorem{formalproof}[thm]{Demostraci\'on Formal}
\newtheorem{solution}[thm]{Soluci\'on}
\makeatother
\makeatletter
%\bibliography{bibliografia}
%\usepackage{polyglossia}
\usepackage{csquotes}
\usepackage{fancyvrb}
\usepackage{listings}
\lstset{ %
    language=C++,                	% choose the language of the code
    basicstyle=\footnotesize,       % the size of the fonts that are used for the code
    numbers=left,                   % where to put the line-numbers
    numberstyle=\footnotesize,      % the size of the fonts that are used for the line-numbers
    stepnumber=1,                   % the step between two line-numbers.
                                    % If it is 1 each line will be numbered
    numbersep=5pt,                  % how far the line-numbers are from
                                    %the code
    backgroundcolor=\color{white},  % choose the background color.
                                    %You must add \usepackage{color}
    showspaces=false,               % show spaces adding
                                    %particular underscores
    showstringspaces=false,         % underline spaces within strings
    showtabs=false,                 % show tabs within strings
                                    %adding particular underscores
    frame=single,           		% adds a frame around the code
    tabsize=2,          			% sets default tabsize to 2 spaces
    captionpos=b,           		% sets the caption-position to bottom
    breaklines=true,        		% sets automatic line breaking
    breakatwhitespace=false,    	% sets if automatic breaks should
                                    %only happen at whitespace
    escapeinside={\%*}{*)}          % if you want to add a
                                    %comment within your code
}

\makeatother

\makeatletter

\begin{document}

\title{
    Clasificaci\'on salvo isomorfismo de los grupos finitos hasta el orden 22\\
    {{\normalsize{T.F.G. del grado en Matem\'aticas de la U.N.E.D.}}}
}

\author{
    Juli\'an Calder\'on Almendros\\
    \textsc{\small{TUTORIZADO POR} \textbf{Emilio Bujalance Garc{\'{\i}}a}}
}

\date{\today}

\maketitle


\makeatletter

\hypersetup{pdftitle={Clasificaci\'on de los grupos hasta el orden 22},
pdfauthor={Juli\'an Calder\'on Almendros},
pdfsubject={Teor{\'{\i}}a de Grupos Finitos},
pdfkeywords={\'Algebra, Grupos Finitos, Matem\'aticas Discretas, \'Algebra Computacional}}
\tableofcontents{}
\makeatletter
\listoffigures
\makeatletter
\listoftables
\makeatletter

%\printnomenclature{}
%\makeatletter


\chapter{Introducci\'on}

En este trabajo para el fin de grado de Matem\'aticas, doy una clasificaci\'on, salvo
isomorfismos, de los grupos de los primeros \'ordenes, hasta el orden $22$ incluido.
Adem\'as de cumplir esta petici\'on inicial, requerida por el mismo t{\'{\i}}tulo, he intentado
darle el car\'acter lo m\'as constructivo posible, de forma que, en vez de tratar el orden $9$,
he tratado la clase de \'ordenes $p^2$, siendo $p$ un primo, o en vez de los grupos $15$ y $21$,
los grupos de la clase $pq$ d\'onde $p$ y $q$ son primos diferentes y mayores que $2$
(el caso del primo $2$ parece preferible tratarlo por separado por su numerosa aparici\'on,
adem\'as de haber diferencias si $p\vert q-1$),
as{\'{\i}} los tipos $4p$, en vez de $12$ y $20$ (hay ciertas diferencias entre el caso $12$ y los
dem\'as de este tipo). Y as{\'{\i}} los que se fuesen presentando.

\section{La herramienta de computaci\'on}

Si digo constructivo y no general (que tambi\'en se podr{\'{\i}}a decir), es por la derivaci\'on que
he tomado en este proyecto: hacer el software que me hac{\'{\i}}a falta. El software GAP es, desde luego,
impresionante, con colaboraciones de casi todos los dedicados a la teor{\'{\i}}a de grupos y teor{\'{\i}}as algebraicas
aledañas, tanto en enseñanza como en investigaci\'on. Pero yo llevo alguna desventaja, que en \'ultimo
t\'ermino puedo convertir en una ventaja: conozco desde hace mucho algunos lenguajes de programaci\'on.
Acostumbrado al lenguaje C++ pronto me di cuenta que iba a tener que dedicarle un tiempo {\'{\i}}mprobo a
conocer los entresijos de GAP para poder hacer cosas como las que yo quer{\'{\i}}a hacer (probablemente, de
forma muy sencilla finalmente, cuando hubiese aprendido GAP a suficiente nivel): escribir en un fichero
la informaci\'on m\'as importante de cada clase de equivalencia de un grupo finito.

Esto viene de la mano de utilizar LaTeX para describir todos los grupos de la forma m\'as gr\'afica posible.
Sab{\'{\i}}a que en C++ pod{\'{\i}}a hacer f\'acilmente ese tipo de tareas. Adem\'as, inicialmente, pens\'e que no necesitaba
que el c\'odigo fuese muy eficiente ni tampoco ten{\'{\i}}a porqu\'e buscar los algoritmos m\'as \'optimos conocidos.
Al contrario, lo que necesitaba eran solo las definiciones. Por ejemplo, aunque pudiese utilizar un algoritmo
que probara la asociatividad (o la no asociatividad) de un conjunto con una determinada operaci\'on binaria
interna con un algoritmo de complejidad $n^2$ ($n$ es el n\'umero de elementos del grupo y $n^2$ la
complejidad en tiempo o n\'umero de operaciones), yo prefer{\'{\i}}a la
claridad del algoritmo que la definici\'on de asociatividad induce directamente y tiene complejidad $n^3$.
En \'ultimo t\'ermino, si llego a computar un grupo de orden 64 y quiero comprobar si una determinada
operaci\'on de un producto semidirecto est\'a correctamente desarrollada, $64^3=262400$ operaciones,
que con ser muchas, es un precio que podemos pagar, quiz\'as una fracci\'on grande de segundo. Escribir la tabla
de la operaci\'on en el disco duro seguro que tarda m\'as.

Por otra parte encontr\'e desde el principio un problema para trabajar con GAP: no hab{\'{\i}}a un IDE adecuado
(se puede utilizar \textsc{notepad++}, que detecta elementos del lenguaje, sin ser ni mucho menos un IDE).
Sin embargo C++ tiene un buen n\'umero de ellos, y bastante amigables, como \textsc{Code::Blocks},
\textsc{CodeLite} y otros.

Este software que he escrito, adjunto, tiene que depurarse mucho a\'un y limpiarse m\'as de c\'odigo innecesario.
Con todo ha producido todos los datos del presente trabajo fin de grado.

Adem\'as, hay que decir que las librer{\'{\i}}as Boost (C++) en sus \'ultimas versiones contiene una librer{\'{\i}}a de programaci\'on
funcional y tipado fuerte que se llama \textsc{boost::hana} y que ya contiene el tipo \textsc{boost::hana::group},
y que si no he utilizado es porque  me obligaba a gastar demasiado tiempo en aprender a utilizar unas librer{\'{\i}}as
algo cr{\'{\i}}pticas. Con todo, estas
librer{\'{\i}}as han de encontrarse cerca (por el mismo camino) de lo que he estado codificando. Es muy probable,
que si este software lo sigo desarrollando incorpore ya estas librer{\'{\i}}as. De hecho he utilizado tipado fuerte para
mi software y ``metaprogramaci\'on de plantillas'' (la que hay disponible en C++ y que est\'a avanzando a grandes pasos
en los \'ultimos est\'andares, C++11, C++14, C++17 y que parece que en su versi\'on C++20 har\'a mucho m\'as sencillo
todo este tipo de programaci\'on por la entrada del metatipo \textsc{std::concept}).

Para resumir, cada grupo finito c{\'{\i}}clico est\'a capturado en la clase \textsc{int\_mod\_N}$<N>$, d\'onde
\textsc{N} significa el m\'odulo
de la relaci\'on de equivalencia, que ya utilic\'e para representar n\'umeros en base arbitraria en forma de
\textsc{std::basic\_string}$<$\textsc{int\_mod\_N}$<N>>$, esto es,
con el mismo metatipo de una cadena de caracteres,
d\'onde tipo car\'acter se han cambiado por el tipo mencionado.

Antes de seguir dir\'e que desde el principio ten{\'{\i}}a inter\'es en representar los grupos finitos en productos cartesianos
de tipos modulares como los que he dicho. En otras palabras, quer{\'{\i}}a representar los grupos finitos pequeños sobre
los productos directos de grupos c{\'{\i}}clicos, y que estos grupos estuviesen lo m\'as ajustados posibles. Me explicar\'e,
por ejemplo, un conjunto de matrices cuadradas de dimensi\'on $2^3$ sobre un cuerpo finito como $\mathds{F}_{3}$
puede representar un grupo, por ejemplo de orden $6$, pero esta representaci\'on tiene una capacidad $3^8=6561$
cuando realmente estamos utilizando solo $6$ lugares. Otra posibilidad son las permutaciones, tal y como lo hace GAP,
pero por alguna raz\'on, que es f\'acil de imaginar y de expresar,
siempre he preferido el producto cartesiano a las permutaciones, es
m\'as ``est\'atico'', se deja ``ver'' mejor. Con todo, el realizar unos sistemas (librer{\'{\i}}as C++) de representaci\'on
sobre espacios vectoriales complejos, sobre espacios vectoriales cuyo cuerpo sea finito o sobre permutaciones,
no tiene gran problema, y ser{\'{\i}}a muy bonito que, por ejemplo, el tipo grupo alternado $4$, fuese un tipo que cobre
entidad solo al decirle el tipo de representaci\'on, de forma que se trabajase en este grupo con la representaci\'on m\'as
adecuada (o requerida por la finalidad que se tenga), y que en su tipo no se distinguiese
por la forma de representarse. Esto veo que es casi como parece que
funciona en GAP.


?`Qu\'e clase de informaci\'on sobre un grupo le pido al programa que me d\'e?

Al trabajar
sobre un grupo en el presente proyecto TFG, lo que tengo es b\'asicamente su conjunto subyacente que se concreta
en un tipo
de producto cartesiano con una aritm\'etica modular en cada componente del producto, esto es, una tupla de n\'umeros
fijan parte del tipo del grupo, concretamente,
\textsc{element\_of\_set\_t<$N_n, N_{n-1},\ldots,N_1$>}. Sobre esta estructura
añadimos una operaci\'on, concretamente:

\textsc{group\_binop\_t<group\_name,$N_n, N_{n-1},\ldots,N_1$>},


d\'onde
\textsc{group\_name} es solo un nombre que damos a grupos o familias de grupos en una enumeraci\'on, por ejemplo,
\textsc{group\_name::$\mathsf{C}$},
indica que es un grupo c{\'{\i}}clico, \textsc{group\_name::$\mathsf{Dih}$}, indica que es un grupo
dih\'edrico, y as{\'{\i}} una variedad.

A partir de aqu{\'{\i}} tenemos el elemento de grupo propiamente dicho,
\textsc{element\_of\_group\_t}
$\langle{group\_name},{N_n, N_{n-1}\ldots N_1}\rangle$.
Las operaciones que tenemos sobre estos elementos son las esperadas:
 \textsc{bop}$(a,b)$, \textsc{commutator}$(a,b)$, \textsc{commute}$(a,b)$, $a.$\textsc{conjugate}$(b)$,
$a.$\textsc{elem\_order}$()$. $a.$\textsc{inv}$()$.

El grupo se construye como un contenedor de
la bilioteca est\'andar C++ (\textsc{std::set}), un conjunto de tipo homog\'eneo,
de tipo el elemento de grupo anteriormente dicho.
B\'asicamente hereda todas las propiedades y operaciones anteriores.


El tipo es ahora \textsc{group\_t<group\_name,$N_n,
N_{n-1},\ldots,N_1$>}.


En este punto tenemos ya los test de propiedades, \textsc{is\_associative()}, \textsc{has\_unit\_element()} (neutro),
\textsc{exists\_inverse\_for\_every\_element()},
y con esto ya sabemos si la operaci\'on conforma un grupo, \textsc{is\_abelian()} como
propiedad inmediata. Adem\'as, dada una instancia \textsc{group\_t} del mismo tipo que el grupo completo,
pero que es subconjunto, podemos comprobar adem\'as si este subconjunto tiene el inverso de cada uno
de sus elementos en el propio subconjunto, \textsc{is\_inner\_inverter()}, o si cada pareja de elementos del subconjunto
al operarse tiene el resultado dentro del propio subconjunto, \textsc{is\_inner\_binary\_operation()}: en definitiva
si el subconjunto es subgrupo.

Adem\'as se puede operar con la operaci\'on del grupo, \textsc{bop(a,b)}, obtener el inverso de un elemento
del grupo, \textsc{inv()}, o llenar la instancia del grupo con todos sus elementos, \textsc{fill\_the\_group()}, o vaciarlo,
\textsc{do\_empty()}.

Adem\'as podemos llamar a \textsc{G.L\_derived(n)} (derivaci\'on de la serie central descendente),
 \textsc{G.derived(n)} (la derivaci\'on del subgrupo derivado y sucesivos),
 \textsc{H.centralizer()}, \textsc{G.centre()}, \textsc{G.is\_solvable()},
\textsc{G.is\_nilpotent()}, entre muchas otras utilidades de teor{\'{\i}}a de grupos.

A partir tenemos la clase \textsc{subgroup\_lattice\_t<gn,$N_n$,$N_{n-1}$...>}, que se puede montar f\'acilmente
porque hemos dotado a los grupos (conjuntos en realidad) de relaciones de orden e igualdad, al igual que hemos
hecho con todos los tipos
hasta ahora, con la excepci\'on de esta representaci\'on del ret{\'{\i}}culo de subgrupos y de la representaci\'on de la
operaci\'on binaria interna. Como es de prever, aqu{\'{\i}} podremos obtener los subgrupos de determinado orden, etc.

Todas los tipos anteriores est\'an provistos de un m\'etodo, \textsc{to\_string(abecedario) }, que nos dar\'a en una cadena
de caracteres la representaci\'on gr\'afica, ya pasada a formato LaTeX del objeto correspondiente. La clase
\textsc{group\_binop\_t} no tiene esta posibilidad. Sin embargo, para codificar un grupo solo hay de
suministrar una especializaci\'on de esta clase para el grupo concreto que estemos tratando,
\memph{esto es lo \'unico que tiene que suministrar el usuario de la librer{\'{\i}}a}. En este TFG se utiliza de esta
forma la librer{\'{\i}}a, se estudia el grupo y cuando se tiene la operaci\'on de grupo o clase de grupos se introduce en su
metatipo de operaci\'on de grupo y ya est\'a todo lo dem\'as hecho.

Finalmente tenemos la clase-funci\'on \textsc{get\_group\_info<gn,$N_n$,$N_{n-1}$...>},
que tiene solo una funci\'on miembro,
\textsc{work(std::string,std::string,std::string)}, a la que se le dicen el nombre de fichero, el abecedario y una entrada
en LaTeX para personalizar la salida a archivo.

\section{La procedencia de la informaci\'on}

En estos cap{\'{\i}}tulos que ahora escribo quiero dejar sentado los ``hechos''
que tomar\'e como ya probados, mayormente provenientes de la Unidad
Did\'actica I de la asignatura de \'Algebra de 2"o curso de la antigua
licenciatura en Matem\'aticas, libro dedicado precisamente a la teor{\'{\i}}a
de grupos y escrito por Emilio Bujalance Garc{\'{\i}}a\cite{EBujalance1989}.
Adem\'as me he basado en los cuatro art{\'{\i}}culos siguientes: \cite{Wilde2005} y
\cite{Clausen2012}.

Tambi\'en me he informado y he tomado elementos en varios sitios web, como son,
\cite{GroupProps},
\cite{WIKIPEDIAen},
\cite{WIKIPEDIAes},
\cite{ProofWiki},
\cite{GroupNames}
y con gran profusi\'on \cite{MATHStackExchange}.

Para la parte computacional, \cite{GAP} y numerosas p\'aginas de
programaci\'on en C++, para resumir la m\'as empleada ha sido
\cite{CppProgStackExchange}, \cite{CPPReference}, \cite{CPlusPlus}, y
algunos blogs. Los libros referencia utilizados son ``C++ Language Programming'',
la \'ultima versi\'on que cubre el est\'andard C++11\cite{Stroustrup2013},
``The Boost C++ Libraries'' de Boris Shcäling \cite{Schaeling2011}, y
``C++ Template Metaprogramming'' de David Abrahams y Alesksey Gurtovoy
\cite{DavidAbrahams2005}.

\section{La parte te\'orica}

Esta parte la he trabajado, como ya he dicho sobre todo desde la unidad did\'actica \cite{EBujalance1989}, sin embargo he curioseado por muchos sitios, y en otros libros como \cite{Burnside1897} y \cite{Hall1959}. Para el orden $16$ he utilizado los art{\'{\i}}culos \cite{Clausen2012}, \cite{Wilde2005}.

Mi inter\'es por el producto semidirecto y su generalizaci\'on m\'as directa, el producto de Zappa-Szèp (ZS para abreviar en adelante), me han inclinado a hallar todos los grupos como una construcci\'on sobre los grupos c{\'{\i}}clicos simples, y a pensar en los grupos como extensiones de grupos de \'ordenes m\'as bajos.

Por eso comienzo dando una versi\'on del producto Zappa-Szèp, describiendo sus propiedades, y despu\'es veo el producto semidirecto como un caso del anterior y el producto directo de grupos como el producto directo de acci\'on trivial.

\chapter{Algunas definiciones}

Comenzaremos directamente. Los grupos c{\'{\i}}clicos de un determinado orden
$n$ lo denominaremos $\cyclicgr{n}$. El grupo trivial ser\'a denotado
por $\trivialgr$. Los homomorfismos entre dos grupos, digamos
de $G$ en $L$ ser\'an denotados como $\grhom{G}{L}$. Si estos son biyecciones,
esto es, isomorfismos,
se denotar\'a como $\griso{G}{L}$. Si los homomorfismos son de un grupo
en s{\'{\i}} mismo pondremos $\grendo{G}$ y si son biyecciones,isomorfismos,
$\graut{G}$.
Tambi\'en
ser\'an de inter\'es las biyecciones de un conjunto en s{\'{\i}} mismo $\biy{S}$
puesto que constituye un grupo bajo la composici\'on de funciones como operaci\'on.
Un subgrupo importante de este \'ultimo ser\'a $\biyuno{\grG}$, las biyecciones
de un grupo sobre si mismo que dejan invariante la unidad.


Las operaciones entre grupos que utilizaremos ser\'an $\rtimes$
para el producto semidirecto, d\'onde, si escribimos $N\rtimes H$ siendo
$N \normsubgreq G$ subgrupo normal del grupo en el que est\'an y
$H \subgreq G$ simplemente subgrupo. Esta forma de nombrar el
producto semidirecto es imprecisa, as{\'{\i}} que siendo
$\varphi          \in \grhom{H}{\graut{N}}$,
el simbolismo completo ser\'a $\pactsd{N}{H}{\varphi }$. Cuando ambos subgrupos
sean normales simplemente escribiremos el grupo como $\pd{N}{H}$, coincidiendo
con $\pactsd{N}{H}{id}$, que es el producto
directo (no confundir el grupo producto directo con el conjunto
base del grupo, que es el producto cartesiano, y que escribimos igual).
A su vez el producto semidirecto ser\'a considerado como un
caso particular del producto ZS o ``knit'' o ``bicross'', que solo requiere
que ambos subconjuntos
sean subgrupos, as{\'{\i}} que necesitaremos dos homomorfismos,
$\alpha         \in \grhom {N}{\biy{H}}$
y $\beta         \in \grhom {H}{\biy{N}}$, denot\'andose como
$\pactxd {N}{H}{\alpha }{\beta }$. M\'as adelante desarrollaremos las relaciones
entre estos productos
y algunas de sus propiedades, y veremos las diferencias entre las
versiones externa e interna de ellos. Adem\'as, cuando un subgrupo de
$G$, pongamos $H \subgreq G$, podemos definir dos relaciones
de equivalencia $\mathcal{R}_{H}$ y $\mathcal{R}^{H}$, y el conjunto
de particiones para una de ellas ser\'a
$\grcoc{G}{{\mathcal{R}}_{\mathrm{H}}}
\subgreq\extset {gH}{g          \in G}$ para la primera,
y para la otra
$\grcoc{G}{{\mathcal{R}}^{\mathrm{H}}}\subgreq\extset{Hg}{g        \in G}$.
Cuando $H\normsubgreq G$ lo denotaremos por $\grcoc{G}{H}$
y adem\'as constituir\'a un grupo.

Cuando un grupo sea hamiltoniano usaremos el s{\'{\i}}mbolo $\quaterniongr{}$ (ya
que tiene alguna copia de $\quaterniongr{8}$). Son importantes estos grupos porque
son los primeros que no se pueden construir a trav\'es de los productos semidirectos antes citados
(al menos en
sus formas m\'as elementales). Encontraremos m\'as casos de grupos no constructibles
mediante los productos mencionados.

Hay a\'un dos productos importantes que citar: uno es el producto wreath,
denotado $\pwr{H}{N}$, y que es un producto semidirecto a partir de
$\sigma         \in \grhom {H}{\biy{N}}$, d\'onde, si $h         \in H$
tenemos que $\sigma_{h}         \in \biy{N}$
se forma de una manera especial.

El otro producto que nos queda
es el siguiente. Sean $H$ y $K$ grupos, ($G=HK$ es la forma aproximada del grupo que buscamos construir),
tales que tienen un subgrupo cada uno
de ellos $I \subgreq \centro{H}$ y $J \subgreq \centro{K}$
tales que podemos encontrar un isomorfismo $\delta         \in \griso{I}{J}$
con el que podemos crear $L\normsubgreq\extset{(i,j)} {i         \in I,j         \in J,
\delta (i)*j=\unity}$.
Ahora el producto central es $\grcoc{\pd{H}{K}}{L}$. O m\'as generalmente
$\grcoc{\psd{H}{K}}{L}$ o incluso  $\grcoc{\pxd{H}{K}}{L}$.

De los subgrupos de un grupo dado nos interesan el ya mencionado $\centro{G}$
o centro, en general $\centralizador{G}{H}$, siendo $H \subgreq G$ y
el normalizador del anterior subgrupo en el grupo total,
$\normalizador{G}{H}$. Para definir
los siguientes subgrupos, vamos a definir las principales familias
de subgrupos.
$\grlattice{G}\mathdef\extset {H         \subseteq G} {H \subgreq G}$,
$\normgrlattice{G}\mathdef\extset {H         \in \grlattice{G}}{H\normsubgreq G}$.
Desde aqu{\'{\i}} podemos definir
$\supcoregrlattice{G}{K}\mathdef\extset{H         \in \normgrlattice{G}}{
K         \subseteq H}$,
y la definici\'on de un grupo normal
$\supcore{G}{K}\mathdef\bigcap\supcoregrlattice{G}{K}$
y de otro
$\subcoregrlattice{G}{K}\mathdef\extset{H          \in \normgrlattice{G}}{H         \subseteq K}$
y el correspondiente subgrupo normal
$\subcore{G}{K}
\mathdef
\generadopor{
    \subcoregrlattice{G}{K}
}$.
Tambi\'en nos quedan los subgrupos carácter{\'{\i}}sticos, invariantes ante
isomorfismos cualquiera, $ H \charsubgreq G $. As{\'{\i}} definimos
$\chargrlattice{G}\mathdef{\set{H          \in \normgrlattice{G}\vert{H \charsubgreq G}}}$.
Adem\'as, ser\'a interesante el conjunto de los subgrupos maximales de $G$,
$\maxgrlattice{G} \mathdef
\extset {H          \in \left({\grlattice{G}}\setminus\set{G}\right)}
        {
                     \exists          K          \in {\grlattice{G}}\setminus\set{G}\;
                                {H         \subseteq K         \Rightarrow          H=K}
        }
$.
Para subgrupos minimales tenemos dualmente
$\mingrlattice{G} \mathdef
\extset {H         \in {\grlattice{G}\setminus\set{\trivialgr}}}
{        \exists             K            \in {\grlattice{G}\setminus\set{\trivialgr}}
    \; K         \subseteq H          \Rightarrow            H=K}$.
Llamamos subgrupo de Frattini a $\frattini{G}\mathdef        \cap \maxgrlattice{G}$.

De otro lado, un operador binario com\'un entre elementos de un grupo es el \emph{conmutador},
$\conmutador{a}{b} \mathdef a^{-1}        \cdot b^{-1}        \cdot a         \cdot b$
(si $a$ y $b$ conmutan $\conmutador{a}{b}=\unity$).
As{\'{\i}} $\conmutador{H}{K}\mathdef \extset{\conmutador{h}{k}}
{h         \in H    \ly    k         \in K}$.
El conjunto conmutador as{\'{\i}} definido de dos subgrupos no tiene porqu\'e
ser subgrupo, as{\'{\i}} se suele considerar $\grconmutador{H}{K}$
que es el subgrupo generado por el conjunto conmutador.
A partir de aqu{\'{\i}} tenemos dos operadores
sobre grupos. El primero es el subgrupo \emph{derivado} de diversos \'ordenes,
$\grderivado{\grG}{0}\mathdef \grG$, $\forall i         \in \setN\setminus\set{0}; \grderivado{\grG}{i+1}\mathdef{\grconmutador{\grderivado{\grG}{i}}{\grderivado{\grG}{i}}}$.
Otro operador importante es
$\grLderivado{\grG}{0}\mathdef \grG$,
$\grLderivado{\grG}{1}\mathdef \grderivado{\grG}{1}$,
$
\forall i         \in \setN\setminus\set{0,1}\;\;
\grLderivado{\grG}{i+1}\mathdef
\grconmutador{\grLderivado{\grG}{i}}{\grG}
$.


Se dice que un grupo es \emph{nilpotente} si hay un natural $n          \in \setN$ tal que
$\grLderivado{\grG}{n}=\trivialgr$. Un grupo es \emph{nilpotente} de clase $n$ si
$\grLderivado{\grG}{n-1}          \neq          \trivialgr$ y $\grLderivado{\grG}{n}=\trivialgr$.

Igualmente, se dice de un grupo $\grG$ que es resoluble de clase $n$ si
$\grderivado{\grG}{n-1}         \neq         \trivialgr$ y $\grderivado{\grG}{n}=\trivialgr$.

Un grupo se dice el \emph{abelianizado} del grupo $\grG$ cuando nos referimos al cociente
de $\grG$ por el primer subgrupo derivado $\grderivado{\grG}{1}$, esto es
$\grabelianizado{\grG}\mathdef \grcoc{\grG}{\grderivado{\grG}{1}}$.

As{\'{\i}} un grupo es \emph{perfecto} si su abelianizado es trivial,
$\grabelianizado{\grG} = \trivialgr$, o dicho de otra manera,
$\grG = \grderivado{\grG}{1}$.

Un grupo $\grG$ se dice \emph{cuasisimple} si es perfecto, $\grG = \grderivado{\grG}{1}$, y
$\grcoc{\grG}{\centro{\grG}}$ es simple. Hay que entender que
$\grcoc{\grG}{\centro{\grG}}          \cong          \grinn{\grG}          \subseteq \graut{\grG}$,
luego la condici\'on
anterior equivale a que $\grinn{\grG}$ sea simple.

El subgrupo de Fitting de $\grG$, $\fitting{\grG}$, se define como
$\fitting{\grG}
\mathdef
\generadopor{
    \extset{H         \in \normgrlattice{\grG}}
    {\grH \textsf{ nilpotente}}
}
$.

Ahora definiremos un grupo \emph{semisimple}. Para esto tiene que existir
$r          \in \setN\setminus\set{0}$ tal que existe
$\mathcal{S} \mathdef \extset{G_i       \in \grlattice{\grG}}
{\forall i\; 0<i<r       \Rightarrow          G_i \textsf{ es cuasisimple}}\rbrace$,
y $\generadopor{\bigcup\mathcal{S}}=\grG$. Dicho en palabras, un grupo es semisimple si existe
una subfamilia finita de subgrupos cuasisimples que lo genere.

Un grupo es \emph{soluble} si existe una subfamilia de $\grlattice{G}$ ordenada
como cadena (por ejemplo creciente) por la inclusi\'on, que comience en el
subgrupo trivial y termine en el impropio, tal que ya no quepan m\'as
subgrupos intermedios, m\'as la condici\'on de que cada grupo de la cadena sea
normal al previo y el cociente sea siempre abeliano.
Esta cadena ir{\'{\i}}a desde
$\trivialgr$ hasta $\grG$. Digamos que es
$\grH_0 \mathdef \trivialgr \subgr \grH_1 \subgr \ldots
\subgr \grH_{n-1} \subgr \grH_n \mathdef \grG$, de tal forma que
$\forall i             \in \setN\;0             \neq             i\leq n             \Rightarrow
\grcoc{\grH_i}{\grH_{i-1}} \textsf{ es abeliano}$.

Un subgrupo $\grH$ de $\grG$ es \emph{subnormal} si existe una subfamilia de
$\grlattice{\grG}$ finita y ordenada en una cadena por inclusi\'on, de
tipo ascendente, tal que
$\grH_0 \mathdef \grH\subgr H_1 \subgr \ldots \subgr \grH_{n-1} \subgr \grH_n \mathdef \grG$, $n$ m{\'{\i}}nimo.
Es subnormal de clase $n$.

Un subgrupo $\grH$ de $\grG$ se dice \emph{componente} si es subnormal y cuasisimple.

Llamamos ``layer'' de un grupo $\grG$ al producto (que conmuta) de sus componentes.
Tambi\'en podemos definirlo como el subgrupo normal semisimple m\'as grande que hay en $\grG$.
Digo m\'as grande (o m\'aximo) y no maximal porque es \'unico. Se suele denominar $\mathsf{E}(G)$.

El subgrupo de Fitting generalizado es
$\gfitting{\grG} \mathdef \fitting{\grG} \mathsf{E}(\grG)$.

Sobre el cardinal del grupo finito $\grG$, lo denominamos por
$\textsf{card }\grG = \left| \grG \right| = \mathsf{o}(\grG)$, el orden de un grupo.
Para los subgrupos utilizamos la misma denominaci\'on, y hablaremos
habitualmente de orden de un grupo. Para un grupo $\grH$ cualquiera
de $\grG$, hablaremos del {\'{\i}}ndice de un grupo cuando queramos decir
el cardinal de conjunto de las clases de equivalencia lateral inducidas,
$\indice{\grG}{\grH} = \left| \grcoc{\grG}{\mathcal{R}_H} \right| =
\left| \grcoc{G}{\mathcal{R}^H} \right| $.

Ahora tenemos el teorema de Lagrange, que nos dice que si $\grG$ es finito,
$\grH\subgreq\grG$ entonces ${\mathsf{o}(\grH)} \,\shortmid\, {{\mathsf{o}}(\grG)}$.
Esto nos da en principio un punto de conexi\'on con el teorema fundamental de la aritm\'etica.

De esta forma se hacen importantes los \emph{p-grupos} y los
\emph{p-subgrupos}, que son en general aquellos grupos finitos cuyo
orden es $p^n$ donde $p$ es un n\'umero primo y
$n             \in \setN\setminus\lbrace 0\rbrace$.

Llamamos $
{\psylowlattice{\grG}{p}}
\mathdef
{\extset
	{\grH          \in \grlattice{\grG}}	
	{
		{\left| {\grH} \right|}={p^n}\,\ly\,
		{p^n} ∣{\left| {\grG} \right|}\,\ly\,
		{{p}^{n+1}} \nmid {\left| \grG \right|\,\ly\, p  \in \setP}
	}
}$

As{\'{\i}} llamamos \emph{p-n\'ucleo} de $G$ a
$\mathsf{O}_{p}(\grG) \mathdef \bigcap {\psylowlattice{\grG}{p}}$.

Emplearemos adem\'as los teoremas de Sylow, el teorema de Cauchy, el de Lagrange y
dem\'as hechos que iremos nombrando, aparte de los habituales sobre productos directos,
semidirectos y otros que sirven para extender grupos, ya nombrados.

Como el objetivo del presente trabajo es la clasificaci\'on los grupos de
determinados \'ordenes, ser\'a fundamental en este trabajo el \emph{Teorema de
Clasificaci\'on de los Grupos Finitos Abelianos}, que nos hace el trabajo
 ya completo para los grupos abelianos,
 dej\'andonos concentrarnos en los grupos no-abelianos.

\section{Enumeraci\'on de los distintos resultados}

El primero ser\'a el teorema de Lagrange.
\begin{thm}{Teorema de Lagrange\\}{\label{ThLagrange}}
    Si $\grH\subgr\grG$ entonces $\mathbf{o}\left(\grH\right) ∣\mathbf{o}\left(\grG\right)$
\end{thm}

Seguiremos por el teorema de Cauchy.
\begin{thm}{Teorema de Cauchy\\}{\label{ThCauchy}}
    Si $         \exists          p         \in \setP\; : \; p ∣\mathbf{o}\left(\grG\right)$ entonces
    $\left(         \exists          g         \in G\; : \; g^p=u \right)\ly \left(0<i<p \then g^i         \neq          u\right)$
\end{thm}

El siguiente teorema de correspondencia no es el de Cayley, esto es, no es de los grupos como correspondientes a subgrupos de permutaciones.
\begin{thm}{Teorema de Correspondencia entre los ret{\'{\i}}culos de un grupo y los de un grupo cociente\\}{\label{ThCorresp}}
   Sea $\grH \normsubgreq \grG$, entonces existe una correspondencia biyectiva $\varphi: \grlattice{\grcoc{\grG}{\grH}}\To \mathcal{G}_\grH\mathdef\extset{\grK         \in \grlattice{G}}{\grH         \subseteq \grK}$, con la forma concreta, $\varphi \entreparentesis{\grL}= \generadopor{\grL\grH}$
\end{thm}

Teoremas de Sylow\\
\begin{thm}{Teoremas de Sylow\\}{\label{ThSilow}}
    Sea $\grG$ un grupo cuyo cardinal u orden es $\left|{\grG}\right|=p^nm$ d\'onde $p         \in \setP$ y $p\nmid m$.
    Entonces son verdaderas las siguientes afirmaciones:
    \begin{itemize}
        \item[1] $  \exists          \grH         \in \grlattice{\grG} \left|\grH\right|=p^n$
        \item[2] $\forall \grK,\grH          \in \psylowlattice{\grG}{p} \quad          \exists          g         \in \grG\quad H^g=K$
        \item[3] Sea $n_p\mathdef{\left|\psylowlattice{\grG}{p}\right|}$. Entonces permanecen las siguientes afirmaciones:
        \begin{itemize}
            \item[3.1] $\forall \grP         \in \psylowlattice{\grG}{p}\quad n_p ∣m=\indice{\grG}{\grP}$
            \item[3.2] $n_p {         \equiv         }_{p} 1$
            \item[3.3] $\forall \grP         \in \psylowlattice{\grG}{p}\quad n_p =\indice{\grG}{\normalizador{\grG}{\grP}}$
        \end{itemize}
    \end{itemize}
\end{thm}


Teorema de clasificaci\'on de los grupos abelianos finitos.\\
\begin{thm}{Teorema de Clasificaci\'on de los Grupos Abelianos Finitos\\}{\label{ThClGAF}}
    Sea $\grG$ un grupo abeliano finito cuyo orden es $\left|\grG\right|=n$. Entonces todos los grupos abelianos no isomorfos son productos directos de grupos c{\'{\i}}clicos, y adem\'as se pueden enumerar sin repetici\'on en funci\'on de determinadas factorizaciones de $n$. Sean $c_i^j$ una cadena de n\'umeros naturales mayores que $1$, ordenados de forma decreciente, d\'onde $i$ establece la longitud de la cadena, y $j$ es un {\'{\i}}ndice dentro de las factorizaciones de $n$ en $i$ factores, de forma que las $j$ $i-tuplas$ son diferentes entre s{\'{\i}}. Estas cadenas tienen la siguiente representaci\'on:
    \[
        c_i^j = (\nu_i,\nu_{i-1},\ldots,\nu_k,\ldots,\nu_2,\nu_1)
    \]
    \begin{itemize}
        \item[1] $\prod         _{k=1}^{k=i} \nu_k=n$
        \item[2] $\nu_{i-1}∣\nu_i$ y $\nu_1∣\nu_2$
        \item[3] $\nu_{k-1}∣\nu_k$ si $1<k<i-1$
    \end{itemize}
\end{thm}


Teoremas sobre $p$-grupos.
\begin{thm}{Los $p$-grupos tienen centro no trivial\\}{\label{CentroP-GrupoNoTrivial}}
    Si $\grG$ es un $p$-grupo d\'onde $p  \in \setP$, entonces el centro es no-trivial y $\left|\centro{\grG}\right|\geq p$. Si adem\'as $\grG$ es no-abeliano, $\left| \grG\right|=p^n \then \left|\centro{\grG}\right|  \in \set{p,p^2,\ldots,p^{n-2}}$.
\end{thm}


Sobre el grupo cociente del grupo total por el centro.
\begin{thm}{El cociente por el centro es c{\'{\i}}clico implica que el grupo es abeliano\\}{\label{CocientePorCentroCiclicoYAbelianidad}}
    Si $\grcoc{\grG}{\centro{\grG}}$ es c{\'{\i}}clico, $\grG$ es abeliano.
\end{thm}


Condiciones simples de normalidad.
\begin{thm}{El {\'{\i}}ndice de un grupo por un subgrupo es el m{\'{\i}}nimo divisor del total implica que el subgrupo es normal\\}{\label{IndiceYNormalidad}}
    Si $\grH\subgreq\grG$ y $\indice{\grG}{\grH}=p$ es divisor m{\'{\i}}nimo del orden de $\grG$, entonces $\grH\normsubgreq\grG$.
\end{thm}


Sobre el grupo derivado o conmutador.
\begin{thm}{El cociente por el derivado es abeliano\\}{\label{DerivadoYAbelianidad}}
    El grupo $\grcoc{\grG}{\grderivado{\grG}{1}}$ es abeliano. Adem\'as, si $\grH  \subsetneq {\grderivado{\grG}{1}}$, $\grH\unlhd\grG$,
    entonces  $\grcoc{\grG}{\grH}$ no es abeliano.
\end{thm}


\chapter{Las clases no isomorfas de grupos para cada orden dado}

\section{{\'O}rdenes 1 y primos}

Estos casos son elementales, ya que por el teorema de Lagrange, los \'unicos subgrupos que pueden tener estos grupos son ellos mismos (el subgrupo impropio) y el subgrupo trivial. El orden $1$ es d\'onde coinciden ambos y su importancia radica en que para todo grupo, el subgrupo m{\'{\i}}nimo en su ret{\'{\i}}culo de inclusiones es siempre el $\trivialgr$ y de ah{\'{\i}} la tendencia a considerarlo o pensarlo como una entidad \'unica, un poco al estilo del conjunto vac{\'{\i}}o en el ret{\'{\i}}culo de subconjuntos de cualquier conjunto (pero este es necesariamente \'unico).

Veamos que dado un n\'umero primo $p$, todos los grupos de ese grado son isomorfos. Lo primero, no tienen subgrupos propios no triviales. Por \'ultimo sean $\grG$ y $\grH$ grupos de orden $p$. Si $p=1$ (no est\'a en los primos pero lo admitiremos en este contexto), ambos grupos est\'an un{\'{\i}}vocamente determinados $\grG=\set{\unity_{\grG}}$, $\grH=\set{\unity_{\grH}}$. Solo cabe una funci\'on de uno en otro $f:\grG\To\grH::\unity_{\grG}         \mapsto \unity_{\grH}$ y $f^{-1}$ que son isomorfismos de grupos, de hecho los \'unicos endomorfismos posibles. Son el isomórficamente el mismo.
Sea ahora $p         \in \setP$. Sean $\grG$ y $\grH$ del orden $p$. Sean $g         \in \grG\setminus\set{\unity}$ y $h         \in \grH\setminus{\unity}$, entonces $g         \mapsto h$ determina un{\'{\i}}vocamente un isomorfismo entre ambos grupos ya que el orden de $g$ y de $h$ es el mismo, $p$. Luego son isomorfos. A estos grupos los denominamos $\trivialgr$ y $\cyclicgr{p}$ para $p         \in \setP$. Son isomorfos al conjunto de los restos de divisi\'on m\'odulo $p$ y as{\'{\i}} los implementamos en la programaci\'on.

\section{\'Ordenes p{\^{}}2 d\'onde p es un primo}

En estos casos el procedimiento es casi inmediato. ?`Puede haber grupos no-abelianos?

El centro no puede ser otro que, existir\'a un elemento de orden $p$ del grupo $\grG$, digamos $g$, que sabemos que existe por el teorema de Cauchy, que pongamos que genera $\centro{\grG}$. Entonces $\indice{\grG}{\generadopor{g}}=p$ y por lo tanto, al ser el menor divisor del orden del grupo, resulta que $\generadopor{g}\normsubgr\grG$ y $\grcoc{\grG}{\generadopor{g}         \cong         \cyclicgr{p}}$ y por lo tanto $\grG$ es abeliano y $\grG=\centro{\grG}$: absurdo. Por lo tanto $\grG$ de orden $p^2$ ser\'a siempre abeliano.

Por el teorema de clasificaci\'on de los grupos abelianos finitos solo existen dos posibilidades:

\begin{thm}
\begin{flalign*}
&\Egr{p}{2}\mathdef\pdcyclics{p}{p}\\
&\left|\grG\right|=p^2 \ly p         \in \setP \text{ entonces}\\
&\grG         \cong
\begin{cases}
&\cyclicgr{p^2}\quad \lo\\
&\gKleingr{p}{2}{0} \mathdef \Egr{p}{2}         \equiv          \pdcyclics{p}{p}
\end{cases}
\end{flalign*}
\end{thm}


\section{\'Ordenes pq d\'onde p,q  est\'an en P considerando p<q}

Si $\grG$ es abeliano tiene que ser $\cyclicgr{pq}$ ya que $mcd(p,q)=1$, y no existen factorizaciones no triviales de las que se requerir{\'{\i}}an para obtener alg\'un otro grupo no abeliano.

Tambi\'en podemos darnos cuenta que existen, por Cauchy, elementos de orden $p$ (digamos $g$) y elementos de orden $q$ (digamos $f$), que conforman los \'unicos grupos posibles y que de hecho se tienen que dar.

Tambi\'en podemos ver que el grupo generado por el elemento de orden $q$, $f$, debe ser normal, pues $\indice{\grG}{\generadopor{f}}=p$ que es el menor divisor no trivial de $\left|\grG\right|=pq$.

Podemos ver que $g\generadopor{f}         \in \grcoc{\grG}{\generadopor{f}}$ hacen una partici\'on de $\grG$, adem\'as que $\generadopor{g}         \cap \generadopor{f}=\set{\unity}$. As{\'{\i}} que todos los elementos de $grG$ han de generarse por una multiplicaci\'on de un elemento de $\generadopor{g}$ y otro de $\generadopor{h}$. En un diagrama de factorizaci\'on queda:
\begin{flalign*}
&1\quad
         \longrightarrow
\generadopor{f}
\quad
\stackrel
{\alpha (f^i)\mathdef f^i}
{         \longrightarrow         }
\quad\grG\quad
\underset
    {\gamma (f^k)\mathdef f^k}
    {\overset
        {{\beta (g^i f^j)}\mathdef f^{k(i,j)}}
        {\rightleftarrows}
    }
\quad
{\generadopor{f}}
         \longrightarrow          \quad
1\\
&\beta         \circ \gamma=\mathtt{id}_{\generadopor{f}}\\
&\qquad
\beta         \circ \gamma (f^k)=\beta (f^k)=f^k\\
&{{\gamma         \circ \beta }|}_{\generadopor{f}} = \mathtt{id}_{\generadopor{f}}\\
&\qquad
\gamma         \circ \beta (g^if^j)=\gamma (f^{k(i,j)})=f^{k(i,j)}\\
\end{flalign*}

?`Cu\'ando la secuencia anterior determina el producto directo? Siempre que $\forall i,j \; k(i,j)=0$.

?`Qu\'e posibilidad hay de obtener un diagrama que no nos de el producto directo (que saldr{\'{\i}}a necesariamente $\grG$ abeliano)?

Como es un homomorfismo de grupos, solo si hay una acci\'on $\generadopor{g}      \times      \generadopor{f}\To\generadopor{f}$, esto es, cuando haya un homomorfismo de $\generadopor{g}\To{\graut{\generadopor{f}}}$, y para que no sea el homomorfismo trivial tiene que dividir el orden de $g$, $p$, al orden del grupo de automorfismos de $\generadopor{f}$, que sabemos que su orden es $q-1$. En caso contrario solo cabe la acci\'on trivial y por tanto el grupo $\grG$ ser\'a el c{\'{\i}}clico antes hallado.

A partir de aqu{\'{\i}} vamos a calcular el grupo completo como un producto semidirecto de $\psdcyclics{q}{p}$, esto es en \'ultimo t\'ermino, hallar el anterior $k(i,j)$ que nos sal{\'{\i}}a en el homomorfismo $\beta$.

\subsection{Detalle del producto semidirecto cuando p divide a q-1}

Primero vemos el grupo de automorfismos del grupo c{\'{\i}}clico mayor:

\begin{flalign*}
&f          \in {{\mathsf{C}}_{q}} \; f^q=1\\
&\graut{\cyclicgr{q}}=
\set{\generadopor{f         \mapsto f},\ldots,\generadopor{f         \mapsto f^{q-1}}}
\end{flalign*}

Ahora los homomorfismos posibles que buscamos son :

\begin{flalign*}
&\cyclicgr{q}            \cong            \generadopor{ f ∣f^q=1}\\
&\cyclicgr{p}            \cong            \generadopor{ g ∣g^p=1}\\
& p ∣q-1\\
&{{\varphi }_{i}}:\cyclicgr{p}            \longrightarrow             {\graut{\cyclicgr{q}}\quad i=0,1,2,\ldots,q-1}\\
&{{\varphi }_{i}}\generadopor{f}\mathdef
\begin{cases}
\generadopor{f          \mapsto f}            \Leftarrow             i=0\\
\generadopor{f          \mapsto f^2}            \Leftarrow             i=1\\
\qquad\ldots\ldots\\
\generadopor{f          \mapsto f^{q-1}}            \Leftarrow             i=q-1\\
\end{cases}
\end{flalign*}

Pero de todos estos tenemos que escoger a lo sumo $p$ distintos.
Si fijamos un $i=r$, veremos que se genera un homomorfismo
${\psi } : \cyclicgr{p}         \rightarrow \graut{\cyclicgr{q}}$
que debe ser no trivial y cumplir alguna otra condici\'on que
enseguida veremos. La no trivialidad y otro requisito recaer\'a
en la elecci\'on de $r$. Por ahora no tenemos en cuenta estos
dos importantes detalles.

\[
{\psi }{\left({f}\right)}\mathdef {f^r}
\]
\[
{\psi }{\left({f^j}\right)}={ {\left({f^r}\right)}^j}=f^{rj}
\]

Es homomorfismo, isomorfismo de hecho, por ser $\generadopor{f}$ abeliano (por c{\'{\i}}clico).
Si hacemos, por ejemplo, $r=\frac{q-1}{p}$, se cumplir\'a:

\[
{\psi }{\left({f^p}\right)}=f^{rp}=f^{q-1}=f^{-1}
\]

Introducimos ahora el par\'ametro debido a $g^k$:

\[
{{\theta }_{g=g^1}}{\left({f^j}\right)} \mathdef f^{rj}
\]

A partir de aqu{\'{\i}}, para conformar una acci\'on, vemos para aclararnos el caso $k=2$ :

\[
{{\theta }_{g^2}}{\left({f^j}\right)}=
{{\theta }_{g}}            \circ {{\theta }_{g}}{\left({f^j}\right)}=
{{\theta }_{g}}{\left({f^{rj}}\right)}=f^{r^2 j}
\]

Y as{\'{\i}},  en general:

\[
{{\theta }_{g^k}}{\left({f^j}\right)} =
{{\theta }_{g^{k-2}}}            \circ {{\theta }_{g^2}}{\left({f^j}\right)}=
{{\theta }_{g^{k-2}}}{\left({f^{r^2j}}\right)}=
f^{r^kj}
\]

El caso especial $k=0$ se resuelve f\'acilmente en ${\theta }_{1}{\left({f^j}\right)}=f^j$, esto es, la identidad.

Tomando solo los {\'{\i}}ndices nos queda:

\[
{{\vartheta }_{k}}{\left({j}\right)} \mathdef r^kj
\]

La acci\'on de un elemento sobre otro (tomando solo exponentes), no solo componentes:

\[
(i,j)\blacktriangleleft(k,m)\mathdef{(i,{{\vartheta }_{k}}{\left({j}\right)})}
\]

Deshaciendo definiciones:

\[
(i,j)\blacktriangleleft(k,m)=(i,r^k j);
\]

Nos queda ver los requisitos de $r$:
\begin{flalign*}
& \mbox{\textbf{ Hip\'otesis :  }} r  {         \equiv         }_q 1\\
& g f g^{-1} = f^r = f\\
& gf = fg\\
& \psdcyclics{q}{p}          \cong          \pdcyclics{q}{p} \mbox{\textbf{absurdo}}\\
& \bm{ r {\not         \equiv         }_{q} 1 } \mbox{\textbf{conclusi\'on}}
\end{flalign*}

El otro requisito para $r$ es sencillo de ver y tiene que ver tambi\'en con la aritm\'etica modular:
\begin{flalign*}
&  g^{p} f g^{-p} = 1 f 1 = f = f^1\\
&  g^{p} f g^{-p} = f^{r^p} \\
& \bm{r^p {         \equiv         }_{q} 1}
\end{flalign*}

As{\'{\i}} que $r$ ser\'a tal que
\begin{flalign*}
&           \nexists            s           \in \setN \boldsymbol{ r-1 = sq }\\
&           \exists            t           \in \setN \boldsymbol{ r^p -1 = tq }
\end{flalign*}

Seguimos ya con todas las condiciones para que la acci\'on conforme un grupo no abeliano.

Ahora el producto con solo los exponentes queda (el signo $+$ hace referencia a la operaci\'on propia del producto directo):
\[
(i,j)          \cdot (k,m)\mathdef (i,r^kj)+(k,m)=(i+k,r^kj+m)
\]

\begin{center}
    \begin{flalign*}
    &\bm{(i,j)          \cdot (k,m) \qquad=\qquad (i+k,r^kj+m)}\\
    &           \nexists           s           \in \setN \boldsymbol{ r-1 = sq }\\
    &           \exists            t           \in \setN \boldsymbol{ r^p -1 = tq }
    \end{flalign*}
\end{center}

(Cada componente vive en su respectivo $\grcoc{\setZ}{n\setZ}$)

\begin{thm}
    Dado un grupo $\grG$ de orden $pq$, d\'onde $p         \in \setP$ y $q         \in \setP$, siendo $p<q$, entonces solo caben los siguientes grupos:
    \begin{flalign*}
    &\grG          \cong
    \begin{cases}
    & \cyclicgr{pq} \\
    & \psdcyclics{q}{p}          \Leftarrow          p∣q-1\\
    \end{cases}\\
    \end{flalign*}
    La operaci\'on que queda en el caso no-abeliano es:
    \[
        g^if^j          \cdot g^k f^m = g^{i+k}f^{r^kj+m}
    \]
    D\'onde $r$ ha de cumplir:
    \begin{center}
        \begin{flalign*}
        &           \nexists          s           \in \setN \bm{ r-1 = sq }\\
        &           \exists            t           \in \setN \bm{ r^p -1 = tq }
        \end{flalign*}
    \end{center}
\end{thm}

\section{\'Ordenes p{\^{}}3 d\'onde p es un primo}

Estos \'ordenes tiene por lo pronto $3$ clases de grupos abelianos (por el teorema de clasificaci\'on de los grupos abelianos finitos):

\begin{flalign*}
&\Egr{p}{3}\mathdef\pd{\pdcyclics{p}{p}}{\cyclicgr{p}}\\
&\gKleingr{p}{2}{1}\mathdef\pdcyclics{p^2}{p}\\
&\cyclicgr{p^3}
\end{flalign*}


\begin{proof}[Clasificaci\'on de grupos no abelianos de orden $p^3$]
\begin{itemize}
\item[1]
Para los grupos no abelianos, tenemos que el centro del grupo $\grG$ ha de ser necesariamente de orden $p$. Adem\'as el cociente ha de ser no c{\'{\i}}clico, por lo que solo cabe $\grcoc{\grG}{\centro{\grG}}         \cong         \pdcyclics{p}{p}$.
\item[2]
Sea el subgrupo normal maximal propio $\grN\normsubgr\grG$, que debe existir por ser $\grG$ finito y ser cualquier subgrupo de orden $p^2$ (que existe por los teoremas de Sylow) normal por ser el {\'{\i}}ndice $\indice{\grG}{\grN}$ de orden $p$ que es el menor divisor no trivial del orden del grupo.
\item[3]
$\grN$ ser\'a isomorfo bien a $\cyclicgr{p^2}$, bien a $\pdcyclics{p}{p}$, siendo en cualquier caso que $\centro{\grG}\subsetneqq\grN$, esto es, est\'a dentro de este grupo normal de forma estricta.
\item[3.1]
Por el lado del orden, si es subgrupo, es de orden $p$ mientras que el orden de $\grN$ es $p^2$. $\centro{\grG}         \neq         \grN$.
\item[3.2]
Que $\centro{\grN}         \subset \grN$ hay que verlo algo mejor. Supongamos que $\grN         \cap \centro{\grG}=\set{unity}$. Sea $c         \in \centro{\grG}$ tal que $c         \neq         \unity$ y por lo tanto $c         \notin         \grN$. Ahora $\generadopor{c,\grN}=\grG$, pues el orden es $p^3$, y por lo tanto $\grG         \cong         \centro{\grG}         \subset \grN$, contra lo inicialmente supuesto. Luego $\grN         \cap \centro{\grG}         \neq         \set{unity}$. Como el orden del centro es $p$ primo, cualquier elemento no trivial del centro lo genera, as{\'{\i}} que cualquier elemento no trivial del centro que est\'e en $grN$, genera todo $\centro{\grG}$ y por lo tanto $\centro{\grG}\normsubgr\grN$.
\item[4]
Desde que $\grcoc{\grG}{\centro{\grG}}$ es abeliano (pues tiene orden $p^2$), $\grderivado{\grG}{1}=\generadopor{\conmutador{\grG}{\grG}}         \subseteq \centro{\grG}$. Como el centro tiene orden primo, los \'unicos \'ordenes posibles para el subgrupo derivado son $1$ y $p$. El grupo derivado no puede ser el trivial pues $\grG$ es no-abeliano, y de aqu{\'{\i}}, que debe coincidir con el centro, $\grderivado{\grG}{1}=\centro{\grG}$.
\item[5]
$\grcoc{\grG}{\centro{\grG}}$ es abeliano elemental de orden $p^2$, ya sea $\grN         \cong         \pdcyclics{p}{p}$, ya sea $\grN         \cong         \cyclicgr{p^2}$. De aqu{\'{\i}} veremos que est\'a generado por un conjunto de cardinal $2$, siendo elementos distintos de la identidad, digamos $\set{a\centro{\grG},b\centro{\grG}}$, d\'onde $a         \notin          \centro{\grG}$, $b         \notin          \centro{\grG}$ y $a\centro{\grG}         \cup          b\centro{\grG}=\emptyset$. Entonces $\grG=\generadopor{a,b,\centro{\grG}}$.
\item[6]
Lo siguiente es ver que $ab         \neq          ba$. $\centro{\grG}$ conmuta con todo $\grG$. Si $ab=ba$, tenemos que tambi\'en $\centralizador{\grG}{a}=\generadopor{\centro{\grG},b}$ y $\centralizador{\grG}{b}=\generadopor{\centro{\grG},a}$ siendo entonces $\centralizador{\grG}{a}=\centralizador{\grG}{b}=\generadopor{\centro{\grG},a,b}=\grG$ de d\'onde $\set{a,b}         \subset \centro{\grG}$. Contra lo que hab{\'{\i}}amos ya demostrado: $a$ y $b$ no conmutan.
\item[7]
Como hemos visto que el centralizador y el subgrupo derivado son lo mismo, sea $z\mathdef\conmutador{a}{b}         \in \grderivado{\grG}{1}=\centro{\grG}$. $z         \neq         \unity$ (porque no conmutan), y $generadopor{z}=\centro{\grG}$ (porque tiene orden primo el centro y el elemento $z$ est\'a dentro del centro sin ser el neutro).
\item[8]
El orden de $a$ tiene que dividir al orden de $grG$, que no es c{\'{\i}}clico. Las \'unicas posibilidades son $p$ y $p^2$.
\item[9]
A la vez, la proyecci\'on de $a$ in $\grcoc{\grG}{\centro{\grG}}$ tiene orden $p$ porque $\grcoc{\grG}{\centro{\grG}}$ es, como ya hemos visto, abeliano elemental, por eso $a^p          \in \centro{\grG}$. Lo mismo podemos decir de $b$, esto es $b^p          \in \centro{\grG}$.
\item[10]
Ahora procedemos a un estudio por casos:
\begin{itemize}
    \item[10.1]
    \textbf{$\mathsf{o}(a)=\mathsf{o}(b)=p$}.

    Podemos suponer sin p\'erdida de generalidad que $a\in \grN\setminus\centro{\grG}$, ya que el centro tiene orden $p$ y $\grN$ tiene orden $p^2$, y este caso, $b\in \grG\setminus\grN$. Tenemos de otra parte que $z\in \grderivado{\grG}{1}=\centro{\grG}$, tambi\'en tiene orden $p$. De esta forma $\grG=\generadopor{z,a,b}$.
    Queda la siguiente presentaci\'on:
    \[
    \grG=\generadopor{a,b,z∣a^p=b^p=z^p=\unity, az=za, bz=zb,\conmutador{a}{b}=z}
    \]
    \item[10.2]
    \textbf{$\mathsf{o}(a)=p^2\ly\mathsf{o}(b)=p$}.

    Lo primero a tener en cuenta en este caso es que $a^p          \in \centro{\grG} = \generadopor{z}$. Desde el momento que $a^p$ no es el elemento neutro (por el orden de $a$), existe un n\'umero $r\mod p$ que cumple $a^p = z^r$.
    Sea $c \mathdef b^r$. Tenemos:

    \[ \conmutador{a}{c} = \conmutador{a}{b^r} = {\conmutador{a}{b}}^r = z^r = a^p \]

    La presentaci\'on que queda es:

    \[\grG=\generadopor{ a,c ∣a^{p^2} = c^p = e, \conmutador{a}{c} = a^p }\]
    \item[10.3]
    \textbf{$\mathsf{o}(a)=p^2=\mathsf{o}(b)$}.

    $\generadopor{a}         \cap \generadopor{b}         \neq         \set{\unity}$ (el grupo tendr{\'{\i}}a orden $p^4$). Por lo pronto el orden de la intersecci\'on ha de ser $p$ o $p^2$. Si fuese $p^2$, $\generadopor{a}=\generadopor{b}$ solo tendr{\'{\i}}amos la posibilidad de $p=2$, puesto que $a\centro{\grG}=b\centro{\grG}$, y $\grG=a\centro{\grG}         \cup         \centro{\grG}$ de d\'onde $\entreparentesis{a\centro{\grG}}^2=\centro{\grG}$.

    Pero tampoco es v\'alido para el caso $p=2$, puesto que entonces tendr{\'{\i}}amos un orden como mucho $2^2$, siendo el grupo $\grG         \cong         \cyclicgr{4}$. Absurdo, luego:
    \[
    \left|\generadopor{a}         \cap \generadopor{b}\right|=p
    \]

    Definimos $
    \grH\mathdef\generadopor{x         \in \generadopor{a}         \cap \generadopor{b}\setminus\set{\unity}}$. Vamos a ver que en el caso $p=2$, $\centro{\grG}=\generadopor{a}         \cap \generadopor{b}$.
    Si no fuese as{\'{\i}}, solo queda $\centro{\grG}         \cap \entreparentesis{\generadopor{a}         \cap \generadopor{b}}=\set{\unity}$ (porque los \'ordenes son primos). Por lo tanto, $\grG={\entreparentesis{\generadopor{a}         \cap \generadopor{b}}\centro{\grG}}$ y el orden de $\grG$ es $p^2$ y no $p^3$.

    Ahora nos queda, $\centro{\grG}         \cong
    {\grcoc{\generadopor{a}}{\centro{\grG}}}
    {\grcoc{\generadopor{b}}{\centro{\grG}}}
    $. Como $a$ y $b$ no conmutan, $\set{
        \grcoc{\generadopor{a}}{\centro{\grG}},
        \grcoc{\generadopor{b}}{\centro{\grG}},
        \grcoc{\generadopor{ab}}{\centro{\grG}}
    }
    $ constituye una partici\'on de $\grcoc{\grG}{\centro{\grG}}$.

    La presentaci\'on que queda es, llamando a $\overline{\unity}\mathdef x$:

    \[
    \grG=\generadopor{a,b ∣a^p=b^p=(ab)^b=\overline{\unity},{\overline{\unity}}^p=\unity,\generadopor{\overline{\unity}}=\centro{\grG}}
    \]

\end{itemize}
\end{itemize}
\end{proof}

\subsection{Explorando el producto semidirecto de CpxCp por Cp}

Veamos los automorfismos de $\pdcyclics{p}{p}$. Podemos considerar dos elementos $a$ y $b$ de igual orden, $p$. Los elementos orden $p$ ser\'an
\[\set{a^{p-1},\dots,a,a^{p-1}b,\ldots,ab,\ldots,a^{p-1}b^{p-1},\ldots,ab^{p-1},b^{p-1},\ldots,b}\]
, en total son $p^2-1$ elementos de orden $p$.As{\'{\i}} para cada elecci\'on $a         \mapsto a^ib^j$, $b$ puede enviarse a cualquier otro que no sea el ocupado por la imagen de $a$. Son posibles $(p^2-1)(p^2-2)$ automorfismos. En el caso de orden $2^2$ son posibles $(2^2-1)(2^2-2)=6$, tal como sabemos que los automorfismos del $V$ grupo de Klein tiene un grupo de automorfismos isomorfo al $\dihgr{3}$.

Aunque el n\'umero de automorfismos es muy grande, las posibilidades de grupos con $p^3$ elementos son reducidas. As{\'{\i}} que comenzaremos por encontrar alg\'un homomorfismo de $\cyclicgr{p}$ en $\graut{\pdcyclics{p}{p}}$. De hecho, por este camino, solo hay oportunidad (para grupos no-abelianos) si $p|(p^2-2)$.

\begin{lem}
Si el grupo $\grG$ es de orden $p^3$ siendo $p         \in \setP$ y no es abeliano, tiene la posibilidad de desarrollarse como producto semidirecto $\psd{(\pdcyclics{p}{p})}{\cyclicgr{p}}$ si $p|p^2-2$.
\end{lem}

Sea
${{\varphi }_{c}}\mathdef{\generadopor{\gfrac{a         \mapsto a}{b         \mapsto ab}}}$.
Entonces
\[\varphi_{c^2} = \generadopor{\gfrac{a         \mapsto a}{b         \mapsto ab}}         \circ \generadopor{\gfrac{a         \mapsto a}{b         \mapsto ab}} = \generadopor{\gfrac{a         \mapsto a}{b         \mapsto a^2b}}\].
As{\'{\i}} obtendremos $\varphi_{c^k} = \generadopor{\gfrac{a         \mapsto a}{b         \mapsto a^kb}}$. El homomorfismo es:

\begin{flalign*}
\varphi : &\cyclicgr{p} \To \graut{\pdcyclics{p}{p}}\\
& c^k          \mapsto \generadopor{\gfrac{a          \mapsto a}{b          \mapsto a^kb}}
\end{flalign*}

Podemos formar el producto semidirecto as{\'{\i}}:

\begin{flalign*}
\mathbf{a^ib^jc^k         \cdot a^lb^mc^n\mathdef a^{i+kl}b^{j+m}c^{k+n}}
\end{flalign*}

As{\'{\i}} el inverso de un elemento queda:
\begin{flalign*}
&a^ib^jc^k         \cdot a^lb^mc^n=a^{0}b^{0}c^{0}=\unity\\
&i+kl\;{         \equiv         }_{p}\;0 \Leftrightarrow l = ({{(p-i)}{k}^{-1}} \mod p) \quad k^{-1}\;\text{ en el cuerpo }\nicefrac{\setZ}{p\setZ}\\
&j+m\;{         \equiv         }_{p}\;0 \Leftrightarrow m = (p-j \mod p)\\
&k+n\;{         \equiv         }_{p}\;0 \Leftrightarrow n = (p-k \mod p)\\
&\mathbf{(a^ib^jc^k)^{-1}={a^{{(p-i)}{k}^{-1}}}{b^{p-j}}c^{p-k}}
\end{flalign*}

Hay que darse cuenta que $ab=ba$, pero no conmutan los elementos de $\Egr{p}{2}         \cong         \generadopor{a,b}$ con los de $\cyclicgr{p}         \cong         \generadopor{c}$.


\subsection{Producto semidirecto Cp{\^{}}2 por Cp}

El grupo automorfismos de $\cyclicgr{p^2}$ es isomorfo a $\cyclicgr{p(p-1)}$. Podemos considerar un elemento $f$ de orden $p^2$. Los elementos de orden $p$ ser\'an aquellos de la forma $f^p$. Enumeramos los elementos de orden $p^2$,
\[\set{f,\dots,f^{p-1},f^{p+1},\ldots,f^{2p-1},f^{2p+1},\ldots,f^{p(p-1)-1},f^{p(p-1)+1},\ldots,f^{p^2-1}}\]
, en total son $p(p-1)$ elementos. As{\'{\i}} que tenemos $f         \mapsto f^i$ con $i$ con $p(p-1)$ elecciones posibles.

\begin{lem}
Si el grupo $\grG$ es de orden $p^3$ siendo $p         \in \setP$ y no es abeliano, tiene la posibilidad de desarrollarse como producto semidirecto $\psd{\cyclicgr{p^2}}{\cyclicgr{p}}$.
\end{lem}

Sea
${{\varphi }_{g}}\mathdef{\generadopor{f         \mapsto f^i}}$.
Entonces
\[\varphi_{g^2} = \generadopor{f         \mapsto f^i}         \circ \generadopor{f         \mapsto f^i} = \generadopor{f         \mapsto f^{i^2}}\]. As{\'{\i}} obtendremos $\varphi_{g^k} = \generadopor{f         \mapsto f^{i^k}}$. El valor de $i         \equiv         _{p^2}1$ nos da el homomorfismo trivial por lo que no nos vale. El homomorfismo es:

\begin{flalign*}
&r{\not         \equiv         }_{p^2}1\\
\varphi : &\cyclicgr{p} \To \graut{\cyclicgr{p^2}}\\
& g^k          \mapsto \generadopor{f          \mapsto f^{r^k}}
\end{flalign*}

Podemos formar el producto semidirecto as{\'{\i}}:

\begin{flalign*}
\mathbf{g^if^j         \cdot g^kf^l\mathdef g^{i+k}f^{r^kj+l}}
\end{flalign*}

As{\'{\i}} el inverso de un elemento queda:
\begin{flalign*}
&g^if^j         \cdot g^kf^l=g^{0}f^{0}=\unity\\
&i+k\;{         \equiv         }_{p}\;0 \Leftrightarrow k = p-i \mod p\\
&r^kj+l\;{         \equiv         }_{p^2}\;0 \Leftrightarrow l = (p^2-r^k)+(p^2-l) \mod p^2\\
&\mathbf{{(g^if^j)}^{-1}={g^{p-i}}{f^{p^2-r^k-l}}}
\end{flalign*}

\subsection{Una factorizaci\'on m\'as all\'a del producto semidirecto}


Dado que el producto semidirecto dado en la anterior subsecci\'on es posible para todos los \'ordenes primos, seguiremos trabajando a partir de los resultados conseguidos en ese producto.

Resumidamente:

\begin{flalign*}
&\mathbf{r{\not         \equiv         }_{p}1}\\
&\mathbf{g^if^j         \cdot g^kf^l\mathdef g^{i+k}f^{r^kj+l}}\\
&\mathbf{{(g^if^j)}^{-1}={g^{p-i}}{f^{p^2-r^k-l}}}
\end{flalign*}

Ahora podemos darnos cuenta que, dados los dos argumentos de un producto, digamos $g^if^j$ como operando izquierdo y $g^kf^l$ como derecho, podemos coger las proyecciones $g^i$ y $g^k$ y darnos cuenta que ambas viven en $\pdcyclics{p}{p}$. Si logr\'aramos una acci\'on del tipo:

\begin{flalign*}
\varphi : &\pdcyclics{p}{p}          \times          \cyclicgr{p^2} \To \cyclicgr{p^2}\\
& ((g^i,g^k),f^l)\qquad         \mapsto \qquad f^{\zeta (i,k,l)}
\end{flalign*}

Podr{\'{\i}}amos intentar obtener una nueva operaci\'on no isomorfa a las anteriores. En alg\'un caso es probable que, incluso, ambos grupos fuesen normales (sin ser el grupo total abeliano). El requisito es que la operaci\'on ha de seguir siendo asociativa, pues lo dem\'as vendr\'a dado.

Veamos la acci\'on. La m\'as sencilla posible que no sea la trivial, puede ser una multiplicaci\'on por la derecha por un factor elegido por $(g^i,g^k)$. Entonces las acci\'on quedar{\'{\i}}a:

\begin{flalign*}
\varphi : &\pdcyclics{p}{p}          \times          \cyclicgr{p^2} \To \cyclicgr{p^2}\\
& ((g^i,g^k),f^l)\qquad         \mapsto \qquad f^{l+\zeta (i,k)}
\end{flalign*}

Veamos que es acci\'on:

\begin{flalign*}
& (g^i,g^k)\blacktriangleright f^l = f^{l+\zeta (i,k)}\\
& (g^0,g^0)\blacktriangleright f^l = f^{l+\zeta (0,0)}\\
&\qquad\qquad 1\blacktriangleright f^l = f^{l+\zeta (0,0)}\\
&\qquad\qquad\qquad\qquad f^l = f^{l+\zeta (0,0)}\\
&\qquad\qquad\qquad\qquad\qquad\qquad \zeta (0,0){         \equiv         }_{p^2}0\\
& ((g^i,g^k)(g^r,g^s))\blacktriangleright f^l = (g^i,g^k)\blacktriangleright((g^r,g^s)\blacktriangleright f^l) :\\
&\qquad\qquad (g^{i+r},g^{k+s})\blacktriangleright f^l = (g^i,g^k)\blacktriangleright f^{l+\zeta (r,s)}\\
&\qquad\qquad\qquad\qquad f^{l+\zeta (i+r,k+s)} = f^{l+\zeta (r,s)+\zeta (i,k)}\\
&\qquad\qquad\qquad\qquad\qquad\qquad \zeta (i+r,k+s) {         \equiv         }_{p^2} \zeta (r,s)+\zeta (i,k)
\end{flalign*}

Las condiciones para que sea acci\'on $\varphi$ (con la condici\'on puesta) es que la funci\'on $\zeta$ cumpla:

\begin{flalign*}
&\mathbf{(g^i,g^k)\blacktriangleright f^l = f^{l+\zeta (i,k)}}\\
&\mathbf{\qquad\qquad\qquad \zeta (0,0){         \equiv         }_{p^2}0}\\
&\mathbf{\qquad\qquad\qquad \zeta (i+r,k+s) {         \equiv         }_{p^2} \zeta (r,s)+\zeta (i,k)}
\end{flalign*}

Veamos como quedar{\'{\i}}a la operaci\'on, haciendo el producto semidirecto anterior y aplicando la acci\'on $\varphi$:

\begin{flalign*}
&\mathbf{r{\not         \equiv         }_{p}1}\\
&\mathbf{g^if^j         \cdot g^kf^l\mathdef g^{i+k}f^{r^kj+l+\zeta (i,k)}}\\
\end{flalign*}

La operaci\'on es a\'un dif{\'{\i}}cil de analizar, as{\'{\i}} que veamos un tipo concreto:

\begin{flalign*}
&s         \in \nicefrac{\setZ}{p^2\setZ}\;\textbf{  constante fijada}\\
&\zeta (i,k)\mathdef
\begin{cases}
&0\mod p^2          \Leftarrow          i+k{         \equiv         }_{p}0\\
&s\mod p^2          \Leftarrow          i+k{\not         \equiv         }_{p}0\\
\end{cases}\\
\end{flalign*}

Veamos qu\'e tiene que cumplir $s$ para que sea acci\'on:

\begin{flalign*}
&(g^i,g^k)\blacktriangleright f^l = f^{l+\zeta (i,k)}\\
&\zeta (0,0){         \equiv         }_{p^2}0 \\
&\qquad\qquad \zeta (0,0)= 0          \Leftarrow          0+0{         \equiv         }_{p}0 \textbf{  cumple este requisito}\\
&\zeta (i+t,k+v) {         \equiv         }_{p^2} \zeta (i,k)+\zeta (t,v)\\
&
\begin{cases}
0 =\zeta (0,0) {         \equiv         }_{p^2} \zeta (i,p-i)+\zeta (t,p-t)          \Leftarrow          i+k{         \equiv         }_{p}0 \ly t+v{         \equiv         }_{p}0\\
s= 0+s =\zeta (0,0)+\zeta (t,v) {         \equiv         }_{p^2} \zeta (i,p-i)+\zeta (t,v)          \Leftarrow          i+k{         \equiv         }_{p}0 \ly t+v{\not         \equiv         }_{p}0\\
s= s+0 =\zeta (i,k)+0 {         \equiv         }_{p^2} \zeta (i,k)+\zeta (t,p-t)          \Leftarrow          i+k{\not         \equiv         }_{p}0 \ly t+v{         \equiv         }_{p}0\\
2s= s+s =\zeta (i,k)+\zeta (t,v) {         \equiv         }_{p^2} \zeta (i,k)+\zeta (t,v)          \Leftarrow          i+k{\not         \equiv         }_{p}0 \ly t+v{\not         \equiv         }_{p}0\\
\end{cases}\\
&s         \equiv         _{p^2} 2s \Leftrightarrow s         \equiv         _p 0          \vee p=2
\end{flalign*}

Con esto s\'olo en el caso en que $p=2$ sea par vale esta acci\'on, concretamente, siendo $s=2$. Luego este camino nos revela los grupos posibles para $\grG$ de orden $2^3=8$.

\chapter{Detalle de los grupos de los diferentes \'ordenes}
Partiremos de los que hemos llamado grupos c{\'{\i}}clicos (que incluyen el trivial),
pero considerando solo los de orden primo. Sea :
\begin{flalign*}
\mathds{P} \triangleq
{\left\lbrace{
{
p              \in \mathds{N}\setminus\lbrace 0,1\rbrace
}
\mid
{
p>1
\,
            \wedge
\,
\forall n
              \in \mathds{N}\setminus{\left\lbrace{0,1,p}\right\rbrace}\; n\nmid p
}
}\right\rbrace}
\end{flalign*}

Para comenzar, los grupos de un orden primo $p$ son todos isomorfos entre s{\'{\i}}.
Sean G y H grupos de orden $p             \in \mathds{P}$. La prueba es clara.
Si $p             \in \mathds{P}$, cualquier subgrupo
$G^\prime             \preccurlyeq             G$ y $H^\prime             \preccurlyeq             H$
tendr\'an, por el teorema de Lagrange, orden $1$ o $p$. El subgrupo de orden $1$
solo puede ser
$G^\prime = \left\lbrace \unity\right\rbrace                      \cong                      \cyclicgr{1}$,
o bien $G^\prime = \left\lbrace g                     \in G∣g                     \neq                      1
\right\rbrace                      \cong                      \cyclicgr{p}$. Id\'enticamente en
$H$, $H^\prime = \left\lbrace 1\right\rbrace                      \cong                      \mathsf{C}_1$,
o bien $H^\prime = \left\lbrace h                     \in G∣h                     \neq                      1\right\rbrace
                  \cong                      \cyclicgr{p}$. Como el orden de $G$ y de $H$ es $p>1$,
$                     \exists                      g                     \in G\setminus\{1\}                     \wedge
                  \exists                      g                     \in G\setminus\{1\}$. Forzosamente esos $g$ y $h$ tienen orden $p$,
ya que si tuviesen orden $1$ ser{\'{\i}}an $g=1$ y $h=1$, que por elecci\'on no es as{\'{\i}}.
La aplicaci\'on, siempre que $g             \in G\setminus\{1\}$ y que $h             \in H\setminus\{1\}$,
$\varphi \,:\,G             \longrightarrow              H\,::\, g^n             \mapsto h^n \; n             \in \mathds{N}$ es
isomorfismo de grupos.

De esta forma los $\cyclicgr{p}\;p             \in \mathds{P}$, son los grupos que de forma
inicial es claro que no tienen estructura interna alguna, esto es,
$\mathcal{S}_{G}{(\cyclicgr{p})} = \mathcal{N}_{G}{(\cyclicgr{p})} =
\mathcal{C}_{G}{(\cyclicgr{p})} = \left\lbrace{\{1\},G}\right\rbrace$.
No tienen composici\'on alguna. Aunque posteriormente encontremos grupos simples
no c{\'{\i}}clicos, ser\'an no abelianos y tienen cierta estructura interna ya que por el
teorema de Cauchy, existir\'a al menos un elemento de orden $p             \in \mathds{P}$
para cada $p∣\left|G\right|$ y por lo tanto tendr\'a subgrupos de esos \'ordenes
(de hecho subgrupos c{\'{\i}}clicos de orden primo, para cada primo que divida al orden).

En la clasificaci\'on que nos ocupa, esto nos da que salvo isomorfismo de grupos,
todos los grupos de orden $p             \in \mathds{P}$ y tambi\'en el de orden $1$, son
\'unicos. Adem\'as todo grupo c{\'{\i}}clico es abeliano, caracter{\'{\i}}stica
que es propia de la suma de exponentes naturales con los que se pueden representar
los grupos c{\'{\i}}clicos, suma m\'odulo el orden del grupo c{\'{\i}}clico. El intento ser\'a
construir todos los grupos con un producto cartesiano de grupos c{\'{\i}}clicos simples
como conjunto subyacente del grupo.

Para cada grupo c{\'{\i}}clico, el grupo de automorfismos internos es trivial,
y en general $\mathtt{Aut}(\cyclicgr{p})             \cong             \cyclicgr{p-1}$. Es f\'acil ver que
para orden $p^n , p             \in \mathds{P} , n              \in \mathds{N}\setminus\{0\}$,
$\mathtt{Aut}(\cyclicgr{p^n})             \cong             \cyclicgr{p^{n-1}             \cdot (p-1)}$.
Esta funci\'on que asigna a cada primo su predecesor, a cada potencia de un primo
le asigna el producto de la potencia anterior por el predecesor del primo, y que
adem\'as preserva la multiplicaci\'on cuando los factores son coprimos, es la funci\'on
$\phi_{E}$ de Euler, que nos da la cantidad de divisores de un n\'umero dado,
menores que ese n\'umero (que es el argumento). As{\'{\i}}, si
$n             \in \mathds{N}\setminus\{0\}$, por el teorema fundamental de la aritm\'etica,
existen una cantidad finita de primos,
$\left\lbrace p_1 , p_2 , \ldots , p_r \right\rbrace             \subset \mathds{P}$ con
$r             \in \mathds{N}\,            \wedge             \, r>1$ y una cantidad similar de naturales
estrictamente positivos, $\left\lbrace n_1, n_2 , \ldots , n_r \right\rbrace
            \subset \mathds{N}\setminus\{0\}$ tales que $n = {\prod}_{k=1}^{r} p_k^{n_k}$.
As{\'{\i}} $Aut(\cyclicgr{n})             \cong             \cyclicgr{{\phi_E}(n)}$.

Para cualquier grupo c{\'{\i}}clico de orden primo $p$, el diagrama de Hasse de
los subgrupos es:

\begin{center}
\begin{tikzpicture}
    \label{tab:GruposTrivial}
    % First, locate each of the nodes and name them
    \node (1) at (0,0) {$\cyclicgr{1}$};
    % Now draw the lines:
    \draw [black, thick] (1) -- (1);
\end{tikzpicture}
\end{center}


\begin{center}
\begin{tikzpicture}
    \label{tab:GruposCiclicosSimplesLatticeDiagram}
    % First, locate each of the nodes and name them
    \node (top) at (0,0) {$\cyclicgr{p}$ con $p        \in \setP$};
    \node (bottom) [below of=top] {$\set{{1}_{\cyclicgr{p}}}            \cong            \trivialgr$};
    % Now draw the lines:
    \draw [black, thick] (top) -- (bottom);
\end{tikzpicture}
\end{center}

Adem\'as tenemos una tabla de grupos y \'ordenes con algunas propiedades:

\begin{table}[h!]
  \begin{center}
    \caption{Grupos c{\'{\i}}clicos de orden primo o trivial}
    \label{tab:GruposCiclicosSimplesCayleyTable}
    \begin{tabular}{|c|c|c|c|c|c|c|}
    \hline
      \tiny{\textbf{Orden}} & \tiny{\textbf{n\'um grupos}} & \tiny{\textbf{numeraci\'on}} & \tiny{\textbf{abeliano}} & \tiny{\textbf{resoluble}} & \tiny{\textbf{nilpotente}}\\
      \hline
      1&$\cyclicgr{1}$ & 1 & 1& Si & Si & Si\\
      \hline
      2&$\cyclicgr{2}$ & 1 & 1& Si & Si  & Si\\
      \hline
      3&$\cyclicgr{3}$ & 1 & 1& Si & Si  & Si\\
      \hline
      5&$\cyclicgr{5}$ & 1 & 1 & Si   & Si & Si\\
      \hline
      7&$\cyclicgr{7}$ & 1 & 1 & Si  &  Si & Si\\
      \hline
      11&$\cyclicgr{11}$ & 1 & 1 & Si  &  Si & Si\\
      \hline
      13&$\cyclicgr{13}$ & 1 & 1 & Si  &  Si & Si\\
      \hline
      17&$\cyclicgr{17}$ & 1 & 1 & Si  &  Si & Si\\
      \hline
      19&$\cyclicgr{19}$ & 1 & 1 & Si  &  Si & Si\\
      \hline
      23&$\cyclicgr{23}$ & 1 & 1 & Si  &  Si & Si\\
      \hline
      29&$\cyclicgr{29}$ & 1 & 1 & Si &  Si & Si\\
      \hline
      31&$\cyclicgr{31}$ & 1 & 1 & Si  & Si & Si\\
      \hline
      37&$\cyclicgr{37}$ & 1 & 1 & Si  &  Si & Si\\
      \hline
      41&$\cyclicgr{41}$ & 1 & 1 & Si  & Si & Si\\
      \hline
      43&$\cyclicgr{43}$ & 1 & 1 & Si  & Si & Si\\
      \hline
      47&$\cyclicgr{47}$ & 1 & 1 & Si  &  Si & Si\\
      \hline
      53&$\cyclicgr{53}$ & 1 & 1 & Si  &  Si & Si\\
      \hline
      59&$\cyclicgr{59}$ & 1 & 1 & Si  &  Si & Si\\
      \hline
      61&$\cyclicgr{61}$ & 1 & 1 & Si  &  Si & Si\\
      \hline
    \end{tabular}
  \end{center}
\end{table}

\chapter{Primeros grupos compuestos: orden 4}

Ahora nos ocupa el orden $4$. Directamente caben al menos dos posibilidades, $\cyclicgr{2^2}$
que tiene un solo grupo normal de orden $2$, un elemento de orden $4$ y que es evidentemente
abeliano por ser c{\'{\i}}clico, y el grupo $\cyclicgr{2}      \times      \cyclicgr{2}$ que es tambi\'en
abeliano por serlo cada subgrupo que forma el producto. Adem\'as, tiene $3$ grupos normales
de orden $2$ y ning\'un elemento de orden $4$.

En cuanto a si puede haber m\'as grupos no isomorfos a los anteriores, tenemos que ver que,
o bien tiene elemento de orden $4$ o no, y necesariamente, por Cauchy, tienen elementos de
orden $2$. Si tiene elemento de orden $4$ solo puede ser el grupo c{\'{\i}}clico, por definici\'on
de grupo c{\'{\i}}clico, y si no, ha de tener m\'as de un elemento de orden $2$. Y tienen que ser
exactamente $3$. Si tuviese m\'as elementos (necesariamente de orden $2$), tendr{\'{\i}}amos m\'as
de $4$ elementos y el orden ya no ser{\'{\i}}a $4$. En caso de ensayar alguna forma no abeliana con
elementos de orden $2$, tomando $2$ generadores (menos ser{\'{\i}}a c{\'{\i}}clico), obtenemos
siempre m\'as elementos que en el caso abeliano, por lo que incluso con solo $2$ generadores solo
cabe el caso abeliano ($\left|\left(\langle a,b ∣a^2=b^2=1, ab        \neq        ba \rangle\right)\right|\geqslant 5$).

El diagrama de Hasse para el grupo c{\'{\i}}clico:

\begin{center}
\begin{tikzpicture}
    % First, locate each of the nodes and name them
    \node (arriba) at (0,0) {$\cyclicgr{2^2}=\set{1,x,x^2,x^3}$};
    \node[below of=arriba] (medio) {$\cyclicgr{2}                    \cong                    \set{1,x^2}$};
    \node[below of=medio] (abajo) {$\cyclicgr{1}                    \cong                    \set{{1}_{\cyclicgr{2^2}}}$};

    % Now draw the lines:
    \draw [black, thick] (arriba) -- (medio);
    \draw [black, thick] (medio) -- (abajo);
\end{tikzpicture}
\end{center}

Y para el producto directo $\pdcyclics{2}{2}       \equiv       \mathsf{V}       \equiv       \grK_{4}$ grupo de Klein:

\begin{center}
\begin{tikzpicture}
    % First, locate each of the nodes and name them
    \node (arriba) at (0,0) {$\cyclicgr{2}       \times       \cyclicgr{2}        \cong        \set{1,a,b,a·b=b·a}$};
    \node (mediocentro) at (0,-2) {$\cyclicgr{2}                    \cong                    \set{1,a·b}$};
    \node (medioderecha) at (-3,-2) {$\cyclicgr{2}                    \cong                    \set{1,a}$};
    \node (medioizquierda) at (3,-2) {$\cyclicgr{2}                    \cong                    \set{1,b}$};
    \node (abajo) at (0,-4) {$\cyclicgr{1}                    \cong                    \set{1_{\cyclicgr{2}       \times       \cyclicgr{2}}}$};
    % Now draw the lines:
    \draw [black, thick] (arriba) -- (medioderecha);
    \draw [black, thick] (arriba) -- (mediocentro);
    \draw [black, thick] (arriba) -- (medioizquierda);
    \draw [black, thick] (medioderecha) -- (abajo);
    \draw [black, thick] (mediocentro) -- (abajo);
    \draw [black, thick] (medioizquierda) -- (abajo);
\end{tikzpicture}
\end{center}


\begin{table}[h!]
  \begin{center}
    \caption{Grupos de orden $4$}
    \label{tab:GruposOrden4}
    \begin{tabular}{c|c|c|c|c|c} % <-- Alignments: 1st column left, 2nd middle and 3rd right, with vertical lines in between
      \textbf{} & \textbf{n\'umero} &\textbf{numeraci\'on} & \textbf{Nombre} & \textbf{} & \textbf{}\\
      \textbf{} & \textbf{de} &\textbf{dentro} & \textbf{del} & \textbf{} & \textbf{}\\
      \textbf{orden} & \textbf{grupos} &\textbf{del} & \textbf{grupo} & \textbf{Simple} & \textbf{abeliano}\\
      \textbf{} & \textbf{no} &\textbf{orden} & \textbf{} & \textbf{} & \textbf{}\\
      \textbf{} & \textbf{isomorfos} &\textbf{} & \textbf{} & \textbf{} & \textbf{}\\
      \hline\\
      \hline
      4 & 2 & 1 & $\cyclicgr{4}$ & No & Si\\
      \hline
      4 & 2 & 2  & ${\cyclicgr{2}}        \times        {\cyclicgr{2}}$ & No & Si\\
      \hline
    \end{tabular}
  \end{center}
\end{table}

\chapter{Grupos de orden 2p d\'onde p es un primo mayor que 2}

Colocamos esta clase de grupos en este lugar porque el siguiente orden en el ascenso que nos ocupa es el $6$ y es del tipo dicho, $6=2·3$ y
$3      \in \setP_{>2}$ y toda esta clase de \'ordenes tiene una estructura de grupos muy similar.

Por una parte est\'a el grupo c{\'{\i}}clico que siempre existe, que es abeliano, y que adem\'as es \'unico por el teorema de clasificaci\'on
de grupos abelianos finitos. Este grupo tiene $2$ subgrupos normales, caracter{\'{\i}}sticos, uno de orden $2$ y otro de orden $p$. As{\'{\i}}
estos \'ordenes tienen un solo grupo abeliano y es, entonces, necesariamente, el c{\'{\i}}clico. Por el teorema de Cauchy tienen un elemento
$a$ de orden $2$ y otro $b$ de orden $p$. Entonces $ab$ tendr\'a orden $mcm(2,p)=2p$ y as{\'{\i}} es c{\'{\i}}clico.

De otro lado est\'an los grupos no abelianos si existiesen. Pero podemos comprobar que el homomorfismo inyectivo
$\varphi : \cyclicgr{2}       \longrightarrow       \graut{\cyclicgr{p}}$ tal que $\varphi (1)={\mathtt{id}}_{\cyclicgr{p}}$ y siendo
$x       \in {\cyclicgr{p}}\setminus\{1\}$ el valor $\varphi (x)$ es el automorfismo que asigna $x       \mapsto x^{p-1}$ y de ah{\'{\i}} asigna a
$x^k       \mapsto x^{k(p-1)}$ obtenemos el producto semidirecto $\cyclicgr{p} {\rtimes}_{\varphi } \cyclicgr{2}$. Aqu{\'{\i}} los $p$ subgrupos
de orden $2$ son no-normales, y el de orden $p$ es \'unico, y por lo tanto no solo normal sino caracter{\'{\i}}stico.
Estos grupos se suelen denominar $\dihgr{p}       \equiv       \cyclicgr{p} {{\rtimes}_{\varphi }} \cyclicgr{2}$, y ya tenemos lo necesario
para confeccionar su tabla de Cayley. Optaremos por denominar a la imagen del homomorfismo $\varphi_x       \equiv       \varphi (x)$ por razones
de mejor legibilidad y escritura m\'as corta.

Sea $(1,a)        \in \left(\cyclicgr{p}      \times      \cyclicgr{2}\right)\setminus\left\lbrace{(1,1)}\right\rbrace$, entonces ha de ser $\varphi_{a}(f^\imath)=f^{(p-1)\imath}$,
 as{\'{\i}}, $(1,a)·(b^\imath,a^\jmath)=(1·\varphi_a (b^\imath),a       \cdot a^\jmath)=(b^{(p-1)\imath},a^{\jmath + 1})$.
 De otro lado, $(b^\imath,a^\jmath)             \cdot (1,a)=(b^\imath ·\varphi_a (1),a^\jmath ·a)=(b^{\imath},a^{\jmath + 1})$,
 y como $(1,a)       \cdot (b^\imath,a^\jmath)=(b^{(p-1)\imath},a^{\jmath + 1})       \neq      (b^{\imath},a^{\jmath + 1}) = (b^\imath,a^\jmath)·(1,a)$,
 al menos para alg\'un valor de $\imath$, $\dihgr{p}$ es necesariamente no-conmutativo.

\[
    (b^i,a^j)*(b^k,a^l) =
    \begin{cases}
        (b^{i+k},a^{l}) &                   \Leftarrow                   j = 0 \\
        (b^{i-k},a^{1+l})&                   \Leftarrow                   j = 1
    \end{cases}
\]

Y partiendo del resultado anterior hallamos el inverso de $(b^i,a^j)$:

\[
(b^i,a^j)^{-1} =
\begin{cases}
(b^{-i},1) &                   \Leftarrow                   j = 0 \\
(b^{i},a)&                   \Leftarrow                   j = 1
\end{cases}
\]

Renombrando $\mathfrak{f}\triangleq (b,1)$ y $\mathfrak{g}\triangleq (1,a)$, tenemos una forma de escribir los elementos del grupo de forma m\'as sencilla y corta. Podemos ver que el elemento $(b^i,a^j)=(b^i,1)*(1,a^j)=(b,1)^i * (1,a)^j$ se puede escribir (es igual a) como $\mathfrak{f}^i \mathfrak{g}^j$. De la misma manera $(1,a^j)*(b^i,1)=(\varphi_{a^j}(b^i),a^j)$ y dependiendo del valor de $j$ podremos ponerla como:

\[
(1,a^j)*(b^i,1) =
\begin{cases}
(b^{i},1) &                   \Leftarrow                   j = 0 \\
(b^{-i},a)&                   \Leftarrow                   j = 1
\end{cases}
\]
\[
(b^i,a^j)=
(b^i,1)*(1,a^j)=
 \mathfrak{f}^i \mathfrak{g}^j
\]
Y ahora podemos traducir
\[
\mathfrak{g}^{j} \mathfrak{f}^{i} =
\begin{cases}
\mathfrak{f}^{i} &                   \Leftarrow                   j = 0 \\
\mathfrak{f}^{-i}\mathfrak{g}&                   \Leftarrow                   j = 1
\end{cases}
\]
En definitiva tenemos una regla de la multiplicaci\'on dih\'edrica:
\[
\mathfrak{g} \mathfrak{f}^{i} =
\mathfrak{f}^{-i}\mathfrak{g}
\]
Para los inversos tenemos:
\[
{(\mathfrak{g}^{i}\mathfrak{f}^{j})}^{-1} =
\begin{cases}
\mathfrak{f}^{-j}&                   \Leftarrow                   i = 0 \\
\mathfrak{g} \mathfrak{f}^{j}&                   \Leftarrow                   i = 1
\end{cases}
\]

Vemos por lo tanto que $\mathfrak{f}^{i} \mathfrak{g} = \mathfrak{g} \mathfrak{f}^{-i}$ es elemento de orden $2$ para todo $i$.

La correspondiente tabla de Cayley es:
\begin{flalign*}
\begin{array}{|c|c|c|c|c|c|c|c|}
\hline
1&\mathfrak{f}&\mathfrak{f}^2&\mathfrak{g}&\mathfrak{gf}&\mathfrak{gf}^2\\
\hline
\mathfrak{f}&\mathfrak{f}^2&1&\mathfrak{gf}^2&\mathfrak{g}&\mathfrak{gf}\\
\hline
\mathfrak{f}^2&1&\mathfrak{f}&\mathfrak{gf}&\mathfrak{gf}^2&\mathfrak{g}\\
\hline
\mathfrak{g}&\mathfrak{gf}&\mathfrak{gf}^2&1&\mathfrak{f}&\mathfrak{gf}^2\\
\hline
\mathfrak{gf}&\mathfrak{gf}^2&\mathfrak{g}&\mathfrak{f}^2&1&\mathfrak{f}\\
\hline
\mathfrak{gf}^2&\mathfrak{g}&\mathfrak{gf}&\mathfrak{f}&\mathfrak{f}^2&1\\
\hline
\end{array}
\end{flalign*}

Podemos ver ahora como es la conjugaci\'on de un elemento:

\begin{align}
(\mathfrak{f}^i\mathfrak{g})(\mathfrak{f}^k\mathfrak{g}^l)(\mathfrak{f}^i\mathfrak{g})^{-1} &=&\\
&=\mathfrak{f}^i\mathfrak{g}\mathfrak{f}^k\mathfrak{g}^l(\mathfrak{f}^i\mathfrak{g})^{-1}=&\\
&=\mathfrak{f}^i\mathfrak{f}^{-k}\mathfrak{g}^{1+l}(\mathfrak{f}^i\mathfrak{g})^{-1}=&\\
&= \mathfrak{f}^{i-k}\mathfrak{g}^{1+l}\mathfrak{g} \mathfrak{f}^{-i} =& \\
&= \mathfrak{f}^{i-k}\mathfrak{g}^{2+l} \mathfrak{f}^{-i} =& \\
&= \mathfrak{f}^{i-k}\mathfrak{g}^{l} \mathfrak{f}^{-i} =& \\
&&=\begin{cases}
\mathfrak{f}^{2i-k} \mathfrak{g} &                   \Leftarrow                   l = 1\\
\mathfrak{f}^{-k} &                   \Leftarrow                   l = 0
\end{cases}
\end{align}

Y desde aqu{\'{\i}} vemos que $\left\langle \mathfrak f \right\rangle \lhd \mathsf{Dih}_{p}$.

Y como de hecho no hemos utilizado la primalidad de $p$ en ning\'un sitio, los grupos $\mathsf{Dih}_{n} \triangleq \cyclicgr{n}\rtimes_{\varphi }\cyclicgr{2}$, $n                   \in \mathds{N}_{{}>2}$, siendo $\varphi : \cyclicgr{2}                    \longrightarrow                    \mathtt{Aut}(\cyclicgr{n})$, de forma que $\varphi_1 = \mathtt{id}_{\cyclicgr{n}}$ y $\varphi_{a                   \neq                    1}(b^{k \mod n}) = b^{(n-1)k \mod n}$. El orden de $\mathsf{o}(\mathsf{Dih}_{n})=\left|\cyclicgr{n}                   \times                    \cyclicgr{2}\right|=2n$, y ser\'a no-conmutativo.
Podemos ver por c\'alculo directo en la f\'ormula general, que los elementos $\mathfrak{g},\mathfrak{gf},\mathfrak{gf}^2, \ldots , \mathfrak{gf}^{n-1}$ son de orden $2$, que son conjugados entre si, y por lo tanto generan subgrupos no-normales y que el subgrupo generado por $\mathfrak{f}$ es caracter{\'{\i}}stico. La \'unica diferencia entre que sea primo o no, es que cuando no sea primo aparecer\'an otros subgrupos, que ser\'an subgrupos de $\left\langle\mathfrak{f}\right\rangle$, am\'en del grupo de automorfismos, que diferir\'a en su tamaño, ya que si $n$ es primo hay un total de $n-1$ automorfismos de $\mathtt{Aut}(\left\langle\mathfrak{f}\right\rangle) = \left\lbrace{\mathfrak{f}                  \mapsto \mathfrak{f},\mathfrak{f}             \mapsto \mathfrak{f}^2,\mathfrak{f}             \mapsto \mathfrak{f}^3,\ldots ,\mathfrak{f}             \mapsto \mathfrak{f}^{n-1} }\right\rbrace$. De esta forma $\left|\mathtt{Aut}(\left\langle\mathfrak{f}\right\rangle)\right|=n-1$ si $n             \in \mathds{P}\setminus\{2\}$. De forma m\'as general, el orden ser\'a $\phi_{E}(n)$, la funci\'on de Euler que cuenta los coprimos con $n$ que son menores que $n$.
Para los elementos de orden $2$, podemos buscar los siguientes isomorfismos: $\left\{{\mathfrak{g}                  \mapsto \mathfrak{g},\mathfrak{g}             \mapsto \mathfrak{g}\mathfrak{f},\mathfrak{g}             \mapsto \mathfrak{g}\mathfrak{f}^2,\ldots,,\mathfrak{g}                  \mapsto \mathfrak{g}\mathfrak{f}^{n-1}}\right\}$ que nos da un total de $n$ elecciones. As{\'{\i}} $Aut(\mathsf{Dih}_{n})$ tiene $n\phi_{E}(n)$ elecciones posibles. En el caso que $n=p$ primo, el orden del grupo de automorfismos es $p(p-1)$.

El subgrupo de Frattini de $\dihgr{p}$ lo confeccionaremos explicitando primero la familia de maximales:
${\mathcal{M}}_{\dihgr{p}} = \left\{ { \langle \mathfrak{f} \rangle,\langle \mathfrak{g} \rangle,\langle
\mathfrak{gf}\rangle,\ldots,\langle \mathfrak{gf}^{p-1}\rangle } \right\} $. As{\'{\i}} queda $\Phi_{\dihgr{p}} =
\bigcap\mathcal{M}_{\dihgr{p}} = \{1\}$.

Adem\'as, el centro de este grupo es claramente $\mathsf{Z}\left({\dihgr{p}}\right)=\{1\}$, de d\'onde $Int\left({\dihgr{p}}\right)                  \cong                  \nicefrac{{\dihgr{p}}}{\mathsf{Z}({\dihgr{p}})}                  \cong                  {\dihgr{p}}$, con lo que ordinal de los automorfismos interiores es $2p$, de los $(p-1)^2$ del total de automorfismos.

El subgrupo de Fitting es el generado por los subgrupos normales y nilpotentes, por lo tanto solo tenemos que ver si $\langle \mathfrak{f} \rangle$ es nilpotente. Pero esto es claro porque es un grupo c{\'{\i}}clico, y por lo tanto abeliano, por lo que el subgrupo derivado ${\langle \mathfrak{f} \rangle}^{(1)} = \{1\}$. Como no hay otros normales, $\mathsf{F}(\dihgr{p})=\langle \mathfrak{f} \rangle$.

El subgrupo derivado primero del grupo dihedral es
$\left\{ \left[ \mathfrak{g}^{i}\mathfrak{f}^{j},\mathfrak{g}^{k}\mathfrak{f}^{l} \right]=\mathfrak{f}^{2(j-l)} \right\} = \langle \mathfrak{f} \rangle$.
 As{\'{\i}} ${\dihgr{p}}^{(1)} = \langle \mathfrak{f} \rangle$. Sin embargo ${\mathsf{L}}_{2}\left( \dihgr{p}
\right) = {\mathsf{L}}_{1}\left( \dihgr{p} \right) = {\dihgr{p}}^{(1)} = \langle \mathfrak{f} \rangle$.
 Sin embargo ${\dihgr{p}}^{(2)} = \left\{1\right\}$.


Para el caso del grupo c{\'{\i}}clico $\graut{\cyclicgr{2p}}       \cong        \cyclicgr{p-1}$.


El diagrama de Hasse de los subgrupos para $\cyclicgr{2p}$ ser\'a:

\begin{center}
    \begin{tikzpicture}
    % First, locate each of the nodes and name them
    \node (arriba) at (0,0) {$\cyclicgr{2p}                     \cong                     \{1,a,c\triangleq a^2,\ldots,b\triangleq a^p,\ldots,a^{2p-1}\}$};
    \node (medioarriba) at (3,-2) {$\cyclicgr{p}                    \cong                    \{1,c,\ldots,c^{p-1}\}$};
    \node (medioabajo) at (-3,-4) {$\cyclicgr{2}                    \cong                    \{1,b\}$};
    \node (abajo) at (0,-6) {$\cyclicgr{1}                    \cong                    \{1_{\cyclicgr{2p}}\}$};

    % Now draw the lines:
    \draw [black, thick] (arriba) -- (medioarriba);
    \draw [black, thick] (arriba) -- (medioabajo);
    \draw [black, thick] (medioarriba) -- (abajo);
    \draw [black, thick] (medioabajo) -- (abajo);
    \end{tikzpicture}
\end{center}

El diagrama de Hasse de los subgrupos para $\dihgr{p}$ ser\'a:


\begin{center}
    \begin{tikzpicture}
    \label{tab:GruposDih\'edricosDeOrden2p}
    % First, locate each of the nodes and name them
    \node (arriba) at (0,0) {$\mathsf{Dih}_{p}                     \cong                     \left\langle \mathfrak{g},\mathfrak{f} \right\rangle$};
    \node (medioarriba) at (2,-2) {$\cyclicgr{p}                    \cong                    \{1,\mathfrak{f},\ldots,{\mathfrak{f}}^{p-1}\}$};
    \node (medioabajo1) at (-2,-4) {$                  \cong                  \{1,\mathfrak{g}\}$};
    \node (medioabajo2) at (-4,-4) {$                  \cong                  \{1,\mathfrak{gf}\}$};
    \node (medioabajo3) at (-6,-4) {$\ldots$};
    \node (medioabajo4) at (-8,-4) {$\cyclicgr{2}                    \cong                    \{1,{\mathfrak{gf}}^{p-1}\}                  \cong                  $};
    \node (abajo) at (0,-6) {$\cyclicgr{1}                    \cong                    \{1_{\mathsf{Dih}_{p}}\}$};

    % Now draw the lines:
    \draw [black, thick] (arriba) -- (medioarriba);
    \draw [blue, thick] (arriba) -- (medioabajo1);
    \draw [blue, thick] (arriba) -- (medioabajo2);
    \draw [blue, thick] (arriba) -- (medioabajo4);
    \draw [black, thick] (medioarriba) -- (abajo);
    \draw [black, thick] (medioabajo1) -- (abajo);
    \draw [black, thick] (medioabajo2) -- (abajo);
    \draw [black, thick] (medioabajo4) -- (abajo);
    \end{tikzpicture}
\end{center}





La tabla con la que ya podemos aumentar las tablas de grupos por \'ordenes es:





\begin{table}[h!]
    \begin{center}
        \caption{Grupos de orden $2p$}
        \label{tab:GruposOrden2p}
        \begin{tabular}{c|c|c|c|c|c|c} % <-- Alignments: 1st column left, 2nd middle and 3rd right, with vertical lines in between
            \textbf{} & \tiny\textbf{n\'umero} &\tiny\textbf{numeraci\'on} & \tiny\textbf{nombre} & \textbf{} & {} & {} \\
            \textbf{} & \tiny\textbf{de} &\tiny\textbf{dentro} & \tiny\textbf{del} & \tiny\textbf{} & {} & {} \\
            \tiny\textbf{orden} & \tiny\textbf{grupos} & \tiny\textbf{del} & \tiny\textbf{grupo} & \tiny\textbf{abeliano}  & \tiny\textbf{nilpotente} & \tiny\textbf{resoluble} \\
            \textbf{} & \tiny\textbf{no} & \tiny\textbf{orden} & \textbf{} & \textbf{} & {} & {} \\
            \textbf{} & \tiny\textbf{isomorfos} &\textbf{} & \textbf{} & \textbf{} & {} & {} \\
            \hline
            6 & 2 & 1 & ${\textsf{C}}_{2*3}$ & Si & Si & Si \\
            \hline
            6 & 2 & 2  & ${{\textsf{Dih}}_{3}}                    \cong                    {{\textsf{C}}_{3}}\rtimes{{{\textsf{C}}_{2}}}$ & No & No & Si \\
            \hline
            10 & 2 & 1 & ${\textsf{C}}_{2*5}$ & Si & Si & Si \\
            \hline
            10 & 2 & 2  & ${{\textsf{Dih}}_{5}}                    \cong                    {{\textsf{C}}_{5}}\rtimes{{{\textsf{C}}_{2}}}$ & No & No & Si \\
            \hline
            14 & 2 & 1 & ${\textsf{C}}_{2*7}$ & Si & Si & Si \\
            \hline
            14 & 2 & 2  & ${{\textsf{Dih}}_{7}}                    \cong                    {{\textsf{C}}_{7}}\rtimes{{{\textsf{C}}_{2}}}$ & No & No & Si \\
            \hline
            22 & 2 & 1 & ${\textsf{C}}_{2*11}$ & Si & Si & Si \\
            \hline
            22 & 2 & 2  & ${{\textsf{Dih}}_{11}}                    \cong                    {{\textsf{C}}_{11}}\rtimes{{{\textsf{C}}_{2}}}$ & No & No & Si \\
            \hline
            26 & 2 & 1 & ${\textsf{C}}_{2*13}$ & Si & Si & Si \\
            \hline
            26 & 2 & 2  & ${{\textsf{Dih}}_{13}}                    \cong                    {{\textsf{C}}_{13}}\rtimes{{{\textsf{C}}_{2}}}$ & No & No & Si \\
            \hline
            34 & 2 & 1 & ${\textsf{C}}_{2*17}$ & Si & Si & Si \\
            \hline
            34 & 2 & 2  & ${{\textsf{Dih}}_{17}}                    \cong                    {{\textsf{C}}_{17}}\rtimes{{{\textsf{C}}_{2}}}$ & No & No & Si \\
            \hline
            38 & 2 & 1 & ${\textsf{C}}_{2*19}$ & Si & Si & Si \\
            \hline
            38 & 2 & 2  & ${{\textsf{Dih}}_{19}}                    \cong                    {{\textsf{C}}_{19}}\rtimes{{{\textsf{C}}_{2}}}$ & No & No & Si \\
            \hline
            46 & 2 & 1 & ${\textsf{C}}_{2*23}$ & Si & Si & Si \\
            \hline
            46 & 2 & 2  & ${{\textsf{Dih}}_{23}}                    \cong                    {{\textsf{C}}_{23}}\rtimes{{{\textsf{C}}_{2}}}$ & No & No & Si \\
            \hline
            58 & 2 & 1 & ${\textsf{C}}_{2*29}$ & Si & Si & Si \\
            \hline
            58 & 2 & 2  & ${{\textsf{Dih}}_{29}}                    \cong                    {{\textsf{C}}_{29}}\rtimes{{{\textsf{C}}_{2}}}$ & No & No & Si \\
            \hline
            62 & 2 & 1 & ${\textsf{C}}_{2*31}$ & Si & Si & Si \\
            \hline
            62 & 2 & 2  & ${{\textsf{Dih}}_{31}}                    \cong                    {{\textsf{C}}_{31}}\rtimes{{{\textsf{C}}_{2}}}$ & No & No & Si \\
            \hline
        \end{tabular}
    \end{center}
\end{table}








    Queda a\'un una pregunta por responder: ?`Hay m\'as grupos que los dichos del orden $2p$?

    De tipo abeliano claramente no, ya se dio la explicaci\'on, pero tambi\'en se puede aludir al teorema de
    clasificaci\'on de los grupos abelianos y finitos.

    Para el caso no-abeliano, no cabe que existan m\'as subgrupos normales, pues en ese caso el grupo pasar{\'{\i}}a a
    ser abeliano por ser ambas partes del producto abelianas.

    Pensando en p-subgrupos y p-subgrupos de Sylow, siempre tendremos $2$ familias de subgrupos de Sylow,
    $\mathcal{S}yl_2 (\grG)$ y $\mathcal{S}yl_p (\grG)$. Los p-subgrupos ser\'an necesariamente normales porque $[\grG:\grH_p]=2$,
    de d\'onde los elementos de $\mathcal{S}yl_2 (\grG)$ ser\'an no-normales. As{\'{\i}} que lo que queda por dilucidar
    es si puede ser $\left|\mathcal{S}yl_{p}(\grG)\right|>1$. Pongamos que sea as{\'{\i}}: entonces existir\'an
    $\grH_p^1,\grH_p^2$ p-subgrupos. Como son diferentes, por hip\'otesis, $\grH_p^1       \cap \grH_p^2 = \{1\}$, y as{\'{\i}},
    si hay $r$ p-subgrupos, tendremos un total de $r(p-1)+1$ elementos debido a los p-subgrupos, y $r       \le 2$,
    de d\'onde $r=2$ significa que solo habr{\'{\i}}a un 2-subgrupo, pero entonces el 2-subgrupo habr{\'{\i}}a
    de ser normal y de ah{\'{\i}} el grupo abeliano. Luego no hay m\'as grupos no-abelianos que ${\dihgr{p}}$.

    \chapter{Grupos de orden 8}

    Los grupos de orden $8$ podr{\'{\i}}an ser tratados como grupos de orden $p^3$ con $p       \in \setP$, de la misma forma que
    los grupos de orden $4$ podr{\'{\i}}an haber sido tratados como los grupos de orden $p^2$. No lo hago as{\'{\i}} por
    algunas diferencias entre lo grupos con orden potencia de un primo impar y los grupos con orden potencia de $2$.

    Los grupos abelianos de orden $8$ son f\'aciles de enumerar gracias al teorema de clasificaci\'on de los grupos finitos
    abelianos. Salvo isomorfismo solo pueden ser $\cyclicgr{2}      \times      \cyclicgr{2}      \times      \cyclicgr{2}$, $\cyclicgr{4}      \times      \cyclicgr{2}$ y $\cyclicgr{8}$.
    Al ser abelianos y por su construcci\'on mediante producto directo podemos reconstruir f\'acilmente su tablas de Cayley,
    su ret{\'{\i}}culo de subgrupos (se invierte el teorema de Lagrange) y todos son normales.

    Los diagramas de Hasse correspondientes son:



\resizebox{\linewidth}{!}{
\begin{tikzpicture}
        % First, locate each of the nodes and name them
        \node (C8) at (3,0)
        {$\cyclicgr{8}                  \cong                  \langle a\vert a^8\rangle$};
        \node (C4xC2) at (-1,0)
        {$\mathsf{C}_4            \times            \mathsf{C}_2                   \cong                    \left\langle a,b \vert a^4=b^2 \right\rangle$};
        \node (C2xC2xC2) at (-7,0)
        {$\mathsf{C}_2                   \times                  \mathsf{C}_2                   \times                  \mathsf{C}_2                   \cong                  \left\langle a,b,c \vert a^2=b^2=c^2\right\rangle $};
        \node (C8C4) at (3,-2)
        {$                  \cong                  \left\langle a^2\right\rangle $};
        \node (C8C4C2) at (3,-4)
        {$                  \cong                  \left\langle{a^4}\right\rangle$};
        \node (C8abajo) at (3,-6)
        {$                  \cong                  \left\langle{1_{\mathsf{C}_8}}\right\rangle$};
        \node (C4xC2C4) at (-2,-2)
        {$\mathsf{C}_4                   \cong                  \left\langle{a}\right\rangle$};
        \node (C4xC2C4C2) at (-2,-4)
        {$                  \cong                  \left\langle{a^2}\right\rangle$};
        \node (C4xC2C2) at (1,-4)
        {$                  \cong                  \left\langle{b}\right\rangle$};
        \node (C4xC2abajo) at (-1,-6)
        {$                  \cong                  \left\langle{1_{\mathsf{C}_4                   \times                   \mathsf{C}_2}}\right\rangle$};
        \node (C2xC2xC21) at (-5,-2)
        {$                  \cong                  \left\langle{a,b}\right\rangle$};
        \node (C2xC2xC22) at (-7,-2)
        {$                  \cong                  \left\langle{a,c}\right\rangle$};
        \node (C2xC2xC23) at (-9,-2)
        {$\cyclicgr{2}                  \times                  \cyclicgr{2}                  \cong                  \left\langle{b,c}\right\rangle$};
        \node (C2xC2xC212) at (-5,-4)
        {$                  \cong                  \left\langle{a}\right\rangle$};
        \node (C2xC2xC222) at (-7,-4)
        {$                  \cong                  \left\langle{b}\right\rangle$};
        \node (C2xC2xC232) at (-9,-4)
        {$\cyclicgr{2}                  \cong                  \left\langle{c}\right\rangle$};
        \node (C2xC2xC2abajo) at (-7,-6)
        {$\cyclicgr{1}                  \cong                  \left\langle{1_{\mathsf{C}_2                   \times                  \mathsf{C}_2                   \times                  \mathsf{C}_2}}\right\rangle$};
        % Now draw the lines:
        \draw [black, thick] (C8) -- (C8C4);
        \draw [black, thick] (C8C4) -- (C8C4C2);
        \draw [black, thick] (C8C4C2) -- (C8abajo);
        \draw [black, thick] (C4xC2) -- (C4xC2C4);
        \draw [black, thick] (C4xC2) -- (C4xC2C2);
        \draw [black, thick] (C4xC2C4) -- (C4xC2C4C2);
        \draw [black, thick] (C4xC2C2) -- (C4xC2abajo);
        \draw [black, thick] (C4xC2C4C2) -- (C4xC2abajo);
        \draw [black, thick] (C2xC2xC2) -- (C2xC2xC21);%ab
        \draw [black, thick] (C2xC2xC2) -- (C2xC2xC22);%ac
        \draw [black, thick] (C2xC2xC2) -- (C2xC2xC23);%bc
        \draw [black, thick] (C2xC2xC21) -- (C2xC2xC212);%ab->a
        \draw [black, thick] (C2xC2xC21) -- (C2xC2xC222);%ab->b
        \draw [black, thick] (C2xC2xC22) -- (C2xC2xC212);%ac->a
        \draw [black, thick] (C2xC2xC22) -- (C2xC2xC232);%ac->c
        \draw [black, thick] (C2xC2xC23) -- (C2xC2xC222);%bc->b
        \draw [black, thick] (C2xC2xC23) -- (C2xC2xC232);%bc->c
        \draw [black, thick] (C2xC2xC212) -- (C2xC2xC2abajo);
        \draw [black, thick] (C2xC2xC222) -- (C2xC2xC2abajo);
        \draw [black, thick] (C2xC2xC232) -- (C2xC2xC2abajo);
\end{tikzpicture}
}



    Los grupos no-abelianos han de tener $\mathsf{o}\left(\centro{\grG}\right)       \in \{2,4\}$, por ser un 2-grupo.
    Si el orden del centro fuese $2$, $\nicefrac{\grG}{\centro{\grG}}       \cong       \cyclicgr{4}$ \'o
    $\nicefrac{\grG}{\centro{\grG}}       \cong       \cyclicgr{2}       \times      \cyclicgr{2}$. Pero el caso $\nicefrac{\grG}{\centro{\grG}}       \cong       \cyclicgr{4}$
    hemos de rechazarlo, por que si ese cociente es c{\'{\i}}clico entonces $\grG$ es abeliano contra la hip\'otesis.
    Si el orden del centro fuese $4$, $\nicefrac{\grG}{\centro{\grG}}       \cong       \cyclicgr{2}$, ser{\'{\i}}a c{\'{\i}}clico y
    por lo tanto abeliano contra la hip\'otesis.

    Luego $\grG$ no-abeliano, $\left|\grG\right|=8$, entonces $\centro{\grG}       \cong       \cyclicgr{2}$, y
    $\nicefrac{\grG}{\centro{\grG}}       \cong       \cyclicgr{2}       \times       \cyclicgr{2}$.

    Ahora falta ver si $\mathcal{M}ax_{2}(\grG)=\left\lbrace \grH        \preccurlyeq       \grG \vert \mathsf{o}(\grH)=4 \right\rbrace$
    contiene grupos normales o no-normales, pero en cualquier caso isomorfos entre s{\'{\i}}. Caben dos posibilidades,
    que $\left| \mathcal{M}ax_2 (\grG)\right|=1$ y que $\left| {{\mathcal{M}ax}_{2}}(\grG)\right|>1$.

    Exploremos los subgrupos maximales propios que llamaremos (que no son propiamente los subgrupos de Sylow de
    $\grG$ sino sus maximales propios de orden $4$) $\left| {{\mathcal{M}ax}_{2}}(\grG)\right|>1$:

    El primer caso es $\left| {{\mathcal{M}ax}_{2}}(\grG)\right|=2$,
tomando que la intersecci\'on es $\{1\}$, hacen un total de $2(4-1)+1=7$ elementos
en los 2-subgrupos de Sylow, de d\'onde necesitar{\'{\i}}amos un elemento extra que ser{\'{\i}}a
de orden $2$. Esto hace que el n\'umero de elementos final se duplicara y excediera
el orden buscado. Pongamos por lo tanto que la intersecci\'on de los 2-subgrupos maximales propios
es un subgrupo de orden $2$, de d\'onde $2(4-2)+2=6$, y necesitar{\'{\i}}amos de nuevo
alg\'un elemento extra de orden $2$, que har{\'{\i}}a que el orden de $\grG$ fuese mayor que $8$.
Por lo tanto este caso se puede desechar.

    El segundo caso es $\left| {{\mathcal{M}ax}_{2}}(\grG)\right|=3$, pero $3*4=12 > 8$,
por lo que habremos de tomar que todos los subgrupos maximales propios tienen una intersecci\'on
com\'un de orden $2$. Salen $3(4-2)+2=8$ elementos exactamente. De aqu{\'{\i}} sabemos
que ha de existir en cada subgrupo maximal propio $2$ elementos que s\'olo est\'an en ese subgrupo.
Llamemos $ {{\mathcal{M}ax}_{2}}\left(\grG\right)=\left\{ \grN_1,\grN_2,\grN_3 \right\} $,
llamemos
$\grH = \grN_1          \cap \grN_2 = \grN_1       \cap \grN_3 = \grN_2        \cap \grN_3 =
\grN_1        \cap \grN_2       \cap \grN_3 = \left\{ \overline{1},1 \right\}$.
Sea $\grN_1=\left\{ \overline{a},a\right\}       \cup       \grH$, $\grN_2=\left\{\overline{b},b\right\}       \cup       \grH$ y
$\grN_3=\left\{ \overline{c},c\right\}      \cup       \grH$. Ahora la pregunta es ?`cuales pueden ser los
ordenes de $a$? Pueden ser $2$ o $4$. Si fuese $2$, $\grK_1 = \{1,a\}$, $\grK_2 = \{1,\overline{a}\}$,
y as{\'{\i}} con $b$ y $c$. Pero entonces todos los subgrupos de orden $2$ ser{\'{\i}}an isomorfos,
y los subgrupos maximales propios ser{\'{\i}}an el $4$-grupo de Klein, pero
tambi\'en lo ser{\'{\i}}an $\{1,a\}        \times        \{1,b\}$ y saldr{\'{\i}}an muchos m\'as de $3$ grupos maximales propios.
As{\'{\i}} que solo nos queda que los
grupos maximales propios sean isomorfos a $\cyclicgr{4}$.
Por lo tanto el orden de los elementos $\{a,\overline{a},b,\overline{b},c,\overline{c}\}$
ser\'a $4$. Sabemos que $a^2       \neq       \overline{a}$ pues, si no, su orden ya
ser{\'{\i}}a mayor que $4$. Por lo tanto solo queda
$a^2=\overline{a}^2=b^2=\overline{b}^2=c^2=\overline{c}^2=\overline{1}$ y
adem\'as $\overline{1}^2=1$. S\'olo queda ver c\'omo son el resto de productos. Por lo pronto
$a\overline{1}=\overline{1}a=\overline{a}$, $\overline{a}\overline{1}=\overline{1}\overline{a}=a$,
e igualmente con $b$ y $c$. A todo esto ha de quedar ya claro que $a^{-1}=\overline{a}$
y as{\'{\i}} con $b$ y $c$, y que $\overline{1}^{-1}=\overline{1}$.

    Queda ver solamente como es el producto $ab$ etc. Sabemos que $ab       \neq       \overline{1}$
porque entonces $b=\overline{a}$, y con el mismo tipo de argumento llegamos a que
$ab       \in \{c,\overline{c}\}$. Haremos la prueba con $ab=c$, $bc=a$, $ca=b$, y por la
no-conmutatividad $ba=\overline{c}$, $cb=\overline{a}$, $ac=\overline{b}$.
Si hacemos la prueba con $ab=\overline{c}$, $bc=a$, $ca=b$, y por la no-conmutatividad
$ba=c$, $cb=a$, $ac=b$, y obtenemos un grupo isomorfo al anterior.

    Nos quedamos con $ab=c$, $bc=a$, $ca=b$, $ba=\overline{c}$, $cb=\overline{a}$,
$ac=\overline{b}$, $a\overline{1}=\overline{1}a=\overline{a}$,
$\overline{a}\overline{1}=\overline{1}\overline{a}=a$,
$b\overline{1}=\overline{1}b=\overline{b}$,
$\overline{b}\overline{1}=\overline{1}\overline{b}=b$,
$c\overline{1}=\overline{1}c=\overline{c}$,
$\overline{c}\overline{1}=\overline{1}\overline{c}=c$,
$a^2=\overline{a}^2=b^2=\overline{b}^2=c^2=\overline{c}^2=1$,
$a\overline{a}=b\overline{b}=c\overline{c}=\overline{1}^2=1$.

    Definimos al grupo del p\'arrafo anterior como $\grQ_8 = \langle a,b | a^4=b^4=(ab)^4=1, a^2=b^2=(ab)^2, ab=a^{-1}b^{-1} \rangle$. Es f\'acil ver que $\langle a\rangle \lhd \grQ_8$, $\langle b\rangle \lhd \grQ_8$, $\langle ab\rangle \lhd \grQ_8$ (son todos los subgrupos de orden $4$) y que $\centro{\grQ_8}=\{1,\overline{1}\}\lhd\grQ_8$ que es \'unico subgrupo de orden $2$. Es un grupo hamiltoniano (Dedekind no-abeliano).

    El diagrama de Hasse de $\grQ_8$ ser\'a:


    \begin{center}
        \begin{tikzpicture}
        \label{tab:DiagramaDeHasseParaElGrupoQuaternion8}
        % First, locate each of the nodes and name them
        \node (arriba) at (0,0) {$\mathsf{Q}_{8}                     \cong                     \left\langle a,b \right\rangle$};
        \node (medioarribaderecha) at (2,-2) {$                  \cong                    \left\langle ab \right\rangle$};
        \node (medioarriba) at (0,-2) {$                  \cong                  \left\langle b\right\rangle $};
        \node (medioarribaizquierda) at (-2,-2) {$\cyclicgr{4}                    \cong                  \left\langle a\right\rangle $};
        \node (medioabajo) at (0,-4) {$\cyclicgr{2}                  \cong                  \left\langle{a^2=b^2=(ab)^2}\right\rangle$};
        \node (abajo) at (0,-6) {$\cyclicgr{1}                    \cong                    \{1_{\mathsf{Q}_{8}}\}$};

        % Now draw the lines:
        \draw [black, thick] (arriba) -- (medioarribaderecha);
        \draw [black, thick] (arriba) -- (medioarriba);
        \draw [black, thick] (arriba) -- (medioarribaizquierda);
        \draw [black, thick] (medioarribaderecha) -- (medioabajo);
        \draw [black, thick] (medioarriba) -- (medioabajo);
        \draw [black, thick] (medioarribaizquierda) -- (medioabajo);
        \draw [black, thick] (medioabajo) -- (abajo);
        \end{tikzpicture}
    \end{center}

    Antes de pasar a los subgrupos caracter{\'{\i}}sticos y con nombre propio o a los grupos de automorfismos
    a\'un podemos ver algo interesante. Veamos como se extiende el grupo $\mathsf{Q}_{8}$ desde $\cyclicgr{4}$
    por $\cyclicgr{2}$, o lo que es lo mismo, como existen secuencias exactas cortas desde $\cyclicgr{4}$
    hasta $\cyclicgr{2}$:



\begin{center}
\begin{flalign*}
& \bold{1}
                  \longrightarrow
\langle a \rangle
\underset{{\substack{{a                  \mapsto a}\\{b                  \mapsto b}}}}{\overset{\alpha }{                  \longrightarrow                  }}
\mathsf{Q}_{8}
\underset{{\substack{{a                  \mapsto 1}\\{b                  \mapsto \overline{1}}}}}{\overset{\beta_a }{                  \longrightarrow                  }}
\langle \overline{1} \rangle
                  \longrightarrow
\bold{1}\\
& \bold{1}
                  \longrightarrow
\langle b \rangle
\underset{{\substack{{a                  \mapsto a}\\{b                  \mapsto b}}}}{\overset{\alpha }{                  \longrightarrow                  }}
\mathsf{Q}_{8}
\underset{{\substack{{a                  \mapsto \overline{1}}\\{b                  \mapsto 1}}}}{\overset{\beta_b }{                  \longrightarrow                  }}
\langle \overline{1} \rangle
                  \longrightarrow
\mathbf{1}
\end{flalign*}
\end{center}


	De aqu{\'{\i}} podemos pensar en modelar una acci\'on de $\cyclicgr{4}$ en $\cyclicgr{2}$ que nos de
la operaci\'on propia de $\grQ_{8}          \equiv          \dicgr{2}$. Expongo primero la
dos acciones, (la operaci\'on tiene que variar del producto semidirecto, y resulta haber una segunda
acci\'on de $\cyclicgr{2}       \times       \cyclicgr{2}$ sobre $\cyclicgr{4}$ y que
depende de las dos entradas de orden $2$ en el par ordenado, consistiendo la acci\'on en la
multiplicaci\'on a derechas por un determinado elemento de $\cyclicgr{4}$ determinado por la
pareja de ${\cyclicgr{2}}{      \times      }{\cyclicgr{2}}$):



\begin{flalign*}
&{\phi }\; : \;
{
{\cyclicgr{2}}
         \longrightarrow
\mathtt{Aut}
\left(
{\cyclicgr{4}}
\right)
}
\; ::\;
{
{y}
\mapsto
{\langle{{x^k}         \mapsto {{x}^{4-k}}}\rangle}
}
\end{flalign*}

La siguiente funci\'on elige un elemento de
$\cyclicgr{4}$ en funci\'on de $\cyclicgr{2}      \times      \cyclicgr{2}$
(constituye un homomorfismo):

\begin{flalign*}
&\zeta \;:\;
{{{\mathsf{C}}_{2}}{                  \times                  }{{\mathsf{C}}_{2}}}
                  \longrightarrow
{{\mathsf{C}}_{4}}\;::\;\left(y,z\right) \mapsto
\begin{cases}
b^2                   \Leftarrow                   {y=z                  \neq                   1}\\
1                   \Leftarrow                   (y                  \neq                   z)                  \vee (y=z=1)
\end{cases}
\end{flalign*}

	Ahora la acci\'on que buscamos es la multiplicaci\'on por la derecha del segundo componente del par por el valor devuelto por $\zeta$.

\begin{flalign*}
&\theta \;:\;
{{{\mathsf{C}}_{2}}{                  \times                  }{{\mathsf{C}}_{2}}{                  \times                  }{{\mathsf{C}}_{4}}}
                  \longrightarrow
{{\mathsf{C}}_{4}}\;::\;\left(y,z,w\right)                   \mapsto w\zeta \left(y,z\right)
\end{flalign*}


	Desempaquetada la acci\'on queda:


\begin{flalign*}
&\theta \;:\;
{{{\mathsf{C}}_{2}}{                  \times                  }{{\mathsf{C}}_{2}}{                  \times                  }{{\mathsf{C}}_{4}}}
                  \longrightarrow
{{\mathsf{C}}_{4}}\;::\;\left(y,z,w\right) \mapsto
\begin{cases}
wf^2                   \Leftarrow                   {y=z                  \neq                   1}\\
w                   \Leftarrow                   (y                  \neq                   z)                  \vee (y=z=1)
\end{cases}
\end{flalign*}


	La operaci\'on quedar{\'{\i}}a:

\begin{flalign*}
&(x_1,y_1),(x_2,y_2)         \in \cyclicgr{2}         \times         \cyclicgr{4}\\
&\{a,1\}         \cong         \cyclicgr{2}\\
&\{f,f^2,f^3,1\}         \cong         \cyclicgr{4}\\
&\left({x_1},{y_1}\right)\left({x_2},{y_2}\right)
=
\left(
{x_1}{x_2},
{\phi }{\left({x_2}\right)}{\left({y_1}\right)}{y_2}
{\zeta }{\left({x_1},{x_2}\right)}
\right)\\
&{\left({x_1},{y_1}\right)}{\left({x_2},{y_2}\right)}=
\begin{cases}
&(x_1,y_1 y_2\zeta (x_1,1))                   \Leftarrow                   x_2 = 1\\
&(x_1 a,{y_1}^{-1} {y_2}\zeta (x_1,a))           \Leftarrow                   x_2 = a
\end{cases}
\end{flalign*}

	Hasta el momento, la operaci\'on ser{\'{\i}}a la del grupo $\mathsf{Dih}_{4}$ si no fuese por el aditamento de la acci\'on
$\theta : \cyclicgr{2}                   \times                   \cyclicgr{2}                  \times                   \cyclicgr{4}                   \longleftarrow                   \cyclicgr{4}$, que act\'ua sobre
la segunda componente del operando derecho. Hagamos ya las sustituci\'on:
	
\begin{flalign*}
&(x_1,y_1),(x_2,y_2)                  \in \cyclicgr{2}                  \times                  \cyclicgr{4}\\
&\{a,1\}                  \cong                  \cyclicgr{2}\\
&\{f,f^2,f^3,1\}                  \cong                  \cyclicgr{4}\\
&(x_1,y_1) (x_2,y_2)=
\begin{cases}
&(x_1,y_1 y_2 \zeta (x_1,1)                   \Leftarrow                   x_2 = 1\\
&(x_1 a,y_1^{-1} y_2 \zeta (x_1,a)                   \Leftarrow                   x_2 = a
\end{cases}
\\
&(x_1,y_1) (x_2,y_2)=
\begin{cases}
&(1,y_1 y_2)                   \Leftarrow                   \left(x_2 = 1                   \wedge                   x_1 = 1\right)\\
&(a,{y_1}^{-1} y_2)                   \Leftarrow                   \left(x_2 = a \ly x_2                   \neq                   x_1 \right)\\
&(a,y_1 y_2)                   \Leftarrow                   \left(x_2 = 1 \ly x_2                   \neq                   x_1 \right)\\
&(1,y_1^{-1} y_2 f^2)                   \Leftarrow                   \left(x_2 = a                   \wedge                   x_1 = a\right)
\end{cases}
\end{flalign*}

	Resumiendo, la operaci\'on binaria de grupo de $\mathsf{Q}_{8}$ quedar{\'{\i}}a:

\begin{flalign*}
&(x_1,y_1),(x_2,y_2)                  \in \cyclicgr{2}                  \times                  \cyclicgr{4}\\
&\{a,1\}                  \cong                  \cyclicgr{2}\\
&\{f,f^2,f^3,1\}                  \cong                  \cyclicgr{4}\\
&(x_1,y_1) (x_2,y_2)=
\begin{cases}
&(1,y_1 y_2 )                   \Leftarrow                   \left(x_2 = x_1 = 1\right)\\
&(a,{y_1}^{-1} y_2 )                   \Leftarrow                   \left(x_2 =a \ly x_1=1\right)\\
&(a,y_1 y_2 )                   \Leftarrow                   \left(x_2 =1 \ly x_1=a\right)\\
&(1,y_1^{-1} y_2 f^2)                   \Leftarrow                   \left(x_2 = x_1 = a\right)
\end{cases}
\end{flalign*}

	Concretando la ecuaci\'on con los valores que pueden tomar las variables $y_1 , y_2$ que no son otra cosa que potencias de $f$ :
	
\begin{flalign*}
&\{a,1\}                  \cong                  \cyclicgr{2}\\
&\{f,f^2,f^3,1\}                  \cong                  \cyclicgr{4}\\
&(x_1,f^k),(x_2,f^l)                  \in \cyclicgr{2}                  \times                  \cyclicgr{4}\\
&(x_1,f^k) (x_2,f^l)=
\begin{cases}
&(1,f^{k+l})                   \Leftarrow                   \left(x_2 = x_1 = 1\right)\\
&(a,f^{4-k+l})                   \Leftarrow                   \left(x_2=a \ly x_1=1 \right)\\
&(a,f^{k+l})                   \Leftarrow                   \left(x_2=1 \ly x_1=a \right)\\
&(1,f^{2l-k})                   \Leftarrow                   \left(x_2 = x_1 = a\right)
\end{cases}
\end{flalign*}

	Para la implementaci\'on de esta operaci\'on con aritm\'etica modular, nos quedamos exclusivamente con los exponentes, consiguiendo as{\'{\i}} el algoritmo a implementar en el programa inform\'atico que nos dar\'a la tabla, los subgrupos y otros tratamientos del grupo:
	
\begin{flalign*}
&\{a,1\}                  \cong                  \cyclicgr{2}\\
&\{f,f^2,f^3,1\}                  \cong                  \cyclicgr{4}\\
&(a^i,f^k),(a^j,f^l)                  \in \cyclicgr{2}                  \times                  \cyclicgr{4}
\end{flalign*}


\begin{flalign*}
&(i,k)(j,l)=
\begin{cases}
(0,k+l) &                  \Leftarrow                   {\left(j = i = 0\right)}\\
(1,4-k+l) &                  \Leftarrow                   {\left( j = 1                   \wedge                   i=0\right)}\\
(1,k+l) &                  \Leftarrow                   {\left( j = 0                   \wedge                   i=1\right)}\\
(0,2l-k) &                  \Leftarrow                   {\left(j = i = 1\right)}
\end{cases}
\end{flalign*}


	La tabla de la operaci\'on de grupo resultante (quitado la fila y la columna de argumentos), o tabla de Cayley es, una vez cambiados todas las apariciones de $f^2$ por ${-1}$, y si no aparece solo el elemento ${-1}$, por ejemplo en $x{-1}$, cambiamos la expresión por ${-x}$. Hacemos otra trasformaci\'on, cambiamos $a$ por $\imath$, $f$ por $\jmath$ y $af$ por $\kappa$, y vemos si se parece m\'as al grupo cuaterni\'on en su forma habitual:


\begin{flalign*}
\begin{array}{|c|c|c|c|c|c|c|c|}
\hline
1&{\jmath}&{-1}&{-{\jmath}}&{\imath}&{\kappa }&{-{\imath}}&{-{\kappa }}\\
\hline
{\jmath}&{-1}&{-{\jmath}}&1&{-{\kappa }}&{\imath}&{\kappa }&{-{\imath}}\\
\hline
{-1}&{-{\jmath}}&1&{\jmath}&{-{\imath}}&{-{\kappa }}&{\imath}&{\kappa }\\
\hline
{-{\jmath}}&1&{\jmath}&{-1}&{\kappa }&{-{\imath}}&{-{\kappa }}&{\imath}\\
\hline
{\imath}&{\kappa }&{-{\imath}}&{-{\kappa }}&{-1}&{-{\jmath}}&1&{\jmath}\\
\hline
{\kappa }&{-{\imath}}&{-{\kappa }}&{\imath}&{\jmath}&{-1}&{-{\jmath}}&1\\
\hline
{-{\imath}}&{-{\kappa }}&{\imath}&{\kappa }&1&{\jmath}&{-1}&{-{\jmath}}\\
\hline
{-{\kappa }}&{\imath}&{\kappa }&{-{\imath}}&{-{\jmath}}&1&{\jmath}&{-1}\\
\hline
\end{array}
\end{flalign*}

	\begin{quote}
		Es importante darse cuenta que este m\'etodo extiende $\cyclicgr{4}                  \times                  \cyclicgr{2}$, de forma distinta a como lo hace el producto directo y el semidirecto. As{\'{\i}}, para los grupos de orden $4n$, aplicando este m\'etodo obtendremos los grupos $\mathsf{Dic}_{n}$, e incluso podemos pensar que en general tendremos una generalizaci\'on para los grupos de la forma $p^2 n$.
	\end{quote}	
	
	
		
	Seguimos con el estudio habitual de este grupo.	

    Como $\cyclicgr{4}                  \cong                  \langle a\rangle                  \cong                  \langle b\rangle                  \cong                  \langle ab\rangle$ son nilpotentes y normales, al igual que $\cyclicgr{2}                  \cong                  \langle\overline{1}\rangle$, el subgrupo de Fitting es $\mathsf{F}(\mathsf{Q}_8)=\mathsf{Q}_8$, y el grupo de Frattini, $\Phi_{\mathsf{Q}_8}=\langle \overline{1}\rangle$ y coincide con el centro $\Phi_{\mathsf{Q}_8}=\mathsf{Z}(\mathsf{Q}_8)$.

    En cuanto a $D_{\mathsf{Q}_8}^{(1)}=\left[{\mathsf{Q}_8,\mathsf{Q}_8}\right]=\langle \overline{1}\rangle$. Por \'ultimo $D_{\mathsf{Q}_8}^{(2)}=\{1\}$, luego es resoluble.

    Para ver si es nilpotente, tenemos $\mathsf{L}_{1}(\mathsf{Q}_8)=D_{\mathsf{Q}_8}^{(1)}=\Phi_{\mathsf{Q}_8}=\mathsf{Z}(\mathsf{Q}_8)=\langle{\overline{1}}\rangle$. Ahora $\mathsf{L}_{2}(\mathsf{Q}_8)=\left[\mathsf{Q}_8,\mathsf{L}_1(\mathsf{Q}_8)\right]=\left[\langle{a,b}\rangle,\langle{\overline{1}}\rangle\right]$. Esto nos da $\mathsf{L}_{2}(\mathsf{Q}_8)=\{1\}$, de d\'onde es nilpotente.

    En cuanto a los automorfismos internos $\mathtt{Int}(\mathsf{Q}_8)                  \cong                   \nicefrac{\mathsf{Q}_8}{\mathsf{Z}(\mathsf{Q}_8)}                  \cong                  \mathsf{C}_4$. Para los automorfismos en general habr{\'{\i}}a que mantener fijo el centro (es caracter{\'{\i}}stico), y solo quedan $a                  \mapsto a,a                  \mapsto \overline{a},a                  \mapsto b,a                  \mapsto \overline{b},a                  \mapsto {ab},a                  \mapsto \overline{ab}$ e id\'enticamente con $b$ en el origen, aunque no todas las elecciones son compatibles: da un m\'aximo de $36$ y tiene que ser menor. Tiene que ser m\'ultiplo de $3$ y de $4$, de d\'onde sale $\mathsf{o}\left(\mathtt{Aut}(\mathsf{Q}_8)\right)=24$. $\mathtt{Int}(\mathtt{Q}_8)                  \cong                  \mathsf{C}_4$ y sin prueba vemos que $\mathtt{Aut}(\mathtt{Q}_8)                  \cong                  \mathsf{S}_4$.

    Ahora nos quedamos con el caso que queda: $\mathcal{S}yl_2 (G)=\{N\}$. Sabemos que $\mathsf{Z}(G)$ es orden $2$, d\'onde $\nicefrac{G}{\mathsf{Z}(\mathsf{G})}                  \cong                   {{\mathsf{C}}_{2}}                  \times                  {{\mathsf{C}}_{2}}$. Sabemos que $N\lhd G$ porque $[G:N]=2$, ya que $\mathsf{o}(N)=4$. Ha de existir un grupo de orden $2$ por Cauchy, pero buscamos alguno que no sea normal en $G$. Sea ese grupo $H=\{1,g\}$ y $N=\{1,f,f^2,f^3\}$, nos proponemos hacer $\cyclicgr{4}\rtimes_{\varphi }\cyclicgr{2}$ d\'onde $\forall k \;\varphi_g (f^k)=f^{4-k}$ y $\forall k\;\varphi_1 (f^k)=f^k$. Obtenemos el grupo $G                  \cong                  \mathsf{Dih}_{4}$. En este caso queda claro que $K \triangleq\mathsf{Z}\left(\mathsf{Dih}_4\right)={1,f^2}$. Como todos los subgrupos son abelianos todos los subgrupos son nilpotentes. Normales solo ser\'an $N$ y $K$. ${\mathsf{D}}_{\mathsf{Dih}_4}^{(1)} = N$ y ${\mathsf{D}}_{\mathsf{Dih}_4}^{(2)} = \{1\}$ luego es resoluble. Adem\'as, $\mathsf{L}_{1}(\mathsf{Dih}_4)=N$, y de aqu{\'{\i}} $\mathsf{L}_{2}(\mathsf{Dih}_4)=\left[\mathsf{Dih}_4,N\right]=\{1\}$ luego es nilpotente. El subgrupo de Fitting es $\mathsf{F}(\mathsf{Dih}_4)=\langle f\rangle$ y el de Frattini $\Phi_{\mathsf{Dih}_4} = \{1\}$.

    La tabla de Cayley es:

\begin{flalign*}
\begin{array}{|c|c|c|c|c|c|c|c|}
\hline
1&f&f^2&f^3&g&gf&gf^2&gf^3\\
\hline
f&f^2&f^3&1&gf^3&g&gf&gf^2\\
\hline
f^2&f^3&1&f&gf^2&gf^3&g&gf\\
\hline
f^3&1&f&f^2&gf&gf^2&gf^3&g\\
\hline
g&gf&gf^2&gf^3&1&f&f^2&f^3\\
\hline
gf&gf^2&gf^3&g&f^3&1&f&f^2\\
\hline
gf^2&gf^3&g&gf&f^2&f^3&1&f\\
\hline
gf^3&g&gf&gf^2&f&f^2&f^3&1\\
\hline
\end{array}
\end{flalign*}

    El diagrama de Hasse de este grupo es:



    \begin{center}
        \begin{tikzpicture}
        % First, locate each of the nodes and name them
        \node (arriba) at (0,0) {$\mathsf{Dih}_{4}                     \cong                     \left\langle g,f \right\rangle$};
        \node (N) at (4,-2) {$\mathsf{C}_4                   \cong                    \left\langle f \right\rangle$};
        \node (Z) at (4,-4) {$                  \cong                  \left\langle f^2\right\rangle $};
        \node (H) at (0,-4) {$                  \cong                  \left\langle g\right\rangle $};
        \node (H1) at (-2,-4) {$                  \cong                  \left\langle{gf}\right\rangle$};
        \node (H2) at (-4,-4) {$                  \cong                  \left\langle{gf^2}\right\rangle$};
        \node (H3) at (-6,-4) {$\cyclicgr{2}                  \cong                  \left\langle{gf^3}\right\rangle$};
        \node (abajo) at (0,-6) {$\cyclicgr{1}                    \cong                    \{1_{\mathsf{Dih}_{4}}\}$};

        % Now draw the lines:
        \draw [black, thick] (arriba) -- (N);
        \draw [blue, thick] (arriba) -- (H);
        \draw [blue, thick] (arriba) -- (H1);
        \draw [blue, thick] (arriba) -- (H2);
        \draw [blue, thick] (arriba) -- (H3);
        \draw [black, thick] (N) -- (Z);
        \draw [black, thick] (H) -- (abajo);
        \draw [black, thick] (H1) -- (abajo);
        \draw [black, thick] (H2) -- (abajo);
        \draw [black, thick] (H3) -- (abajo);
        \draw [black, thick] (Z) -- (abajo);
        \end{tikzpicture}
    \end{center}








\begin{table}[h!]
    \begin{center}
        \caption{Grupos de orden $8$}
        \label{tab:GruposOrden8}
        \begin{tabular}{c|c|c|c|c|c|c} % <-- Alignments: 1st column left, 2nd middle and 3rd right, with vertical lines in between
            \textbf{} & \tiny\textbf{n\'umero} &\tiny\textbf{numeraci\'on} & \tiny\textbf{nombre} & \textbf{} & {} & {} \\
            \textbf{} & \tiny\textbf{de} &\tiny\textbf{dentro} & \tiny\textbf{del} & \tiny\textbf{} & {} & {} \\
            \tiny\textbf{orden} & \tiny\textbf{grupos} & \tiny\textbf{del} & \tiny\textbf{grupo} & \tiny\textbf{abeliano}  & \tiny\textbf{nilpotente} & \tiny\textbf{resoluble} \\
            \textbf{} & \tiny\textbf{no} & \tiny\textbf{orden} & \textbf{} & \textbf{} & {} & {} \\
            \textbf{} & \tiny\textbf{isomorfos} &\textbf{} & \textbf{} & \textbf{} & {} & {} \\
            \hline
            8 & 5 & 1 & ${\textsf{C}}_{8}$ & Si & Si & Si \\
            \hline
            8 & 5 & 2  & ${{\textsf{C}}_{4}}                  \times                  {{\textsf{C}}_{2}}$ & Si & Si & Si \\
            \hline
            8 & 5 & 3 & ${{\textsf{C}}_{2}}                  \times                  {{\textsf{C}}_{2}}                  \times                  {{\textsf{C}}_{2}}$ & Si & Si & Si \\
            \hline
            8 & 5 & 4  & ${{\textsf{Q}}_{8}}$ & No & Si & Si \\
            \hline
            8 & 5 & 5 & ${{\textsf{Dih}}_{4}}                    \cong                    {{\textsf{C}}_{4}}\rtimes{{{\textsf{C}}_{2}}}$ & No & Si & Si \\
            \hline
        \end{tabular}
    \end{center}
\end{table}



\chapter{Grupos de orden p{\^{}}2 d\'onde p es un primo mayor que 2}


Para abordar los grupos de orden $9$, dado que $9=3^2$, ser\'a igualmente f\'acil tratar el caso general.

Comenzaremos por los grupos abelianos: o bien tiene un elemento $a$ de orden $p^2$ y entonces es c{\'{\i}}clico y tenemos $\cyclicgr{p^2}\triangleq\langle{a\vert a^{p^2}=1}\rangle$, o bien no tiene ese elemento y tendr\'a elementos de orden $p$, digamos $b$ y $c$, y tenemos (a la manera del grupo 4 de Klein) $\cyclicgr{p}                  \times                  \cyclicgr{p}\triangleq\langle{b,c\vert b^{p}=c^{p}=1}\rangle$. El diagrama de su ret{\'{\i}}culo de subgrupos es:






\begin{center}
    \begin{tikzpicture}
    % First, locate each of the nodes and name them
    \node (Cp2) at (3,0) {$\cyclicgr{p^2}                  \cong                  \langle a\vert a^{(p^2)}\rangle$};
    \node (Cp2Cp) at (3,-2) {$                  \cong                  \left\langle a^p\right\rangle $};
    \node (Cp2abajo) at (3,-4) {$                  \cong                  \left\{1_{\cyclicgr{p^2}}\right\}$};
    \node (CpxCp) at (-3,0) {$\cyclicgr{p}                  \times                  \cyclicgr{p}                   \cong                    \left\langle b,c \vert b^p=c^p \right\rangle$};
    \node (CpxCpCp1) at (-5,-2) {$\cyclicgr{p}                   \cong                  \left\langle{b}\right\rangle$};
    \node (CpxCpCp2) at (-3,-2) {$                  \cong                  \left\langle{c}\right\rangle$};
    \node (CpxCpCp3) at (-1,-2) {$                  \cong                  \left\langle{bc}\right\rangle$};
    \node (CpxCpabajo) at (-3,-4) {$\textsf{C}_{1}                  \cong                  \left\{1_{\mathsf{C}_2                   \times                   \mathsf{C}_2}\right\}$};
% Now draw the lines:
    \draw [black, thick] (Cp2) -- (Cp2Cp);
    \draw [black, thick] (Cp2Cp) -- (Cp2abajo);
    \draw [black, thick] (CpxCp) -- (CpxCpCp1);
    \draw [black, thick] (CpxCp) -- (CpxCpCp2);
    \draw [black, thick] (CpxCp) -- (CpxCpCp3);
    \draw [black, thick] (CpxCpCp1) -- (CpxCpabajo);
    \draw [black, thick] (CpxCpCp2) -- (CpxCpabajo);
    \draw [black, thick] (CpxCpCp3) -- (CpxCpabajo);
    \end{tikzpicture}
\end{center}


A partir de aqu{\'{\i}} tenemos que buscar los grupos no abelianos. Pero sabemos que al se un p-grupo, $\mathsf{Z}\left(G\right)                  \neq                  \{1\}$, y al ser no-abeliano solo queda $\mathsf{Z}\left(G\right)                  \cong                  \cyclicgr{p}$,
de d\'onde $\nicefrac{G}{\mathsf{Z}\left(G\right)}                   \cong                   \cyclicgr{p}$ que es c{\'{\i}}clico, y por lo tanto $G$ es abeliano, contra la hip\'otesis. No existen grupos no abelianos de orden $p^2$.




\begin{table}[h!]
    \begin{center}
        \caption{Grupos de orden $p^2$}
        \label{tab:GruposDeOrden$p^2$}
        \begin{tabular}{c|c|c|c|c|c|c} % <-- Alignments: 1st column left, 2nd middle and 3rd right, with vertical lines in between
            \textbf{} & \tiny\textbf{n\'umero} &\tiny\textbf{numeraci\'on} & \tiny\textbf{nombre} & \textbf{} & {} & {} \\
            \textbf{} & \tiny\textbf{de} &\tiny\textbf{dentro} & \tiny\textbf{del} & \tiny\textbf{} & {} & {} \\
            \tiny\textbf{orden} & \tiny\textbf{grupos} & \tiny\textbf{del} & \tiny\textbf{grupo} & \tiny\textbf{abeliano}  & \tiny\textbf{nilpotente} & \tiny\textbf{resoluble} \\
            \textbf{} & \tiny\textbf{no} & \tiny\textbf{orden} & \textbf{} & \textbf{} & {} & {} \\
            \textbf{} & \tiny\textbf{isomorfos} &\textbf{} & \textbf{} & \textbf{} & {} & {} \\
            \hline
            9 & 2 & 1 & ${\textsf{C}}_{9}$ & Si & Si & Si \\
            \hline
            9 & 2 & 2  & ${{\textsf{C}}_{3}}                  \times                  {{\textsf{C}}_{3}}$ & Si & Si & Si \\
            \hline
            25 & 2 & 1 & ${{\textsf{C}}_{25}}$ & Si & Si & Si \\
            \hline
            25 & 2 & 2 & ${{\textsf{C}}_{5}}                  \times                  {{\textsf{C}}_{5}}$ & Si & Si & Si \\
            \hline
            49 & 2 & 1 & ${{\textsf{C}}_{49}}$ & Si & Si & Si \\
            \hline
            49 & 2 & 2 & ${{\textsf{C}}_{7}}                  \times                  {{\textsf{C}}_{7}}$ & Si & Si & Si \\
            \hline
        \end{tabular}
    \end{center}
\end{table}


\chapter{Grupos de orden 12}

Veremos estos grupos de forma espec{\'{\i}}fica, porque el orden $12$ tiene cierta peculiaridad (el grupo alternante) que no tiene an\'alogos en los \'ordenes $4p$ d\'onde hay dos familias, $p=1 \mod 4$ y $p=3\mod 4$ (esta \'ultima ser{\'{\i}}a la correspondiente al orden $12$).

Primero enumeramos los grupos abelianos de ese orden: $\cyclicgr{2}                  \times                  \cyclicgr{6} \triangleq \langle a,b\vert a^2=b^6=1 \rangle$ y $\cyclicgr{12} \triangleq \langle c \vert c^{12}=1 \rangle$.

Sus diagramas de Hasse del ret{\'{\i}}culo de sus subgrupos:



\begin{center}
    \begin{tikzpicture}
    % First, locate each of the nodes and name them
    \node (C12) at (5,15) {$\cyclicgr{12}                  \cong                  \langle a\vert a^{12}=1\rangle$};
    \node (C12C6) at (12,12) {$\cyclicgr{6}                  \cong                  \left\langle a^2\right\rangle $};
    \node (C12C4) at (0,9) {$\cyclicgr{4}                  \cong                  \left\langle a^3\right\rangle $};
    \node (C12C3) at (12,6) {$\cyclicgr{3}                  \cong                  \left\langle a^4\right\rangle $};
    \node (C12C4C2) at (0,3) {$\cyclicgr{2}                  \cong                  \left\langle a^6\right\rangle $};
    \node (C12abajo) at (5,0) {$\cyclicgr{1}                  \cong                  \left\{1_{\cyclicgr{12}}\right\}$};
    % Now draw the lines:
    \draw [black, thick] (C12) -- (C12C6);
    \draw [black, thick] (C12) -- (C12C4);
    \draw [black, thick] (C12C6) -- (C12C3);
    \draw [black, thick] (C12C6) -- (C12C4C2);
    \draw [black, thick] (C12C4) -- (C12C4C2);
    \draw [black, thick] (C12C4C2) -- (C12abajo);
    \draw [black, thick] (C12C3) -- (C12abajo);
    \end{tikzpicture}
\end{center}



\begin{center}
    \begin{tikzpicture}
    % First, locate each of the nodes and name them
    \node (a-b) at (6,12)
    {$\cyclicgr{6}                  \times                  \cyclicgr{2}                  \cong                  \left\langle{a,b\vert a^6=b^2=1}\right\rangle$};
    \node (a2b) at (3,9)
    {$                  \cong                   \left\langle{a^2 b}\right\rangle$};
    \node (a) at (0,9)
    {$\cyclicgr{6}                  \cong                  \left\langle{a}\right\rangle$};
    \node (ab) at (12,9)
    {$                  \cong                  \left\langle{ab}\right\rangle$};
    \node (a3-b) at (9,6)
    {$\cyclicgr{2}                  \times                  \cyclicgr{2}                  \cong                  \left\langle{a^3,b}\right\rangle$};
    \node (a2) at (3,3)
    {$\cyclicgr{3}                  \cong                  \left\langle{a^2}\right\rangle$};
    \node (b) at (9,0)
    {$                  \cong                  \left\langle{b}\right\rangle$};
    \node (a3) at (0,0)
    {$\cyclicgr{2}                  \cong                  \left\langle{a^3}\right\rangle$};
    \node (a3b) at (12,0)
    {$                  \cong                  \left\langle{a^3 b}\right\rangle$};
    \node (1) at (6,-3)
    {$\cyclicgr{1}                  \cong                  \left\langle{1_{\cyclicgr{6}                  \times                  \cyclicgr{2}}}\right\rangle$};
   % Now draw the lines:
    \draw [black, thick] (a-b) -- (ab);
    \draw [black, thick] (a-b) -- (a);
    \draw [black, thick] (a-b) -- (a2b);
    \draw [black, thick] (a-b) -- (a3-b);
    \draw [black, thick] (ab) -- (a3b);
    \draw [black, thick] (ab) -- (a2);
    \draw [black, thick] (a) -- (a2);
    \draw [black, thick] (a) -- (a3);
    \draw [black, thick] (a2b) -- (a2);
    \draw [black, thick] (a2b) -- (b);
    \draw [black, thick] (a3-b) -- (a3b);
    \draw [black, thick] (a3-b) -- (a3);
    \draw [black, thick] (a3-b) -- (b);
    \draw [black, thick] (a2) -- (1);
    \draw [black, thick] (a3b) -- (1);
    \draw [black, thick] (a3) -- (1);
    \draw [black, thick] (b) -- (1);
    \end{tikzpicture}
\end{center}

Ahora comenzaremos con los grupos no abelianos.

Para comenzar veamos las posibilidades que tenemos en cuanto al centro: podr{\'{\i}}a ser $\mathsf{Z}\left(G\right)$ el subgrupo
trivial. Tambi\'en podr{\'{\i}}a ser de \'ordenes $6$, $4$, $3$ y $2$. Como el centro es abeliano solo puede ser
$\mathsf{Z}\left(G\right)                  \cong                   \cyclicgr{6} \;\acute{o}\;                   \cong                   \cyclicgr{4} \;\acute{o}\;
                  \cong                   \cyclicgr{2}                  \times                  \cyclicgr{2} \;\acute{o}\;                   \cong                   \cyclicgr{3} \;\acute{o}\;                   \cong                   \cyclicgr{2}
\;\acute{o}\;                   \cong                   \cyclicgr{1}$. Podemos descartar como centro $\cyclicgr{6}$ pues el cociente ser{\'{\i}}a
c{\'{\i}}clico, y por lo tanto, $G$ ser{\'{\i}}a abeliano. Podemos utilizar el mismo argumento con los grupos de orden $4$.

 Quedar{\'{\i}}a las siguientes posibilidades:
\begin{align}
\mathsf{Z}\left(G\right) &                  \cong                   \cyclicgr{3} &\text{ si }\quad\quad\quad\quad&
\nicefrac{G}{\mathsf{Z}\left(G\right)}                  \cong                   \cyclicgr{2}                  \times                   \cyclicgr{2}\\
&                  \cong                   \cyclicgr{2}  &\text{ si }\quad\quad\quad\quad& \nicefrac{G}{\mathsf{Z}\left(G\right)}
                  \cong                   \cyclicgr{3}                  \times                   \cyclicgr{2}\\
&                  \cong                   \cyclicgr{1}&&
\end{align}

Ahora podemos pensar en grupos seg\'un tengan un grupos de determinados \'ordenes normales o no.

Podemos pensar que tenga subgrupos normales de orden $6$, llam\'emos a uno de ellos $N_6$, o bien que alg\'un grupo de
orden $6$ sea normal. Este \'ultimo caso no se puede dar, porque en el caso que hubiese uno solo, digamos $H_6$, entonces $[G:H_6]=2$ y autom\'aticamente $H_6$ ser{\'{\i}}a normal. Luego si hay subgrupos de orden $6$ han de ser normales.

Pero siempre tendremos elementos de orden $3$ y $2$ por el teorema de Cauchy. Subgrupos de orden $4$ existir\'an por el primer teorema de Sylow. Supondremos que este subgrupo maximal es normal, $N_4$, y el de orden $3$, digamos $H_3$ no lo puede ser, por que sino tenemos que el grupo es abeliano. Estos casos ser\'an del tipo $N_4\rtimes H_3$, o bien a la inversa, y tendr{\'{\i}}amos $N_3\rtimes H_4$ teniendo en cuenta que los grupos de orden $4$ pueden ser bien $\cyclicgr{4}$, bien $\cyclicgr{2}                   \times                   \cyclicgr{2}$.

Enumerando las opciones:
\begin{align}
\cyclicgr{4}\rtimes \cyclicgr{3}\\
\cyclicgr{3}\rtimes \cyclicgr{4}\\
\left(\cyclicgr{2}                  \times                  \cyclicgr{2}\right)\rtimes \cyclicgr{3}\\
\cyclicgr{3}\rtimes \left(\cyclicgr{2}                  \times                  \cyclicgr{2}\right)
\end{align}

En el caso que consideremos que existe grupo de orden $6$ (por lo tanto normal), tendremos las siguientes posibilidades:

\begin{align}
\mathsf{Dih}_{6}                  \cong                  \cyclicgr{6}\rtimes_{\langle{g                  \mapsto g\,,\,f                  \mapsto f^{-1}}\rangle} \cyclicgr{2}\\
\mathsf{Dih}_{3}                   \times                   \cyclicgr{2}\\
\mathsf{Dih}_{3} \rtimes \cyclicgr{2}
\end{align}

Para hacer los c\'alculos necesarios debemos tener los grupos de automorfismos.

Los automorfismos del grupo dih\'edrico de orden $6$ son:
\begin{align}
&\texttt{Aut}(\mathsf{Dih}_{3})=\nonumber\\
&=\quad
    \left\lbrace \right.
        \varphi_0\triangleq\binom{g                  \mapsto g}{f                  \mapsto f},
        \varphi_1\triangleq\binom{g                  \mapsto g}{f                  \mapsto f^2},
        \varphi_2\triangleq\binom{g                  \mapsto gf}{f                  \mapsto f},\nonumber\\
&\phantom{=}\quad\varphi_3\triangleq\binom{g                  \mapsto gf}{f                  \mapsto f^2}, 			
	\varphi_4\triangleq\binom{g                  \mapsto gf^2}{f
	                  \mapsto f}, \varphi_5\triangleq\binom{g                  \mapsto gf^2}{f                  \mapsto f^2}
    \left. \right\rbrace\nonumber                   \cong                   \mathsf{Dih}_{3}
\end{align}

Los automorfismos del grupo $4$ de Klein son:
\begin{align}
&\texttt{Aut}(\cyclicgr{2}                  \times                  \cyclicgr{2})=\nonumber\\
&=\quad
	\left\lbrace \right.
	\phi_0\triangleq\binom{a                  \mapsto a}{b                  \mapsto b},
	\phi_1\triangleq\binom{a                  \mapsto a}{b                  \mapsto ab},
	\phi_2\triangleq\binom{a                  \mapsto ab}{b                  \mapsto b},\nonumber\\
&\phantom{=}\quad
	\phi_3\triangleq\binom{a                  \mapsto ab}{b                  \mapsto a}, 			
	\phi_4\triangleq\binom{a                  \mapsto b}{b                  \mapsto ab},
	\phi_5\triangleq\binom{a                  \mapsto b}{b                  \mapsto a}
	\left. \right\rbrace\nonumber                   \cong                   \mathsf{Dih}_{3}
\end{align}

Los automorfismos del grupo $\cyclicgr{4}$ son:

\begin{align}
&\texttt{Aut}(\cyclicgr{4})=
	\left\lbrace
		\psi_0\triangleq\left({a                  \mapsto a}\right),
		\psi_1\triangleq\left({a                  \mapsto a^3}\right)
	\right\rbrace
	                  \cong                   \cyclicgr{2}
\end{align}

Los automorfismos del grupo $\cyclicgr{3}$ son:

\begin{align}
&\texttt{Aut}(\cyclicgr{3})=
\left\lbrace
	\vartheta_0\triangleq\left({a                  \mapsto a}\right),
	\vartheta_1\triangleq\left({a                  \mapsto a^2}\right)
\right\rbrace
                  \cong                   \cyclicgr{2}
\end{align}

Una vez enumerados los distintos grupos de automorfismos plantearemos las acciones
que se utilizar\'an en los productos semidirectos, mediante homomorfismos de un grupo
(el no necesariamente normal) al grupo de automorfismos del otro grupo
(el que ser\'a normal).

Para $\cyclicgr{4}\rtimes\cyclicgr{3}$ tenemos solo el homomorfismo identidad, pues si
$\langle a\rangle=\cyclicgr{3}$, y teniendo en cuenta $\psi_1$ es de orden
$2$, $a                   \mapsto \psi_1$ entonces, por ser homomorfismo,
$a^2                   \mapsto \psi_1                  \circ \psi_1=\psi_0=Id$. Ahora $a^3=1                   \mapsto \psi_1$
y llegamos a que no puede ser homomorfismo, ya que esto exigir{\'{\i}}a $a^3=1                   \mapsto \psi_0 = Id$.
Por lo tanto, por este camino podemos desechar
$\cyclicgr{4}\rtimes_{Id}\cyclicgr{3} =
\cyclicgr{4}                  \times                  \cyclicgr{3}                  \cong                   \cyclicgr{12}$
que es abeliano y ya hemos dado cuenta de \'el.

Para $\cyclicgr{3}\rtimes\cyclicgr{4}$ tenemos una \'unica acci\'on:

\begin{flalign*}
& \Psi : \cyclicgr{4}                   \longrightarrow                   \texttt{Aut}\left( \cyclicgr{3} \right) \\
& \qquad        1				\mapsto		{\vartheta }_0 = Id \\
& \qquad		a				\mapsto		{\vartheta }_1 \\
& \qquad		a^2				\mapsto		{\vartheta }_0 \\
& \qquad		a^3				\mapsto		{\vartheta }_1
\end{flalign*}

As{\'{\i}}, tomando el conjunto base $\cyclicgr{4}                  \times                  \cyclicgr{3}$ establecemos la
operaci\'on de grupo para $\cyclicgr{3}\rtimes_{\Psi } \cyclicgr{4}$:

\begin{flalign*}
& a^4=1 \\
& b^3=1 \\
& (a^i,b^j)(a^k,b^m) \triangleq (a^{i+k} ,{\Psi }_{a^k}(b^j) {b^{m}}) \\
&
(a^i,b^j)(a^k,b^m) =
	\begin{cases}
		\left( a^{i  } , \vartheta_0(b^{j}){b^m} \right)   &                  \Leftarrow                   k=0\\
		\left( a^{i+1} , \vartheta_1(b^{j}){b^m} \right)   &                  \Leftarrow                   k=1\\
		\left( a^{i+2} , \vartheta_0(b^{j}){b^m} \right)   &                  \Leftarrow                   k=2\\
		\left( a^{i+3} , \vartheta_1(b^{j}){b^m} \right)   &                  \Leftarrow                   k=3
	\end{cases}
\end{flalign*}

Ahora sustituimos por el valor de $\vartheta_t$:

\begin{flalign*}
& a^4=1 \\
& b^3=1 \\
& (a^i,b^j)(a^k,b^m) \triangleq (a^{i+k} ,{\Psi }_{a^k}(b^j) {b^{j}}) \\
&
(a^i,b^j)(a^k,b^m) =
	\begin{cases}
		\left( a^{i  } , {b^{j}} {b^m} \right)   &                  \Leftarrow                   k=0\\
		\left( a^{i+1} , {b^{2j}}{b^m} \right)   &                  \Leftarrow                   k=1\\
		\left( a^{i+2} , {b^{j}} {b^m} \right)   &                  \Leftarrow                   k=2\\
		\left( a^{i+3} , {b^{2j}}{b^m} \right)   &                  \Leftarrow                   k=3
	\end{cases}
\end{flalign*}


Y obtenemos ya la expresi\'on final:


\begin{flalign*}
& a^4=1 \\
& b^3=1 \\
& (a^i,b^j)(a^k,b^m) \triangleq (a^{i+k} ,{\Psi }_{a^k}(b^j) {b^{j}}) \\
&
(a^i,b^j)(a^k,b^m) =
	\begin{cases}
		\left( a^{i  } , {b^{m+j}}  \right)   &                  \Leftarrow                   k=0\\
		\left( a^{i+1} , {b^{m+2j}} \right)   &                  \Leftarrow                   k=1\\
		\left( a^{i+2} , {b^{m+j}}  \right)   &                  \Leftarrow                   k=2\\
		\left( a^{i+3} , {b^{m+2j}} \right)   &                  \Leftarrow                   k=3
	\end{cases}
\end{flalign*}


En la siguiente tabla de Cayley, en vez de poner $(a^k,b^m)$
pondr\'e solo $(k,m)$. La multiplicaci\'on cogiendo solo los
exponentes quedar{\'{\i}}a as{\'{\i}}:


\begin{flalign*}
& a^4=1 \\
& b^3=1 \\
& (a^i,b^j)(a^k,b^m) \triangleq (a^{i+k} ,{\Psi }_{a^k}(b^j) {b^{m}})                   \cong                  \\
&                   \cong                    (i,j)(k,m) \triangleq ({i+k}\mod 4 ,\exp{({\Psi }_{a^k}(b^j))}+m \mod 3) \\
&
(i,j)(k,m) =
	\begin{cases}
		\left( {i  }\mod 4 , {m+j}  \mod 3 \right)   &                  \Leftarrow                   k=0\\
		\left( {i+1}\mod 4 , {m+2j} \mod 3 \right)   &                  \Leftarrow                   k=1\\
		\left( {i+2}\mod 4 , {m+j}  \mod 3 \right)   &                  \Leftarrow                   k=2\\
		\left( {i+3}\mod 4 , {m+2j} \mod 3 \right)   &                  \Leftarrow                   k=3
	\end{cases}
\end{flalign*}

Ahora llamando $1 \triangleq (0,0)$, $a^k \triangleq (k,0)$ y $b^m \triangleq (0,m)$, y sabiendo que $a^k b^m = (k,0)(0,m)$:
\\

\begin{flushleft}
\begin{tabular}{|c|c|c|c|c|c|c|c|c|c|c|c|}
\hline
 $1$ & $b$ & $b^2$ & $a$ & $ab$ & $ab^2$ & $a^2$ & $a^2b$ & $a^2b^2$ & $a^3$ & $a^3b$ & $a^3b^2$ \\
\hline
 $b$ & $b^2$ & $1$ & $ab^2$ & $a$ & $ab$ & $a^2b$ & $a^2b^2$ & $a^2$ & $a^3b^2$ & $a^3$ & $a^3b$ \\
\hline
 $b^2$ & $1$ & $b$ & $ab$ & $ab^2$ & $a$ & $a^2b^2$ & $a^2$ & $a^2b$ & $a^3b$ & $a^3b^2$ & $a^3$ \\
\hline
 $a$ & $ab$ & $ab^2$ & $a^2$ & $a^2b$ & $a^2b^2$ & $a^3$ & $a^3b$ & $a^3b^2$ & $1$ & $b$ & $b^2$ \\
\hline
 $ab$ & $ab^2$ & $a$ & $a^2b^2$ & $a^2$ & $a^2b$ & $a^3b$ & $a^3b^2$ & $a^3$ & $b^2$ & $1$ & $b$ \\
\hline
 $ab^2$ & $a$ & $ab$ & $a^2b$ & $a^2b^2$ & $a^2$ & $a^3b^2$ & $a^3$ & $a^3b$ & $b$ & $b^2$ & $1$ \\
\hline
 $a^2$ & $a^2b$ & $a^2b^2$ & $a^3$ & $a^3b$ & $a^3b^2$ & $1$ & $b$ & $b^2$ & $a$ & $ab$ & $ab^2$ \\
\hline
 $a^2b$ & $a^2b^2$ & $a^2$ & $a^3b^2$ & $a^3$ & $a^3b$ & $b$ & $b^2$ & $1$ & $ab^2$ & $a$ & $ab$ \\
\hline
 $a^2b^2$ & $a^2$ & $a^2b$ & $a^3b$ & $a^3b^2$ & $a^3$ & $b^2$ & $1$ & $b$ & $ab$ & $ab^2$ & $a$ \\
\hline
 $a^3$ & $a^3b$ & $a^3b^2$ & $1$ & $b$ & $b^2$ & $a$ & $ab$ & $ab^2$ & $a^2$ & $a^2b$ & $a^2b^2$ \\
\hline
 $a^3b$ & $a^3b^2$ & $a^3$ & $b^2$ & $1$ & $b$ & $ab$ & $ab^2$ & $a$ & $a^2b^2$ & $a^2$ & $a^2b$ \\
\hline
 $a^3b^2$ & $a^3$ & $a^3b$ & $b$ & $b^2$ & $1$ & $ab^2$ & $a$ & $ab$ & $a^2b$ & $a^2b^2$ & $a^2$ \\
\hline
\end{tabular}
\end{flushleft}

Los \'ordenes de cada uno de los elementos es:

\begin{flalign*}
&\mathsf{o}\left(1\right)=1\\
&\mathsf{o}\left(a^2\right)=2\\
&\mathsf{o}\left(b\right)=3\\
&\mathsf{o}\left(b^2\right)=3\\
&\mathsf{o}\left(a\right)=4\\
&\mathsf{o}\left(ab\right)=4\\
&\mathsf{o}\left(ab^2\right)=4\\
&\mathsf{o}\left(a^3\right)=4\\
&\mathsf{o}\left(a^3b\right)=4\\
&\mathsf{o}\left(a^3b^2\right)=4\\
&\mathsf{o}\left(a^2b\right)=6\\
&\mathsf{o}\left(a^2b^2\right)=6
\end{flalign*}

Por lo pronto podemos ver que se invierte el teorema de Lagrange en este grupo.

Como $a^2$ tiene orden $2$, y es el \'unico elemento de ese orden, es normal luego
$\langle a^2\rangle\lhd \cyclicgr{3}\rtimes_{\Psi } \cyclicgr{4}$,
entonces ha de estar en el centro:

\[a^2                   \in \mathsf{Z}(\cyclicgr{3}\rtimes_{\Psi } \cyclicgr{4})\]

Adem\'as los subgrupos c{\'{\i}}clicos son los \'unicos subgrupos propios (siendo el grupo entero
el \'unico que no es c{\'{\i}}clico) :

\begin{flalign*}
{\mathcal{L}}\left({{{\mathsf{C}}_{3}}\rtimes{{\mathsf{C}}_{4}}}\right) =
\begin{cases}
&{\left\langle{1}\right\rangle}={\left\lbrace{1}\right\rbrace}\\
&\qquad {\mathsf{o}}{\left({\left\langle{1}\right\rangle}\right)}=1\\
&{\left\langle{a^2}\right\rangle}={\left\lbrace{1,a^2}\right\rbrace}\\
&\qquad {\mathsf{o}}{\left({\left\langle{a^2}\right\rangle}\right)}=2\\
&{\left\langle{b}\right\rangle}={\left\lbrace{1,b,b^2}\right\rbrace}\\
&\qquad {\mathsf{o}}{\left({\left\langle{b}\right\rangle}\right)}=3\\
&{\left\langle{a}\right\rangle}={\left\lbrace{1,a,a^2,a^3}\right\rbrace}\\
&\qquad {\mathsf{o}}{\left({\left\langle{a}\right\rangle}\right)}=4\\
&{\left\langle{a b}\right\rangle}={\left\lbrace{1,ab,a^2,a^3b}\right\rbrace}\\
&\qquad {\mathsf{o}}{\left({\left\langle{a b}\right\rangle}\right)}=4\\
&{\left\langle{a b^2}\right\rangle}={\left\lbrace{1,ab^2,a^2,a^3b^2}\right\rbrace}\\
&\qquad {\mathsf{o}}{\left({\left\langle{a b^2}\right\rangle}\right)}=4\\
&{\left\langle{a^2 b}\right\rangle}={\left\lbrace{1,b,b^2,a^2,a^2b,a^2b^2}\right\rbrace}\\
&\qquad {\mathsf{o}}{\left({\left\langle{a^2 b}\right\rangle}\right)}=6\\
&{\left\langle{a , b}\right\rangle}={{{\mathsf{C}}_{3}}\rtimes{{\mathsf{C}}_{4}}}\\
&\qquad {\mathsf{o}}{\left({\left\langle{a , b}\right\rangle}\right)}=12
\end{cases}
\end{flalign*}


El subgrupo centro (caracter{\'{\i}}stico):

\begin{flalign*}
&\mathsf{Z}\left({{{\mathsf{C}}_{3}}\rtimes{{\mathsf{C}}_{4}}}\right)
= {\left\langle{a^2}\right\rangle}
\end{flalign*}

Los subgrupos normalizadores de cada subgrupo en el grupo de estudio,
nos dar\'a tambi\'en la informaci\'on de qu\'e subgrupos son normales:

\begin{flalign*}
&{\mathsf{N}}_{{{\mathsf{C}}_{3}}\rtimes{{\mathsf{C}}_{4}}}
\left({\left\lbrace{1}\right\rbrace}\right) =
{{{\mathsf{C}}_{3}}\rtimes{{\mathsf{C}}_{4}}}\\
&{\mathsf{N}}_{{{\mathsf{C}}_{3}}\rtimes{{\mathsf{C}}_{4}}}
\left({\left\langle{a^2}\right\rangle}\right) =
{{{\mathsf{C}}_{3}}\rtimes{{\mathsf{C}}_{4}}}\\
&{\mathsf{N}}_{{{\mathsf{C}}_{3}}\rtimes{{\mathsf{C}}_{4}}}
\left({\left\lbrace{1,b,b^2}\right\rbrace}\right) =
{{{\mathsf{C}}_{3}}\rtimes{{\mathsf{C}}_{4}}}\\
&{\mathsf{N}}_{{{\mathsf{C}}_{3}}\rtimes{{\mathsf{C}}_{4}}}
\left({\left\langle{a}\right\rangle}\right) =
{\left\langle{a}\right\rangle}\\
&{\mathsf{N}}_{{{\mathsf{C}}_{3}}\rtimes{{\mathsf{C}}_{4}}}
\left({\left\langle{ab}\right\rangle}\right) =
{\left\langle{ab}\right\rangle}\\
&{\mathsf{N}}_{{{\mathsf{C}}_{3}}\rtimes{{\mathsf{C}}_{4}}}
\left({\left\langle{ab^2}\right\rangle}\right) =
{\left\langle{ab^2}\right\rangle}\\
&{\mathsf{N}}_{{{\mathsf{C}}_{3}}\rtimes{{\mathsf{C}}_{4}}}
\left({\left\langle{a^2b}\right\rangle}\right) =
{{{\mathsf{C}}_{3}}\rtimes{{\mathsf{C}}_{4}}}\\
\end{flalign*}

Los centralizadores de cada uno de los subgrupos en el subgrupo de estudio o ambiente son:

\begin{flalign*}
&{\cyclicgr{{{{\mathsf{C}}_{3}}\rtimes{{\mathsf{C}}_{4}}}}}
\left({\left\lbrace{1}\right\rbrace}\right) =
{{{\mathsf{C}}_{3}}\rtimes{{\mathsf{C}}_{4}}}\\
&{\cyclicgr{{{{\mathsf{C}}_{3}}\rtimes{{\mathsf{C}}_{4}}}}}
\left({\left\langle{a^2}\right\rangle}\right) =
{{{\mathsf{C}}_{3}}\rtimes{{\mathsf{C}}_{4}}}\\
&{\cyclicgr{{{{\mathsf{C}}_{3}}\rtimes{{\mathsf{C}}_{4}}}}}
\left({\left\langle{b}\right\rangle}\right) =
{\left\langle{a^2b}\right\rangle}\\
&{\cyclicgr{{{{\mathsf{C}}_{3}}\rtimes{{\mathsf{C}}_{4}}}}}
\left({\left\langle{a}\right\rangle}\right) =
{\left\langle{a}\right\rangle}\\
&{\cyclicgr{{{{\mathsf{C}}_{3}}\rtimes{{\mathsf{C}}_{4}}}}}
\left({\left\langle{ab}\right\rangle}\right) =
{\left\langle{ab}\right\rangle}\\
&{\cyclicgr{{{{\mathsf{C}}_{3}}\rtimes{{\mathsf{C}}_{4}}}}}
\left({\left\langle{ab^2}\right\rangle}\right) =
{\left\langle{ab^2}\right\rangle}\\
&{\cyclicgr{{{{\mathsf{C}}_{3}}\rtimes{{\mathsf{C}}_{4}}}}}
\left({\left\langle{a^2b}\right\rangle}\right) =
{\left\langle{a^2b}\right\rangle}
\end{flalign*}

Los subgrupos derivados son:

\begin{flalign*}
&{\left\langle{a^2}\right\rangle}^{(1)} = {\left\lbrace{1}\right\rbrace}\\
&{\left\langle{b}\right\rangle}^{(1)} = {\left\lbrace{1}\right\rbrace}\\
&{\left\langle{a}\right\rangle}^{(1)} = {\left\lbrace{1}\right\rbrace}\\
&{\left\langle{ab}\right\rangle}^{(1)} = {\left\lbrace{1}\right\rbrace}\\
&{\left\langle{ab^2}\right\rangle}^{(1)} = {\left\lbrace{1}\right\rbrace}\\
&{\left\langle{a^2b}\right\rangle}^{(1)} = {\left\lbrace{1}\right\rbrace}\\
&{{{\mathsf{C}}_{3}}\rtimes{{\mathsf{C}}_{4}}}^{(1)} ={\left\langle{a^2}\right\rangle}\\
&{{{\mathsf{C}}_{3}}\rtimes{{\mathsf{C}}_{4}}}^{(2)} ={\left\lbrace{1}\right\rbrace}
\end{flalign*}

De aqu{\'{\i}} que el grupo que nos ocupa es resoluble.

Los subgrupos caracter{\'{\i}}sticos que nos dicen si el grupo nilpotente:

\begin{flalign*}
&\mathsf{L}_{1}\left({\left\langle{a^2}\right\rangle}\right)
=
{\left\lbrace{1}\right\rbrace}\\
&\mathsf{L}_{1}\left({\left\langle{b}\right\rangle}\right)
=
{\left\lbrace{1}\right\rbrace}\\
&\mathsf{L}_{1}\left({\left\langle{a}\right\rangle}\right)
=
{\left\lbrace{1}\right\rbrace}\\
&\mathsf{L}_{1}\left({\left\langle{ab}\right\rangle}\right)
=
{\left\lbrace{1}\right\rbrace}\\
&\mathsf{L}_{1}\left({\left\langle{ab^2}\right\rangle}\right)
=
{\left\lbrace{1}\right\rbrace}\\
&\mathsf{L}_{1}\left({\left\langle{a^2b}\right\rangle}\right)
=
{\left\lbrace{1}\right\rbrace}\\
&{\mathsf{L}}_{1,2 ...}
\left(
{{{\mathsf{C}}_{3}}\rtimes{{\mathsf{C}}_{4}}}
\right)
=
{\left\langle{b}\right\rangle}                   \Rightarrow
{{{\mathsf{C}}_{3}}\rtimes{{\mathsf{C}}_{4}}}
\mbox{ NO ES NILPOTENTE}
\end{flalign*}

El subgrupo de Frattini de ${{\mathsf{C}}_{3}}\rtimes{{\mathsf{C}}_{4}}$ es
$
\left\langle{ a^2 b}\right\rangle
\cap
\left\langle{a}\right\rangle
\cap
\left\langle{a b^2}\right\rangle
\cap
\left\langle{a b}\right\rangle
=
{
\left\langle
{a^2}
\right\rangle
}$.
As{\'{\i}}:
\[
{\Phi }_{{{{\mathsf{C}}_{3}}\rtimes{{\mathsf{C}}_{4}}}} =
{\mathsf{Z}}{\left({{{\mathsf{C}}_{3}}\rtimes{{\mathsf{C}}_{4}}}\right)}=
{\left\langle{a^2}\right\rangle}
\]

Para el grupo de Fitting solo tenemos que coger los subgrupos nilpotentes normales, esto es,
los subgrupos normales y dar con el subgrupo generado por estos: ${\mathsf{F}}
{\left({{{\mathsf{C}}_{3}}\rtimes{{\mathsf{C}}_{4}}}\right)}=
{\left\langle{a^2 b}\right\rangle}$ ya que
${\left\langle{a^2}\right\rangle}                  \subset {\left\langle{a^2 b}\right\rangle}$
y ${\left\langle{b}\right\rangle}                  \subset {\left\langle{a^2 b}\right\rangle}$,
de forma que ${\left\langle{a^2 b}\right\rangle}={\left\langle{a^2 , b}\right\rangle}$.


Faltar{\'{\i}}a ver a qu\'e equivale
$\nicefrac{{{{\mathsf{C}}_{3}}\rtimes{{\mathsf{C}}_{4}}}}{\left\langle{a^2}\right\rangle}$
, de la misma forma que ser{\'{\i}}a interesante ver
$\nicefrac{{{{\mathsf{C}}_{3}}\rtimes{{\mathsf{C}}_{4}}}}{\left\langle{b}\right\rangle}$
y $\nicefrac{{{{\mathsf{C}}_{3}}\rtimes{{\mathsf{C}}_{4}}}}{\left\langle{a^2 b}\right\rangle}$.
Como el grupo del denominador es el centro cualquier conjugado constituye un clase de
equivalencia que incluye $x,y,z$ no necesariamente distintas, tal que, $xy=1                   \wedge                   xz=a^2$:

\begin{flalign*}
&\overline{1}
&\triangleq&\quad
{1}                  \cdot {\left\langle{a^2}\right\rangle}
&                  \equiv                  &\quad
{\left\lbrace{1,a^2}\right\rbrace}\\
&\overline{a}
&\triangleq&\quad
{a}                  \cdot {\left\langle{a^2}\right\rangle}
&                  \equiv                  &\quad
{\left\lbrace{a,a^3}\right\rbrace}\\
&\overline{ab}
&\triangleq&\quad
{ab}                  \cdot {\left\langle{a^2}\right\rangle}
&                  \equiv                  &\quad
{\left\lbrace{ab,a^3b}\right\rbrace}\\
&\overline{ab^2}
&\triangleq&\quad
{ab^2}                  \cdot {\left\langle{a^2}\right\rangle}
&                  \equiv                  &\quad
{\left\lbrace{ab^2,a^3b^2}\right\rbrace}
\\
&\overline{b}
&\triangleq&\quad
{b}                  \cdot {\left\langle{a^2}\right\rangle}
&                  \equiv                  &\quad
{\left\lbrace{b,a^2b}\right\rbrace}\\
&\overline{b^2}
&\triangleq&\quad
{b^2}                  \cdot {\left\langle{a^2}\right\rangle}
&                  \equiv                  &\quad
{\left\lbrace{b^2,a^2b^2}\right\rbrace}
\end{flalign*}

En cuanto a la operaci\'on heredada, tenemos:

\begin{flalign*}
&\overline{1}                  \cdot \overline{x}=\overline{x}                  \cdot \overline{1}=\overline{x}\quad \mbox{conforma el grupo trivial}\\
&{\overline{b}}^{-1}={\overline{b^2}}\quad \mbox{conforman un grupo c{\'{\i}}clico de orden 3}\\
&{\overline{a}}^{-1}={\overline{a}}\quad \mbox{conforman un grupo c{\'{\i}}clico de orden 2}\\
&{\overline{ab}}^{-1}={\overline{ab}}\quad \mbox{conforman un grupo c{\'{\i}}clico de orden 2}\\
&{\overline{ab^2}}^{-1}={\overline{ab^2}}\quad \mbox{conforman un grupo c{\'{\i}}clico de orden 2}\\
&{\overline{b}}                  \cdot {\overline{a}}=\overline{ab^2}\\
&{\overline{a}}                  \cdot {\overline{b}}=\overline{ab}\\
&{\overline{b}}                  \cdot {\overline{a}}=\overline{ab^2}                  \neq                  \overline{ab}={\overline{a}}                  \cdot {\overline{b}}\\
&\nicefrac{{{\mathsf{C}}_{3}}\rtimes{{\mathsf{C}}_{4}}}{\left\langle a^2 \right\rangle}                   \cong                  {\dihgr{3}}                  \cong                   {\mathtt{Int}\left({{{\mathsf{C}}_{3}}\rtimes{{\mathsf{C}}_{4}}}\right)}
\end{flalign*}

En cuanto a todos los automorfismos sabemos que tienen que cumplir que $1                   \mapsto 1 , a^2                   \mapsto a^2$ y tenemos las posibilidades $a                   \mapsto x , \mathsf{o}(x)=4, ab                   \mapsto y, \mathsf{o}(y)=4$ y $b                  \mapsto b,b                  \mapsto b^2$ de forma que sean compatibles, que da un m\'aximo de $12$ posibilidades. Como medida de control, todo automorfismo en el grupo total, cuando se restringe a $\langle a^2 b\rangle$ ha de ser tambi\'en un automorfismo de $\mathtt{Aut}\left(\langle a^2 b\rangle\right)                  \cong                  \mathtt{Aut}\left({{\mathsf{C}}_{6}}\right)                  \cong                   {{\mathsf{C}}_{2}}$.

Por \'ultimo tambi\'en tenemos $\nicefrac{{{{\mathsf{C}}_{3}}\rtimes{{\mathsf{C}}_{4}}}}{\left\langle b \right\rangle}                   \cong                   {{\mathsf{C}}_{4}}$, y muy claramente , $\nicefrac{{{{\mathsf{C}}_{3}}\rtimes{{\mathsf{C}}_{4}}}}{\left\langle a^2 b \right\rangle}                   \cong                   {{\mathsf{C}}_{2}}$.

El diagrama de Hasse quedar{\'{\i}}a:


\begin{center}
    \begin{tikzpicture}
    % First, locate each of the nodes and name them
    \node (a-b) at (0,12)
    {$\cyclicgr{3}\rtimes\cyclicgr{4}                  \cong                  \left\langle{a,b\vert a^4=b^3=1,xa^2=a^2x,ab=b^2a}\right\rangle$};
    \node (a2b) at (3,10)
    {$\cyclicgr{6}                  \cong                   \left\langle{a^2 b}\right\rangle$};
    \node (a) at (1,8)
    {$                  \cong                  \left\langle{a}\right\rangle$};
    \node (ab) at (-1,8)
    {$                  \cong                  \left\langle{ab}\right\rangle$};
    \node (ab2) at (-3,8)
    {$\cyclicgr{4}                  \cong                  \left\langle{ab^2}\right\rangle$};
    \node (b) at (3,6)
    {$\cyclicgr{3}                  \cong                  \left\langle{b}\right\rangle$};
    \node (a2) at (0,4)
    {$\cyclicgr{2}                  \cong                  \left\langle{a^2}\right\rangle$};
    \node (1) at (0,0)
    {${\cyclicgr{1}}                  \cong                  {\left\langle{1_{\cyclicgr{3}\rtimes\cyclicgr{4}}}\right\rangle}$};
   % Now draw the lines:
    \draw [black, thick] (a-b) -- (a2b);
    \draw [blue, thick] (a-b) -- (a);
    \draw [blue, thick] (a-b) -- (ab);
    \draw [blue, thick] (a-b) -- (ab2);
    \draw [black, thick] (a2b) -- (b);
    \draw [black, thick] (a2b) -- (a2);
    \draw [black, thick] (a) -- (a2);
    \draw [black, thick] (ab) -- (a2);
    \draw [black, thick] (ab2) -- (a2);
    \draw [black, thick] (b) -- (1);
    \draw [black, thick] (a2) -- (1);
    \end{tikzpicture}
\end{center}

El siguiente grupo no abeliano lo buscaremos tambi\'en por los productos semidirectos de grupos de orden $4$ y de orden $3$. Se nos plantea el mismo problema que antes, ver si son posibles ${\left({{{\mathsf{C}}_{2}}                  \times                  {{\mathsf{C}}_{2}}}\right)}\rtimes{{\mathsf{C}}_{3}}$ y/o ${{\mathsf{C}}_{3}}\rtimes{\left({{{\mathsf{C}}_{2}}                  \times                  {{\mathsf{C}}_{2}}}\right)}$. Comenzamos por el \'ultimo, esto es, si es posible construir un homomorfismo $\Psi : {{{\mathsf{C}}_{2}}                  \times                  {{\mathsf{C}}_{2}}}                   \longrightarrow                   {\mathtt{Aut}}{\left({{\mathsf{C}}_{3}}\right)}                  \cong                  {{\mathsf{C}}_{2}}$.

Veamos:

\begin{flalign*}
&{{\mathsf{C}}_{3}} &\triangleq \qquad&{\left\langle{c \shortmid c^3=1}\right\rangle}&&\\
&{{\mathsf{C}}_{2}} &\triangleq \qquad&{\left\langle{a,b\shortmid a^2=b^2=1, ab=ba}\right\rangle}&&\\
&{\mathtt{Aut}}{\left({{\mathsf{C}}_{3}}\right)} &\triangleq
\qquad&\left\lbrace{\left\langle c         \mapsto c \right\rangle,
\left\langle c         \mapsto c^2 \right\rangle}\right\rbrace&&
\end{flalign*}

\begin{flalign*}
&\Psi_0 : &{{{\mathsf{C}}_{2}}                  \times                  {{\mathsf{C}}_{2}}}\qquad &                  \longrightarrow                   &{\mathtt{Aut}}{\left({{\mathsf{C}}_{3}}\right)}                  \cong                  {{\mathsf{C}}_{2}}&\\
&&a&                  \mapsto \qquad & \left\langle c                  \mapsto c \right\rangle&\\
&&b&                  \mapsto \qquad & \left\langle c                  \mapsto c \right\rangle&\\
&&1&                  \mapsto \qquad & \left\langle c                  \mapsto c \right\rangle &\mbox{ en todo homomorfismo}\\
&&ab&                  \mapsto \qquad & \left\langle c                  \mapsto c \right\rangle &\mbox{ conclusi\'on}\\
&\Psi_1 : &{{{\mathsf{C}}_{2}}                  \times                  {{\mathsf{C}}_{2}}}\qquad &                  \longrightarrow                   &{\mathtt{Aut}}{\left({{\mathsf{C}}_{3}}\right)}                  \cong                  {{\mathsf{C}}_{2}}&\\
&&a&                  \mapsto \qquad & \left\langle c                  \mapsto c^2 \right\rangle&\\
&&b&                  \mapsto \qquad & \left\langle c                  \mapsto c^2 \right\rangle&\\
&&1&                  \mapsto \qquad & \left\langle c                  \mapsto c \right\rangle &\mbox{ en todo homomorfismo}\\
&&ab&                  \mapsto \qquad & \left\langle c                  \mapsto c \right\rangle &\mbox{ conclusi\'on}
\end{flalign*}

Tomando $\Phi_1$ como nuestro homomorfismo, puesto que $\Phi_0$ no es mas que la identidad, que nos dar{\'{\i}}a un producto directo y por lo tanto abeliano. La acci\'on quedar{\'{\i}}a as{\'{\i}}:

\begin{flushright}
\begin{flalign*}
&\Psi_1 : &{{{\mathsf{C}}_{2}}                  \times                  {{\mathsf{C}}_{2}}}                  \times                  {{\mathsf{C}}_{3}} &                  \longrightarrow                   &{{\mathsf{C}}_{3}}\\
&&a^ib^jc^k&                  \mapsto \qquad & c^l
\end{flalign*}


\begin{flalign*}
&l=
\begin{cases}
-k                   \Rightarrow                   (i                  \neq                   j)\\
k \;                  \Rightarrow                   (i=j)
\end{cases}
\\
& i,j                   \in \left\lbrace 0,1 \right\rbrace\\
& k,l                   \in \left\lbrace 0,1,2 \right\rbrace
\end{flalign*}
\end{flushright}

	Aplicando esta acci\'on al producto directo sale la siguiente tabla de multiplicar:


\begin{flalign*}
\begin{array}{|c|c|c|c|c|c|c|c|c|c|c|c|}
\hline
1&g&g^2&b&bg&bg^2&a&ag&ag^2&ab&abg&abg^2\\
\hline
g&g^2&1&bg^2&b&bg&ag^2&a&ag&abg&abg^2&ab\\
\hline
g^2&1&g&bg&bg^2&b&ag&ag^2&a&abg^2&ab&abg\\
\hline
b&bg&bg^2&1&g&g^2&ab&abg&abg^2&a&ag&ag^2\\
\hline
bg&bg^2&b&g^2&1&g&abg^2&ab&abg&ag&ag^2&a\\
\hline
bg^2&b&bg&g&g^2&1&abg&abg^2&ab&ag^2&a&ag\\
\hline
a&ag&ag^2&ab&abg&abg^2&1&g&g^2&b&bg&bg^2\\
\hline
ag&ag^2&a&abg^2&ab&abg&g^2&1&g&bg&bg^2&b\\
\hline
ag^2&a&ag&abg&abg^2&ab&g&g^2&1&bg^2&b&bg\\
\hline
ab&abg&abg^2&a&ag&ag^2&b&bg&bg^2&1&g&g^2\\
\hline
abg&abg^2&ab&ag^2&a&ag&bg^2&b&bg&g&g^2&1\\
\hline
abg^2&ab&abg&ag&ag^2&a&bg&bg^2&b&g^2&1&g\\
\hline
\end{array}
\end{flalign*}

Primero vemos que el orden de los elementos es:

\begin{flalign*}
&\mathsf{o}{\left({1}\right)} = 1\quad &&\mathsf{o}{\left({g}\right)} = 3\quad &&\mathsf{o}{\left({g^2}\right)} = 3\\
&\mathsf{o}{\left({b}\right)} = 2\quad &&\mathsf{o}{\left({bg}\right)} = 2\quad &&\mathsf{o}{\left({bg^2}\right)} = 2\\
&\mathsf{o}{\left({a}\right)} = 2\quad &&\mathsf{o}{\left({ag}\right)} = 2\quad &&\mathsf{o}{\left({ag^2}\right)} = 2\\
&\mathsf{o}{\left({ab}\right)} = 2\quad &&\mathsf{o}{\left({abg}\right)} = 6\quad &&\mathsf{o}{\left({abg^2}\right)} = 6
\end{flalign*}

De aqu{\'{\i}} ya podemos decir que subgrupos c{\'{\i}}clicos tiene nuestro grupo:

\begin{flalign*}
&{\mathcal{L}_{1}}{\left({G}\right)}\triangleq{\left\lbrace{\langle x \rangle\vert x                  \in G}\right\rbrace}\\
&\\
&{\mathcal{L}_{1}}{\left({{\cyclicgr{3}}\rtimes{\cyclicgr{2}                  \times                  \cyclicgr{2}}}\right)} = {\left\lbrace\right.}
{\left\lbrace{1}\right\rbrace},{\left\lbrace{1,b}\right\rbrace},\\
&\qquad\qquad\qquad\qquad{\left\lbrace{1,bg}\right\rbrace},
{\left\lbrace{1,bg^2}\right\rbrace},\\
&\qquad\qquad\qquad\qquad{\left\lbrace{1,a}\right\rbrace},
{\left\lbrace{1,ag}\right\rbrace},
{\left\lbrace{1,ag^2}\right\rbrace},
{\left\lbrace{1,ag}\right\rbrace},\\
&\qquad\qquad\qquad\qquad{\left\lbrace{1,g,g^2}\right\rbrace},
{\left\lbrace{1,g,g^2,ab,abg,abg^2}\right\rbrace}{\left.\right\rbrace}
\end{flalign*}

Los subgrupos generados estrictamente por $2$ elementos son:

\begin{flalign*}
&{\mathcal{L}_{2}}{\left({G}\right)}\triangleq
{\left\lbrace{
\langle x,y \rangle
\vert
(x,y                  \in G)                   \wedge                   (                  \nexists                   z                   \in G \;\langle x,y \rangle = \langle z\rangle)
}\right\rbrace}\\
&\\
&{\mathcal{L}_{2}}{\left({{\mathsf{C}_3}\rtimes{({\cyclicgr{2}}                  \times                  {\cyclicgr{2}})}}\right)}=\left\lbrace\right.
{\left\lbrace{1,b,a,ab}\right\rbrace},
{\left\lbrace{1,bg,ag,ab}\right\rbrace},
{\left\lbrace{1,bg^2,ag^2,ab}\right\rbrace},\\
&\qquad\qquad\qquad\qquad\qquad\;\,{\left\lbrace{1,g,g^2,b,bg,bg^2,a,ag,ag^2,ab,abg,abg^2}\right\rbrace}
\left.\right\rbrace
\end{flalign*}

Con esto queda claro que no hay subgrupos propiamente generados por $3$ elementos de forma minimal, pues con $2$ elementos se genera todo el grupo.

El centro del grupo es:
\[
\mathsf{Z}\left({{\cyclicgr{3}}\rtimes({\cyclicgr{2}}                  \times                  {\cyclicgr{2}})}\right) = {\left\lbrace{1,ab}\right\rbrace}
\]


Los subgrupos que son " absolutamente no normales ", esto es, aquellos que coinciden con su normalizador:


\begin{flalign*}
&
\mathsf{N}_{{\cyclicgr{3}}\rtimes({\cyclicgr{2}}                  \times                  {\cyclicgr{2}})}
\left(
{\left\lbrace
{1,b,a,ab}
\right\rbrace}
\right) =
{\left\lbrace
{1,b,a,ab}
\right\rbrace}\\
&
\mathsf{N}_{{\cyclicgr{3}}\rtimes({\cyclicgr{2}}                  \times                  {\cyclicgr{2}})}
\left(
{\left\lbrace
{1,bg,ag,ab}
\right\rbrace}
\right) =
{\left\lbrace
{1,bg,ag,ab}
\right\rbrace}\\
&
\mathsf{N}_{{\cyclicgr{3}}\rtimes({\cyclicgr{2}}                  \times                  {\cyclicgr{2}})}
\left(
{\left\lbrace
{1,bg^2,ag^2,ab}
\right\rbrace}
\right) =
{\left\lbrace
{1,bg^2,ag^2,ab}
\right\rbrace}
\end{flalign*}

A continuaci\'on, los que tienen un normalizador estrictamente mayor que ellos mismos, pero sin ser normales, son:

\begin{flalign*}
&\mathsf{N}_{{\cyclicgr{3}}\rtimes({\cyclicgr{2}}                  \times                  {\cyclicgr{2}})}
\left({\left\lbrace{1,b}\right\rbrace}\right) =
\mathsf{N}_{{\cyclicgr{3}}\rtimes({\cyclicgr{2}}                  \times                  {\cyclicgr{2}})}
\left({\left\lbrace{1,a}\right\rbrace}\right) =
\left\lbrace{1,b,a,ab}\right\rbrace
\\
&\mathsf{N}_{{\cyclicgr{3}}\rtimes({\cyclicgr{2}}                  \times                  {\cyclicgr{2}})}
\left({\left\lbrace{1,bg}\right\rbrace}\right) =
\mathsf{N}_{{\cyclicgr{3}}\rtimes({\cyclicgr{2}}                  \times                  {\cyclicgr{2}})}
\left({\left\lbrace{1,ag}\right\rbrace}\right) =
{\left\lbrace{1,bg,ag,ab}\right\rbrace}
\\
&\mathsf{N}_{{\cyclicgr{3}}\rtimes({\cyclicgr{2}}                  \times                  {\cyclicgr{2}})}
\left({\left\lbrace{1,bg^2}\right\rbrace}\right) =
\mathsf{N}_{{\cyclicgr{3}}\rtimes({\cyclicgr{2}}                  \times                  {\cyclicgr{2}})}
\left({\left\lbrace{1,ag^2}\right\rbrace}\right) =
{\left\lbrace{1,bg^2,ag^2,ab}\right\rbrace}
\end{flalign*}

Los subgrupos normales quedan:

\begin{flalign*}
&{\mathcal{L}}_{{{\mathsf{C}}_{3}}\rtimes
{\left({
{{\mathsf{C}}_{2}}                  \times                  {{\mathsf{C}}_{2}}
}\right)}}=
{\left\lbrace\right.}
{\left\lbrace{1}\right\rbrace},
{\left\lbrace{1,ab}\right\rbrace},
{\left\lbrace{1,g,g^2}\right\rbrace},
{\left\lbrace{1,g,g^2,b,bg,bg^2}\right\rbrace},\\
&{\left\lbrace{1,g,g^2,a,ag,ag^2}\right\rbrace},
{\left\lbrace{1,g,g^2,ab,abg,abg^2}\right\rbrace},
{{\mathsf{C}}_{3}}
\rtimes
{\left({
{{\mathsf{C}}_{2}}                  \times                  {{\mathsf{C}}_{2}}
}\right)}
{\left.\right\rbrace}
\end{flalign*}

De estos hay $2$ subgrupos no abelianos de orden $6$ que no son nilpotentes, y de hecho son isomorfos a $\mathsf{Dih}_{3}$, a saber son ${\left\lbrace{1,g,g^2,b,bg,bg^2}\right\rbrace}$y ${\left\lbrace{1,g,g^2,a,ag,ag^2}\right\rbrace}$. Por lo tanto el ret{\'{\i}}culo de subgrupos nilpotentes normales es:

\begin{flalign*}
&{\mathcal{L}}_{{{\mathsf{C}}_{3}}\rtimes
{\left({
{{\mathsf{C}}_{2}}                  \times                  {{\mathsf{C}}_{2}}
}\right)}}^{\mathcal{NIL}}=
\left\lbrace
{\left\lbrace{1}\right\rbrace},
{\left\lbrace{1,ab}\right\rbrace},
{\left\lbrace{1,g,g^2}\right\rbrace},
{\left\lbrace{1,g,g^2,ab,abg,abg^2}\right\rbrace}
\right\rbrace\\
&\left\langle{\bigcup{\mathcal{L}}_{{{\mathsf{C}}_{3}}\rtimes
{\left({
{{\mathsf{C}}_{2}}                  \times                  {{\mathsf{C}}_{2}}
}\right)}}^{\mathcal{NIL}}}\right\rangle = {\left\lbrace{1,g,g^2,ab,abg,abg^2}\right\rbrace} = \left\langle abg \right\rangle
\end{flalign*}

Buscamos ahora los centralizadores de cada subgrupo en el grupo ambiente, quitando el trivial, el impropio y el centro que ya se ha hallado.
Primero buscamos los subgrupos propios abelianos (lo son todos los propios) pero que su centralizador es el mismo subgrupo:

\begin{flalign*}
&\cyclicgr{{\cyclicgr{3}}\rtimes({\cyclicgr{2}}                  \times                  {\cyclicgr{2}})}\left({\left\lbrace{1,b,a,ab}\right\rbrace}\right) = {\left\lbrace{1,b,a,ab}\right\rbrace}\\
&\cyclicgr{{\cyclicgr{3}}\rtimes({\cyclicgr{2}}                  \times                  {\cyclicgr{2}})}\left({\left\lbrace{1,bg,ag,ab}\right\rbrace}\right) = {\left\lbrace{1,bg,ag,ab}\right\rbrace}\\
&\cyclicgr{{\cyclicgr{3}}\rtimes({\cyclicgr{2}}                  \times                  {\cyclicgr{2}})}\left({\left\lbrace{1,bg^2,ag^2,ab}\right\rbrace}\right) = {\left\lbrace{1,bg^2,ag^2,ab}\right\rbrace}\\
&\cyclicgr{{\cyclicgr{3}}\rtimes({\cyclicgr{2}}                  \times                  {\cyclicgr{2}})}\left({\left\lbrace{1,g,g^2,ab,abg,abg^2}\right\rbrace}\right) = {\left\lbrace{1,g,g^2,ab,abg,abg^2}\right\rbrace}
\end{flalign*}

A continuaci\'on aquellos subgrupos que tienen centralizador mayor que ellos mismos:

\begin{flalign*}
&\cyclicgr{{\cyclicgr{3}}\rtimes
({\cyclicgr{2}}                  \times                  {\cyclicgr{2}})}\left({\left\lbrace{1,b}\right\rbrace}\right) =
\cyclicgr{{\cyclicgr{3}}\rtimes
({\cyclicgr{2}}                  \times                  {\cyclicgr{2}})}\left({\left\lbrace{1,a}\right\rbrace}\right) =
{\left\lbrace{1,b,a,ab}\right\rbrace}\\
&\cyclicgr{{\cyclicgr{3}}\rtimes
({\cyclicgr{2}}                  \times                  {\cyclicgr{2}})}\left({\left\lbrace{1,bg}\right\rbrace}\right) =
\cyclicgr{{\cyclicgr{3}}
\rtimes({\cyclicgr{2}}                  \times                  {\cyclicgr{2}})}\left({\left\lbrace{1,ag}\right\rbrace}\right) =
{\left\lbrace{1,bg,ag,ab}\right\rbrace}\\
&\cyclicgr{{\cyclicgr{3}}
\rtimes
({\cyclicgr{2}}                  \times                  {\cyclicgr{2}})}\left({\left\lbrace{1,bg^2}\right\rbrace}\right) =
\cyclicgr{{\cyclicgr{3}}
\rtimes
({\cyclicgr{2}}                  \times                  {\cyclicgr{2}})}\left({\left\lbrace{1,ag^2}\right\rbrace}\right) =
 {\left\lbrace{1,bg^2,ag^2,ab}\right\rbrace}\\
&\cyclicgr{{\cyclicgr{3}}\rtimes({\cyclicgr{2}}                  \times
{\cyclicgr{2}})}\left({\left\lbrace{1,g,g^2,abg}\right\rbrace}\right) = {\left\lbrace{1,g,g^2,ab,abg,abg^2}\right\rbrace}
\end{flalign*}


Los subgrupos derivados vienen:

\begin{flalign*}
&\mathsf{D^{(1)}}{\left({\cyclicgr{3}}\rtimes({\cyclicgr{2}}                  \times                  {\cyclicgr{2}})\right)} = {\left\lbrace{1,g,g^2}\right\rbrace}\\
&\mathsf{D^{(2)}}{\left({\cyclicgr{3}}\rtimes({\cyclicgr{2}}                  \times                  {\cyclicgr{2}})\right)} = {\left\lbrace{1}\right\rbrace}
\end{flalign*}

De d\'onde vemos que el grupo ${\cyclicgr{3}}\rtimes({\cyclicgr{2}}                  \times                  {\cyclicgr{2}})$ \textsc{ es resoluble}.

Los subgrupos L-derivados del grupo m\'aximo (en el ret{\'{\i}}culo):

\begin{flalign*}
\mathsf{L_{(1)}}{\left({\cyclicgr{3}}\rtimes({\cyclicgr{2}}                  \times                  {\cyclicgr{2}})
\right)} = \mathsf{L_{(2)}}{\left({\cyclicgr{3}}\rtimes({\cyclicgr{2}}                  \times                  {\cyclicgr{2}})
\right)} = {\left\lbrace{1,g,g^2}\right\rbrace}
\end{flalign*}

De d\'onde vemos que el grupo
${\cyclicgr{3}}\rtimes({\cyclicgr{2}}                  \times                  {\cyclicgr{2}})$ \textsc{no es nilpotente}.

Para aclararnos un poco entre tanto elemento y subgrupo veamos como se pueden generar los
subgrupos:

\begin{flalign*}
&{\left\langle{b}\right\rangle} = {\left\lbrace{1,b}\right\rbrace}\\
&{\left\langle{bg}\right\rangle} = {\left\lbrace{1,bg}\right\rbrace}\\
&{\left\langle{bg^2}\right\rangle} = {\left\lbrace{1,bg^2}\right\rbrace}\\
&{\left\langle{a}\right\rangle} = {\left\lbrace{1,a}\right\rbrace}\\
&{\left\langle{ag}\right\rangle} = {\left\lbrace{1,ag}\right\rbrace}\\
&{\left\langle{ag^2}\right\rangle} = {\left\lbrace{1,ag^2}\right\rbrace}\\
&{\left\langle{ab}\right\rangle} = {\left\lbrace{1,ab}\right\rbrace}\\
&{\left\langle{g,b}\right\rangle} = {\left\langle{g^2,b}\right\rangle} =
{\left\langle{g,bg}\right\rangle} = {\left\langle{g,bg^2}\right\rangle} = {\left\lbrace{1,g,g^2,b,bg,bg^2}\right\rbrace}\\
&{\left\langle{g,a}\right\rangle} = {\left\langle{g^2,a}\right\rangle} =
{\left\langle{g,ag}\right\rangle} = {\left\langle{g,ag^2}\right\rangle} = {\left\lbrace{1,g,g^2,a,ag,ag^2}\right\rbrace}\\
&{\left\langle{g}\right\rangle} = {\left\lbrace{1,g,g^2}\right\rbrace}\\
&{\left\langle{abg}\right\rangle} =
{\left\langle{g,ab}\right\rangle} = {\left\langle{g^2,ab}\right\rangle} =
{\left\lbrace{1,g,g^2,ab,abg,abg^2}\right\rbrace}\\
&{\left\langle{b,a}\right\rangle} = {\left\langle{b,ab}\right\rangle} = {\left\lbrace{1,b,a,ab}\right\rbrace}\\
&{\left\langle{bg,ag}\right\rangle} = {\left\langle{bg,ab}\right\rangle} = {\left\lbrace{1,bg,ag,ab}\right\rbrace}\\
&{\left\langle{bg^2,ag^2}\right\rangle} = {\left\langle{bg^2,ab}\right\rangle} = {\left\lbrace{1,bg^2,ag^2,ab}\right\rbrace}\\
&{\left\langle{g,b,a}\right\rangle} = {\left\lbrace{1,g,g^2,b,bg,bg^2,a,ag,ag^2,ab,abg,abg^2}\right\rbrace}
\end{flalign*}

El subgrupo de Frattini, la intersecci\'on de los subgrupos maximales, es el subgrupo trivial,
${\Phi }_{\cyclicgr{3}}\rtimes{\left({\cyclicgr{2}}         \times         {\cyclicgr{2}}\right)}=
\lbrace 1\rbrace$. El subgrupo de Fitting, la suma de todos los subgrupos nilpotentes
(todos excepto el propio) normales, coincidir\'a con el ya hallado anteriormente,
$\mathsf{F}\left(\cyclicgr{3}\rtimes\left({\cyclicgr{2}}         \times         {\cyclicgr{2}}\right)\right)=\langle abg \rangle$.

Solo nos queda ver como se factoriza el grupo, primero a trav\'es del centro:

\begin{flalign*}
&
\nicefrac
{\cyclicgr{3}\rtimes\left({\cyclicgr{2}}                  \times                  {\cyclicgr{2}}\right)}
{\left\langle ab \right\rangle}
=
\left\lbrace\right.
1{\left\langle{ab}\right\rangle},
g{\left\langle{ab}\right\rangle},
g^2{\left\langle{ab}\right\rangle},
b{\left\langle{ab}\right\rangle}=
a{\left\langle{ab}\right\rangle},
bg{\left\langle{ab}\right\rangle}
=
ag{\left\langle{ab}\right\rangle},\\
&
bg^2{\left\langle{ab}\right\rangle}
=
ag^2{\left\langle{ab}\right\rangle}
\left.\right\rbrace
\end{flalign*}

Que tomando representantes:

\begin{flalign*}
&
\nicefrac
{\cyclicgr{3}\rtimes\left({\cyclicgr{2}}                  \times                  {\cyclicgr{2}}\right)}
{\left\langle ab \right\rangle}
=
\left\lbrace
\overline{1},
\overline{g},
\overline{g^2},
\overline{a},
\overline{ag},
\overline{ag^2}
\right\rbrace
\end{flalign*}

Que tiene la estructura del grupo $\mathsf{Dih}_{3}$ y el completo tiene la apariencia de $\mathsf{Dih}_{3}                  \times                  \cyclicgr{2}$.

Para empezar, $ab$ conmuta con todo el grupo, y $\langle a,g \rangle                  \cong                  \mathsf{Dih}_{3}$ es tambi\'en grupo normal. De otro lado $\langle ab\rangle                   \cap \langle a,g \rangle = \lbrace 1 \rbrace$. Por lo tanto:
\[
{\cyclicgr{3}\rtimes\left({\cyclicgr{2}                  \times                  \cyclicgr{2}}\right)}                  \cong                  {\mathsf{Dih}_{3}                  \times                  \cyclicgr{2}}
\]

Si llamamos a $\mathtt{f}\triangleq abg$, $\mathtt{g}\triangleq a$ y $\mathtt{a} \triangleq ab$ explicitamos el isomorfismo antes dicho.

Observamos con facilidad que ${\cyclicgr{3}}\rtimes{\left({{\cyclicgr{2}}                  \times                  {\cyclicgr{2}}}\right)}                  \cong                  \mathsf{Dih}_{6}$. Veamos esto \'ultimo. Sustituimos $abg$ por $\mathfrak{f}\triangleq abg$, siendo $(abg)^2 = g^2$ de d\'onde  $\mathfrak{f}^{2}\triangleq g^2$, ${(abg)}^{3}=ab$, obtenemos $\mathfrak{f}^{3}\triangleq ab$, $(abg)^{4}= g$ que nos da $\mathfrak{f}^{4}\triangleq g$ y por \'ultimo $(abg)^{5}= abg^2$ que se traslada a $\mathfrak{f}^{5}\triangleq abg^2$.  Podemos escoger $b$ como elemento no normal y de orden $2$ del grupo dih\'edrico, escogeremos $\mathfrak{g}\triangleq b$. Como resultado $\mathfrak{gf}=a$, y tenemos solo $2$ generadores como hab{\'{\i}}amos deducido al generar con la herramienta el grupo total. Via las sustituciones hechas, vemos que:

\[{\cyclicgr{3}}\rtimes{\left({{\cyclicgr{2}}                  \times                  {\cyclicgr{2}}}\right)}                  \cong                  {\cyclicgr{6}}\rtimes{\cyclicgr{2}}                  \cong                  \mathsf{Dih}_{6}                  \cong                  {\mathsf{Dih}_{3}}                \times                {\cyclicgr{2}}\]

Vemos como factoriza a trav\'es de $\langle abg\rangle$:

\begin{flalign*}
\overline{a}\triangleq{\left\lbrace{a,b,ag,ag^2,bg,bg^2}\right\rbrace},
\overline{1}\triangleq{\left\lbrace{1,g^2,g,ab,abg,abg^2}\right\rbrace}
\end{flalign*}

Que no es otra cosa que el grupo c{\'{\i}}clico de orden $2$. Por eso no tiene sentido hacer el cociente por los otros subgrupos de orden $6$, puesto que partir\'an el grupo en dos partes, la unidad y un elemento de orden $2$.

El diagrama de Hasse es notablemente m\'as complicado que  los anteriores:

La siguiente versi\'on del diagrama de Hasse a\'un se sale de los m\'argenes, m\'as adelante los retocar\'e:
\begin{center}
    \begin{tikzpicture}
    % First, locate each of the nodes and name them
    \node (a-b-g) at (-8,12)
    {$\mathsf{Dih}_{6}         \cong
    {\left\langle{a,b,g\vert a^2=b^2=g^3=1,xab=abx,ag=g^2a,bg=g^b}\right\rangle}$};
    \node (b-g) at (-14,9)
    {$\left\langle{b,g\vert b^2=g^3=1,bg=g^2b}\right\rangle$};
    \node (a-g) at (-8,9)
    {$\left\langle{a,g\vert a^2=g^3=1,ag=g^2a}\right\rangle$};
    \node (abg) at (-2,9)
    {$         \cong         \left\langle{abg\vert \left(ab\right)^6=1}\right\rangle$};
    \node (a-b) at (-14,6)
    {$\left\langle{a,b\vert a^2=b^2=1}\right\rangle$};
    \node (ag-bg) at (-8,6)
	{$\left\langle{ag,bg\vert (ag)^2=(bg)^2=1}\right\rangle$};
    \node (ag2-bg2) at (-2,6)
    {$\left\langle{ag^2,bg^2\vert (ag^2)^2=(bg^2)^2=1}\right\rangle$};
    \node (g) at (-8,3)
	{$\left\langle{g\vert g^3=1}\right\rangle$};
    \node (ab) at (-14,0)
    {$\left\langle{ab}\right\rangle$};
    \node (b) at (-12,0)
    {$\left\langle{b}\right\rangle$};
    \node (a) at (-10,0)
    {$\left\langle{a}\right\rangle$};
    \node (bg) at (-8,0)
	{$\left\langle{bg}\right\rangle$};
    \node (ag) at (-6,0)
    {$\left\langle{ag}\right\rangle$};
    \node (bg2) at (-4,0)
    {$\left\langle{bg^2}\right\rangle$};
    \node (ag2) at (-2,0)
    {$\left\langle{ag^2}\right\rangle$};
	\node (1) at (-8,-3)
    {$\left\langle{1}\right\rangle$};
   % Now draw the lines:
    \draw [black, thick] (a-b-g) -- (abg);
    \draw [black, thick] (a-b-g) -- (a-g);
    \draw [black, thick] (a-b-g) -- (b-g);
    \draw [black, thick] (a-b-g) -- (a-b);
    \draw [black, thick] (a-b-g) -- (ag-bg);
    \draw [black, thick] (a-b-g) -- (ag2-bg2);
    \draw [black, thick] (abg) -- (g);
    \draw [black, thick] (abg) -- (ab);
    \draw [black, thick] (a-g) -- (a);
    \draw [black, thick] (a-g) -- (ag);
    \draw [black, thick] (a-g) -- (ag2);
    \draw [black, thick] (a-g) -- (g);
    \draw [black, thick] (b-g) -- (b);
    \draw [black, thick] (b-g) -- (bg);
    \draw [black, thick] (b-g) -- (bg2);
    \draw [black, thick] (b-g) -- (g);
    \draw [black, thick] (a-b) -- (a);
    \draw [black, thick] (a-b) -- (ab);
    \draw [black, thick] (a-b) -- (b);
    \draw [black, thick] (ag-bg) -- (ag);
    \draw [black, thick] (ag-bg) -- (bg);
    \draw [black, thick] (ag2-bg2) -- (ag2);
    \draw [black, thick] (ag2-bg2) -- (bg2);
    \draw [black, thick] (g) -- (1);
    \draw [black, thick] (a) -- (1);
    \draw [black, thick] (b) -- (1);
    \draw [black, thick] (ab) -- (1);
    \draw [black, thick] (ag) -- (1);
    \draw [black, thick] (bg) -- (1);
    \draw [black, thick] (ag2) -- (1);
    \draw [black, thick] (bg2) -- (1);
    \end{tikzpicture}
\end{center}

Antes de seguir, nos damos cuenta que al haber encontrado el grupo dih\'edrico, podemos encontrar en este caso, $4\shortmid12$, el grupo dic{\'{\i}}clico correspondiente, tal y como hicimos en el caso $\mathsf{Q}_{8}$. Procedemos:

El grupo $\mathsf{Dih}_{6}$ se puede caracterizar, como ya hicimos en el cap{\'{\i}}tulo de los grupos de orden $2p$, de la siguiente forma:
\begin{flalign*}
&\varphi : \cyclicgr{2}              \longrightarrow              \mathtt{Aut}\left(\cyclicgr{6}\right)\\
&{\varphi }_{a^j}=
\begin{cases}
& \mathtt{id}_{\cyclicgr{6}}              \Leftarrow              a^j=1\\
& \mathtt{inv}_{\cyclicgr{6}}              \Leftarrow              a^j             \neq              1\\
\end{cases}
\end{flalign*}

Ahora le aplicamos la acci\'on al miembro normal (c{\'{\i}}clico de orden $6$), una acci\'on dependiente del $V$ grupo de Klein que forman las dos proyecciones derechas (en nuestra representaci\'on) de los argumentos de la multiplicaci\'on:

\begin{flalign*}
&(a^i,b^k)*(a^j,b^m)=(a^l,b^n) \mbox{ referencia}\\
&\zeta : \cyclicgr{2}             \times             \cyclicgr{2}              \longrightarrow              \cyclicgr{6}\\
&{\zeta }_{(a^i,a^j)}=
\begin{cases}
& b^3              \Leftarrow              a^i             \neq              a^j\\
& 1              \Leftarrow              a^i=a^j\\
\end{cases}\\
&(a^i,b^k)*(a^j,b^m)\triangleq (a^{(i+j)},{\varphi }_{a^j}\left(b^k\right) b^n {\zeta }_{a^i,a^j})
\end{flalign*}

Simplificando la expresi\'on solo a los exponentes:


\begin{flalign*}
&(i,k)*(j,m)=(l,n) \mbox{ referencia}\\
&{\varphi }_{l}^{\prime}\left(k\right)=
\begin{cases}
& k              \Leftarrow              j=0 \mod 2\\
& 6-k              \Leftarrow              j=1 \mod 2\\
\end{cases}\\
&\zeta : \cyclicgr{2}             \times             \cyclicgr{2}              \longrightarrow              \cyclicgr{6}\\
&{\zeta }_{(i,j)}^{\prime}=
\begin{cases}
& 3              \Leftarrow              i             \neq              j \mod 2\\
& 0              \Leftarrow              i=j \mod 2\\
\end{cases}\\
&(i,k)*(j,m)\triangleq (i+j,{\varphi }_{j}^{\prime}\left(k\right)+n+{\zeta }_{a^i,a^j}^{\prime})
\end{flalign*}

A implementar queda:

\begin{flalign*}
&(i,k)*(j,m)\triangleq (i+j,{\varphi }_{j}^{\prime}\left(k\right)+n+{\zeta }_{a^i,a^j}^{\prime})\\
&(i,k)*(j,m)=
\begin{cases}
& (i+j,k+m)              \Leftarrow              i=j=0\\
& (i+j,-k+m+6)              \Leftarrow              i=j=1\\
& (i+j,-k+m+9)              \Leftarrow              j=1             \wedge              i=0\\
& (i+j,k+m+3)              \Leftarrow              j=0             \wedge              i=1\\
\end{cases}
\end{flalign*}

La tabla de Cayley del grupo dic{\'{\i}}clico es:

\begin{flalign*}
\begin{array}{|c|c|c|c|c|c|c|c|c|c|c|c|}
\hline
1&g&g^2&g^3&g^4&g^5&f&fg&fg^2&fg^3&fg^4&fg^5\\
\hline
g&g^2&g^3&g^4&g^5&1&fg^5&f&fg&fg^2&fg^3&fg^4\\
\hline
g^2&g^3&g^4&g^5&1&g&fg^4&fg^5&f&fg&fg^2&fg^3\\
\hline
g^3&g^4&g^5&1&g&g^2&fg^3&fg^4&fg^5&f&fg&fg^2\\
\hline
g^4&g^5&1&g&g^2&g^3&fg^2&fg^3&fg^4&fg^5&f&fg\\
\hline
g^5&1&g&g^2&g^3&g^4&fg&fg^2&fg^3&fg^4&fg^5&f\\
\hline
f&fg&fg^2&fg^3&fg^4&fg^5&g^3&g^4&g^5&1&g&g^2\\
\hline
fg&fg^2&fg^3&fg^4&fg^5&f&g^2&g^3&g^4&g^5&1&g\\
\hline
fg^2&fg^3&fg^4&fg^5&f&fg&g&g^2&g^3&g^4&g^5&1\\
\hline
fg^3&fg^4&fg^5&f&fg&fg^2&1&g&g^2&g^3&g^4&g^5\\
\hline
fg^4&fg^5&f&fg&fg^2&fg^3&g^5&1&g&g^2&g^3&g^4\\
\hline
fg^5&f&fg&fg^2&fg^3&fg^4&g^4&g^5&1&g&g^2&g^3\\
\hline
\end{array}
\end{flalign*}

Desde esta tabla podemos establecer f\'acilmente el isomorfismo que nos falta, que es:
                \[\mathsf{Dic}_{3}                \cong                \cyclicgr{3}\rtimes\cyclicgr{4}\]


Solo queda ver la situaci\'on actual y si son posibles m\'as grupos de orden $12$ de tipo no-abeliano.
Hemos visto el grupo $\mathsf{Dic}_{3}                \cong                \cyclicgr{3}\rtimes\cyclicgr{4}$ y
el $\mathsf{Dih}_{6}$.

El \'ultimo tiene $3$ grupos normales de orden $6$, no isomorfos entre s{\'{\i}}. No son todos los maximales,
aunque su uni\'on es todo el grupo. Tiene un grupo de orden $3$ que es caracter{\'{\i}}stico, y los de orden
$4$ son de nuevo $3$, conjugados y no normales, y $7$ grupos de orden $2$ de los cuales uno de ellos
es normal-caracter{\'{\i}}stico y los otros no son normales.

El primero ($\mathsf{Dic}_{3}$) tiene solo un subgrupo de orden $6$, normal y caracter{\'{\i}}stico, al
igual que los subgrupos \'unicos de orden $2$ y $3$. A su vez tiene como el anterior, $3$ subgrupos de
orden $4$ (c{\'{\i}}clicos a diferencia del dih\'edrico) tambi\'en no normales.

La pregunta es si puede haber un grupos del orden que estudiamos que tenga alg\'un grupo de orden $4$ y
que sea normal. Con subgrupos de orden $4$ c{\'{\i}}clicos solo puede estar $\mathsf{Dic}_{3}$, siendo un
$2-grupo$ de Sylow. Con subgrupos $2-grupo$ de Sylow isomorfos al $V$ grupo de Klein, hemos
encontrado el dih\'edrico, pero podemos buscar uno que factorize sobre el $V$ grupo de Klein, esto es,
que sea normal. Esta es la posibilidad:

\begin{flalign*}
\mathbf{1}                 \hookrightarrow                 \cyclicgr{2}                \times                \cyclicgr{2} \overset{\alpha }{                \hookrightarrow                } \mathbf{G} \overset{\beta }\twoheadrightarrow \cyclicgr{3} \twoheadrightarrow \mathbf{1}
\end{flalign*}

D\'onde la acci\'on no puede ser simple, como las vistas hasta ahora (saldr{\'{\i}}a de nuevo el grupo dih\'edrico o el abeliano). Especifiquemos los homomorfismos $\alpha$ y $\beta$:

\begin{flalign*}
&{\alpha } : {\mathsf{C}}_{2}                \times                {\mathsf{C}}_{2}                 \hookrightarrow                 \mathbf{G}\\
&			\left(a^l,b^m\right)                 \mapsto \left(\overline{a}^l,\overline{b}^{m}\right)\\
& {\alpha }                \circ {{\alpha }^{-1}}
                \equiv
{\mathtt{Id}}_{{{\mathsf{C}}_{2}}                \times                {{\mathsf{C}}_{2}}}\\
&\beta : \mathbf{G} \twoheadrightarrow {\mathsf{C}}_{3}\\
&			\left(\overline{x},\overline{c}^k\right)                 \mapsto c^k\\
& {\alpha }^{-1}\left(\overline{x}\right)=x                \in \cyclicgr{2}                \times                \cyclicgr{2}\\
& \mathsf{o}\left(\overline{c}\right)=\mathsf{o}\left(c\right)=3\\
& {{\beta }^{-1}}{\left(\overline{1}\right)}=
{\alpha }{\left({{\mathsf{C}}_{2}                \times                {\mathsf{C}}_{2}}\right)}=
\ker\left({\beta }\right)\\
& {\beta }^{-1}\left(\overline{c}\right)=\ker\left(\beta \right) c\\
& {\beta }^{-1}\left(\overline{c}^{2}\right)=\ker\left(\beta \right) c^2\\
& \mbox{\textsc{ es secuencia exacta corta }}\\
& {\mathbf{G}}\triangleq {\cyclicgr{2}                \times                \cyclicgr{2}}\bowtie{\cyclicgr{3}}
\end{flalign*}

Todo lo que hemos conocido nos habilita para construir el grupo $\mathbf{G}$ buscado, sobre un conjunto subyacente $\cyclicgr{2}                \times                \cyclicgr{2}                \times                \cyclicgr{3}$ pero con las siguientes acciones que conformar\'an la nueva operaci\'on de grupo, d\'onde primero tomaremos las referencias de exponentes por facilidad al escribir y porque la herramienta lo toma as{\'{\i}}, por exponentes:
\begin{flalign*}
&(a^i,b^j,c^k)*(a^l,b^m,c^n)=(a^r,b^s,c^t)\\
&(i\,j,k)(l,m,n)=(r,s,t)\\
\end{flalign*}
Se trata de averiguar $(r,s,t)$.
Primero los homomorfismos para las acciones:
\begin{flalign*}
&\varphi : \cyclicgr{3}                 \longrightarrow                 \mathtt{Aut}(\cyclicgr{2}               \times               \cyclicgr{2})\\
&\varphi \left( k \right) =
\begin{cases}
&(l,m)                \mapsto (l,m)                 \Leftarrow                 k=0\\
&(l,m)                \mapsto (l+m,l)                \Leftarrow                 k=1\\
&(l,m)                \mapsto (m,l+m)                \Leftarrow                 k=2\\
\end{cases}
\end{flalign*}
\begin{flalign*}
&\varphi \mbox{\textsc{ es homomorfismo de grupos }}
\end{flalign*}
\begin{flalign*}
&\zeta : {\mathsf{C}}_{2}                \times                {\mathsf{C}}_{2}                 \longrightarrow                 {\mathcal{S}ym}
{\left({{\mathsf{C}}_{2}                \times                {\mathsf{C}}_{2}}\right)}
\end{flalign*}
\begin{flalign*}
&{\zeta }_{\left((l,m)\right)}=
\begin{cases}
(i,j)                 \Leftarrow                (l,m)=(0,0)\\
(i,i+1)                 \Leftarrow                 (l,m)=(0,1)\\
(i+1,j)                 \Leftarrow                 (l,m)=(1,0)\\
(i+1,j+1)                 \Leftarrow                 (l,m)=(1,1)
\end{cases}
\end{flalign*}
\begin{flalign*}
& \zeta \mbox{\textsc{ es homomorfismo de grupos }}
\end{flalign*}

Ahora la operaci\'on binaria queda:

\begin{flalign*}
(i,j,k)*(l,m,n)\triangleq
(
{\zeta }_{(l,m)}(i,j),
k
)
(
{\varphi }_{k}(l,m),
n
)
\end{flalign*}

La tabla de multiplicaci\'on de este grupo queda as{\'{\i}}:

\begin{flalign*}
\begin{array}{|c|c|c|c|c|c|c|c|c|c|c|c|}
\hline
1&c&c^2&b&bc&bc^2&a&ac&ac^2&ab&abc&abc^2\\
\hline
c&c^2&1&abc&abc^2&ab&bc&bc^2&b&ac&ac^2&a\\
\hline
c^2&1&c&ac^2&a&ac&abc^2&ab&abc&bc^2&b&bc\\
\hline
b&bc&bc^2&1&c&c^2&ab&abc&abc^2&a&ac&ac^2\\
\hline
bc&bc^2&b&ac&ac^2&a&c&c^2&1&abc&abc^2&ab\\
\hline
bc^2&b&bc&abc^2&ab&abc&ac^2&a&ac&c^2&1&c\\
\hline
a&ac&ac^2&ab&abc&abc^2&1&c&c^2&b&bc&bc^2\\
\hline
ac&ac^2&a&bc&bc^2&b&abc&abc^2&ab&c&c^2&1\\
\hline
ac^2&a&ac&c^2&1&c&bc^2&b&bc&abc^2&ab&abc\\
\hline
ab&abc&abc^2&a&ac&ac^2&b&bc&bc^2&1&c&c^2\\
\hline
abc&abc^2&ab&c&c^2&1&ac&ac^2&a&bc&bc^2&b\\
\hline
abc^2&ab&abc&bc^2&b&bc&c^2&1&c&ac^2&a&ac\\
\hline
\end{array}
\end{flalign*}

Los subgrupos c{\'{\i}}clicos son (como era de esperar solo \'ordenes $2$ y $3$):

\begin{flalign*}
&\mathcal{L}_{1}\left(\mathbf{G}\right)=\left\lbrace\right.
{\left\lbrace{1}\right\rbrace} ,
{\left\lbrace{1,b}\right\rbrace} ,
{\left\lbrace{1,a}\right\rbrace} ,
{\left\lbrace{1,ab}\right\rbrace} ,
{\left\lbrace{1,c,c^2}\right\rbrace} ,\\
&\qquad{\left\lbrace{1,bc,ac^2}\right\rbrace} ,
{\left\lbrace{1,bc^2,abc}\right\rbrace} ,
{\left\lbrace{1,ac,abc^2}\right\rbrace}
\left.\right\rbrace
\end{flalign*}

Los subgrupos generados minimalmente por m\'as de un elemento (por $2$ concretamente):

\begin{flalign*}
&\mathcal{L}_{2}\left(\mathbf{G}\right)=
\left\lbrace\right.{
{\left\lbrace{1,b,a,ab}\right\rbrace} ,
{\mathbf{G}}
}\left.\right\rbrace
\end{flalign*}

Siendo el subgrupo $\left\{1,a,b,ab\right\}$ el \'unico propio normal que no sea el trivial.
Adem\'as es el \'unico de su orden de d\'onde es adem\'as caracter{\'{\i}}stico.
Por otra parte solo los subgrupos de orden $2$ que sean subgrupo del $V$ grupo de Klein tienen normalizador mayor que ellos mismos.

Como vemos no existe subgrupo de orden $6$, por lo que es el primer grupo en el que no se invierte el teorema de Lagrange.

Construyendo el homomorfismo de $\left(\cyclicgr{2}             \times             \cyclicgr{2}\right)\bowtie\cyclicgr{3}$
en el grupo alternado $\mathsf{A}_{4}$, $a             \mapsto (1,2)(3,4)$, $b             \mapsto (1,3)(2,4)$(este se podr{\'{\i}}a deducir),
y $c             \mapsto (1,2,3)$ podemos comprobar que se trata de un isomorfismo, as{\'{\i}} que tenemos:

\[
  	\left(\cyclicgr{2}             \times             \cyclicgr{2}\right)\bowtie\cyclicgr{3} \;             \cong             \;\mathsf{A}_{4}
\]

El centro es el trivial $\mathsf{Z}\left(\mathsf{A}_{4}\right)=\lbrace{1}\rbrace$,
de forma que el grupo de automorfismos internos es isomorfo al propio grupo.

En cuanto a los grupos derivados:

\[
  \mathsf{D^{(1)}}{\left({\cyclicgr{2}             \times             \cyclicgr{2}\bowtie\cyclicgr{3}}\right)} = {\left\lbrace{1,b,a,ab}\right\rbrace}
\]
\[
\mathsf{D^{(2)}}{\left({\cyclicgr{2}             \times             \cyclicgr{2}\bowtie\cyclicgr{3}}\right)} = {\left\lbrace{1}\right\rbrace}
\]

As{\'{\i}} que el grupo alternado $\mathsf{A}_{4}$ \textsc{es resoluble}. Sin embargo:
\[
  \mathsf{L^{(2)}}{\left({\cyclicgr{2}             \times             \cyclicgr{2}\bowtie\cyclicgr{3}}\right)} = {\left\lbrace{1,b,a,ab}\right\rbrace}
\]

Nos indica que $\mathsf{A}_{4}$ \textsc{no es nilpotente}.

El subgrupo de Fitting es inmediato puesto que el \'unico subgrupo normal es nilpotente, as{\'{\i}} que $\mathsf{F}\left(\mathsf{A}_{4}\right)=\left\langle a,b \right\rangle$. El subgrupo de Frattini es el trivial $\Phi_{\mathsf{A}_{4}}=\left\lbrace 1 \right\rbrace$.

Su grafo de Hasse de los subgrupos es :

\begin{center}
    \begin{tikzpicture}
    % First, locate each of the nodes and name them
    \node (a-c) at (-6,12)
    {$\mathsf{A}_{4}$};
    \node (a-b) at (-10,9)
    {$\left\langle{a,b\vert a^2=b^2,ab=ba}\right\rangle$};
    \node (c) at (-2,6)
    {$\left\langle{c\vert c^3}\right\rangle$};
    \node (bc) at (-4,6)
    {$\left\langle{bc\vert (bc)^3}\right\rangle$};
    \node (abc) at (-6,6)
	{$\left\langle{abc\vert (abc)^3}\right\rangle$};
    \node (ac) at (-8,6)
    {$\left\langle{ac\vert(ac)^3}\right\rangle$};
    \node (b) at (-12,3)
    {$\left\langle{b\vert b^2}\right\rangle$};
    \node (ab) at (-10,3)
    {$\left\langle{ab\vert (ab)^2}\right\rangle$};
    \node (a) at (-8,3)
    {$\left\langle{a\vert a^2}\right\rangle$};
	\node (1) at (-6,0)
    {$\left\langle{1}\right\rangle$};
   % Now draw the lines:
    \draw [black, thick] (a-c) -- (a-b);
    \draw [black, thick] (a-c) -- (c);
    \draw [black, thick] (a-c) -- (ac);
    \draw [black, thick] (a-c) -- (abc);
    \draw [black, thick] (a-c) -- (bc);
    \draw [black, thick] (a-b) -- (a);
    \draw [black, thick] (a-b) -- (b);
    \draw [black, thick] (a-b) -- (ab);
    \draw [black, thick] (c) -- (1);
    \draw [black, thick] (ac) -- (1);
    \draw [black, thick] (abc) -- (1);
    \draw [black, thick] (bc) -- (1);
    \draw [black, thick] (a) -- (1);
    \draw [black, thick] (b) -- (1);
    \draw [black, thick] (ab) -- (1);
    \end{tikzpicture}
\end{center}

Nos queda la siguiente tabla resumen del orden $12$ de grupos clasificados salvo isomorfismos:

\begin{table}[h!]
    \begin{center}
        \caption{Grupos de orden $12$}
        \label{tab:GruposOrden12}
        \begin{tabular}{c|c|c|c|c|c|c} % <-- Alignments: 1st column left, 2nd middle and 3rd right, with vertical lines in between
            \textbf{} & \tiny\textbf{n\'umero} &\tiny\textbf{numeraci\'on} & \tiny\textbf{nombre} & \textbf{} & {} & {} \\
            \textbf{} & \tiny\textbf{de} &\tiny\textbf{dentro} & \tiny\textbf{del} & \tiny\textbf{} & {} & {} \\
            \tiny\textbf{orden} & \tiny\textbf{grupos} & \tiny\textbf{del} & \tiny\textbf{grupo} & \tiny\textbf{abeliano}  & \tiny\textbf{nilpotente} & \tiny\textbf{resoluble} \\
            \textbf{} & \tiny\textbf{no} & \tiny\textbf{orden} & \textbf{} & \textbf{} & {} & {} \\
            \textbf{} & \tiny\textbf{isomorfos} &\textbf{} & \textbf{} & \textbf{} & {} & {} \\
            \hline
            12 & 5 & 1 & ${{\mathsf{C}}_{12}}$ & Si & Si & Si \\
            \hline
            12 & 5 & 2  & ${{\textsf{C}}_{6}}             \times             {{\mathsf{C}}_{12}}$ & Si & Si & Si \\
            \hline
            12 & 5 & 3 & ${\textsf{C}}_{3}\rtimes{\textsf{C}}_{4}             \cong             \mathsf{Dic}_{3}$& No & No & Si \\
            \hline
            12 & 5 & 4 & ${{\textsf{Dih}}_{6}}             \cong             {{\mathsf{Dih}_{3}}             \times             }{\cyclicgr{2}}             \cong             {{\cyclicgr{3}}\rtimes}{\cyclicgr{4}}$ & No & No & Si \\
            \hline
            12 & 5 & 5 & ${\left({{\textsf{C}}_{2}}             \times             {{\textsf{C}}_{2}}\right)}\bowtie {{\textsf{C}}_{3}}             \cong             {{\mathsf{A}_{4}}}$ & No & No & Si \\
            \hline
        \end{tabular}
    \end{center}
\end{table}

\chapter{Grupos de orden pq d\'onde p,q son primos mayores que 2 y p<q}

Estos grupos son grupos libres de cuadrados y se podr{\'{\i}}an ver como una generalizadi\'on de los grupos de orden $2p$.

En cuanto a grupos abelianos, solo hay uno, ya que $mcd(p,q)=1$ que indica, por ejemplo recordando el teorema de clasificaci\'on de los grupos abelianos finitos, que solo es posible $\cyclicgr{pq}             \cong             \cyclicgr{p}             \times             \cyclicgr{q}$. Es lo mismo que ocurr{\'{\i}}a con los grupos abelianos en los \'ordenes $2p$.

En cuanto a los grupos no abelianos tenemos realmente pocas posibilidades. De un lado, por el teorema de Cauchy, tendremos elementos de orden $p$ y de orden $q$, por lo tanto, se invierte el teorema de Lagrange. Adem\'as, dado que $p<q$, suponiendo que $a             \in \mathbf{G}$ con $a^p=1$ y $b             \in \mathbf{G}$ con $b^q=1$, $\left[\mathbf{G}:\left\langle b\right\rangle\right]=q$ que es el menor divisor del orden del grupo, tendremos que necesariamente $\left\langle{b}\right\rangle             \vartriangleleft             \mathbf{G}$. Esto nos deja solo dos posibilidades: si $\left\langle{a}\right\rangle             \vartriangleleft             \mathbf{G}$, como ${\left\langle{a}\right\rangle}             \cap {\left\langle{b}\right\rangle}={\left\lbrace{1}\right\rbrace}$ y  ${\left\langle{a,b}\right\rangle}={\mathbf{G}}$, entonces $\mathbf{G}$ es el grupo abeliano isomorfo a ${\left\langle{a}\right\rangle}             \times             {\left\langle{b}\right\rangle}$.

As{\'{\i}} que el grupo no abeliano, en caso de existir, ha de ser con los subgrupos de orden $q$ normales y los de orden $p$ han de ser no normales. El caso que ambos sean no normales es imposible ya que hemos visto que los subgrupos de orden $q$ son necesariamente normales.

Esto sugiere que puede haber algún homomorfismo de $\cyclicgr{p}$ en $\mathtt{Aut}\left(\cyclicgr{q}\right)$ que resulte en una acci\'on que nos de un producto semidirecto con el que podamos construir el grupo $\mathbf{G}$ no abeliano.

Para esto $\mathsf{o}(a) \shortmid \mathsf{o}(\mathtt{Aut}(\langle b \rangle))$. Sabemos que esto significa que $p \shortmid q-1$. Esto es lo que pasaba con los grupos de orden $2p$, d\'onde al ser $p$ impar $p-1$ es m\'ultiplo de $2$, y de ah{\'{\i}} vienen los grupos dih\'edricos (en esos casos es claro, en los que $n$ es cualquiera, par o impar hay que verlo mejor).

Por el primer teorema de Sylow, $n_p$ es el n\'umero de subgrupos conjugados de $p-grupo$ de Sylow, y ha de ser $n_p \shortmid q$ y $n_p {             \cong             }_p 1$. Entonces $n_p$ ha de ser $1$ o $q$. Si $n_p = 1$ tendr{\'{\i}}amos que habr{\'{\i}}an un \'unico grupo de orden $p$, y por lo tanto ser{\'{\i}}a caracter{\'{\i}}stico, y de aqu{\'{\i}} normal y abeliano. Luego estamos buscando solo el caso $n_p = q$ y $q {             \cong             }_p 1$, de d\'onde s\'olo queda $p \shortmid q-1$. Como esta condici\'on se puede solventar con el producto semidirecto, hemos encontrado la \'unica posibilidad de grupo no abeliano.

Veamos como ser{\'{\i}}a:
\[
\psi: {\cyclicgr{p}}              \longrightarrow              {{\mathtt{Aut}}{\left({\cyclicgr{q}}\right)}}             \cong             {{\mathsf{C}}_{q-1}}
\]
Sea $pr=q-1$. Podemos tomar el homomorfismo ${\psi }_{a}\triangleq \langle b             \mapsto b^{r=\nicefrac{q-1}{p}}\rangle $.
Primero vemos el grupo de automorfismos del grupo c{\'{\i}}clico mayor:

\begin{flalign*}
&b            \in {{\mathsf{C}}_{q}} \; b^q=1\\
&{\mathtt{Aut}}{\left({{\mathsf{C}}_{q}}\right)}=
{\left\lbrace{\left\langle{b            \mapsto b}\right\rangle},\ldots,{\left\langle{b            \mapsto b^{q-1}}\right\rangle}\right\rbrace}
\end{flalign*}

Ahora los homomorfismos posibles que buscamos son :

\begin{flalign*}
&{{\mathsf{C}}_{q}}            \cong            {\left\langle{ b \shortmid b^q=1}\right\rangle}\\
&{{\mathsf{C}}_{p}}            \cong            {\left\langle{ a \shortmid a^p=1}\right\rangle}\\
& p \shortmid q-1\\
&{{\varphi }_{i}}:{{\mathsf{C}}_{p}}            \longrightarrow             {{\mathtt{Aut}}{\left({{\mathsf{C}}_{q}}\right)}}\quad i=0,1,2,\ldots,q-1\\
&{{\varphi }_{i}}{\left({b}\right)}\triangleq
\begin{cases}
&{\left\langle{b            \mapsto b}\right\rangle}            \Leftarrow             i=0\\
&{\left\langle{b            \mapsto b^2}\right\rangle}            \Leftarrow             i=1\\
&\qquad\ldots\ldots\\
&{\left\langle{b            \mapsto b^{q-1}}\right\rangle}            \Leftarrow             i=q-1\\
\end{cases}\\
\end{flalign*}

Pero de todos estos tenemos que escoger a lo sumo $p$ distintos. Si fijamos un $i=r$, veremos que se genera un homomorfismo ${\psi } : {{\mathsf{C}}_{p}}            \longrightarrow            {{\mathtt{Aut}}{\left({{\mathsf{C}}_{q}}\right)}}$ que debe ser no trivial y cumplir alguna otra condici\'on que enseguida veremos. La no trivialidad y otro requisito recaer\'a en la elecci\'on de $r$. Por ahora no tenemos en cuenta estos dos importantes detalles.

\[
  {\psi }{\left({b}\right)}\triangleq {b^r}
\]
\[
  {\psi }{\left({b^j}\right)}={ {\left({b^r}\right)}^j}=b^{rj}
\]

Es homomorfismo, isomorfismo de hecho, por ser $\left\langle b \right\rangle$ abeliano (por c{\'{\i}}clico).
Si, por ejemplo, $r=\frac{q-1}{p}$, se cumplir\'a:

\[
  {\psi }{\left({b^p}\right)}=b^{rp}=b^{q-1}=b^{-1}
\]

Introducimos ahora el par\'ametro debido a $a^k$:

\[
  {{\theta }_{a=a^1}}{\left({b^j}\right)} \triangleq b^{rj}
\]

A partir de aqu{\'{\i}}, para conformar una acci\'on, vemos para aclararnos el caso $k=2$ :

\[
  {{\theta }_{a^2}}{\left({b^j}\right)}=
  {{\theta }_{a}}            \circ {{\theta }_{a}}{\left({b^j}\right)}=
  {{\theta }_{a}}{\left({b^{rj}}\right)}={\left({b^{r^2 j}}\right)}
\]

Y as{\'{\i}},  en general:

\[
  {{\theta }_{a^k}}{\left({b^j}\right)} =
  {{\theta }_{a^{k-2}}}            \circ {{\theta }_{a^2}}{\left({b^j}\right)}=
  {{\theta }_{a^{k-2}}}{\left({b^{r^2j}}\right)}=
  {\left({b^{r^kj}}\right)}
\]

El caso especial $k=0$ se resuelve f\'acilmente en ${\theta }_{1}{\left({b^j}\right)}=b^j$, esto es, la identidad.

Tomando solo los {\'{\i}}ndices nos queda:

\[
  {{\vartheta }_{k}}{\left({j}\right)} \triangleq  {\left({r^kj}\right)}
\]

La acci\'on de un elemento sobre otro (tomando solo exponentes), no solo componentes:

\[
(i,j)\blacktriangleleft(k,m)\triangleq{(i,{{\vartheta }_{k}}{\left({j}\right)})}
\]

Deshaciendo definiciones:

\[
(i,j)\blacktriangleleft(k,m)=(i,r^k j);
\]

Nos queda ver los requisitos de $r$:
\begin{flalign*}
& \mbox{\textbf{Hip\'otesis : }} r           \equiv           _q 1\\
&  a b a^{-1} = b^r = b\\
& ab = ba\\
& \cyclicgr{q}\rtimes\cyclicgr{p}            \cong            \cyclicgr{q}           \times           \cyclicgr{p} \mbox{\textbf{  absurdo}}\\
& \boldsymbol{ r {\not           \equiv           }_{q} 1 } \mbox{\textbf{  conclusi\'on}}
\end{flalign*}

El otro requisito para $r$ es sencillo de ver y tiene que ver tambi\'en con la aritm\'etica modular:
\begin{flalign*}
&  a^{p} b a^{-p} = 1 b 1 = b = b^1\\
&  a^{p} b a^{-p} = b^{r^p} \\
& \boldsymbol{r^p {           \equiv           }_{q} 1}
\end{flalign*}

As{\'{\i}} que $r$ ser\'a tal que
\begin{flalign*}
&           \nexists            s           \in \mathds{N} \boldsymbol{ r-1 = sq }\\
&           \exists            t           \in \mathds{N} \boldsymbol{ r^p -1 = tq }
\end{flalign*}

Seguimos ya con todas las condiciones para que la acci\'on conforme un grupo no abeliano.

Ahora el producto con solo los exponentes queda (el signo $+$ hace referencia a la operaci\'on propia del producto directo):
\[
(i,j)            \cdot (k,m)\triangleq (i,r^kj)+(k,m)=(i+k,r^kj+m)
\]

\begin{center}
\begin{flalign*}
&\boldsymbol{(i,j)            \cdot (k,m) \qquad=\qquad (i+k,r^kj+m)}\\
&           \nexists            s           \in \mathds{N} \boldsymbol{ r-1 = sq }\\
&           \exists            t           \in \mathds{N} \boldsymbol{ r^p -1 = tq }
\end{flalign*}
\end{center}

(Cada componente vive en su respectivo $\nicefrac{\mathds{Z}}{n\mathds{Z}}$)



Despu\'es de implementar la operaci\'on en la herramienta que hemos construido, tenemos la siguiente tabla para el grupo no abeliano de orden $21$:



\smallskip\noindent
\resizebox{\linewidth}{!}{
\begin{tabular}{|c|c|c|c|c|c|c|c|c|c|c|c|c|c|c|c|c|c|c|c|c|}
\hline
$1$&$b$&$b^2$&$b^3$&$b^4$&$b^5$&$b^6$&$a$&$ab$&$ab^2$&$ab^3$&$ab^4$&$ab^5$&$ab^6$&$a^2$&$a^2b$&$a^2b^2$&$a^2b^3$&$a^2b^4$&$a^2b^5$&$a^2b^6$\\
\hline
$b$&$b^2$&$b^3$&$b^4$&$b^5$&$b^6$&$1$&$ab^2$&$ab^3$&$ab^4$&$ab^5$&$ab^6$&$a$&$ab$&$a^2b^4$&$a^2b^5$&$a^2b^6$&$a^2$&$a^2b$&$a^2b^2$&$a^2b^3$\\
\hline
$b^2$&$b^3$&$b^4$&$b^5$&$b^6$&$1$&$b$&$ab^4$&$ab^5$&$ab^6$&$a$&$ab$&$ab^2$&$ab^3$&$a^2b$&$a^2b^2$&$a^2b^3$&$a^2b^4$&$a^2b^5$&$a^2b^6$&$a^2$\\
\hline
$b^3$&$b^4$&$b^5$&$b^6$&$1$&$b$&$b^2$&$ab^6$&$a$&$ab$&$ab^2$&$ab^3$&$ab^4$&$ab^5$&$a^2b^5$&$a^2b^6$&$a^2$&$a^2b$&$a^2b^2$&$a^2b^3$&$a^2b^4$\\
\hline
$b^4$&$b^5$&$b^6$&$1$&$b$&$b^2$&$b^3$&$ab$&$ab^2$&$ab^3$&$ab^4$&$ab^5$&$ab^6$&$a$&$a^2b^2$&$a^2b^3$&$a^2b^4$&$a^2b^5$&$a^2b^6$&$a^2$&$a^2b$\\
\hline
$b^5$&$b^6$&$1$&$b$&$b^2$&$b^3$&$b^4$&$ab^3$&$ab^4$&$ab^5$&$ab^6$&$a$&$ab$&$ab^2$&$a^2b^6$&$a^2$&$a^2b$&$a^2b^2$&$a^2b^3$&$a^2b^4$&$a^2b^5$\\
\hline
$b^6$&$1$&$b$&$b^2$&$b^3$&$b^4$&$b^5$&$ab^5$&$ab^6$&$a$&$ab$&$ab^2$&$ab^3$&$ab^4$&$a^2b^3$&$a^2b^4$&$a^2b^5$&$a^2b^6$&$a^2$&$a^2b$&$a^2b^2$\\
\hline
$a$&$ab$&$ab^2$&$ab^3$&$ab^4$&$ab^5$&$ab^6$&$a^2$&$a^2b$&$a^2b^2$&$a^2b^3$&$a^2b^4$&$a^2b^5$&$a^2b^6$&$1$&$b$&$b^2$&$b^3$&$b^4$&$b^5$&$b^6$\\
\hline
$ab$&$ab^2$&$ab^3$&$ab^4$&$ab^5$&$ab^6$&$a$&$a^2b^2$&$a^2b^3$&$a^2b^4$&$a^2b^5$&$a^2b^6$&$a^2$&$a^2b$&$b^4$&$b^5$&$b^6$&$1$&$b$&$b^2$&$b^3$\\
\hline
$ab^2$&$ab^3$&$ab^4$&$ab^5$&$ab^6$&$a$&$ab$&$a^2b^4$&$a^2b^5$&$a^2b^6$&$a^2$&$a^2b$&$a^2b^2$&$a^2b^3$&$b$&$b^2$&$b^3$&$b^4$&$b^5$&$b^6$&$1$\\
\hline
$ab^3$&$ab^4$&$ab^5$&$ab^6$&$a$&$ab$&$ab^2$&$a^2b^6$&$a^2$&$a^2b$&$a^2b^2$&$a^2b^3$&$a^2b^4$&$a^2b^5$&$b^5$&$b^6$&$1$&$b$&$b^2$&$b^3$&$b^4$\\
\hline
$ab^4$&$ab^5$&$ab^6$&$a$&$ab$&$ab^2$&$ab^3$&$a^2b$&$a^2b^2$&$a^2b^3$&$a^2b^4$&$a^2b^5$&$a^2b^6$&$a^2$&$b^2$&$b^3$&$b^4$&$b^5$&$b^6$&$1$&$b$\\
\hline
$ab^5$&$ab^6$&$a$&$ab$&$ab^2$&$ab^3$&$ab^4$&$a^2b^3$&$a^2b^4$&$a^2b^5$&$a^2b^6$&$a^2$&$a^2b$&$a^2b^2$&$b^6$&$1$&$b$&$b^2$&$b^3$&$b^4$&$b^5$\\
\hline
$ab^6$&$a$&$ab$&$ab^2$&$ab^3$&$ab^4$&$ab^5$&$a^2b^5$&$a^2b^6$&$a^2$&$a^2b$&$a^2b^2$&$a^2b^3$&$a^2b^4$&$b^3$&$b^4$&$b^5$&$b^6$&$1$&$b$&$b^2$\\
\hline
$a^2$&$a^2b$&$a^2b^2$&$a^2b^3$&$a^2b^4$&$a^2b^5$&$a^2b^6$&$1$&$b$&$b^2$&$b^3$&$b^4$&$b^5$&$b^6$&$a$&$ab$&$ab^2$&$ab^3$&$ab^4$&$ab^5$&$ab^6$\\
\hline
$a^2b$&$a^2b^2$&$a^2b^3$&$a^2b^4$&$a^2b^5$&$a^2b^6$&$a^2$&$b^2$&$b^3$&$b^4$&$b^5$&$b^6$&$1$&$b$&$ab^4$&$ab^5$&$ab^6$&$a$&$ab$&$ab^2$&$ab^3$\\
\hline
$a^2b^2$&$a^2b^3$&$a^2b^4$&$a^2b^5$&$a^2b^6$&$a^2$&$a^2b$&$b^4$&$b^5$&$b^6$&$1$&$b$&$b^2$&$b^3$&$ab$&$ab^2$&$ab^3$&$ab^4$&$ab^5$&$ab^6$&$a$\\
\hline
$a^2b^3$&$a^2b^4$&$a^2b^5$&$a^2b^6$&$a^2$&$a^2b$&$a^2b^2$&$b^6$&$1$&$b$&$b^2$&$b^3$&$b^4$&$b^5$&$ab^5$&$ab^6$&$a$&$ab$&$ab^2$&$ab^3$&$ab^4$\\
\hline
$a^2b^4$&$a^2b^5$&$a^2b^6$&$a^2$&$a^2b$&$a^2b^2$&$a^2b^3$&$b$&$b^2$&$b^3$&$b^4$&$b^5$&$b^6$&$1$&$ab^2$&$ab^3$&$ab^4$&$ab^5$&$ab^6$&$a$&$ab$\\
\hline
$a^2b^5$&$a^2b^6$&$a^2$&$a^2b$&$a^2b^2$&$a^2b^3$&$a^2b^4$&$b^3$&$b^4$&$b^5$&$b^6$&$1$&$b$&$b^2$&$ab^6$&$a$&$ab$&$ab^2$&$ab^3$&$ab^4$&$ab^5$\\
\hline
$a^2b^6$&$a^2$&$a^2b$&$a^2b^2$&$a^2b^3$&$a^2b^4$&$a^2b^5$&$b^5$&$b^6$&$1$&$b$&$b^2$&$b^3$&$b^4$&$ab^3$&$ab^4$&$ab^5$&$ab^6$&$a$&$ab$&$ab^2$\\
\hline
\end{tabular}
}

De un vistazo a la tabla podemos ver que efectivamente, $\mathbf{G}$\textsc{ no es abeliano}.

El ret{\'{\i}}culo de subgrupos es:


\begin{flalign*}
&\mathcal{L}{\left({\mathbf{G}}\right)}=\left\lbrace\right.
{\left\lbrace{1}\right\rbrace} ,
{\left\lbrace{1,g,g^2}\right\rbrace} ,
{\left\lbrace{1,gf,g^2f^3}\right\rbrace} ,
{\left\lbrace{1,gf^2,g^2f^6}\right\rbrace} ,
{\left\lbrace{1,gf^3,g^2f^2}\right\rbrace} ,\\
&{\left\lbrace{1,gf^4,g^2f^5}\right\rbrace} ,
{\left\lbrace{1,gf^5,g^2f}\right\rbrace} ,
{\left\lbrace{1,gf^6,g^2f^4}\right\rbrace} ,
{\left\lbrace{1,f,f^2,f^3,f^4,f^5,f^6}\right\rbrace} ,\\
&\left\lbrace\right. 1,f,f^2,f^3,f^4,f^5,f^6,g,gf,gf^2,gf^3,gf^4,gf^5,gf^6,\\
&g^2,g^2f,g^2f^2,g^2f^3,g^2f^4,g^2f^5,g^2f^6\left.\right\rbrace
\left.\right\rbrace
\end{flalign*}

Una f\'ormula expl{\'{\i}}cita para el inverso ser{\'{\i}}a:

\begin{flalign*}
&(i,j)            \cdot (k,m) = (i+k,r^kj+m)\\
&{(i,j)}^{-1} = (p-i,{\vartheta }_{p-i}(q-j))\\
&\boldsymbol{{(i,j)}^{-1} = (p-i,{r}^{p-i}(q-j))}\\
&\boldsymbol{(a^i b^j)^{-1} = a^{p-i}b^{{{r}^{p-i}}(q-j)}}
\end{flalign*}

Todos los subgrupos son c{\'{\i}}clicos excepto el completo. Hay un solo grupo de orden
$7$, $\left\lbrace{b}\right\rbrace$ que es normal y caracter{\'{\i}}stico. Hay $7$ subgrupos
conjugados (isomorfos) de orden $3$ y que son no-normales. El subgrupo de Frattini
es trivial (todos los subgrupos del ret{\'{\i}}culo son maximales, quitando el impropio y el trivial,
y su intersecci\'on es claramente el subgrupo trivial) y el subgrupo de Fitting coincide
con el \'unico subgrupo de orden $7$ (es el \'unico normal, y adem\'as es,
por c{\'{\i}}clico-abeliano, nilpotente).

Si un elemento pertenece al centro debe pertenecer a $\langle f\rangle$.
Pero ning\'un elemento $b^l$ conmuta con $a$ en el caso que tenemos a no ser que sea la unidad.
El centro es trivial.

En cuanto a los subgrupos derivado y L-derivado (para el caso $p=3$ y $q=7$):

\[
\mathsf{D^{(1)}}{\left({\mathsf{GDih}_{21}}\right)} = {\left\lbrace{1,b,b^2,b^3,b^4,b^5,b^6}\right\rbrace}
\]
\[
{{\mathsf{D}}^{2}}
{\entreparentesis{{\mathsf{GDih}}_{21}}}
=
{\set{1}}
\]
\[
\mathsf{L^{(2)}}{\left({\mathsf{GDih}_{21}}\right)} = {\left\lbrace{1,b,b^2,b^3,b^4,b^5,b^6}\right\rbrace}
\]


Por lo que $\mathbf{G}$ \textsc{ es resoluble} y \textsc{ no es nilpotente}.


El diagrama de Hasse, de forma general, siempre que $p$ no divida a ning\'un $n<q-1$, es:

\begin{center}
    \begin{tikzpicture}
    % First, locate each of the nodes and name them
    \node (g-f) at (-8,12)
    {${\mathsf{GDih}}_{p,q}
                      \cong
    {\left\langle{a,b\vert a^p=b^q=1,b^rg=gf}\right\rangle}$};
    \node (f) at (-14,8)
    {$\left\langle{b\vert b^q=1}\right\rangle$};
    \node (g) at (-10,6)
    {$\left\langle{a\vert a^p=1}\right\rangle$};
    \node (gfim1) at (-8,6)
    {$\ldots$};
    \node (gfi) at (-6,6)
    {$\left\langle{ab^i\vert (ab^i)^p=1}\right\rangle$};
    \node (gfip1) at (-4,6)
    {$\ldots$};
    \node (gfqm1) at (-2,6)
    {$\left\langle{ab^{q-1}\vert {(ab^{q-1})}^p=1}\right\rangle$};
    \node (1) at (-8,0)
    {$\left\lbrace{1}\right\rbrace$};
   % Now draw the lines:
    \draw [black, thick] (g-f) -- (f);
    \draw [black, thick] (g-f) -- (g);
    \draw [black, thick] (g-f) -- (gfi);
    \draw [black, thick] (g-f) -- (gfqm1);
    \draw [black, thick] (f) -- (1);
    \draw [black, thick] (g) -- (1);
    \draw [black, thick] (gfi) -- (1);
    \draw [black, thick] (gfqm1) -- (1);
    \end{tikzpicture}
\end{center}

Para los \'ordenes $pq$ nos queda la siguiente tabla de grupos distintos salvo isomorfismos:

\begin{table}[h!]
    \begin{center}
        \caption{Grupos de orden $pq$}
        \label{tab:GruposOrdenPQ}
        \begin{tabular}{c|c|c|c|c|c|c} % <-- Alignments: 1st column left, 2nd middle and 3rd right, with vertical lines in between
            \textbf{} & \tiny\textbf{n\'umero} &\tiny\textbf{numeraci\'on} & \tiny\textbf{nombre} & \textbf{} & {} & {} \\
            \textbf{} & \tiny\textbf{de} &\tiny\textbf{dentro} & \tiny\textbf{del} & \tiny\textbf{} & {} & {} \\
            \tiny\textbf{orden} & \tiny\textbf{grupos} & \tiny\textbf{del} & \tiny\textbf{grupo} & \tiny\textbf{abeliano}  &
            \tiny\textbf{nilpotente} & \tiny\textbf{resoluble} \\
            \textbf{} & \tiny\textbf{no} & \tiny\textbf{orden} & \textbf{} & \textbf{} & {} & {} \\
            \textbf{} & \tiny\textbf{isomorfos} &\textbf{} & \textbf{} & \textbf{} & {} & {} \\
            \hline
            15 & 1 & 1 & ${{\mathsf{C}}_{15}}$ & Si & Si & Si \\
            \hline
            21 & 2 & 1  & ${{\mathsf{C}}_{21}}$ & Si & Si & Si \\
            \hline
            21 & 2 & 2  & ${{\cyclicgr{7}}\rtimes{\cyclicgr{3}}}$ & No & No & Si \\
            \hline
            33 & 1 & 1 & ${{\mathsf{C}}_{23}}$ & Si & Si & Si \\
            \hline
            35 & 1 & 1 & ${{\mathsf{C}}_{35}}$ & Si & Si & Si \\
            \hline
            39 & 2 & 1 & ${{\mathsf{C}}_{39}}$ & Si & Si & Si \\
            \hline
            39 & 2 & 2 & ${{\mathsf{C}}_{13}}\rtimes{{\mathsf{C}}_{13}}$ & No & No & Si \\
            \hline
            51 & 1 & 1 & ${{\mathsf{C}}_{41}}$ & Si & Si & Si \\
            \hline
            55 & 2 & 1 & ${{\mathsf{C}}_{55}}$ & Si & Si & Si \\
            \hline
            55 & 2 & 2 & ${{\mathsf{C}}_{11}}\rtimes{{\mathsf{C}}_{5}}$ & No & No & Si \\
            \hline
            57 & 2 & 1 & ${{\mathsf{C}}_{57}}$ & Si & Si & Si \\
            \hline
            57 & 2 & 2 & ${{\mathsf{C}}_{19}}\rtimes{{\mathsf{C}}_{3}}$ & No & No & Si \\
            \hline
        \end{tabular}
    \end{center}
\end{table}

Addenda: grupos $\cyclicgr{13}\rtimes\cyclicgr{3}$ y $\cyclicgr{11}\rtimes\cyclicgr{5}$.

\chapter{Los 2-grupos de orden 16}

Los grupos abelianos no isomorfos de este orden son f\'aciles de enumerar gracias al teorema de clasificaci\'on de los grupos abelianos finitos:

\begin{flalign*}
&{\grG}_{1}\mathdef\cyclicgr{16}\\
&{\grG}_{2}\mathdef\cyclicgr{2}         \times         \cyclicgr{8}\\
&{\grG}_{3}\mathdef\cyclicgr{4}         \times         \cyclicgr{2}\\
&{\grG}_{4}\mathdef\cyclicgr{2}         \times         \cyclicgr{2}         \times         \cyclicgr{4}\\
&{\grG}_{5}\mathdef\cyclicgr{2}         \times         \cyclicgr{2}         \times         \cyclicgr{2}         \times         \cyclicgr{2}
\end{flalign*}

Para hacernos una primera idea en cuanto a los no abelianos, tenemos, como en un primer intento,
$\mathsf{Q}_{8}            \times            \cyclicgr{2}$, $\mathsf{Dih}_{4}            \times            \cyclicgr{2}$.
$\mathsf{Dih}_{8}$,  $\mathsf{Dic}_{8}$ (es seguro que estos primeros no son isomorfos),
tenemos posibles acciones (que dar\'an como soluciones productos semidirectos posiblemente
distintos a los anteriores) como son $\cyclicgr{8}\rtimes\cyclicgr{2}$ (de estos
veremos tres productos semidirectos no isomorfos), $\cyclicgr{4}\rtimes\cyclicgr{4}$
y tantas combinaciones de grupos de orden $2$, $4$ y $8$ como se nos ocurran,
junto a una variedad de formas de producirlos. ?`Ser\'an suficientes estos productos
para generar todos los
grupos no-isomorfos de orden diecis\'eis? Y de los grupos que saquemos por vias distintas
?`ser\'an todos no-isomorfos o al contrario los habr\'a isomorfos?

Por lo pronto ya vimos que el grupo cuaterni\'on no se pod{\'{\i}}a generar mediante
producto directo o semidirecto (si quiera con un producto ZS que exigir{\'{\i}}a que
los grupos con los que formamos el grupo fuesen no normales, cosa que en el grupo
cuaterni\'on no se da).

As{\'{\i}} que en este cap{\'{\i}}tulo vamos a añadir a nuestro arsenal algo de teor{\'\i}a
de la exstensi\'on de grupos. Concretamente la parte m\'as sencilla: la teor{\'\i}a de
extensi\'on por grupos c{\'\i}clicos.

Para esto partimos de un grupo $\grN$, que va a ser normal al grupo extendido. Baer hizo
contribuciones a extensiones no normales, en las que estoy especialmente interesado, pues
quiz\'as nos den la forma de construir también grupos finitos simples, cosa que yo hab{\'\i}a
estado buscando en el producto \texttt{ZS}. A decir verdad, me dan ganas de rehacer todo lo hecho
hasta ahora para ponerlo en la manera que vamos ahora a presentar la clasificaci\'on de los grupos
de orden $16$. Pero ya no queda tiempo. A\'un as{\'\i} he añadido algunos m\'odulos al programa,
concretamente un m\'odulo para establecer los automorfismos de un grupo
$\mathtt{Aut\langle num isomorph\rangle(y         \in \grN)}$,
aunque no como grupos ellos mismos por ahora, sino m\'as bien con funciones objeto bien definidas.
Adem\'as he añadido un m\'odulo que se encarga del algoritmo de multiplicaci\'on inducida por la
extensi\'on, de forma que si queremos hacer una extensi\'on, solo tenemos que llamar al
\emph{constructor del tipo} \texttt{binop\_t}, y pasarle el elemento $v         \in \grN$ que ha
de quedar fijo por el isomorfismo elegido. Las comprobaciones que el isomorfismo
$\tau         \in \mathtt{Aut\left({\grN}\right)}$ elegido es tal que, primero es realmente isomorfismo,
segundo, que $\tau \left({v}\right)=v$ y por \'ultimo, que si estamos extendiendo $\grN$ por el
grupo c{\'\i}clico $\cyclicgr{n}$, entonces ${\tau }^{n}=\mathtt{id}$ no est\'an incorporadas a
la herramienta, pero es claro que son f\'aciles de implementar.

A decir verdad, ha sido tal la revelaci\'on que me ha supuesto conocer esta teor{\'\i}a que me
he decidido a llamar a estos m\'etodos para el orden $16$ como m\'etodos ${\textsf{MW}}_{n}$
d\'onde $n$ es el {\'\i}ndice de Marcel Wilde en su art{\'\i}culo "La clasificaci\'on de los grupos
de orden 16 hecha f\'acil".

Antes de comenzar, veremos el concepto de \emph{tipo de extensi\'on}. Tenemos un grupo $\grN$ y
queremos que por medio de un grupo, en principio abeliano (y en nuestro caso c{\'\i}clico), $\grA$,
construir una extensi\'on $\grE$, tal que como queremos que $\grN \normsubgr \grE$,
$\grA=\grQ=\nicefrac{\grE}{\grN}$. Además queremos que se dé la siguiente secuencia exacta corta:
\par
\begin{defn} Decimos que un grupo $\grE$ es una extensi\'on de $\grN$ por $\grA$ (abeliano)
si y solo si, existen $\alpha         \in {\mathtt{Hom}}_{Gr}\left({\grN,\grE}\right)$, inyectiva,
y $\pi         \in {\mathtt{Hom}}_{Gr}\left({\grE,\grA}\right)$, sobreyectiva, tales que:
\begin{flalign*}
& \qquad 1          \rightarrow \grN \overset{\alpha }{         \hookrightarrow         } \grE \overset{\pi }{\twoheadrightarrow}
\nicefrac{\grE}{\grN}          \rightarrow 1 \\
& \qquad\qquad ker\left(\pi \right)         \subseteq imag\left(\alpha \right)
\end{flalign*}
\end{defn}
\par
As{\'\i} que en todo este c\'alculo lo \'unico que est\'a a priori indeterminado es $\grE$.
Las condiciones para que esto ocurra
cuando el grupo $\grQ         \cong         \cyclicgr{n}$, son las mencionadas sobre el automorfismo $\tau$ de
${\mathtt{Aut}}{\left(\grN\right)}$ y $v         \in \grN$.

\par
\begin{defn} Un \emph{tipo de extensi\'on c{\'\i}clica de grupo} es una cuádrupla que notaremos por
$\left\langle{\grN\;grupo , n         \in \setN , \tau         \in \mathtt{Aut}\left(\grN\right)} ,
v         \in \grN\right\rangle$ tal que, ${\tau }^{n}={\mathtt{id}}_{\grN}$ y $\tau \left(v\right)=v$.
Diremos que el grupo resultante, construido sobre el conjunto subyacente $\grN         \times         \cyclicgr{n}$,
es una extensi\'on de $\grN$ por $\cyclicgr{n}$ al dotarlo de la siguiente multiplicaci\'on:
\begin{flalign*}
&x,y,v          \in \grN\\
&\tau          \in {\mathtt{Aut}}_{Gr}{\left({\grN}\right)}\\
&\left|\cyclicgr{n}\right| = n\\
&{\tau }^n = {\mathtt{id}}{∣}_{\grN}\\
&\tau \left(v\right)=v\\
&\langle a \rangle = \cyclicgr{n}\\
&{\left({x,a^i}\right)}\star{\left({y,a^j}\right)}=
\begin{cases}
\left( {x         \cdot {{\tau }^{i}}{\left(y\right)}}, {a^{i+j}}\right)    	&\;         \Leftarrow         \;  i+j<n\\
\left( {x         \cdot {{\tau }^{i}}{\left(y\right)}         \cdot {v}}, {a^{(i+j)-n}}	\right)		&\;         \Leftarrow         \;  i+j\geq n
\end{cases}
\end{flalign*}

\end{defn}
\par

Veamos algunos casos en los que degenera la multiplicación anteriormente definida. Si $v         \in \grN$ es $v=e$,
$\left\langle{\grN\;grupo , n         \in \setN , \tau         \in \mathtt{Aut}\left(\grN\right)} ,
e         \in \grN\right\rangle$, nos queda necesariamente $\tau \left(e\right)=e$, y la ecuaci\'on de la deifinici\'on decae a:

\begin{flalign*}
&{\left({x,a^i}\right)}\star{\left({y,a^j}\right)}=
\left( {x         \cdot {{\tau }^{i}}{\left(y\right)}}, {a^{(i+j)\mod n}}\right)
\end{flalign*}

Que no es mas que el producto semidirecto de $\grN {\rtimes}_{\tau } \cyclicgr{2}$. Si adem\'as $\tau = \mathtt{id}{∣}_{\grN}$,
podemos ver que no se trata de otra cosa que el producto directo. A\'un no siendo ninguno de los casos dichos,
puede que la extensión se pueda expresar de manera más sencilla.

Antes de comenzar, estableceremos primero algunas definiciones y los siguientes teoremas
como hechos sin demostrar: en
todo caso son sencillos de demostrar a partir de lo que ya sabemos.
\par
\begin{defn}[\textbf{Subgrupos $L$-derivados de orden $i=0,1,2\ldots$}]
  \phantom{    }Dado un grupo $\grG$, decimos que ${L^0}{\left({\grG}\right)}\mathdef {\grG}$,
  ${L^1}{\left({\grG}\right)}\mathdef \left\langle\left[{\grG,\grG}\right]\right\rangle$ y
  $\forall i         \in \setN\setminus\set{0,1}\;\;{L^{(i+1)}}{\left({\grG}\right)}\mathdef
  \left\langle\left[{{L^i}(\grG),\grG}\right]\right\rangle$. En general los llamaremos
  los subgrupos $L$-derivados.
\end{defn}
\par
\begin{defn}[\textbf{Serie central descendente}]
  \phantom{    }Dado un grupo $\grG$, su serie central ascendente ser\'a:
  \[
  \grG = {\mathsf{L}^0}{\left({\grG}\right)} \triangleright {\mathsf{L}^1}{\left({\grG}\right)}={\grG}^{\prime}
  \triangleright {\mathsf{L}^2}{\left({\grG}\right)}\triangleright \ldots
  \]
\end{defn}
\par
\begin{defn}[\textbf{Centros de orden $i=0,1,2\ldots$}]
  \phantom{    }Dado un grupo $\grG$, decimos que ${\mathsf{Z}_0}{\left({\grG}\right)}\mathdef {\trivialgr}$,
  ${\mathsf{Z}_1}{\left({\grG}\right)}\mathdef \centro{\grG}$ y
  $\forall i         \in \setN\setminus\set{0,1}\;\;{\mathsf{Z}_{(i+1)}}{\left({\grG}\right)}\subgreq{\grG}$ lo definiremos por
  la siguiente relaci\'on: $\centro{\nicefrac{\grG}{{\mathsf{Z}_{i}}{\left({\grG}\right)}}}=
  \nicefrac{{\mathsf{Z}_{(i+1)}}{\left({\grG}\right)}}{{\mathsf{Z}_{i}}{\left({\grG}\right)}}$
\end{defn}
\par
\begin{defn}[\textbf{Serie central ascendente}]
  \phantom{    }Dado un grupo $\grG$, su serie central ascendente ser\'a:
  \[
  \trivialgr=\set{1_{\grG}}={\mathsf{Z}_0}{\left({\grG}\right)}
  \normsubgr {\mathsf{Z}_1}{\left({\grG}\right)}={\mathsf{Z}}{\left({\grG}\right)}
  \normsubgr {\mathsf{Z}_2}{\left({\grG}\right)}\normsubgr \ldots
  \]
\end{defn}
\par
\begin{defn}[\textbf{Grupo nilpotente}]
  \phantom{    }Dado un grupo $\grG$, decimos que es \emph{nilpotente} si su serie central
  descendente tiene un n\'umero total de inclusiones (estrictas) finita, digamos $n$, que ser\'a
  la longitud de la serie, y el t\'ermino $L$-derivado de orden $n$ ser\'a $\trivialgr$.
  Del grupo $\grG$ diremos que su clase de nilpotencia es $\mathtt{nil}_n$.
\end{defn}
\par
\begin{thm}[\textbf{Condici\'on dual de ser nilpotente}]
  \phantom{    }Si un grupo $\grG$ es \emph{nilpotente} según la definición anterior,
  entonces su \emph{serie normal ascendente} tendr\'a un número finito de inclusiones,
  digamos $n$, y alcanzará $\grG$ en el término centro de orden $n$. Las longitudes de las
  dos series normales ser\'a la misma.

  Si un grupo $\grG$ tiene una \emph{serie central ascendente} que en un n\'umero de
  inclusiones finito alcanza $\grG$ como su centro de orden $n$, entonces $\grG$ es nilpotente
  de clase $\mathtt{nil}_n$.
\end{thm}
\par
\begin{thm}[\textbf{Si $\grG$ es un $p$-grupo, $\grG$ es nilpotente}]
Un $p$-grupo $\grG$ es siempre $nilpotente$, esto es, la \emph{serie central ascendente}
alcanza desde $\trivialgr$ el grupo completo $\grG$, a través del centro primero $\centro{\grG}$
y posteriores. Igualmente la \emph{serie central descendente} alcanza desde $\grG$ el grupo
m{\'{\i}}nimo $\trivialgr$, a trav\'es del grupo derivado ${\grG}^{\prime}$ y siguientes.
\end{thm}
\par
\begin{thm}
Dado un grupo $\grG$, es condici\'on suficiente para que sea \emph{abeliano}
 que $\nicefrac{\grG}{\centro{\grG}}         \cong         \cyclicgr{n}$ para alg\'un $n         \in \setN$.
\end{thm}
\par
\begin{thm}
Para un $p$-grupo $\grG$ de orden ${p^n}$, el orden del centro
$\left|\mathsf{Z}\left({\grG}\right)\right|         \in \set{p,p^2,\ldots,p^{n-2}}$
\end{thm}
\par

Como los grupos que no est\'an directamente clasificados son los no abelianos, su nilpotencia ser\'a de clase superior a $\mathtt{nil}_{1}$.
Por el teorema de Lagrange no podrá ser superior a $4$. Pero el caso $\mathtt{nil}_{4}$ quedar\'a excluido: en caso contrario,
el centro ser\'a de orden $2$ (el mínimo posible para un $2$-grupo), debemos tener un centro segundo de orden $4$,
y un centro tercero de orden $8$. A su vez hemos de tener un grupo derivado de orden $8$, un grupo
$L_2$-derivado de orden $4$, y un grupo $L_3$-derivado de orden $2$. Ahora tanto en la serie central
descendente como en la ascendente tendremos que los cocientes ser\'an c{\'\i}clicos de
orden primo ($p=2$ para nosotros). Sea $z_1          \in \centro{\grG}$, claramente $o(z_1)=2$.
Ahora $\centro{\grG}\subsetneqq\mathsf{Z_2}\entreparentesis{\grG}$, luego existe
${z_2}          \in {\mathsf{Z_2}\entreparentesis{\grG}}\setminus{\centro{\grG}}$,
y tenemos que ${z_2}^2         \in \centro{\grG}$.
Igualmente podemos escoger
$z_3          \in {\mathsf{Z_3}\entreparentesis{\grG}}\setminus{\mathsf{Z_2}\entreparentesis{\grG}}$
tal que ${z_3}^2         \in {\mathsf{Z_2}\entreparentesis{\grG}}$. Por \'ultimo, existir\'a
$z_4          \in {\grG}\setminus{\mathsf{Z_3}\entreparentesis{\grG}}$ tal que
${z_4}^2         \in {\mathsf{Z_3}\entreparentesis{\grG}}$. Luego pod{\'\i}amos haber elegido
$z_1={z_4}^8         \in \centro{\grG}$,
$z_2={z_4}^4         \in {\mathtt{Z_2}{\grG}}$, $z_3={z_4}^2         \in {\mathtt{Z_3}{\grG}}$. Pero entonces
${\mathtt{Z_3}{\grG}}         \subseteq \generadopor{z_4}$, luego por Lagrange,
$o(\generadopor{z_4})=16$ y $\grG         \cong         \cyclicgr{16}$, contra lo supuesto
(que $\grG$ no es abeliano). Solo queda que un grupo de orden $16$ no puede ser de clase
$\mathtt{nil}_4$. Este resultado es f\'acil de generalizar, el mismo razonamiento vale para
cualquier $p$-grupo con exponente $\nu         \in \setN$.
\par
\begin{prop}
Si ${\left|{\grG}\right|}=p^\nu$, y $\grG          \neq          \centro{\grG}$, $\grG$ es $\mathtt{nil}_2$ o
$\mathtt{nil}_3$ o $\ldots$ o $\mathtt{nil}_{\nu-1}$.
\end{prop}
\par
\begin{cor}
Si ${\left|{\grG}\right|}=2^4$, y $\grG          \neq          \centro{\grG}$, $\grG$ es $\mathtt{nil}_2$ o
$\mathtt{nil}_3$.
\end{cor}
\par
\begin{prop}
Si ${\left|{\grG}\right|}=2^4$, y $\grG          \neq          \centro{\grG}$, $\left|\centro{\grG}\right|         \in \set{2,4}$.
\par
Si $\left|\centro{\grG}\right|=2$, entonces $\centro{\grG}         \cong         \cyclicgr{2}$.
Tambi\'en tenemos las siguientes posibilidades, $\nicefrac{\grG}{\centro{\grG}}         \cong         \mathsf{Dih_4}$,
$\nicefrac{\grG}{\centro{\grG}}         \cong         \mathsf{Q_2}$, $\nicefrac{\grG}{\centro{\grG}}         \cong         {\mathsf{C_4}         \times         \mathsf{C_2}}$
o por \'ultimo, $\nicefrac{\grG}{\centro{\grG}}         \cong         {\mathsf{C_2}         \times         \mathsf{C_2}         \times         \mathsf{C_2}}$.
\par
Si $\left|\centro{\grG}\right|=4$, entonces $\centro{\grG}         \cong         \cyclicgr{4}$ o
$\centro{\grG}         \cong         \cyclicgr{2}         \times         \cyclicgr{2}$. A su vez, esto nos da que
$\nicefrac{\grG}{\centro{\grG}}         \cong         \cyclicgr{2}         \times         \cyclicgr{2}$.
\end{prop}
\par
En el art{\'\i}culo \cite{Wilde2005} se establece una lista de extensiones que
ser\'an las que nos generen los grupos de orden $16$. Estas extensiones no son las extensiones que sugieren
la anterior proposici\'on, pues estas no son necesariamente extensiones c{\'\i}clicas, ni siquiera abelianas
en algunos casos.
Utilizaremos las extensiones del art{\'\i}culo por ser de gran sencillez, todas ser\'an extensiones de un
grupo que ser\'a el normal de la extensi\'on, de orden $8$, por un grupo cociente isomorfo a $\cyclicgr{2}$,
que facilita enormemente la elecci\'on del homomorfismo-acci\'on.
Hay que decir que a partir de aqu\'\i tenemos que los casos numerados son, necesariamente, no isomorfos,
pero no queda establecido que est\'en un{\'\i}vocamente  definidos los grupos numerados con las condiciones
impuestas, o, dicho de otro modo, puede haber m\'as de un grupo no isomorfo en cada condici\'on numerada.

Antes de enumerar las extensiones, para concordar en la terminología expondr\'e los distintos grupos de
automorfismos que utilizaremos.

Los automorfismos de $\cyclicgr{4}$ son:
\begin{flalign*}
&         \exists          x          \in \cyclicgr{4} \;:\; \cyclicgr{4} = \generadopor{x}\\
&{\varphi }_{1} = \generadopor{x          \mapsto x}\\
&{\varphi }_{2} = \generadopor{x          \mapsto x^3}\\
&{\mathsf{Aut}}_{\mathsf{Gr}}\entreparentesis{\cyclicgr{4}}         \cong         {\cyclicgr{2}}
\end{flalign*}

Los automorfismos de $\cyclicgr{8}$ son:
\begin{flalign*}
&         \exists          x          \in \cyclicgr{8} \;:\; \cyclicgr{8} \mathdef \generadopor{x}\\
&{\phi }_{1} \mathdef \generadopor{x          \mapsto x}\\
&{\phi }_{2} \mathdef \generadopor{x          \mapsto x^3}\\
&{\phi }_{3} \mathdef \generadopor{x          \mapsto x^5}\\
&{\phi }_{4} \mathdef \generadopor{x          \mapsto x^7}\\
&{\mathsf{Aut}}_{\mathsf{Gr}}\entreparentesis{\cyclicgr{8}}         \cong         \cyclicgr{2}         \times         \cyclicgr{2}={\grK}_{4}
\end{flalign*}

Los automorfismos de ${\grK}_{8}\mathdef{\cyclicgr{4}         \times         \cyclicgr{2}}$ son:
\begin{flalign*}
&         \exists          (x,y)          \in \cyclicgr{4}         \times         \cyclicgr{2} \;:\; \cyclicgr{4}=\generadopor{x}          \wedge          \cyclicgr{2}=\generadopor{y}\\
&\psi_1 \mathdef \generadopor{\begin{matrix}
{x          \mapsto x}\\{y          \mapsto y}
\end{matrix}}\\
&\psi_2 \mathdef \generadopor{\begin{matrix}
{x          \mapsto x^3y}\\{y          \mapsto x^2y}
\end{matrix}}\\
&\psi_3 \mathdef \generadopor{\begin{matrix}
{x          \mapsto x^3}\\{y          \mapsto y}
\end{matrix}}\\
&\psi_4 \mathdef \generadopor{\begin{matrix}
{x          \mapsto xy}\\{y          \mapsto x^2y}
\end{matrix}}\\
&\psi_5 \mathdef \generadopor{\begin{matrix}
{x          \mapsto xy}\\{y          \mapsto y}
\end{matrix}}\\
&\psi_6 \mathdef \generadopor{\begin{matrix}
{x          \mapsto x^3}\\{y          \mapsto x^2y}
\end{matrix}}\\
&\psi_7 \mathdef \generadopor{\begin{matrix}
{x          \mapsto x^3y}\\{y          \mapsto y}
\end{matrix}}\\
&\psi_8 \mathdef \generadopor{\begin{matrix}
{x          \mapsto x}\\{y          \mapsto x^2y}
\end{matrix}}\\
&{\mathsf{Aut}}_{\mathsf{Gr}}\entreparentesis{\cyclicgr{8}}         \cong         \cyclicgr{2}         \times         \cyclicgr{2}={\grK}_{4}
\end{flalign*}

Las extensiones son (lista de Wilde):
\begin{flalign*}
\MW{1}\mathdef&\entreangulos{\cyclicgr{8},2,\phi_1,\unity}         \cong         \pdcyclics{8}{2}\\
\MW{2}\mathdef&\entreangulos{\cyclicgr{8},2,\phi_2,\unity}         \cong         \sdihgr{8}{2}         \cong         \psdefcyclics{8}{2}{\phi_2}\\
\MW{3}\mathdef&\entreangulos{\cyclicgr{8},2,\phi_3,\unity}         \cong         \psdefcyclics{8}{2}{\phi_3}\\
\MW{4}\mathdef&\entreangulos{\cyclicgr{8},2,\phi_4,\unity}         \cong         \dihgr{8}         \cong         \psdefcyclics{8}{2}{\phi_4}\\
\MW{5}\mathdef&\entreangulos{\cyclicgr{8},2,\phi_4,x^4}         \cong         \quaterniongr{2}\;{}^{\star}\\
\MW{6}\mathdef&\entreangulos{\cyclicgr{8},2,\phi_1,x}         \cong         \cyclicgr{16}\;{}^{\star\,2}\\
\MW{7}\mathdef&\entreangulos{\cyclicgr{8},2,\psi_1,\unity}         \cong         \cyclicgr{4}       \times       {\cyclicgr{2}}^2\\
\MW{8}\mathdef&\entreangulos{\cyclicgr{8},2,\psi_3,\unity}
         \cong         {\psdef{\cyclicgr{4}       \times       \cyclicgr{2}}{\cyclicgr{2}}{\psi_3}}         \cong         \dihgr{4}       \times       \cyclicgr{2}\\
\MW{9}\mathdef&\entreangulos{\cyclicgr{8},2,\psi_5,\unity}
         \cong         {\psdef{\cyclicgr{4}       \times       \cyclicgr{2}}{\cyclicgr{2}}{\psi_5}}         \cong         \psdef{{\cyclicgr{2}}^2}{\cyclicgr{4}}{\psi_5}\\
\MW{10}\mathdef&\entreangulos{\cyclicgr{8},2,\psi_6,\unity}
         \cong         {\psdef{\cyclicgr{4}       \times       \cyclicgr{2}}{\cyclicgr{2}}{\psi_6}}         \cong         \psdef{\quaterniongr{1}}{\cyclicgr{2}}{\psi_6}\\
\MW{11}\mathdef&\entreangulos{\cyclicgr{8},2,\psi_3,x^2}
         \cong         {\psdef{\cyclicgr{4}       \times       \cyclicgr{2}}{\cyclicgr{2}}{\psi_3}}         \cong         \pd{\quaterniongr{1}}{\cyclicgr{2}}\;{}^{\star\,3}\\
\MW{12}\mathdef&\entreangulos{\cyclicgr{8},2,\psi_5,x^2}
         \cong         {\psdef{\cyclicgr{4}       \times       \cyclicgr{2}}{\cyclicgr{2}}{\psi_5}}         \cong         \psdefcyclics{4}{4}{\psi_5}\;{}^{\star\,4}\\
\MW{13}\mathdef&\entreangulos{\cyclicgr{8},2,\psi_1,y}
         \cong         {\psdef{\cyclicgr{4}       \times       \cyclicgr{2}}{\cyclicgr{2}}{\psi_1}}         \cong         \pdcyclics{4}{4}\;{}^{\star\,5}\\
\MW{14}\mathdef&
\entreangulos{{\cyclicgr{2}}^3,2,{\mathtt{id}}{{∣}_{{\cyclicgr{2}}^3}},\unity}
         \cong         {\pdquatrocyclics{2}{2}{2}{2}}
\end{flalign*}

Las ecuaciones de las extensiones marcadas con $\star$ mas numeraci\'on merecen un peque{\~n}o comentario aclaratorio,
puesto que no corresponden a productos directos ni semidirectos, y sin embargo obtienen formas de representaci\'on
notacional simples.

El caso de $\star$ primero (sin numerar), es claramente del grupo de los cuaterniones generalizados, que no se escinden
("split") en dos grupos, esto es, no existe un
$\zeta         \in \biyuno{\grcoc{\quaterniongr{2}}{\cyclicgr{8}},\quaterniongr{2}}$
tal que, $\zeta         \circ \pi \left(t\right)=t\;\;\forall t         \in \quaterniongr{2}$.


El caso $\star\,2$ es curioso porque el grupo obtenido es el m\'as simple de todos, el grupo c{\'\i}clico de $16$ elementos.
La extensi\'on solo ha ampliado el orden del grupo c{\'\i}clico de orden $8$, d\'andole un aditamento en la parte ampliada
que hace de prologanci\'on. Estos casos los veremos m\'as detenidamente al ir poniendo sus tablas y dem\'as informaci\'on.

Comenzamos por orden. El primer grupo ser\'a, ${\MW{1}}$.
\par
\smallskip\noindent
\resizebox{\linewidth}{!}{
\begin{tabular}{|c|c|c|c|c|c|c|c|c|c|c|c|c|c|c|c|}
\hline
$1$&$a$&$a^2$&$a^3$&$a^4$&$a^5$&$a^6$&$a^7$&$b$&$ab$&$a^2b$&$a^3b$&$a^4b$&$a^5b$&$a^6b$&$a^7b$\\
\hline
$a$&$a^2$&$a^3$&$a^4$&$a^5$&$a^6$&$a^7$&$1$&$ab$&$a^2b$&$a^3b$&$a^4b$&$a^5b$&$a^6b$&$a^7b$&$b$\\
\hline
$a^2$&$a^3$&$a^4$&$a^5$&$a^6$&$a^7$&$1$&$a$&$a^2b$&$a^3b$&$a^4b$&$a^5b$&$a^6b$&$a^7b$&$b$&$ab$\\
\hline
$a^3$&$a^4$&$a^5$&$a^6$&$a^7$&$1$&$a$&$a^2$&$a^3b$&$a^4b$&$a^5b$&$a^6b$&$a^7b$&$b$&$ab$&$a^2b$\\
\hline
$a^4$&$a^5$&$a^6$&$a^7$&$1$&$a$&$a^2$&$a^3$&$a^4b$&$a^5b$&$a^6b$&$a^7b$&$b$&$ab$&$a^2b$&$a^3b$\\
\hline
$a^5$&$a^6$&$a^7$&$1$&$a$&$a^2$&$a^3$&$a^4$&$a^5b$&$a^6b$&$a^7b$&$b$&$ab$&$a^2b$&$a^3b$&$a^4b$\\
\hline
$a^6$&$a^7$&$1$&$a$&$a^2$&$a^3$&$a^4$&$a^5$&$a^6b$&$a^7b$&$b$&$ab$&$a^2b$&$a^3b$&$a^4b$&$a^5b$\\
\hline
$a^7$&$1$&$a$&$a^2$&$a^3$&$a^4$&$a^5$&$a^6$&$a^7b$&$b$&$ab$&$a^2b$&$a^3b$&$a^4b$&$a^5b$&$a^6b$\\
\hline
$b$&$ab$&$a^2b$&$a^3b$&$a^4b$&$a^5b$&$a^6b$&$a^7b$&$1$&$a$&$a^2$&$a^3$&$a^4$&$a^5$&$a^6$&$a^7$\\
\hline
$ab$&$a^2b$&$a^3b$&$a^4b$&$a^5b$&$a^6b$&$a^7b$&$b$&$a$&$a^2$&$a^3$&$a^4$&$a^5$&$a^6$&$a^7$&$1$\\
\hline
$a^2b$&$a^3b$&$a^4b$&$a^5b$&$a^6b$&$a^7b$&$b$&$ab$&$a^2$&$a^3$&$a^4$&$a^5$&$a^6$&$a^7$&$1$&$a$\\
\hline
$a^3b$&$a^4b$&$a^5b$&$a^6b$&$a^7b$&$b$&$ab$&$a^2b$&$a^3$&$a^4$&$a^5$&$a^6$&$a^7$&$1$&$a$&$a^2$\\
\hline
$a^4b$&$a^5b$&$a^6b$&$a^7b$&$b$&$ab$&$a^2b$&$a^3b$&$a^4$&$a^5$&$a^6$&$a^7$&$1$&$a$&$a^2$&$a^3$\\
\hline
$a^5b$&$a^6b$&$a^7b$&$b$&$ab$&$a^2b$&$a^3b$&$a^4b$&$a^5$&$a^6$&$a^7$&$1$&$a$&$a^2$&$a^3$&$a^4$\\
\hline
$a^6b$&$a^7b$&$b$&$ab$&$a^2b$&$a^3b$&$a^4b$&$a^5b$&$a^6$&$a^7$&$1$&$a$&$a^2$&$a^3$&$a^4$&$a^5$\\
\hline
$a^7b$&$b$&$ab$&$a^2b$&$a^3b$&$a^4b$&$a^5b$&$a^6b$&$a^7$&$1$&$a$&$a^2$&$a^3$&$a^4$&$a^5$&$a^6$\\
\hline
\end{tabular}
}

Los elementos clasificados por su orden dan el invariante:

\begin{flalign*}
&\left[1,3,4,8,0\right]
\end{flalign*}

Teniendo en cuenta que ${\MW{1}}=\generadopor{a}\generadopor{b}       \cong       \pdcyclics{8}{2}$
ya que tenemos las siguientes relaciones
\begin{flalign*}
& o(a)=8 \ly o(b)=2 \ly ab=ba\\
&\generadopor{a}       \cap \generadopor{b}=\set{1}
\end{flalign*}

Podemos presentar este grupo con las relaciones ya dichas,
${\MW{1}}=\entreangulos{a,b∣a^2=b^2=1,ab=ba}       \cong        \pdcyclics{8}{2}$.

Despu\'es de mostrar el centro del grupo y su derivado, comenzamos por los subgrupos c{\'\i}clicos.

\[
\centro{\MW{1}} = {\MW{1}}         \cong        \pdcyclics{8}{2}
\]

Y el subgrupo derivado es

\[
\grderivado{\MW{1}}{1} = {\left\lbrace{1}\right\rbrace}          \cong          \cyclicgr{1}
\]


Los grupos c{\'\i}clicos son:
\begin{flalign*}
&\mathcal{L}_1\left({{\MW{1}}}\right)=\left\lbrace\right.\\
&\quad {\left\lbrace{1}\right\rbrace} ,
{\left\lbrace{1,a^4}\right\rbrace} ,
{\left\lbrace{1,b}\right\rbrace} ,
{\left\lbrace{1,a^4b}\right\rbrace} ,\\%
&\quad {\left\lbrace{1,a^2,a^4,a^6}\right\rbrace} ,
{\left\lbrace{1,a^4,a^2b,a^6b}\right\rbrace} ,\\%
&\quad {\left\lbrace{1,a,a^2,a^3,a^4,a^5,a^6,a^7}\right\rbrace} ,\\%
&\quad {\left\lbrace{1,a^2,a^4,a^6,ab,a^3b,a^5b,a^7b}\right\rbrace}\\
&\left.\right\rbrace\\
&\left[(1,1),c(3,2),c(2,4),c(2,8)\right]
\end{flalign*}


Ahora los subgrupos generados por dos elementos, siguiendo el mismo patr\'on anterior,
sacaremos un nuevo invariante (vac{\'\i}o en este caso):

\begin{flalign*}
&\mathcal{L}_{2}\left({\MW{1}}\right)=\left\lbrace\right.\\
&\quad {\left\lbrace{1,a^4,b,a^4b}\right\rbrace} ,\\
&\quad {\left\lbrace{1,a^2,a^4,a^6,b,a^2b,a^4b,a^6b}\right\rbrace} ,\\
&\quad {\MW{1}}
&\left.\right\rbrace\\
&\left[c(1,4),c(1,8),c(1,16)\right]
\end{flalign*}

El invariante que anotamos con lo visto hasta aqu{\'\i}:


\begin{flalign*}
&\left[1,3,4,8,0\right]\\
&\left[
\begin{array}{l}
\left[
(1,1),c(3,2),c(2,4),c(2,8)
\right]\\
\left[
c(1,4),c(1,8),c(1,16)
\right]\\
\left[\right]\\
\left[\right]
\end{array}
\right]\\
&\left[
\begin{array} {c|c|c}
\mathtt{nil}_1 & \trivialgr  & {\MW{1}} \\
{} & {\MW{1}}  & \trivialgr
\end{array}
\right]\\
&\left[
\begin{array}{l}
\left[\pdcyclics{8}{2}         \cong         \centro{\MW{1}}\supsetneqq \grderivado{\MW{1}}{1}         \cong         \cyclicgr{1}\right]\\
\left[\right]
\end{array}
\right]
\end{flalign*}


Nos falta ver cu\'ales sean los cocientes por el centro y por el derivado.


\begin{flalign*}
&\graut{{\MW{1}}}         \cong         \nicefrac{{\MW{1}}}{\centro{{\MW{1}}}}         \cong         {\trivialgr}\\
&{{\MW{1}}}^{\mathtt{ab}}=\nicefrac{{{\MW{1}}}}{{{\MW{1}}}^{\prime}}         \cong         {\pdquatrocyclics{2}{2}{2}{2}}
\end{flalign*}


El ret{\'\i}culo de subgrupos maximales (propios) lo da la herramienta desarrollada para el trabajo, al igual que el subgrupo de Frattini:


\begin{flalign*}
&\frattini{{\MW{1}}}={\left\lbrace{1,a^2,a^4,a^6}\right\rbrace}        \cong        \cyclicgr{4}
\end{flalign*}

En cuanto al ret{\'\i}culo de subgrupos normales y nilpotentes tambi\'en lo da la herramienta. El subgrupo de Fitting es:

\begin{flalign*}
&\fitting{{\MW{1}}}={\MW{1}}
\end{flalign*}

As{\'\i} tenemos que $\frattini{{\MW{1}}}=\set{1,a^2,a^4,a^6}$. A su vez el subgrupo de
Fitting es el completo ya que hay varios maximales normales y nilpotentes,
uno de orden $4$ ($\generadopor{a^2,b}$) y otro de orden $8$ ($\generadopor{a}$),
que necesariamente generar\'an el grupo completo.


\begin{flalign*}
&\left[1,3,4,8,0\right]\\
&\left[
\begin{array}{l}
\left[
(1,1),c(3,2),c(2,4),c(2,8)
\right]\\
\left[
c(1,4),c(1,8),c(1,16)
\right]\\
\left[\right]\\
\left[\right]
\end{array}
\right]\\
&\left[
\begin{array} {c|c|c}
\mathtt{nil}_1 & \trivialgr  & {\MW{1}} \\
{} & {\MW{1}}  & \trivialgr
\end{array}
\right]\\
&\left[
\begin{array}{l}
\left[\pdcyclics{8}{2}         \cong         \centro{\MW{1}}\supsetneqq{\grderivado{\MW{1}}{1}}         \cong         \cyclicgr{1}\right]\\
\left[\right]
\end{array}
\right]\\
&\left[
\begin{array}{l}
\left[
\graut{{\MW{1}}}         \cong         \grcoc{\MW{1}}{\centro{\MW{1}}}         \cong         \trivialgr,
{\grabelianizado{\MW{1}}}=\grcoc{\MW{1}}{{\MW{1}}^{\prime}}         \cong         \pdcyclics{8}{2}
\right]\\
\left[\right]
\end{array}
\right]\\
&\left[
\frattini{{\MW{1}}}        \cong        \cyclicgr{4},
\fitting{{\MW{1}}}={{\MW{1}}}
\right]
\end{flalign*}


El siguiente grupo ser\'a, $\MW{2}$.


\smallskip\noindent
\resizebox{\linewidth}{!}{
\begin{tabular}{|c|c|c|c|c|c|c|c|c|c|c|c|c|c|c|c|}
\hline
$1$&$a$&$a^2$&$a^3$&$a^4$&$a^5$&$a^6$&$a^7$&$b$&$ab$&$a^2b$&$a^3b$&$a^4b$&$a^5b$&$a^6b$&$a^7b$\\
\hline
$a$&$a^2$&$a^3$&$a^4$&$a^5$&$a^6$&$a^7$&$1$&$ab$&$a^2b$&$a^3b$&$a^4b$&$a^5b$&$a^6b$&$a^7b$&$b$\\
\hline
$a^2$&$a^3$&$a^4$&$a^5$&$a^6$&$a^7$&$1$&$a$&$a^2b$&$a^3b$&$a^4b$&$a^5b$&$a^6b$&$a^7b$&$b$&$ab$\\
\hline
$a^3$&$a^4$&$a^5$&$a^6$&$a^7$&$1$&$a$&$a^2$&$a^3b$&$a^4b$&$a^5b$&$a^6b$&$a^7b$&$b$&$ab$&$a^2b$\\
\hline
$a^4$&$a^5$&$a^6$&$a^7$&$1$&$a$&$a^2$&$a^3$&$a^4b$&$a^5b$&$a^6b$&$a^7b$&$b$&$ab$&$a^2b$&$a^3b$\\
\hline
$a^5$&$a^6$&$a^7$&$1$&$a$&$a^2$&$a^3$&$a^4$&$a^5b$&$a^6b$&$a^7b$&$b$&$ab$&$a^2b$&$a^3b$&$a^4b$\\
\hline
$a^6$&$a^7$&$1$&$a$&$a^2$&$a^3$&$a^4$&$a^5$&$a^6b$&$a^7b$&$b$&$ab$&$a^2b$&$a^3b$&$a^4b$&$a^5b$\\
\hline
$a^7$&$1$&$a$&$a^2$&$a^3$&$a^4$&$a^5$&$a^6$&$a^7b$&$b$&$ab$&$a^2b$&$a^3b$&$a^4b$&$a^5b$&$a^6b$\\
\hline
$b$&$a^3b$&$a^6b$&$ab$&$a^4b$&$a^7b$&$a^2b$&$a^5b$&$1$&$a^3$&$a^6$&$a$&$a^4$&$a^7$&$a^2$&$a^5$\\
\hline
$ab$&$a^4b$&$a^7b$&$a^2b$&$a^5b$&$b$&$a^3b$&$a^6b$&$a$&$a^4$&$a^7$&$a^2$&$a^5$&$1$&$a^3$&$a^6$\\
\hline
$a^2b$&$a^5b$&$b$&$a^3b$&$a^6b$&$ab$&$a^4b$&$a^7b$&$a^2$&$a^5$&$1$&$a^3$&$a^6$&$a$&$a^4$&$a^7$\\
\hline
$a^3b$&$a^6b$&$ab$&$a^4b$&$a^7b$&$a^2b$&$a^5b$&$b$&$a^3$&$a^6$&$a$&$a^4$&$a^7$&$a^2$&$a^5$&$1$\\
\hline
$a^4b$&$a^7b$&$a^2b$&$a^5b$&$b$&$a^3b$&$a^6b$&$ab$&$a^4$&$a^7$&$a^2$&$a^5$&$1$&$a^3$&$a^6$&$a$\\
\hline
$a^5b$&$b$&$a^3b$&$a^6b$&$ab$&$a^4b$&$a^7b$&$a^2b$&$a^5$&$1$&$a^3$&$a^6$&$a$&$a^4$&$a^7$&$a^2$\\
\hline
$a^6b$&$ab$&$a^4b$&$a^7b$&$a^2b$&$a^5b$&$b$&$a^3b$&$a^6$&$a$&$a^4$&$a^7$&$a^2$&$a^5$&$1$&$a^3$\\
\hline
$a^7b$&$a^2b$&$a^5b$&$b$&$a^3b$&$a^6b$&$ab$&$a^4b$&$a^7$&$a^2$&$a^5$&$1$&$a^3$&$a^6$&$a$&$a^4$\\
\hline
\end{tabular}
}


Este grupo ya no es abeliano.
Como en el caso anterior, el primer invariante ser\'a el n\'umero de elementos de cada orden:


\begin{flalign*}
&\left[1,5,6,4,0\right]
\end{flalign*}


El centro del grupo y el subgrupo derivado son:

\begin{flalign*}
&\centro{\MW{2}}={\left\lbrace{1,a^4}\right\rbrace}       \cong       \cyclicgr{2}\\
&\grderivado{\MW{2}}{1}={\left\lbrace{1,a^2,a^4,a^6}\right\rbrace}       \cong       \cyclicgr{4}\\
&\cyclicgr{2}       \cong       \centro{\MW{2}}\supsetneqq\grderivado{\MW{2}}{1}       \cong       \cyclicgr{4}
\end{flalign*}

En este grupo no todos los subgrupos son normales, y los que vienen a
continuaci\'on son los c{\'{\i}}clicos, por lo tanto abelianos:


\begin{flalign*}
&\mathcal{L}_1\left({\grG_{7}}\right)=\left\lbrace\right.\\
&\quad {\left\lbrace{1}\right\rbrace} ,
{\left\lbrace{1,a^4}\right\rbrace} ,%dc
{\left\lbrace{1,b}\right\rbrace} ,%
{\left\lbrace{1,a^2b}\right\rbrace} ,%
{\left\lbrace{1,a^4b}\right\rbrace} ,%
{\left\lbrace{1,a^6b}\right\rbrace} ,\\%
&\quad {\left\lbrace{1,a^2,a^4,a^6}\right\rbrace} ,%d
{\left\lbrace{1,a^4,ab,a^5b}\right\rbrace} ,%
{\left\lbrace{1,a^4,a^3b,a^7b}\right\rbrace} ,\\%
&\quad {\left\lbrace{1,a,a^2,a^3,a^4,a^5,a^6,a^7}\right\rbrace}\\%n
&\left.\right\rbrace\\
&\left[(1,1),dc(1,2),(4,2),d(1,4),(2,4),n(1,8)\right]
\end{flalign*}


Primero vemos que $ba^2=a^6b=a^{-2}b$. Luego $\generadopor{a^2,b}       \cong       \dihgr{4}$.
A continuci\'on $aba^2=a^7b=a^{-1}b=a^{-2}ab$. Luego $\generadopor{a^2,ab}       \cong       \dihgr{4}$.
Ahora los subgrupos generados por dos elementos, siguiendo el mismo patr\'on anterior,
sacaremos un nuevo invariante:


\begin{flalign*}
&\mathcal{L}_{2}\left({\MW{2}}\right)=\left\lbrace\right.\\
&\quad {\left\lbrace{1,a^4,b,a^4b}\right\rbrace} ,%a
{\left\lbrace{1,a^4,a^2b,a^6b}\right\rbrace} ,\\%a
&\quad {\left\lbrace{1,a^2,a^4,a^6,b,a^2b,a^4b,a^6b}\right\rbrace} ,%n
{\left\lbrace{1,a^2,a^4,a^6,ab,a^3b,a^5b,a^7b}\right\rbrace} ,\\%n
&\quad{\MW{2}}\\
&\left.\right\rbrace\\
&\left[a(2,4),n(2,\dihgr{4}),(1,16)\right]
\end{flalign*}


Serie central ascendente (comienza en el $\trivialgr$,
y en general no tiene porqu\'e alcanzar $\MW{2}$):
\begin{flalign*}
&{\trivialgr} \normsubgr
{\left\lbrace{1,a^4}\right\rbrace} \normsubgr
{\left\lbrace{1,a^2,a^4,a^6}\right\rbrace} \normsubgr
{\MW{2}}
\end{flalign*}

Serie central descendente (comienza en el $\MW{2}$,
y en general no tiene proqu\'e alcanzar el $\trivialgr$):
\begin{flalign*}
&{\MW{2}}\triangleright
{\left\lbrace{1,a^2,a^4,a^6}\right\rbrace}\triangleright
{\left\lbrace{1,a^4}\right\rbrace}\triangleright
{\trivialgr}
\end{flalign*}

El invariante que derivamos:
\begin{flalign*}
&\left[
\begin{array}{c|c|c|c|c}
\nil{3} &
\trivialgr &
\cyclicgr{2} &
\cyclicgr{4} &
\MW{2}\\
{} &
\MW{2} &
\cyclicgr{4} &
\cyclicgr{2} &
\trivialgr
\end{array}
\right]
\end{flalign*}

Podemos tambi\'en a\~nadir la informaci\'on sobre el $L$-derivado primero y segundo y el centro primero y segundo.

Hemos visto que $\centronum{\MW{2}}{2}=\grderivado{\MW{2}}{1}       \cong        \cyclicgr{2}$ y
que $\centro{\MW{2}}=\grLderivado{\MW{2}}{2}        \cong         \cyclicgr{4}$.

Ahora el invariante total queda:


\begin{flalign*}
&\left[1,5,6,4,0\right]\\
&\left[
	\begin{array}{l}
		\left[dc(1,2),(4,2),d(1,4),(2,4),n(1,8)\right]\\
		\left[a(2,4),n(2,\dihgr{4}),(1,16)\right]\\
		\left[\right]\\
		\left[\right]
	\end{array}
\right]\\
&\left[
	\begin{array}{c|c|c|c|c}
		\nil{3} & \trivialgr & \cyclicgr{2} & \cyclicgr{4} & \MW{2}\\
		{} & \MW{2} & \cyclicgr{4} & \cyclicgr{2} & \trivialgr
	\end{array}
\right]\\
&\left[
	\begin{array}{l}
		\cyclicgr{2}       \cong       \centro{\MW{2}}\subsetneqq\grderivado{\MW{2}}{1}       \cong         \cyclicgr{4}\\
		\cyclicgr{4}       \cong       \centronum{\MW{2}}{2}\supsetneqq\grLderivado{\MW{2}}{2}        \cong        \cyclicgr{2}
	\end{array}
\right]
\end{flalign*}


Podemos tambi\'en
calcular las clases de conjugaci\'on y ver los grupos cocientes (las clases de conjugaci\'on las da tambi\'en el programa).

\begin{flalign*}
&\left[
	\begin{array}{l}
		\dihgr{4}       \cong       \grcoc{\MW{2}}{\centro{\MW{2}}}       \cong       \grinn{\MW{2}}\\
		\pdcyclics{2}{2}       \cong       \grcoc{\MW{2}}{\grderivado{\MW{2}}{1}}=\grabelianizado{\MW{2}}
	\end{array}
\right]
\end{flalign*}

Adem\'as el grupo de Frattini y el Fitting son:
\begin{flalign*}
&\left[
	\begin{array}{l}
		\frattini{\MW{2}}={\left\lbrace{1,a^2,a^4,a^6}\right\rbrace}       \cong       \cyclicgr{4}\\
		\fitting{\MW{2}}=\MW{2}
	\end{array}
\right]
\end{flalign*}

Finalmente tenemos toda la informaci\'on recogida en la siguiente ecuaci\'on:


\begin{flalign*}
&\left[1,5,6,4,0\right]\\
&\left[
	\begin{array}{l}
		\left[dc(1,2),(4,2),d(1,4),(2,4),n(1,8)\right]\\
		\left[a(2,4),n(2,\dihgr{4}),(1,16)\right]\\
		\left[\right]\\
		\left[\right]
	\end{array}
\right]\\
&\left[
	\begin{array}{c|c|c|c|c}
		\nil{3} & \trivialgr & \cyclicgr{2} & \cyclicgr{4} & \MW{2}\\
		{} & \MW{2} & \cyclicgr{4} & \cyclicgr{2} & \trivialgr
	\end{array}
\right]\\
&\left[
	\begin{array}{l}
		\cyclicgr{2}       \cong       \centro{\MW{2}}\subsetneqq\grderivado{\MW{2}}{1}       \cong         \cyclicgr{4}\\
		\cyclicgr{4}       \cong       \centronum{\MW{2}}{2}\supsetneqq\grLderivado{\MW{2}}{2}        \cong        \cyclicgr{2}
	\end{array}
\right]\\
&\left[
	\begin{array}{l}
		\dihgr{4}       \cong       \grcoc{\MW{2}}{\centro{\MW{2}}}       \cong       \grinn{\MW{2}}\\
		\pdcyclics{2}{2}       \cong       \grcoc{\MW{2}}{\grderivado{\MW{2}}{1}}=\grabelianizado{\MW{2}}
	\end{array}
\right]\\
&\left[
	\begin{array}{l}
		\frattini{\MW{2}}={\left\lbrace{1,a^2,a^4,a^6}\right\rbrace}       \cong       \cyclicgr{4}\\
		\fitting{\MW{2}}=\MW{2}
	\end{array}
\right]
\end{flalign*}



El ret{\'\i}culo de subgrupos normales y nilpotentes es:

\begin{flalign*}
&\mathcal{L}_{norm,nil}\entreparentesis{\MW{2}}=\left\lbrace\right.\\
&\qquad\set{1},{\left\lbrace{1,a^4}\right\rbrace},
{\left\lbrace{1,b}\right\rbrace},
{\left\lbrace{1,a^2b}\right\rbrace},
{\left\lbrace{1,a^4b}\right\rbrace},
{\left\lbrace{1,a^6b}\right\rbrace},\\
&\qquad{\left\lbrace{1,a^2,a^4,a^6}\right\rbrace},
{\left\lbrace{1,a^4,ab,a^5b}\right\rbrace},
{\left\lbrace{1,a^4,a^3b,a^7b}\right\rbrace},\\
&\qquad{\left\lbrace{1,a^4,a^2b,a^6b}\right\rbrace},
{\left\lbrace{1,a^2,a^4,a^6,b,a^2b,a^4b,a^6b}\right\rbrace},\\
&\qquad{\left\lbrace{1,a^2,a^4,a^6,ab,a^3b,a^5b,a^7b}\right\rbrace},\MW{2}
\left.\right\rbrace
\end{flalign*}

De aqu{\'\i} resulta que el grupo de Frattini $\left\lbrace{1,a^2,a^4,a^6}\right\rbrace=$ y el de Fitting maximal.



El siguiente grupo ser\'a, $\mathsf{MW}_3$.



\smallskip\noindent
\resizebox{\linewidth}{!}{
\begin{tabular}{|c|c|c|c|c|c|c|c|c|c|c|c|c|c|c|c|}
\hline
$1$&$a$&$a^2$&$a^3$&$a^4$&$a^5$&$a^6$&$a^7$&$b$&$ab$&$a^2b$&$a^3b$&$a^4b$&$a^5b$&$a^6b$&$a^7b$\\
\hline
$a$&$a^2$&$a^3$&$a^4$&$a^5$&$a^6$&$a^7$&$1$&$ab$&$a^2b$&$a^3b$&$a^4b$&$a^5b$&$a^6b$&$a^7b$&$b$\\
\hline
$a^2$&$a^3$&$a^4$&$a^5$&$a^6$&$a^7$&$1$&$a$&$a^2b$&$a^3b$&$a^4b$&$a^5b$&$a^6b$&$a^7b$&$b$&$ab$\\
\hline
$a^3$&$a^4$&$a^5$&$a^6$&$a^7$&$1$&$a$&$a^2$&$a^3b$&$a^4b$&$a^5b$&$a^6b$&$a^7b$&$b$&$ab$&$a^2b$\\
\hline
$a^4$&$a^5$&$a^6$&$a^7$&$1$&$a$&$a^2$&$a^3$&$a^4b$&$a^5b$&$a^6b$&$a^7b$&$b$&$ab$&$a^2b$&$a^3b$\\
\hline
$a^5$&$a^6$&$a^7$&$1$&$a$&$a^2$&$a^3$&$a^4$&$a^5b$&$a^6b$&$a^7b$&$b$&$ab$&$a^2b$&$a^3b$&$a^4b$\\
\hline
$a^6$&$a^7$&$1$&$a$&$a^2$&$a^3$&$a^4$&$a^5$&$a^6b$&$a^7b$&$b$&$ab$&$a^2b$&$a^3b$&$a^4b$&$a^5b$\\
\hline
$a^7$&$1$&$a$&$a^2$&$a^3$&$a^4$&$a^5$&$a^6$&$a^7b$&$b$&$ab$&$a^2b$&$a^3b$&$a^4b$&$a^5b$&$a^6b$\\
\hline
$b$&$a^5b$&$a^2b$&$a^7b$&$a^4b$&$ab$&$a^6b$&$a^3b$&$1$&$a^5$&$a^2$&$a^7$&$a^4$&$a$&$a^6$&$a^3$\\
\hline
$ab$&$a^6b$&$a^3b$&$b$&$a^5b$&$a^2b$&$a^7b$&$a^4b$&$a$&$a^6$&$a^3$&$1$&$a^5$&$a^2$&$a^7$&$a^4$\\
\hline
$a^2b$&$a^7b$&$a^4b$&$ab$&$a^6b$&$a^3b$&$b$&$a^5b$&$a^2$&$a^7$&$a^4$&$a$&$a^6$&$a^3$&$1$&$a^5$\\
\hline
$a^3b$&$b$&$a^5b$&$a^2b$&$a^7b$&$a^4b$&$ab$&$a^6b$&$a^3$&$1$&$a^5$&$a^2$&$a^7$&$a^4$&$a$&$a^6$\\
\hline
$a^4b$&$ab$&$a^6b$&$a^3b$&$b$&$a^5b$&$a^2b$&$a^7b$&$a^4$&$a$&$a^6$&$a^3$&$1$&$a^5$&$a^2$&$a^7$\\
\hline
$a^5b$&$a^2b$&$a^7b$&$a^4b$&$ab$&$a^6b$&$a^3b$&$b$&$a^5$&$a^2$&$a^7$&$a^4$&$a$&$a^6$&$a^3$&$1$\\
\hline
$a^6b$&$a^3b$&$b$&$a^5b$&$a^2b$&$a^7b$&$a^4b$&$ab$&$a^6$&$a^3$&$1$&$a^5$&$a^2$&$a^7$&$a^4$&$a$\\
\hline
$a^7b$&$a^4b$&$ab$&$a^6b$&$a^3b$&$b$&$a^5b$&$a^2b$&$a^7$&$a^4$&$a$&$a^6$&$a^3$&$1$&$a^5$&$a^2$\\
\hline
\end{tabular}
}


Este grupo no es abeliano.
El primer invariante ser\'a el n\'umero de elementos de cada orden:


\begin{flalign*}
&\left[1,3,4,8,0\right]
\end{flalign*}


En este grupo no todos los subgrupos c{\'\i}clicos son normales, y sabiendo que el centro del grupo es:

\[
\centro{\mathsf{MW}_3} = {\left\lbrace{1,a^2,a^4,a^6}\right\rbrace}          \cong          \cyclicgr{4}
\]

y el subgrupo derivado es

\[
{\mathsf{MW}_3}^{\prime} = {\left\lbrace{1,a^4}\right\rbrace}          \cong          \cyclicgr{2}
\]

Estamos en condiciones de clasificar los subgrupos c{\'\i}clicos y en relaci\'on con el centro y el derivado:

\begin{flalign*}
&\mathcal{L}_1\left({\mathsf{MW}_{3}}\right)=\left\lbrace\right.\\
&\quad {\left\lbrace{1}\right\rbrace} ,
{\left\lbrace{1,a^4}\right\rbrace} ,%cd
{\left\lbrace{1,b}\right\rbrace} ,%
{\left\lbrace{1,a^4b}\right\rbrace} ,%
{\left\lbrace{1,a^2,a^4,a^6}\right\rbrace} ,\\%c
&\quad {\left\lbrace{1,a^4,a^2b,a^6b}\right\rbrace} ,%n
{\left\lbrace{1,a,a^2,a^3,a^4,a^5,a^6,a^7}\right\rbrace} ,\\%n
&\quad {\left\lbrace{1,a^2,a^4,a^6,ab,a^3b,a^5b,a^7b}\right\rbrace}\\%n
&\left.\right\rbrace\\
&\left[cd(1,2),(2,2),c(1,4),n(1,4),n(2,8)\right]
\end{flalign*}

Ahora los subgrupos generados por dos elementos, siguiendo el mismo patr\'on anterior,
sacaremos un nuevo invariante:

\begin{flalign*}
&\mathcal{L}_{2}\left({\mathsf{MW}_3}\right)=\left\lbrace\right.\\
&\quad {\left\lbrace{1,a^4,b,a^4b}\right\rbrace} ,\\%na
&\quad {\left\lbrace{1,a^2,a^4,a^6,b,a^2b,a^4b,a^6b}\right\rbrace} ,\\%na
&\quad {\mathsf{MW}_3}\\
&\left.\right\rbrace\\
&\left[na(1,4),na(1,\pdcyclics{4}{2})\right]
\end{flalign*}


Serie central ascendente (comienza en el $\trivialgr$,
y en general no tiene porqu\'e alcanzar $\mathsf{MW}_3$):
\begin{flalign*}
&{\trivialgr} \normsubgr
{\left\lbrace{1,a^2,a^4,a^6}\right\rbrace}  \normsubgr
{\mathsf{MW}_3}
\end{flalign*}

Serie central descendente (comienza en el $\mathsf{MW}_3$,
y en general no tiene proqu\'e alcanzar el $\trivialgr$):
\begin{flalign*}
&{\mathsf{MW}_3}\triangleright
{\left\lbrace{1,a^4}\right\rbrace}\triangleright
{\trivialgr}
\end{flalign*}

El invariante que derivamos:
\begin{flalign*}
&\left[
\mathtt{nil}_2,
\cyclicgr{2},
\cyclicgr{4}
\right]
\end{flalign*}



Ahora el invariante total queda:


\begin{flalign*}
&\left[1,3,4,8,0\right]\\
&\left[cd(1,2),(2,2),c(1,4),n(1,4),n(2,8)\right]\\
&\left[na(1,4),na(1,\pdcyclics{4}{2})\right]\\
&\left[
\mathtt{nil}_2,
\cyclicgr{2},
\cyclicgr{4}
\right]
\end{flalign*}

Hemos visto que $\cyclicgr{2}         \cong         {\mathsf{MW}_{3}}^{\prime}\subsetneqq\mathsf{Z}\left(\mathsf{MW}_{3}\right)         \cong         \cyclicgr{4}$.
Ahora podemos calcular las clases de conjugaci\'on y ver los grupos cocientes.

El grupo de automorfismos internos es isomorfo a $\nicefrac{\mathsf{MW}_3}{\centro{\mathsf{MW}_3}}         \cong         \pdcyclics{2}{2}$.

A su vez ${\mathsf{MW}_3}^{\mathtt{ab}}=\nicefrac{\mathsf{MW}_3}{{\mathsf{MW}_3}^{\prime}}         \cong          \pdcyclics{4}{2}$.

El ret{\'\i}culo de subgrupos normales y nilpotentes es:

\begin{flalign*}
&\mathcal{L}_{norm,nil}\entreparentesis{\MW{2}}=\left\lbrace\right.\\
&\qquad \set{1},{\left\lbrace{1,a^4}\right\rbrace},
{\left\lbrace{1,a^2,a^4,a^6}\right\rbrace},
{\left\lbrace{1,a^4,a^2b,a^6b}\right\rbrace},\\
&\qquad {\left\lbrace{1,a,a^2,a^3,a^4,a^5,a^6,a^7}\right\rbrace},\\
&\qquad {\left\lbrace{1,a^2,a^4,a^6,b,a^2b,a^4b,a^6b}\right\rbrace},\\
&\qquad \MW{2}
\left.\right\rbrace
\end{flalign*}

De aqu{\'\i} resulta que el grupo de Frattini sea el subgrupo derivado y el de Fitting maximal.




El siguiente grupo ser\'a, $\MW{4}$.



\smallskip\noindent
\resizebox{\linewidth}{!}{
\begin{tabular}{|c|c|c|c|c|c|c|c|c|c|c|c|c|c|c|c|}
\hline
$1$&$a$&$a^2$&$a^3$&$a^4$&$a^5$&$a^6$&$a^7$&$b$&$ab$&$a^2b$&$a^3b$&$a^4b$&$a^5b$&$a^6b$&$a^7b$\\
\hline
$a$&$a^2$&$a^3$&$a^4$&$a^5$&$a^6$&$a^7$&$1$&$ab$&$a^2b$&$a^3b$&$a^4b$&$a^5b$&$a^6b$&$a^7b$&$b$\\
\hline
$a^2$&$a^3$&$a^4$&$a^5$&$a^6$&$a^7$&$1$&$a$&$a^2b$&$a^3b$&$a^4b$&$a^5b$&$a^6b$&$a^7b$&$b$&$ab$\\
\hline
$a^3$&$a^4$&$a^5$&$a^6$&$a^7$&$1$&$a$&$a^2$&$a^3b$&$a^4b$&$a^5b$&$a^6b$&$a^7b$&$b$&$ab$&$a^2b$\\
\hline
$a^4$&$a^5$&$a^6$&$a^7$&$1$&$a$&$a^2$&$a^3$&$a^4b$&$a^5b$&$a^6b$&$a^7b$&$b$&$ab$&$a^2b$&$a^3b$\\
\hline
$a^5$&$a^6$&$a^7$&$1$&$a$&$a^2$&$a^3$&$a^4$&$a^5b$&$a^6b$&$a^7b$&$b$&$ab$&$a^2b$&$a^3b$&$a^4b$\\
\hline
$a^6$&$a^7$&$1$&$a$&$a^2$&$a^3$&$a^4$&$a^5$&$a^6b$&$a^7b$&$b$&$ab$&$a^2b$&$a^3b$&$a^4b$&$a^5b$\\
\hline
$a^7$&$1$&$a$&$a^2$&$a^3$&$a^4$&$a^5$&$a^6$&$a^7b$&$b$&$ab$&$a^2b$&$a^3b$&$a^4b$&$a^5b$&$a^6b$\\
\hline
$b$&$a^7b$&$a^6b$&$a^5b$&$a^4b$&$a^3b$&$a^2b$&$ab$&$1$&$a^7$&$a^6$&$a^5$&$a^4$&$a^3$&$a^2$&$a$\\
\hline
$ab$&$b$&$a^7b$&$a^6b$&$a^5b$&$a^4b$&$a^3b$&$a^2b$&$a$&$1$&$a^7$&$a^6$&$a^5$&$a^4$&$a^3$&$a^2$\\
\hline
$a^2b$&$ab$&$b$&$a^7b$&$a^6b$&$a^5b$&$a^4b$&$a^3b$&$a^2$&$a$&$1$&$a^7$&$a^6$&$a^5$&$a^4$&$a^3$\\
\hline
$a^3b$&$a^2b$&$ab$&$b$&$a^7b$&$a^6b$&$a^5b$&$a^4b$&$a^3$&$a^2$&$a$&$1$&$a^7$&$a^6$&$a^5$&$a^4$\\
\hline
$a^4b$&$a^3b$&$a^2b$&$ab$&$b$&$a^7b$&$a^6b$&$a^5b$&$a^4$&$a^3$&$a^2$&$a$&$1$&$a^7$&$a^6$&$a^5$\\
\hline
$a^5b$&$a^4b$&$a^3b$&$a^2b$&$ab$&$b$&$a^7b$&$a^6b$&$a^5$&$a^4$&$a^3$&$a^2$&$a$&$1$&$a^7$&$a^6$\\
\hline
$a^6b$&$a^5b$&$a^4b$&$a^3b$&$a^2b$&$ab$&$b$&$a^7b$&$a^6$&$a^5$&$a^4$&$a^3$&$a^2$&$a$&$1$&$a^7$\\
\hline
$a^7b$&$a^6b$&$a^5b$&$a^4b$&$a^3b$&$a^2b$&$ab$&$b$&$a^7$&$a^6$&$a^5$&$a^4$&$a^3$&$a^2$&$a$&$1$\\
\hline
\end{tabular}
}


Este grupo no es abeliano.
El primer invariante ser\'a el n\'umero de elementos de cada orden:


\begin{flalign*}
&\left[1,9,2,4,0\right]
\end{flalign*}


En este grupo no todos los subgrupos c{\'\i}clicos son normales, y sabiendo que el centro del grupo es:

\[
\centro{\MW{4}} = {\left\lbrace{1,a^4}\right\rbrace}          \cong          \cyclicgr{2}
\]

y el subgrupo derivado es

\[
{\MW{4}}^{\prime} = {\left\lbrace{1,a^2,a^4,a^6}\right\rbrace}          \cong          \cyclicgr{4}
\]

Estamos en condiciones de clasificar los subgrupos c{\'\i}clicos y en relaci\'on con el centro y el derivado:

\begin{flalign*}
&\mathcal{L}_1\left({\MW{4}}\right)=\left\lbrace\right.\\
&\quad {\left\lbrace{1}\right\rbrace} ,
{\left\lbrace{1,a^4}\right\rbrace} ,%c
{\left\lbrace{1,b}\right\rbrace} ,%
{\left\lbrace{1,ab}\right\rbrace} ,%
{\left\lbrace{1,a^2b}\right\rbrace} ,%
{\left\lbrace{1,a^3b}\right\rbrace} ,\\%
&\quad {\left\lbrace{1,a^4b}\right\rbrace} ,%
{\left\lbrace{1,a^5b}\right\rbrace} ,%
{\left\lbrace{1,a^6b}\right\rbrace} ,%
{\left\lbrace{1,a^7b}\right\rbrace} ,%
{\left\lbrace{1,a^2,a^4,a^6}\right\rbrace} ,\\%d
&\quad {\left\lbrace{1,a,a^2,a^3,a^4,a^5,a^6,a^7}\right\rbrace}\\%n
&\left.\right\rbrace\\
&\left[c(1,2),(8,2),d(1,4),n(1,8)\right]
\end{flalign*}

Ahora los subgrupos generados por dos elementos, siguiendo el mismo patr\'on anterior,
sacaremos un nuevo invariante:


\begin{flalign*}
&\mathcal{L}_{2}\left({\MW{4}}\right)=\left\lbrace\right.\\
&\quad {\left\lbrace{1,a^4,b,a^4b}\right\rbrace} ,%a
{\left\lbrace{1,a^4,ab,a^5b}\right\rbrace} ,%a
{\left\lbrace{1,a^4,a^2b,a^6b}\right\rbrace} ,\\%a
&\quad {\left\lbrace{1,a^4,a^3b,a^7b}\right\rbrace} ,%a
{\left\lbrace{1,a^2,a^4,a^6,b,a^2b,a^4b,a^6b}\right\rbrace} ,%n
{\left\lbrace{1,a^2,a^4,a^6,ab,a^3b,a^5b,a^7b}\right\rbrace} ,\\%n
&\quad {\MW{4}}\\
&\left.\right\rbrace\\
&\left[a(4,4),n(2,8)\right]
\end{flalign*}


Serie central ascendente (comienza en el $\trivialgr$,
y en general no tiene porqu\'e alcanzar $\MW{4}$):


\begin{flalign*}
&{\trivialgr} \normsubgr
{\left\lbrace{1,a^4}\right\rbrace} \normsubgr
{\left\lbrace{1,a^2,a^4,a^6}\right\rbrace} \normsubgr
{\MW{4}}
\end{flalign*}


Serie central descendente (comienza en el $\MW{4}$,
y en general no tiene proqu\'e alcanzar el $\trivialgr$):


\begin{flalign*}
&{\MW{4}}\triangleright
{\left\lbrace{1,a^2,a^4,a^6}\right\rbrace}\triangleright
{\left\lbrace{1,a^4}\right\rbrace}\triangleright
{\trivialgr}
\end{flalign*}


El invariante que derivamos:


\begin{flalign*}
&\left[
\mathtt{nil}_3,
\left[\cyclicgr{2},\cyclicgr{4}\right],
\left[\cyclicgr{4},\cyclicgr{2}\right]
\right]
\end{flalign*}



Ahora el invariante total queda:


\begin{flalign*}
&\left[1,9,2,4,0\right]\\
&\left[c(1,2),(8,2),d(1,4),n(1,8)\right]\\
&\left[a(4,4),n(2,8)\right]\\
&\left[
\mathtt{nil}_3,
\left[\cyclicgr{2},\cyclicgr{4}\right],
\left[\cyclicgr{4},\cyclicgr{2}\right]
\right]
\end{flalign*}


Hemos visto que $\cyclicgr{2}         \cong         {\MW{4}}^{\prime}\supsetneqq\mathsf{Z}\left(\MW{4}\right)         \cong         \cyclicgr{4}$.
Ahora podemos calcular las clases de conjugaci\'on y ver los grupos cocientes.


El grupo de automorfismos internos es isomorfo a una extensi\'on del grupo $\cyclicgr{8}$ por el grupo $\cyclicgr{2}$, necesariamente no abeliana.


A su vez ${\MW{4}}^{\mathtt{ab}}=\nicefrac{\MW{4}}{{\MW{4}}^{\prime}}         \cong          \pdcyclics{2}{2}$.


El ret{\'\i}culo de subgrupos normales y nilpotentes es:


\begin{flalign*}
&\mathcal{L}_{norm,nil}\entreparentesis{\MW{2}}=\left\lbrace\right.\\
&\qquad \set{1},{\left\lbrace{1,a^4}\right\rbrace},
{\left\lbrace{1,a^2,a^4,a^6}\right\rbrace},
{\left\lbrace{1,a,a^2,a^3,a^4,a^5,a^6,a^7}\right\rbrace},\\
&\qquad {\left\lbrace{1,a^2,a^4,a^6,b,a^2b,a^4b,a^6b}\right\rbrace},\\
&\qquad {\left\lbrace{1,a^2,a^4,a^6,ab,a^3b,a^5b,a^7b}\right\rbrace},\\
&\qquad {\MW{4}}\\
\left.\right\rbrace
\end{flalign*}

De aqu{\'\i} resulta que el grupo de Frattini sea el subgrupo central y el de Fitting maximal.


El siguiente grupo ser\'a, ${\mathsf{MW}}_{5}$.



\smallskip\noindent
\resizebox{\linewidth}{!}{
\begin{tabular}{|c|c|c|c|c|c|c|c|c|c|c|c|c|c|c|c|}
\hline
$1$&$a$&$a^2$&$a^3$&$a^4$&$a^5$&$a^6$&$a^7$&$b$&$ab$&$a^2b$&$a^3b$&$a^4b$&$a^5b$&$a^6b$&$a^7b$\\
\hline
$a$&$a^2$&$a^3$&$a^4$&$a^5$&$a^6$&$a^7$&$1$&$ab$&$a^2b$&$a^3b$&$a^4b$&$a^5b$&$a^6b$&$a^7b$&$b$\\
\hline
$a^2$&$a^3$&$a^4$&$a^5$&$a^6$&$a^7$&$1$&$a$&$a^2b$&$a^3b$&$a^4b$&$a^5b$&$a^6b$&$a^7b$&$b$&$ab$\\
\hline
$a^3$&$a^4$&$a^5$&$a^6$&$a^7$&$1$&$a$&$a^2$&$a^3b$&$a^4b$&$a^5b$&$a^6b$&$a^7b$&$b$&$ab$&$a^2b$\\
\hline
$a^4$&$a^5$&$a^6$&$a^7$&$1$&$a$&$a^2$&$a^3$&$a^4b$&$a^5b$&$a^6b$&$a^7b$&$b$&$ab$&$a^2b$&$a^3b$\\
\hline
$a^5$&$a^6$&$a^7$&$1$&$a$&$a^2$&$a^3$&$a^4$&$a^5b$&$a^6b$&$a^7b$&$b$&$ab$&$a^2b$&$a^3b$&$a^4b$\\
\hline
$a^6$&$a^7$&$1$&$a$&$a^2$&$a^3$&$a^4$&$a^5$&$a^6b$&$a^7b$&$b$&$ab$&$a^2b$&$a^3b$&$a^4b$&$a^5b$\\
\hline
$a^7$&$1$&$a$&$a^2$&$a^3$&$a^4$&$a^5$&$a^6$&$a^7b$&$b$&$ab$&$a^2b$&$a^3b$&$a^4b$&$a^5b$&$a^6b$\\
\hline
$b$&$a^7b$&$a^6b$&$a^5b$&$a^4b$&$a^3b$&$a^2b$&$ab$&$a^4$&$a^3$&$a^2$&$a$&$1$&$a^7$&$a^6$&$a^5$\\
\hline
$ab$&$b$&$a^7b$&$a^6b$&$a^5b$&$a^4b$&$a^3b$&$a^2b$&$a^5$&$a^4$&$a^3$&$a^2$&$a$&$1$&$a^7$&$a^6$\\
\hline
$a^2b$&$ab$&$b$&$a^7b$&$a^6b$&$a^5b$&$a^4b$&$a^3b$&$a^6$&$a^5$&$a^4$&$a^3$&$a^2$&$a$&$1$&$a^7$\\
\hline
$a^3b$&$a^2b$&$ab$&$b$&$a^7b$&$a^6b$&$a^5b$&$a^4b$&$a^7$&$a^6$&$a^5$&$a^4$&$a^3$&$a^2$&$a$&$1$\\
\hline
$a^4b$&$a^3b$&$a^2b$&$ab$&$b$&$a^7b$&$a^6b$&$a^5b$&$1$&$a^7$&$a^6$&$a^5$&$a^4$&$a^3$&$a^2$&$a$\\
\hline
$a^5b$&$a^4b$&$a^3b$&$a^2b$&$ab$&$b$&$a^7b$&$a^6b$&$a$&$1$&$a^7$&$a^6$&$a^5$&$a^4$&$a^3$&$a^2$\\
\hline
$a^6b$&$a^5b$&$a^4b$&$a^3b$&$a^2b$&$ab$&$b$&$a^7b$&$a^2$&$a$&$1$&$a^7$&$a^6$&$a^5$&$a^4$&$a^3$\\
\hline
$a^7b$&$a^6b$&$a^5b$&$a^4b$&$a^3b$&$a^2b$&$ab$&$b$&$a^3$&$a^2$&$a$&$1$&$a^7$&$a^6$&$a^5$&$a^4$\\
\hline
\end{tabular}
}


Este grupo no es abeliano.
El primer invariante ser\'a el n\'umero de elementos de cada orden:


\begin{flalign*}
&\left[1,1,10,4,0\right]
\end{flalign*}


En este grupo no todos los subgrupos c{\'\i}clicos son normales, y sabiendo que el centro del grupo es:

\[
\centro{{\mathsf{MW}}_{5}} = {\left\lbrace{1,a^4}\right\rbrace}          \cong          \cyclicgr{2}
\]

y el subgrupo derivado es

\[
{{\mathsf{MW}}_{5}}^{\prime} = {\left\lbrace{1,a^2,a^4,a^6}\right\rbrace}          \cong          \cyclicgr{4}
\]

Estamos en condiciones de clasificar los subgrupos c{\'\i}clicos y en relaci\'on con el centro y el derivado:

\begin{flalign*}
&\mathcal{L}_1\left({{\mathsf{MW}}_{5}}\right)=\left\lbrace\right.\\
&\quad {\left\lbrace{1}\right\rbrace} ,
{\left\lbrace{1,a^4}\right\rbrace} ,%c
{\left\lbrace{1,a^2,a^4,a^6}\right\rbrace} ,%d
{\left\lbrace{1,a^4,b,a^4b}\right\rbrace} ,%
{\left\lbrace{1,a^4,ab,a^5b}\right\rbrace} ,%
{\left\lbrace{1,a^4,a^2b,a^6b}\right\rbrace} ,\\%
&\quad {\left\lbrace{1,a^4,a^3b,a^7b}\right\rbrace} ,%
{\left\lbrace{1,a,a^2,a^3,a^4,a^5,a^6,a^7}\right\rbrace}\\%n
&\left.\right\rbrace\\
&\left[c(1,2),d(1,4),(4,4),n(1,8)\right]
\end{flalign*}

Ahora los subgrupos generados por dos elementos, siguiendo el mismo patr\'on anterior,
sacaremos un nuevo invariante:

\begin{flalign*}
&\mathcal{L}_{2}\left({{\mathsf{MW}}_{5}}\right)=\left\lbrace\right.\\
&\quad {\left\lbrace{1,a^2,a^4,a^6,b,a^2b,a^4b,a^6b}\right\rbrace} ,
{\left\lbrace{1,a^2,a^4,a^6,ab,a^3b,a^5b,a^7b}\right\rbrace} ,\\
&\quad {{\mathsf{MW}}_{5}}\\
&\left.\right\rbrace\\
&\left[n(2,8)\right]
\end{flalign*}


Serie central ascendente (comienza en el $\trivialgr$,
y en general no tiene porqu\'e alcanzar ${\mathsf{MW}}_{5}$):


\begin{flalign*}
&{\trivialgr} \normsubgr
{\left\lbrace{1,a^4}\right\rbrace} \normsubgr
{\left\lbrace{1,a^2,a^4,a^6}\right\rbrace} \normsubgr
{{\mathsf{MW}}_{5}}
\end{flalign*}


Serie central descendente (comienza en el ${\mathsf{MW}}_{5}$,
y en general no tiene proqu\'e alcanzar el $\trivialgr$):


\begin{flalign*}
&{{\mathsf{MW}}_{5}}\triangleright
{\left\lbrace{1,a^2,a^4,a^6}\right\rbrace}\triangleright
{\left\lbrace{1,a^4}\right\rbrace}\triangleright
{\trivialgr}
\end{flalign*}


El invariante que derivamos:


\begin{flalign*}
&\left[
\mathtt{nil}_3,
\left[\cyclicgr{2},\cyclicgr{4}\right],
\left[\cyclicgr{4},\cyclicgr{2}\right]
\right]
\end{flalign*}



Ahora el invariante total queda:


\begin{flalign*}
&\left[1,1,10,4,0\right]\\
&\left[c(1,2),d(1,4),(4,4),n(1,8)\right]\\
&\left[n(2,8)\right]\\
&\left[
\mathtt{nil}_3,
\left[\cyclicgr{2},\cyclicgr{4}\right],
\left[\cyclicgr{4},\cyclicgr{2}\right]
\right]\\
&\left[\cyclicgr{2}         \cong         \centro{\mathsf{MW}_{5}}\subsetneqq{\mathsf{MW}_{5}}^{\prime}         \cong         \cyclicgr{4}\right]
\end{flalign*}


Hemos visto que $\cyclicgr{2}         \cong         \centro{\mathsf{MW}_5}\subsetneqq{{\mathsf{MW}}_{5}}^{\prime}         \cong         \cyclicgr{4}$.
Ahora podemos calcular las clases de conjugaci\'on y ver los grupos cocientes.


El grupo de automorfismos internos es isomorfo al grupo dihedral de $8$ elementos,
${\mathtt{Int}}\entreparentesis{{\mathsf{MW}}_{5}}
         \cong         \nicefrac{{\mathsf{MW}}_{5}}{\centro{{\mathsf{MW}}_{5}}}
         \cong         \dihgr{4}$, no abeliano.


A su vez ${
{{\mathsf{MW}}_{5}}^{\mathtt{ab}}
=
\nicefrac
{{\mathsf{MW}}_{5}}
{{{\mathsf{MW}}_{5}}^{\prime}}
         \cong
\pdcyclics{2}{2}}$.

A{\~{n}}adiendo esta informaci\'on al conjunto de invariantes:

\begin{flalign*}
&\left[1,1,10,4,0\right]\\
&\left[c(1,2),d(1,4),(4,4),n(1,8)\right]\\
&\left[n(2,8)\right]\\
&\left[
\mathtt{nil}_3,
\left[\cyclicgr{2},\cyclicgr{4}\right],
\left[\cyclicgr{4},\cyclicgr{2}\right]
\right]\\
&\left[\cyclicgr{2}         \cong         \centro{\mathsf{MW}_{5}}\subsetneqq{\mathsf{MW}_{5}}^{\prime}         \cong         \cyclicgr{4}\right]\\
&\left[\left[\nicefrac{{\mathsf{MW}}_{5}}{\centro{{\mathsf{MW}}_{5}}}
         \cong         \dihgr{4},\nicefrac{{\mathsf{MW}}_{5}}{{{\mathsf{MW}}_{5}}^{\prime}}         \cong          \pdcyclics{2}{2}\right],\left[\right]\right]
\end{flalign*}

El grupo de corchetes vac{\'\i}os es para el caso $\mathtt{nil}_3$, d\'onde tendremos el centro segundo y el subgrupo derivado segundo $\mathsf{L}^2$, y as{\'\i} tendremos sus cocientes.

El ret{\'\i}culo de subgrupos normales y nilpotentes es:

\begin{flalign*}
&\mathcal{L}_{norm,nil}\entreparentesis{{\mathsf{MW}}_{5}}=\left\lbrace\right.\\
&\qquad \set{1},{\left\lbrace{1,a^4}\right\rbrace},
{\left\lbrace{1,a^2,a^4,a^6}\right\rbrace},
{\left\lbrace{1,a,a^2,a^3,a^4,a^5,a^6,a^7}\right\rbrace},\\
&\qquad {\left\lbrace{1,a^2,a^4,a^6,b,a^2b,a^4b,a^6b}\right\rbrace},\\
&\qquad {\left\lbrace{1,a^2,a^4,a^6,ab,a^3b,a^5b,a^7b}\right\rbrace},\\
&\qquad {{\mathsf{MW}}_{5}}\\
\left.\right\rbrace
\end{flalign*}

De aqu{\'\i} resulta que el grupo de Frattini sea el subgrupo central y el de Fitting maximal.

\begin{flalign*}
&\left[1,1,10,4,0\right]\\
&\left[c(1,2),d(1,4),(4,4),n(1,8)\right]\\
&\left[n(2,8)\right]\\
&\left[
\mathtt{nil}_3,
\left[\cyclicgr{2},\cyclicgr{4}\right],
\left[\cyclicgr{4},\cyclicgr{2}\right]
\right]\\
&\left[\cyclicgr{2}         \cong         \centro{\mathsf{MW}_{5}}\subsetneqq{\mathsf{MW}_{5}}^{\prime}         \cong         \cyclicgr{4}\right]\\
&\left[\left[\nicefrac{{\mathsf{MW}}_{5}}{\centro{{\mathsf{MW}}_{5}}}
         \cong         \dihgr{4},\nicefrac{{\mathsf{MW}}_{5}}{{{\mathsf{MW}}_{5}}^{\prime}}         \cong          \pdcyclics{2}{2}\right],\left[\right]\right]\\
&\left[\mathbf{\Phi }\entreparentesis{{\mathsf{MW}}_{5}}=\centro{{\mathsf{MW}}_{5}},\mathbf{\mathsf{F}}\entreparentesis{{\mathsf{MW}}_{5}}={\mathsf{MW}}_{5}\right]
\end{flalign*}




El siguiente grupo ser\'a, ${\mathsf{MW}}_{6}$.



\smallskip\noindent
\resizebox{\linewidth}{!}{
\begin{tabular}{|c|c|c|c|c|c|c|c|c|c|c|c|c|c|c|c|}
\hline
$1$&$a$&$a^2$&$a^3$&$a^4$&$a^5$&$a^6$&$a^7$&$b$&$ab$&$a^2b$&$a^3b$&$a^4b$&$a^5b$&$a^6b$&$a^7b$\\
\hline
$a$&$a^2$&$a^3$&$a^4$&$a^5$&$a^6$&$a^7$&$1$&$ab$&$a^2b$&$a^3b$&$a^4b$&$a^5b$&$a^6b$&$a^7b$&$b$\\
\hline
$a^2$&$a^3$&$a^4$&$a^5$&$a^6$&$a^7$&$1$&$a$&$a^2b$&$a^3b$&$a^4b$&$a^5b$&$a^6b$&$a^7b$&$b$&$ab$\\
\hline
$a^3$&$a^4$&$a^5$&$a^6$&$a^7$&$1$&$a$&$a^2$&$a^3b$&$a^4b$&$a^5b$&$a^6b$&$a^7b$&$b$&$ab$&$a^2b$\\
\hline
$a^4$&$a^5$&$a^6$&$a^7$&$1$&$a$&$a^2$&$a^3$&$a^4b$&$a^5b$&$a^6b$&$a^7b$&$b$&$ab$&$a^2b$&$a^3b$\\
\hline
$a^5$&$a^6$&$a^7$&$1$&$a$&$a^2$&$a^3$&$a^4$&$a^5b$&$a^6b$&$a^7b$&$b$&$ab$&$a^2b$&$a^3b$&$a^4b$\\
\hline
$a^6$&$a^7$&$1$&$a$&$a^2$&$a^3$&$a^4$&$a^5$&$a^6b$&$a^7b$&$b$&$ab$&$a^2b$&$a^3b$&$a^4b$&$a^5b$\\
\hline
$a^7$&$1$&$a$&$a^2$&$a^3$&$a^4$&$a^5$&$a^6$&$a^7b$&$b$&$ab$&$a^2b$&$a^3b$&$a^4b$&$a^5b$&$a^6b$\\
\hline
$b$&$ab$&$a^2b$&$a^3b$&$a^4b$&$a^5b$&$a^6b$&$a^7b$&$a$&$a^2$&$a^3$&$a^4$&$a^5$&$a^6$&$a^7$&$1$\\
\hline
$ab$&$a^2b$&$a^3b$&$a^4b$&$a^5b$&$a^6b$&$a^7b$&$b$&$a^2$&$a^3$&$a^4$&$a^5$&$a^6$&$a^7$&$1$&$a$\\
\hline
$a^2b$&$a^3b$&$a^4b$&$a^5b$&$a^6b$&$a^7b$&$b$&$ab$&$a^3$&$a^4$&$a^5$&$a^6$&$a^7$&$1$&$a$&$a^2$\\
\hline
$a^3b$&$a^4b$&$a^5b$&$a^6b$&$a^7b$&$b$&$ab$&$a^2b$&$a^4$&$a^5$&$a^6$&$a^7$&$1$&$a$&$a^2$&$a^3$\\
\hline
$a^4b$&$a^5b$&$a^6b$&$a^7b$&$b$&$ab$&$a^2b$&$a^3b$&$a^5$&$a^6$&$a^7$&$1$&$a$&$a^2$&$a^3$&$a^4$\\
\hline
$a^5b$&$a^6b$&$a^7b$&$b$&$ab$&$a^2b$&$a^3b$&$a^4b$&$a^6$&$a^7$&$1$&$a$&$a^2$&$a^3$&$a^4$&$a^5$\\
\hline
$a^6b$&$a^7b$&$b$&$ab$&$a^2b$&$a^3b$&$a^4b$&$a^5b$&$a^7$&$1$&$a$&$a^2$&$a^3$&$a^4$&$a^5$&$a^6$\\
\hline
$a^7b$&$b$&$ab$&$a^2b$&$a^3b$&$a^4b$&$a^5b$&$a^6b$&$1$&$a$&$a^2$&$a^3$&$a^4$&$a^5$&$a^6$&$a^7$\\
\hline
\end{tabular}
}




Este grupo es abeliano. De hecho se trata del grupo ${\cyclicgr{16}}$.
\[
\generadopor{b}={\mathsf{MW}}_{6}         \cong          \cyclicgr{16}
\]
Los distintos elementos pueden ser generados {as\'\i}:
\begin{flalign*}
&b\\
&b^2=a\\
&b^3=ab\\
&b^4=a^2\\
&b^5=a^2b\\
&b^6=a^3\\
&b^7=a^3b\\
&b^8=a^4\\
&b^9=a^4b\\
&b^{10}=a^5\\
&b^{11}=a^5b\\
&b^{12}=a^6\\
&b^{13}=a^6b\\
&b^{14}=a^7\\
&b^{15}=a^7b\\
&b^{16}=1
\end{flalign*}
El primer invariante ser\'a el n\'umero de elementos de cada orden:


\begin{flalign*}
&\left[1,1,2,4,8\right]
\end{flalign*}


En este grupo todos los subgrupos c{\'\i}clicos son normales, y sabiendo que el centro del grupo es el propio grupo, todos son centrales, sabiendo que el subgrupo derivado es el minimal, ninguno no trivial est\'a contenido en el derivado:

\[
\centro{{\mathsf{MW}}_{6}} = {\mathsf{MW}}_{6}
\]

y el subgrupo derivado es

\[
{{\mathsf{MW}}_{6}}^{\prime} = \trivialgr
\]

Estamos en condiciones de clasificar los subgrupos c{\'\i}clicos y en relaci\'on con el centro y el derivado:


\begin{flalign*}
&\mathcal{L}_1\left({{\mathsf{MW}}_{6}}\right)=\left\lbrace\right.\\
&\quad {\left\lbrace{1}\right\rbrace} ,
{\left\lbrace{1,a^4}\right\rbrace} ,
{\left\lbrace{1,a^2,a^4,a^6}\right\rbrace} ,
{\left\lbrace{1,a,a^2,a^3,a^4,a^5,a^6,a^7}\right\rbrace} ,\\
&\quad {\mathsf{MW}}_{6}\\%n
&\left.\right\rbrace\\
&\left[c(1,1),c(1,4),c(1,8),c(1,16)\right]
\end{flalign*}

Ahora los subgrupos generados por dos elementos, siguiendo el mismo patr\'on anterior,
sacaremos un nuevo invariante (vac{\'\i}o en este caso):

\begin{flalign*}
&\mathcal{L}_{2}\left({{\mathsf{MW}}_{6}}\right)=\left\lbrace\right\rbrace\\
&\left[\right]
\end{flalign*}


Serie central ascendente (comienza en el $\trivialgr$,
y en general no tiene porqu\'e alcanzar ${\mathsf{MW}}_{6}$):


\begin{flalign*}
&{\trivialgr} \normsubgr
{{\mathsf{MW}}_{6}}
\end{flalign*}


Serie central descendente (comienza en el ${\mathsf{MW}}_{6}$,
y en general no tiene proqu\'e alcanzar el $\trivialgr$):


\begin{flalign*}
&{{\mathsf{MW}}_{6}}\triangleright
{\trivialgr}
\end{flalign*}


El invariante que derivamos:


\begin{flalign*}
&\left[
\mathtt{nil}_1,
\left[\right],
\left[\right]
\right]
\end{flalign*}



Ahora el invariante total queda:


\begin{flalign*}
&\left[1,1,1,4,8\right]\\
&\left[c(1,1),c(1,4),c(1,8),c(1,16)\right]\\
&\left[\right]\\
&\left[
\mathtt{nil}_1,
\left[\right],
\left[\right]
\right]\\
&\left[{\mathsf{MW}}_{6}=\centro{{\mathsf{MW}}_{6}}\supsetneqq{{\mathsf{MW}}_{6}}^{\prime}         \cong         \trivialgr\right]
\end{flalign*}


En este caso, los grupos cocientes son triviales, por lo que el \'ultimo invariante tendr\'a los dos corchetes vac{\'\i}os. Esto ocurrir\'a siempre que el grupo sea $\mathtt{nil}_{mathbf{1}}$.


\begin{flalign*}
&\left[1,1,1,4,8\right]\\
&\left[c(1,1),c(1,4),c(1,8),c(1,16)\right]\\
&\left[\right]\\
&\left[
\mathtt{nil}_1,
\left[\right],
\left[\right]
\right]\\
&\left[{\mathsf{MW}}_{6}=\centro{{\mathsf{MW}}_{6}}\supsetneqq{{\mathsf{MW}}_{6}}^{\prime}         \cong         \trivialgr\right]\\
&\left[\left[\right],\left[\right]\right]
\end{flalign*}



En cuanto al ret{\'\i}culo de subgrupos normales y nilpotentes es claro que son todos los que son propios y no triviales:

\begin{flalign*}
&\mathcal{L}_{norm,nil}\entreparentesis{{\mathsf{MW}}_{6}}=\left\lbrace\right.\\
&\qquad \set{1},{\left\lbrace{1,a^4}\right\rbrace},
{\left\lbrace{1,a^2,a^4,a^6}\right\rbrace},
{\left\lbrace{1,a,a^2,a^3,a^4,a^5,a^6,a^7}\right\rbrace},\\
&\qquad {\left\lbrace{1,a^2,a^4,a^6,b,a^2b,a^4b,a^6b}\right\rbrace},\\
&\qquad {\left\lbrace{1,a^2,a^4,a^6,ab,a^3b,a^5b,a^7b}\right\rbrace},\\
&\qquad {{\mathsf{MW}}_{6}}\\
\left.\right\rbrace
\end{flalign*}


De aqu{\'\i} resulta que el grupo de Frattini sea el subgrupo central ser\'a el subgrupo tal que si $\generadopor{x}={\mathsf{MW}}_{6}$  ser\'a $\mathbf{\Phi }\entreparentesis{{\mathsf{MW}}_{6}}= \generadopor{x^8}$ y el de Fitting ser\'a de orden $8$, de forma que $\mathbf{\mathsf{F}}\entreparentesis{{\mathsf{MW}}_{6}}=\generadopor{x^2}$.



\begin{flalign*}
&\left[1,1,1,4,8\right]\\
&\left[
\begin{array}{l}
\left[
c(1,1),c(1,4),c(1,8),c(1,16)
\right]\\
\left[\right]\\
\left[\right]\\
\left[\right]
\end{array}
\right]\\
&\left[
\begin{array} {c|c|c}
\mathtt{nil}_1 & \trivialgr & {\mathsf{MW}}_{6} \\
{} & {\mathsf{MW}}_{6} & \trivialgr
\end{array}
\right]\\
&\left[
\begin{array}{l}
\left[{\mathsf{MW}}_{6}=\centro{{\mathsf{MW}}_{6}}\supsetneqq{{\mathsf{MW}}_{6}}^{\prime}         \cong         \trivialgr\right]\\
\left[\right]
\end{array}
\right]\\
&\left[
\begin{array}{l}
\left[
\nicefrac{{\mathsf{MW}}_{6}}{\centro{{\mathsf{MW}}_{6}}}
         \cong         \dihgr{4},\nicefrac{{\mathsf{MW}}_{6}}{{{\mathsf{MW}}_{6}}^{\prime}}         \cong          \pdcyclics{2}{2}
\right]\\
\left[\right]
\end{array}
\right]\\
&\left[
\mathbf{\Phi }\entreparentesis{{\mathsf{MW}}_{6}}
= \generadopor{b^8}         \cong         \cyclicgr{2},\mathbf{\mathsf{F}}\entreparentesis{{\mathsf{MW}}_{6}}=
\generadopor{b^2}         \cong         \cyclicgr{8}
\right]
\end{flalign*}


El siguiente grupo ser\'a, ${\mathsf{MW}}_{7}$.



\smallskip\noindent
\resizebox{\linewidth}{!}{
\begin{tabular}{|c|c|c|c|c|c|c|c|c|c|c|c|c|c|c|c|}
\hline
$1$&$a$&$a^2$&$a^3$&$b$&$ab$&$a^2b$&$a^3b$&$c$&$ac$&$a^2c$&$a^3c$&$bc$&$abc$&$a^2bc$&$a^3bc$\\
\hline
$a$&$a^2$&$a^3$&$1$&$ab$&$a^2b$&$a^3b$&$b$&$ac$&$a^2c$&$a^3c$&$c$&$abc$&$a^2bc$&$a^3bc$&$bc$\\
\hline
$a^2$&$a^3$&$1$&$a$&$a^2b$&$a^3b$&$b$&$ab$&$a^2c$&$a^3c$&$c$&$ac$&$a^2bc$&$a^3bc$&$bc$&$abc$\\
\hline
$a^3$&$1$&$a$&$a^2$&$a^3b$&$b$&$ab$&$a^2b$&$a^3c$&$c$&$ac$&$a^2c$&$a^3bc$&$bc$&$abc$&$a^2bc$\\
\hline
$b$&$ab$&$a^2b$&$a^3b$&$1$&$a$&$a^2$&$a^3$&$bc$&$abc$&$a^2bc$&$a^3bc$&$c$&$ac$&$a^2c$&$a^3c$\\
\hline
$ab$&$a^2b$&$a^3b$&$b$&$a$&$a^2$&$a^3$&$1$&$abc$&$a^2bc$&$a^3bc$&$bc$&$ac$&$a^2c$&$a^3c$&$c$\\
\hline
$a^2b$&$a^3b$&$b$&$ab$&$a^2$&$a^3$&$1$&$a$&$a^2bc$&$a^3bc$&$bc$&$abc$&$a^2c$&$a^3c$&$c$&$ac$\\
\hline
$a^3b$&$b$&$ab$&$a^2b$&$a^3$&$1$&$a$&$a^2$&$a^3bc$&$bc$&$abc$&$a^2bc$&$a^3c$&$c$&$ac$&$a^2c$\\
\hline
$c$&$ac$&$a^2c$&$a^3c$&$bc$&$abc$&$a^2bc$&$a^3bc$&$1$&$a$&$a^2$&$a^3$&$b$&$ab$&$a^2b$&$a^3b$\\
\hline
$ac$&$a^2c$&$a^3c$&$c$&$abc$&$a^2bc$&$a^3bc$&$bc$&$a$&$a^2$&$a^3$&$1$&$ab$&$a^2b$&$a^3b$&$b$\\
\hline
$a^2c$&$a^3c$&$c$&$ac$&$a^2bc$&$a^3bc$&$bc$&$abc$&$a^2$&$a^3$&$1$&$a$&$a^2b$&$a^3b$&$b$&$ab$\\
\hline
$a^3c$&$c$&$ac$&$a^2c$&$a^3bc$&$bc$&$abc$&$a^2bc$&$a^3$&$1$&$a$&$a^2$&$a^3b$&$b$&$ab$&$a^2b$\\
\hline
$bc$&$abc$&$a^2bc$&$a^3bc$&$c$&$ac$&$a^2c$&$a^3c$&$b$&$ab$&$a^2b$&$a^3b$&$1$&$a$&$a^2$&$a^3$\\
\hline
$abc$&$a^2bc$&$a^3bc$&$bc$&$ac$&$a^2c$&$a^3c$&$c$&$ab$&$a^2b$&$a^3b$&$b$&$a$&$a^2$&$a^3$&$1$\\
\hline
$a^2bc$&$a^3bc$&$bc$&$abc$&$a^2c$&$a^3c$&$c$&$ac$&$a^2b$&$a^3b$&$b$&$ab$&$a^2$&$a^3$&$1$&$a$\\
\hline
$a^3bc$&$bc$&$abc$&$a^2bc$&$a^3c$&$c$&$ac$&$a^2c$&$a^3b$&$b$&$ab$&$a^2b$&$a^3$&$1$&$a$&$a^2$\\
\hline
\end{tabular}
}




Este grupo es abeliano. Observando los distintos elementos y sus \'ordenes veremos qu\'e grupo abeliano es.

\begin{flalign*}
&\left[1,7,8,0,0\right]
\end{flalign*}

Esto reduce las posibilidades a ${\mathsf{MW}}_{7}         \cong         \pdcyclics{4}{4}$ o a ${\mathsf{MW}}_{7}         \cong         \pdtrescyclics{4}{2}{2}$. Probar\'e con el \'ultimo caso por la apariencia de la tabla. Veamos qu\'e $3$ elementos pueden generar el grupo.

\begin{flalign*}
  a&         \neq          1\\
  a^2&         \neq          1\\
  a^3&         \neq          1\\
  a^4&=1\\
  b&         \neq          1\\
  b^2&=1\\
  c&         \neq          1\\
  c^2&=1
\end{flalign*}

As{\'\i} que $\generadopor{a}         \cong         \cyclicgr{4}$, $\generadopor{b}         \cong         \cyclicgr{2}         \cong         \generadopor{c}$ y como adem\'as
$\generadopor{a}         \cap \generadopor{b}=\trivialgr$, $\generadopor{a}         \cap \generadopor{c}=\trivialgr$ y $\generadopor{b}         \cap \generadopor{c}=\trivialgr$,
queda claro que
${\mathsf{MW}}_{7}=\generadopor{a}\generadopor{b}\generadopor{c}         \cong         \generadopor{a}         \times         \generadopor{b}         \times         \generadopor{c}         \cong         \pdtrescyclics{4}{2}{2}$

En este grupo todos los subgrupos son normales, y sabiendo que el centro del grupo es el propio grupo, todos son centrales, sabiendo que el subgrupo derivado es el minimal, ninguno no trivial est\'a contenido en el derivado:

\[
\centro{{\mathsf{MW}}_{7}} = {\mathsf{MW}}_{7}
\]

y el subgrupo derivado es

\[
{{\mathsf{MW}}_{7}}^{\prime} = \trivialgr
\]

Estamos en condiciones de clasificar los subgrupos c{\'\i}clicos y en relaci\'on con el centro y el derivado:


\begin{flalign*}
&\mathcal{L}_1\left({{\mathsf{MW}}_{7}}\right)=\left\lbrace\right.\\
&\quad {\left\lbrace{1}\right\rbrace} ,
{\left\lbrace{1,a^2}\right\rbrace} ,
{\left\lbrace{1,b}\right\rbrace} ,\\
&\quad {\left\lbrace{1,a^2b}\right\rbrace} ,
{\left\lbrace{1,c}\right\rbrace} ,
{\left\lbrace{1,a^2c}\right\rbrace} ,\\
&\quad {\left\lbrace{1,bc}\right\rbrace} ,
{\left\lbrace{1,a^2bc}\right\rbrace} ,\\
&\quad {\left\lbrace{1,a,a^2,a^3}\right\rbrace} ,
{\left\lbrace{1,a^2,ab,a^3b}\right\rbrace} ,\\
&\quad {\left\lbrace{1,a^2,ac,a^3c}\right\rbrace} ,
{\left\lbrace{1,a^2,abc,a^3bc}\right\rbrace}\\
&\left.\right\rbrace\\
&\left[c(1,7),c(4,4)\right]
\end{flalign*}

Ahora los subgrupos generados por dos elementos, siguiendo el mismo patr\'on anterior,
sacaremos un nuevo invariante (vac{\'\i}o en este caso):

\begin{flalign*}
&\mathcal{L}_{2}\left({{\mathsf{MW}}_{7}}\right)=\left\lbrace\right.\\
&\quad {\left\lbrace{1,a^2,b,a^2b}\right\rbrace} ,
{\left\lbrace{1,a^2,c,a^2c}\right\rbrace} ,\\
&\quad {\left\lbrace{1,a^2,bc,a^2bc}\right\rbrace} ,
{\left\lbrace{1,b,c,bc}\right\rbrace} ,\\
&\quad {\left\lbrace{1,b,a^2c,a^2bc}\right\rbrace} ,
{\left\lbrace{1,a^2b,c,a^2bc}\right\rbrace} ,\\
&\quad {\left\lbrace{1,a^2b,a^2c,bc}\right\rbrace} ,\\
&\quad {\left\lbrace{1,a,a^2,a^3,b,ab,a^2b,a^3b}\right\rbrace} ,\\
&\quad {\left\lbrace{1,a,a^2,a^3,c,ac,a^2c,a^3c}\right\rbrace} ,\\
&\quad {\left\lbrace{1,a,a^2,a^3,bc,abc,a^2bc,a^3bc}\right\rbrace} ,\\
&\quad {\left\lbrace{1,a^2,b,a^2b,ac,a^3c,abc,a^3bc}\right\rbrace} ,\\
&\quad {\left\lbrace{1,a^2,ab,a^3b,c,a^2c,abc,a^3bc}\right\rbrace} ,\\
&\quad {\left\lbrace{1,a^2,ab,a^3b,ac,a^3c,bc,a^2bc}\right\rbrace}\\
\left.\right\rbrace\\
&\left[c(7,4),c(6,\pdcyclics{4}{2})\right]
\end{flalign*}

Queda a\'un un subgrupo que es generado por los tres elementos: el propio grupo completo.

\begin{flalign*}
&\mathcal{L}_{2}\left({{\mathsf{MW}}_{7}}\right)=\set{{\mathsf{MW}}_{7}}\\
&\left[c(1,16)\right]
\end{flalign*}

Serie central ascendente (comienza en el $\trivialgr$,
y en general no tiene porqu\'e alcanzar ${\mathsf{MW}}_{7}$):


\begin{flalign*}
&{\trivialgr} \normsubgr
{{\mathsf{MW}}_{7}}
\end{flalign*}


Serie central descendente (comienza en el ${\mathsf{MW}}_{7}$,
y en general no tiene proqu\'e alcanzar el $\trivialgr$):


\begin{flalign*}
&{{\mathsf{MW}}_{7}}\triangleright
{\trivialgr}
\end{flalign*}


El invariante que derivamos con lo visto hasta aq{\'\i}:

\begin{flalign*}
&\left[1,7,8,0,0\right]\\
&\left[
\begin{array}{l}
\left[
c(1,7),c(4,4)
\right]\\
\left[
c(7,4),c(6,\pdcyclics{4}{2})
\right]\\
\left[
c(1,16)
\right]\\
\left[\right]
\end{array}
\right]\\
&\left[
\begin{array} {c|c|c}
\mathtt{nil}_1 & \trivialgr & {\mathsf{MW}}_{7} \\
{} & {\mathsf{MW}}_{7} & \trivialgr
\end{array}
\right]\\
&\left[
\begin{array}{l}
\left[{\mathsf{MW}}_{7}=\centro{{\mathsf{MW}}_{7}}\supsetneqq{{\mathsf{MW}}_{7}}^{\prime}         \cong         \trivialgr\right]\\
\left[\right]
\end{array}
\right]
\end{flalign*}



En cuanto al ret{\'\i}culo de subgrupos normales y nilpotentes es claro que son todos los que son propios y
no triviales, y como hay tres subgrupos minimales que no tienen en com\'un m\'as que la unidad,
el subgrupo de Frattini se hace minimal, y como hay m\'as de un maximal que generan el grupo completo,
el subgrupo de Fitting es el completo.



\begin{flalign*}
&\left[1,7,8,0,0\right]\\
&\left[
\begin{array}{l}
\left[
c(1,7),c(4,4)
\right]\\
\left[
c(7,4),c(6,\pdcyclics{4}{2})
\right]\\
\left[
c(1,16)
\right]\\
\left[\right]
\end{array}
\right]\\
&\left[
\begin{array} {c|c|c}
\mathtt{nil}_1 & \trivialgr & {\mathsf{MW}}_{7} \\
{} & {\mathsf{MW}}_{7} & \trivialgr
\end{array}
\right]\\
&\left[
\begin{array}{l}
\left[{\mathsf{MW}}_{7}=\centro{{\mathsf{MW}}_{7}}\supsetneqq{{\mathsf{MW}}_{7}}^{\prime}         \cong         \trivialgr\right]\\
\left[\right]
\end{array}
\right]\\
&\left[
\begin{array}{l}
\left[
\nicefrac{{\mathsf{MW}}_{7}}{\centro{{\mathsf{MW}}_{7}}}
         \cong         {\trivialgr},\nicefrac{{\mathsf{MW}}_{7}}{{{\mathsf{MW}}_{7}}^{\prime}}         \cong         \pdtrescyclics{4}{2}{2}
\right]\\
\left[\right]
\end{array}
\right]\\
&\left[
\mathbf{\Phi }\entreparentesis{{\mathsf{MW}}_{7}}
= \set{1},\mathbf{\mathsf{F}}\entreparentesis{{\mathsf{MW}}_{7}}=
{{\mathsf{MW}}_{7}}
\right]
\end{flalign*}



El siguiente grupo ser\'a, ${\mathsf{MW}}_{8}$.



\smallskip\noindent
\resizebox{\linewidth}{!}{
\begin{tabular}{|c|c|c|c|c|c|c|c|c|c|c|c|c|c|c|c|}
\hline
$1$&$a$&$a^2$&$a^3$&$b$&$ab$&$a^2b$&$a^3b$&$c$&$ac$&$a^2c$&$a^3c$&$bc$&$abc$&$a^2bc$&$a^3bc$\\
\hline
$a$&$a^2$&$a^3$&$1$&$ab$&$a^2b$&$a^3b$&$b$&$ac$&$a^2c$&$a^3c$&$c$&$abc$&$a^2bc$&$a^3bc$&$bc$\\
\hline
$a^2$&$a^3$&$1$&$a$&$a^2b$&$a^3b$&$b$&$ab$&$a^2c$&$a^3c$&$c$&$ac$&$a^2bc$&$a^3bc$&$bc$&$abc$\\
\hline
$a^3$&$1$&$a$&$a^2$&$a^3b$&$b$&$ab$&$a^2b$&$a^3c$&$c$&$ac$&$a^2c$&$a^3bc$&$bc$&$abc$&$a^2bc$\\
\hline
$b$&$ab$&$a^2b$&$a^3b$&$1$&$a$&$a^2$&$a^3$&$bc$&$abc$&$a^2bc$&$a^3bc$&$c$&$ac$&$a^2c$&$a^3c$\\
\hline
$ab$&$a^2b$&$a^3b$&$b$&$a$&$a^2$&$a^3$&$1$&$abc$&$a^2bc$&$a^3bc$&$bc$&$ac$&$a^2c$&$a^3c$&$c$\\
\hline
$a^2b$&$a^3b$&$b$&$ab$&$a^2$&$a^3$&$1$&$a$&$a^2bc$&$a^3bc$&$bc$&$abc$&$a^2c$&$a^3c$&$c$&$ac$\\
\hline
$a^3b$&$b$&$ab$&$a^2b$&$a^3$&$1$&$a$&$a^2$&$a^3bc$&$bc$&$abc$&$a^2bc$&$a^3c$&$c$&$ac$&$a^2c$\\
\hline
$c$&$a^3c$&$a^2c$&$ac$&$bc$&$a^3bc$&$a^2bc$&$abc$&$1$&$a^3$&$a^2$&$a$&$b$&$a^3b$&$a^2b$&$ab$\\
\hline
$ac$&$c$&$a^3c$&$a^2c$&$abc$&$bc$&$a^3bc$&$a^2bc$&$a$&$1$&$a^3$&$a^2$&$ab$&$b$&$a^3b$&$a^2b$\\
\hline
$a^2c$&$ac$&$c$&$a^3c$&$a^2bc$&$abc$&$bc$&$a^3bc$&$a^2$&$a$&$1$&$a^3$&$a^2b$&$ab$&$b$&$a^3b$\\
\hline
$a^3c$&$a^2c$&$ac$&$c$&$a^3bc$&$a^2bc$&$abc$&$bc$&$a^3$&$a^2$&$a$&$1$&$a^3b$&$a^2b$&$ab$&$b$\\
\hline
$bc$&$a^3bc$&$a^2bc$&$abc$&$c$&$a^3c$&$a^2c$&$ac$&$b$&$a^3b$&$a^2b$&$ab$&$1$&$a^3$&$a^2$&$a$\\
\hline
$abc$&$bc$&$a^3bc$&$a^2bc$&$ac$&$c$&$a^3c$&$a^2c$&$ab$&$b$&$a^3b$&$a^2b$&$a$&$1$&$a^3$&$a^2$\\
\hline
$a^2bc$&$abc$&$bc$&$a^3bc$&$a^2c$&$ac$&$c$&$a^3c$&$a^2b$&$ab$&$b$&$a^3b$&$a^2$&$a$&$1$&$a^3$\\
\hline
$a^3bc$&$a^2bc$&$abc$&$bc$&$a^3c$&$a^2c$&$ac$&$c$&$a^3b$&$a^2b$&$ab$&$b$&$a^3$&$a^2$&$a$&$1$\\
\hline
\end{tabular}
}


Este grupo no es abeliano.

\begin{flalign*}
&\left[1,11,4,0,0\right]
\end{flalign*}

Esto reduce las posibilidades a indica alg\'un producto parecido al anterior, esto es, $\pdtrescyclics{4}{2}{2}$, aunque al no ser abeliano, no puede serlo. Veamos qu\'e nos va apareciendo. Aunque el parecido es el dicho, veamos si se puede generar por dos elementos en vez de tres observando sus distintos subgrupos. Despu\'es de mostrar el centro del grupo y su derivado, comenzamos por los subgrupos c{\'\i}clicos.


\[
\centro{{\mathsf{MW}}_{8}} = {\left\lbrace{1,a^2,b,a^2b}\right\rbrace}
\]

Y el subgrupo derivado es


\[
{{\mathsf{MW}}_{8}}^{\prime} = {\left\lbrace{1,a^2}\right\rbrace}
\]


Estamos en condiciones de clasificar los subgrupos c{\'\i}clicos y en relaci\'on con el centro y el derivado:


\begin{flalign*}
&\mathcal{L}_1\left({{\mathsf{MW}}_{8}}\right)=\left\lbrace\right.\\
&\quad {\left\lbrace{1}\right\rbrace} ,
{\left\lbrace{1,a^2}\right\rbrace} ,%dc
{\left\lbrace{1,b}\right\rbrace} ,%c
{\left\lbrace{1,a^2b}\right\rbrace} ,\\%c
&\quad {\left\lbrace{1,c}\right\rbrace} ,%
{\left\lbrace{1,ac}\right\rbrace} ,%
{\left\lbrace{1,a^2c}\right\rbrace} ,%
{\left\lbrace{1,a^3c}\right\rbrace} ,\\%
&\quad {\left\lbrace{1,bc}\right\rbrace} ,%
{\left\lbrace{1,abc}\right\rbrace} ,%
{\left\lbrace{1,a^2bc}\right\rbrace} ,%
{\left\lbrace{1,a^3bc}\right\rbrace} ,\\%
&\quad {\left\lbrace{1,a,a^2,a^3}\right\rbrace} ,%n
{\left\lbrace{1,a^2,ab,a^3b}\right\rbrace}\\%n
&\left.\right\rbrace\\
&\left[(1,1),dc(1,2),c(2,2),(8,2),n(2,4)\right]
\end{flalign*}


Ahora los subgrupos generados por dos elementos, siguiendo el mismo patr\'on anterior,
sacaremos un nuevo invariante (vac{\'\i}o en este caso):



\begin{flalign*}
&\mathcal{L}_{2}\left({{\mathsf{MW}}_{8}}\right)=\left\lbrace\right.\\
&\quad {\left\lbrace{1,a^2,b,a^2b}\right\rbrace} ,
{\left\lbrace{1,a^2,c,a^2c}\right\rbrace} ,
{\left\lbrace{1,a^2,ac,a^3c}\right\rbrace} ,
{\left\lbrace{1,a^2,bc,a^2bc}\right\rbrace} ,\\
&\quad {\left\lbrace{1,a^2,b,a^2b}\right\rbrace} ,
{\left\lbrace{1,a^2,c,a^2c}\right\rbrace} ,
{\left\lbrace{1,a^2,ac,a^3c}\right\rbrace} ,
{\left\lbrace{1,a^2,bc,a^2bc}\right\rbrace} ,\\
&\quad {\left\lbrace{1,a^2,abc,a^3bc}\right\rbrace} ,
{\left\lbrace{1,b,c,bc}\right\rbrace} ,
{\left\lbrace{1,b,ac,abc}\right\rbrace} ,
{\left\lbrace{1,b,a^2c,a^2bc}\right\rbrace} ,\\
&\quad {\left\lbrace{1,b,a^3c,a^3bc}\right\rbrace} ,
{\left\lbrace{1,a^2b,c,a^2bc}\right\rbrace} ,
{\left\lbrace{1,a^2b,ac,a^3bc}\right\rbrace} ,
{\left\lbrace{1,a^2b,a^2c,bc}\right\rbrace} ,\\
&\quad {\left\lbrace{1,b,a^3c,a^3bc}\right\rbrace} ,
{\left\lbrace{1,a^2b,c,a^2bc}\right\rbrace} ,
{\left\lbrace{1,a^2b,ac,a^3bc}\right\rbrace} ,
{\left\lbrace{1,a^2b,a^2c,bc}\right\rbrace} ,\\
&\quad {\left\lbrace{1,a^2b,a^3c,abc}\right\rbrace} ,
{\left\lbrace{1,a,a^2,a^3,b,ab,a^2b,a^3b}\right\rbrace} ,\\
&\quad {\left\lbrace{1,a,a^2,a^3,c,ac,a^2c,a^3c}\right\rbrace} ,
{\left\lbrace{1,a,a^2,a^3,bc,abc,a^2bc,a^3bc}\right\rbrace} ,\\
&\quad {\left\lbrace{1,a^2,b,a^2b,c,a^2c,bc,a^2bc}\right\rbrace} ,
{\left\lbrace{1,a^2,b,a^2b,ac,a^3c,abc,a^3bc}\right\rbrace} ,\\
&\quad {\left\lbrace{1,a^2,ab,a^3b,c,a^2c,abc,a^3bc}\right\rbrace} ,
{\left\lbrace{1,a^2,ab,a^3b,ac,a^3c,bc,a^2bc}\right\rbrace} ,\\
&\quad {\mathsf{MW}}_{8}\left.\right\rbrace\\
&\left[na(5,4),a(7,4),na(3,8),n(5,8),(1,16)\right]
\end{flalign*}

Los \'ordenes de los generadores de la tabla son $o(a)=4$, $o(b)=o(c)=2$.
Pero el hecho que todo halla salido con dos generadores indica que se puede generar
mediante $2$ elementos del grupo solo.
Veamos que $ab=ba$, $ca=a^3c$, $bc=cb$. De aqu{\'\i} $\generadopor{a,c}         \cong         \dihgr{4}$.
Luego el grupo es describible como ${\mathsf{MW}}_{8}         \cong         \pd{\dihgr{4}}{\cyclicgr{2}}$

Serie central ascendente (comienza en el $\trivialgr$,
y en general no tiene porqu\'e alcanzar ${\mathsf{MW}}_{8}$):


\begin{flalign*}
&{\trivialgr} \normsubgr
{\left\lbrace{1,a^2,b,a^2b}\right\rbrace} \normsubgr
{{\mathsf{MW}}_{8}}
\end{flalign*}


Serie central descendente (comienza en el ${\mathsf{MW}}_{8}$,
y en general no tiene proqu\'e alcanzar el $\trivialgr$):


\begin{flalign*}
&{{\mathsf{MW}}_{8}}\triangleright
{\left\lbrace{1,a^2}\right\rbrace} \triangleright
{\trivialgr}
\end{flalign*}


El invariante que derivamos con lo visto hasta aq{\'\i}:


\begin{flalign*}
&\left[1,11,4,0,0\right]\\
&\left[
\begin{array}{l}
\left[
(1,1),dc(1,2),c(2,2),(8,2),n(2,4)
\right]\\
\left[
na(5,4),a(7,4),na(3,8),n(5,8),(1,16)
\right]\\
\left[\right]\\
\left[\right]
\end{array}
\right]\\
&\left[
\begin{array} {c|c|c|c}
\mathtt{nil}_2 & \trivialgr & {\left\lbrace{1,a^2,b,a^2b}\right\rbrace} & {\mathsf{MW}}_{8} \\
{} & {\mathsf{MW}}_{8} & {\left\lbrace{1,a^2}\right\rbrace} & \trivialgr
\end{array}
\right]\\
&\left[
\begin{array}{l}
\left[\pdcyclics{2}{2}         \cong         \centro{{\mathsf{MW}}_{8}}\supsetneqq{{\mathsf{MW}}_{8}}^{\prime}         \cong         \cyclicgr{2}\right]\\
\left[\right]
\end{array}
\right]
\end{flalign*}

Nos falta ver cu\'ales sean los cocientes por el centro y por el derivado.

\begin{flalign*}
&{{\mathtt{Aut}}_{\mathtt{Gr}}\entreparentesis{{\mathsf{MW}}_{8}}}         \cong         \nicefrac{{\mathsf{MW}}_{8}}{\centro{{\mathsf{MW}}_{8}}}         \cong         {\pdcyclics{2}{2}}\\
&{{\mathsf{MW}}_{\mathsf{8}}}^{\mathtt{ab}}=\nicefrac{{{\mathsf{MW}}_{\mathsf{8}}}}{{{\mathsf{MW}}_{8}}^{\prime}}         \cong         {\pdtrescyclics{2}{2}{2}}
\end{flalign*}

En cuanto al ret{\'\i}culo de subgrupos normales y nilpotentes, lo mostramos a continuaci\'on:

\begin{flalign*}
&\mathcal{L}_{norm,nil} = \left\lbrace\right.\\
&\quad {\left\lbrace{1}\right\rbrace},
{\left\lbrace{1,a^2}\right\rbrace},
{\left\lbrace{1,b}\right\rbrace},
{\left\lbrace{1,a^2b}\right\rbrace},
{\left\lbrace{1,a,a^2,a^3}\right\rbrace},\\
&\quad {\left\lbrace{1,a^2,ab,a^3b}\right\rbrace},
{\left\lbrace{1,a^2,b,a^2b}\right\rbrace},
{\left\lbrace{1,a^2,c,a^2c}\right\rbrace},
{\left\lbrace{1,a^2,ac,a^3c}\right\rbrace},\\
&\quad {\left\lbrace{1,a^2,bc,a^2bc}\right\rbrace},
{\left\lbrace{1,a^2,abc,a^3bc}\right\rbrace},\\
&\quad {\left\lbrace{1,a,a^2,a^3,b,ab,a^2b,a^3b}\right\rbrace},
{\left\lbrace{1,a,a^2,a^3,c,ac,a^2c,a^3c}\right\rbrace},\\
&\quad {\left\lbrace{1,a,a^2,a^3,bc,abc,a^2bc,a^3bc}\right\rbrace},
{\left\lbrace{1,a^2,b,a^2b,c,a^2c,bc,a^2bc}\right\rbrace},\\
&\quad {\left\lbrace{1,a^2,b,a^2b,ac,a^3c,abc,a^3bc}\right\rbrace},
{\left\lbrace{1,a^2,ab,a^3b,c,a^2c,abc,a^3bc}\right\rbrace},\\
&\quad {\left\lbrace{1,a^2,ab,a^3b,ac,a^3c,bc,a^2bc}\right\rbrace},
{{\mathsf{MW}}_{8}}\\
\left.\right\rbrace
\end{flalign*}

As{\'\i} tenemos que $\mathbf{\Phi }\entreparentesis{{\mathsf{MW}}_{8}}=\set{1}$
al haber varios minimales que solo comparten la unidad. A su vez el subgrupo de
Fitting es el completo ya que hay varios maximales, uno incluye $a$, otro $b$ y otro $c$.


\begin{flalign*}
&\left[1,11,4,0,0\right]\\
&\left[
\begin{array}{l}
\left[
(1,1),dc(1,2),c(2,2),(8,2),n(2,4)
\right]\\
\left[
na(5,4),a(7,4),na(3,8),n(5,8),(1,16)
\right]\\
\left[\right]\\
\left[\right]
\end{array}
\right]\\
&\left[
\begin{array} {c|c|c|c}
\mathtt{nil}_2 & \trivialgr & {\left\lbrace{1,a^2,b,a^2b}\right\rbrace} & {\mathsf{MW}}_{8} \\
{} & {\mathsf{MW}}_{8} & {\left\lbrace{1,a^2}\right\rbrace} & \trivialgr
\end{array}
\right]\\
&\left[
\begin{array}{l}
\left[\pdcyclics{2}{2}         \cong         \centro{{\mathsf{MW}}_{8}}\supsetneqq{{\mathsf{MW}}_{8}}^{\prime}         \cong         \cyclicgr{2}\right]\\
\left[\right]
\end{array}
\right]\\
&\left[
\begin{array}{l}
\left[
{{\mathtt{Aut}}_{\mathtt{Gr}}\entreparentesis{{\mathsf{MW}}_{8}}}         \cong         \nicefrac{{\mathsf{MW}}_{8}}{\centro{{\mathsf{MW}}_{8}}}         \cong         {\pdcyclics{2}{2}},
{{\mathsf{MW}}_{\mathsf{8}}}^{\mathtt{ab}}=\nicefrac{{{\mathsf{MW}}_{\mathsf{8}}}}{{{\mathsf{MW}}_{8}}^{\prime}}         \cong         {\pdtrescyclics{2}{2}{2}}
\right]\\
\left[\right]
\end{array}
\right]\\
&\left[
\mathbf{\Phi }\entreparentesis{{\mathsf{MW}}_{8}}
= \set{1},\mathbf{\mathsf{F}}\entreparentesis{{\mathsf{MW}}_{8}}=
{{\mathsf{MW}}_{8}}
\right]
\end{flalign*}




El siguiente grupo ser\'a, ${\mathsf{MW}}_{9}$.



\smallskip\noindent
\resizebox{\linewidth}{!}{
\begin{tabular}{|c|c|c|c|c|c|c|c|c|c|c|c|c|c|c|c|}
\hline
$1$&$a$&$a^2$&$a^3$&$b$&$ab$&$a^2b$&$a^3b$&$c$&$ac$&$a^2c$&$a^3c$&$bc$&$abc$&$a^2bc$&$a^3bc$\\
\hline
$a$&$a^2$&$a^3$&$1$&$ab$&$a^2b$&$a^3b$&$b$&$ac$&$a^2c$&$a^3c$&$c$&$abc$&$a^2bc$&$a^3bc$&$bc$\\
\hline
$a^2$&$a^3$&$1$&$a$&$a^2b$&$a^3b$&$b$&$ab$&$a^2c$&$a^3c$&$c$&$ac$&$a^2bc$&$a^3bc$&$bc$&$abc$\\
\hline
$a^3$&$1$&$a$&$a^2$&$a^3b$&$b$&$ab$&$a^2b$&$a^3c$&$c$&$ac$&$a^2c$&$a^3bc$&$bc$&$abc$&$a^2bc$\\
\hline
$b$&$ab$&$a^2b$&$a^3b$&$1$&$a$&$a^2$&$a^3$&$bc$&$abc$&$a^2bc$&$a^3bc$&$c$&$ac$&$a^2c$&$a^3c$\\
\hline
$ab$&$a^2b$&$a^3b$&$b$&$a$&$a^2$&$a^3$&$1$&$abc$&$a^2bc$&$a^3bc$&$bc$&$ac$&$a^2c$&$a^3c$&$c$\\
\hline
$a^2b$&$a^3b$&$b$&$ab$&$a^2$&$a^3$&$1$&$a$&$a^2bc$&$a^3bc$&$bc$&$abc$&$a^2c$&$a^3c$&$c$&$ac$\\
\hline
$a^3b$&$b$&$ab$&$a^2b$&$a^3$&$1$&$a$&$a^2$&$a^3bc$&$bc$&$abc$&$a^2bc$&$a^3c$&$c$&$ac$&$a^2c$\\
\hline
$c$&$abc$&$a^2c$&$a^3bc$&$bc$&$ac$&$a^2bc$&$a^3c$&$1$&$ab$&$a^2$&$a^3b$&$b$&$a$&$a^2b$&$a^3$\\
\hline
$ac$&$a^2bc$&$a^3c$&$bc$&$abc$&$a^2c$&$a^3bc$&$c$&$a$&$a^2b$&$a^3$&$b$&$ab$&$a^2$&$a^3b$&$1$\\
\hline
$a^2c$&$a^3bc$&$c$&$abc$&$a^2bc$&$a^3c$&$bc$&$ac$&$a^2$&$a^3b$&$1$&$ab$&$a^2b$&$a^3$&$b$&$a$\\
\hline
$a^3c$&$bc$&$ac$&$a^2bc$&$a^3bc$&$c$&$abc$&$a^2c$&$a^3$&$b$&$a$&$a^2b$&$a^3b$&$1$&$ab$&$a^2$\\
\hline
$bc$&$ac$&$a^2bc$&$a^3c$&$c$&$abc$&$a^2c$&$a^3bc$&$b$&$a$&$a^2b$&$a^3$&$1$&$ab$&$a^2$&$a^3b$\\
\hline
$abc$&$a^2c$&$a^3bc$&$c$&$ac$&$a^2bc$&$a^3c$&$bc$&$ab$&$a^2$&$a^3b$&$1$&$a$&$a^2b$&$a^3$&$b$\\
\hline
$a^2bc$&$a^3c$&$bc$&$ac$&$a^2c$&$a^3bc$&$c$&$abc$&$a^2b$&$a^3$&$b$&$a$&$a^2$&$a^3b$&$1$&$ab$\\
\hline
$a^3bc$&$c$&$abc$&$a^2c$&$a^3c$&$bc$&$ac$&$a^2bc$&$a^3b$&$1$&$ab$&$a^2$&$a^3$&$b$&$a$&$a^2b$\\
\hline
\end{tabular}
}


Este grupo no es abeliano.

\begin{flalign*}
&\left[1,7,8,0,0\right]
\end{flalign*}

Despu\'es de mostrar el centro del grupo y su derivado, comenzamos por los subgrupos c{\'\i}clicos.


\[
\centro{{\mathsf{MW}}_{9}} = {\left\lbrace{1,a^2,b,a^2b}\right\rbrace}
\]

Y el subgrupo derivado es


\[
{{\mathsf{MW}}_{9}}^{\prime} = {\left\lbrace{1,b}\right\rbrace}
\]


Estamos en condiciones de clasificar los subgrupos c{\'\i}clicos y en relaci\'on con el centro y el derivado:


\begin{flalign*}
&\mathcal{L}_1\left({{\mathsf{MW}}_{9}}\right)=\left\lbrace\right.\\
&\quad {\left\lbrace{1}\right\rbrace} ,
{\left\lbrace{1,a^2}\right\rbrace} ,%c
{\left\lbrace{1,b}\right\rbrace} ,%dc
{\left\lbrace{1,a^2b}\right\rbrace} ,%c
{\left\lbrace{1,c}\right\rbrace} ,%
{\left\lbrace{1,a^2c}\right\rbrace} ,\\%
&\quad {\left\lbrace{1,bc}\right\rbrace} ,%
{\left\lbrace{1,a^2bc}\right\rbrace} ,%
{\left\lbrace{1,a,a^2,a^3}\right\rbrace} ,%
{\left\lbrace{1,a^2,ab,a^3b}\right\rbrace} ,\\%
&\quad {\left\lbrace{1,a^2b,ac,a^3bc}\right\rbrace} ,%
{\left\lbrace{1,a^2b,a^3c,abc}\right\rbrace}\\%
&\left.\right\rbrace\\
&\left[(1,1),dc(1,2),c(2,2),(4,2),(4,4)\right]
\end{flalign*}


Ahora los subgrupos generados por dos elementos, siguiendo el mismo patr\'on anterior,
sacaremos un nuevo invariante (vac{\'\i}o en este caso):


\begin{flalign*}
&\mathcal{L}_{2}\left({{\mathsf{MW}}_{9}}\right)=\left\lbrace\right.\\
&\quad {\left\lbrace{1,a^2,b,a^2b}\right\rbrace} ,%c
{\left\lbrace{1,a^2,c,a^2c}\right\rbrace} ,%
{\left\lbrace{1,a^2,bc,a^2bc}\right\rbrace} ,%
{\left\lbrace{1,b,c,bc}\right\rbrace} ,\\%n
&\quad {\left\lbrace{1,b,a^2c,a^2bc}\right\rbrace} ,%n
{\left\lbrace{1,a^2b,c,a^2bc}\right\rbrace} ,%
{\left\lbrace{1,a^2b,a^2c,bc}\right\rbrace} ,\\%
&\quad {\left\lbrace{1,a,a^2,a^3,b,ab,a^2b,a^3b}\right\rbrace} ,\\%na
&\quad {\left\lbrace{1,a^2,b,a^2b,c,a^2c,bc,a^2bc}\right\rbrace} ,\\%na
&\quad {\left\lbrace{1,a^2,b,a^2b,ac,a^3c,abc,a^3bc}\right\rbrace} ,\\%na
&\quad {\mathsf{MW}}_{9}\\
&\left.\right\rbrace\\
&\left[c(1,4),na(1,4),a(4,4),na(5,\pdcyclics{4}{2}),(1,16)\right]
\end{flalign*}

Los \'ordenes de los generadores de la tabla son $o(a)=4$, $o(b)=o(c)=2$.
Podemos observar que todos sus subgrupo propios son abelianos.
Pero el hecho que todo halla salido con dos generadores indica que se puede generar
mediante $2$ elementos del grupo solo.
Vemos que $ab=ba$, $ca=abc$, $bc=cb$. De esto, $ab=cac$.
Luego el grupo es describible como ${\mathsf{MW}}_{9}         \cong         \generadopor{a,b,c\;∣\;bx=xb,a^2x=xa^2,ab=cac}$

Serie central ascendente (comienza en el $\trivialgr$,
y en general
 no tiene porqu\'e alcanzar ${\mathsf{MW}}_{9}$):


\begin{flalign*}
&{\trivialgr} \normsubgr
{\left\lbrace{1,a^2,b,a^2b}\right\rbrace} \normsubgr
{{\mathsf{MW}}_{9}}
\end{flalign*}


Serie central descendente (comienza en el ${\mathsf{MW}}_{9}$,
y en general no tiene proqu\'e alcanzar el $\trivialgr$):


\begin{flalign*}
&{{\mathsf{MW}}_{9}}\triangleright
{\left\lbrace{1,b}\right\rbrace} \triangleright
{\trivialgr}
\end{flalign*}


El invariante que anotamos con lo visto hasta aqu{\'\i}:


\begin{flalign*}
&\left[1,7,8,0,0\right]\\
&\left[
\begin{array}{l}
\left[
(1,1),dc(1,2),c(2,2),(4,2),(4,4)
\right]\\
\left[
c(1,4),na(1,4),a(4,4),na(5,\pdcyclics{4}{2}),(1,16)
\right]\\
\left[\right]\\
\left[\right]
\end{array}
\right]\\
&\left[
\begin{array} {c|c|c|c}
\mathtt{nil}_2 & \trivialgr & {\left\lbrace{1,a^2,b,a^2b}\right\rbrace} & {\mathsf{MW}}_{9} \\
{} & {\mathsf{MW}}_{9} & {\left\lbrace{1,b}\right\rbrace} & \trivialgr
\end{array}
\right]\\
&\left[
\begin{array}{l}
\left[\pdcyclics{2}{2}         \cong         \centro{{\mathsf{MW}}_{9}}\supsetneqq{{\mathsf{MW}}_{9}}^{\prime}         \cong         \cyclicgr{2}\right]\\
\left[\right]
\end{array}
\right]
\end{flalign*}

Nos falta ver cu\'ales sean los cocientes por el centro y por el derivado.

\begin{flalign*}
&{{\mathtt{Aut}}_{\mathtt{Gr}}\entreparentesis{{\mathsf{MW}}_{9}}}         \cong         \nicefrac{{\mathsf{MW}}_{9}}{\centro{{\mathsf{MW}}_{9}}}         \cong         {\pdcyclics{2}{2}}\\
&{{\mathsf{MW}}_{9}}^{\mathtt{ab}}=\nicefrac{{{\mathsf{MW}}_{9}}}{{{\mathsf{MW}}_{9}}^{\prime}}         \cong         {\pdcyclics{4}{2}}
\end{flalign*}

En cuanto al ret{\'\i}culo de subgrupos normales y nilpotentes, lo mostramos a continuaci\'on:

\begin{flalign*}
&\mathcal{L}_{norm,nil} = \left\lbrace\right.\\
&\quad \set{1},{\left\lbrace{1,a^2}\right\rbrace},
{\left\lbrace{1,b}\right\rbrace},
{\left\lbrace{1,a^2b}\right\rbrace},
{\left\lbrace{1,a^2,b,a^2b}\right\rbrace},\\
&\quad {\left\lbrace{1,b,c,bc}\right\rbrace},
{\left\lbrace{1,b,a^2c,a^2bc}\right\rbrace},
{\left\lbrace{1,a,a^2,a^3,b,ab,a^2b,a^3b}\right\rbrace},\\
&\quad {\left\lbrace{1,a^2,b,a^2b,c,a^2c,bc,a^2bc}\right\rbrace},\\
&\quad {\left\lbrace{1,a^2,b,a^2b,ac,a^3c,abc,a^3bc}\right\rbrace},\\
&\quad {{\mathsf{MW}}_{9}}\\
\left.\right\rbrace
\end{flalign*}

As{\'\i} tenemos que $\mathbf{\Phi }\entreparentesis{{\mathsf{MW}}_{9}}=\set{1}$
al haber varios minimales que solo comparten la unidad. A su vez el subgrupo de
Fitting es el completo ya que hay varios maximales, que incluyen $11$ elementos
distintos parte de los cuales producen un subgrupo de orden $8$ y la otra parte
produce otro subgrupo de orden $8$ que aunque no son disjuntos no son la uni\'on.
Como han de producir un subgrupo mayor que ambos, solo queda que genere el completo.


\begin{flalign*}
&\left[1,7,8,0,0\right]\\
&\left[
\begin{array}{l}
\left[
(1,1),dc(1,2),c(2,2),(4,2),(4,4)
\right]\\
\left[
c(1,4),na(1,4),a(4,4),na(5,\pdcyclics{4}{2}),(1,16)
\right]\\
\left[\right]\\
\left[\right]
\end{array}
\right]\\
&\left[
\begin{array} {c|c|c|c}
\mathtt{nil}_2 & \trivialgr & {\left\lbrace{1,a^2,b,a^2b}\right\rbrace} & {\mathsf{MW}}_{9} \\
{} & {\mathsf{MW}}_{9} & {\left\lbrace{1,b}\right\rbrace} & \trivialgr
\end{array}
\right]\\
&\left[
\begin{array}{l}
\left[\pdcyclics{2}{2}         \cong         \centro{{\mathsf{MW}}_{9}}\supsetneqq{{\mathsf{MW}}_{9}}^{\prime}         \cong         \cyclicgr{2}\right]\\
\left[\right]
\end{array}
\right]\\
&\left[
\begin{array}{l}
\left[
{{\mathtt{Aut}}_{\mathtt{Gr}}\entreparentesis{{\mathsf{MW}}_{9}}}         \cong         \nicefrac{{\mathsf{MW}}_{9}}{\centro{{\mathsf{MW}}_{9}}}         \cong         {\pdcyclics{2}{2}},
{{\mathsf{MW}}_{9}}^{\mathtt{ab}}=\nicefrac{{{\mathsf{MW}}_{9}}}{{{\mathsf{MW}}_{9}}^{\prime}}         \cong         {\pdcyclics{4}{2}}
\right]\\
\left[\right]
\end{array}
\right]\\
&\left[
\mathbf{\Phi }\entreparentesis{{\mathsf{MW}}_{9}}
= \set{1},\mathbf{\mathsf{F}}\entreparentesis{{\mathsf{MW}}_{9}}=
{{\mathsf{MW}}_{9}}
\right]
\end{flalign*}


El siguiente grupo ser\'a, ${\mathsf{MW}}_{10}$.



\smallskip\noindent
\resizebox{\linewidth}{!}{
\begin{tabular}{|c|c|c|c|c|c|c|c|c|c|c|c|c|c|c|c|}
\hline
$1$&$a$&$a^2$&$a^3$&$b$&$ab$&$a^2b$&$a^3b$&$c$&$ac$&$a^2c$&$a^3c$&$bc$&$abc$&$a^2bc$&$a^3bc$\\
\hline
$a$&$a^2$&$a^3$&$1$&$ab$&$a^2b$&$a^3b$&$b$&$ac$&$a^2c$&$a^3c$&$c$&$abc$&$a^2bc$&$a^3bc$&$bc$\\
\hline
$a^2$&$a^3$&$1$&$a$&$a^2b$&$a^3b$&$b$&$ab$&$a^2c$&$a^3c$&$c$&$ac$&$a^2bc$&$a^3bc$&$bc$&$abc$\\
\hline
$a^3$&$1$&$a$&$a^2$&$a^3b$&$b$&$ab$&$a^2b$&$a^3c$&$c$&$ac$&$a^2c$&$a^3bc$&$bc$&$abc$&$a^2bc$\\
\hline
$b$&$ab$&$a^2b$&$a^3b$&$1$&$a$&$a^2$&$a^3$&$bc$&$abc$&$a^2bc$&$a^3bc$&$c$&$ac$&$a^2c$&$a^3c$\\
\hline
$ab$&$a^2b$&$a^3b$&$b$&$a$&$a^2$&$a^3$&$1$&$abc$&$a^2bc$&$a^3bc$&$bc$&$ac$&$a^2c$&$a^3c$&$c$\\
\hline
$a^2b$&$a^3b$&$b$&$ab$&$a^2$&$a^3$&$1$&$a$&$a^2bc$&$a^3bc$&$bc$&$abc$&$a^2c$&$a^3c$&$c$&$ac$\\
\hline
$a^3b$&$b$&$ab$&$a^2b$&$a^3$&$1$&$a$&$a^2$&$a^3bc$&$bc$&$abc$&$a^2bc$&$a^3c$&$c$&$ac$&$a^2c$\\
\hline
$c$&$a^3c$&$a^2c$&$ac$&$a^2bc$&$abc$&$bc$&$a^3bc$&$1$&$a^3$&$a^2$&$a$&$a^2b$&$ab$&$b$&$a^3b$\\
\hline
$ac$&$c$&$a^3c$&$a^2c$&$a^3bc$&$a^2bc$&$abc$&$bc$&$a$&$1$&$a^3$&$a^2$&$a^3b$&$a^2b$&$ab$&$b$\\
\hline
$a^2c$&$ac$&$c$&$a^3c$&$bc$&$a^3bc$&$a^2bc$&$abc$&$a^2$&$a$&$1$&$a^3$&$b$&$a^3b$&$a^2b$&$ab$\\
\hline
$a^3c$&$a^2c$&$ac$&$c$&$abc$&$bc$&$a^3bc$&$a^2bc$&$a^3$&$a^2$&$a$&$1$&$ab$&$b$&$a^3b$&$a^2b$\\
\hline
$bc$&$a^3bc$&$a^2bc$&$abc$&$a^2c$&$ac$&$c$&$a^3c$&$b$&$a^3b$&$a^2b$&$ab$&$a^2$&$a$&$1$&$a^3$\\
\hline
$abc$&$bc$&$a^3bc$&$a^2bc$&$a^3c$&$a^2c$&$ac$&$c$&$ab$&$b$&$a^3b$&$a^2b$&$a^3$&$a^2$&$a$&$1$\\
\hline
$a^2bc$&$abc$&$bc$&$a^3bc$&$c$&$a^3c$&$a^2c$&$ac$&$a^2b$&$ab$&$b$&$a^3b$&$1$&$a^3$&$a^2$&$a$\\
\hline
$a^3bc$&$a^2bc$&$abc$&$bc$&$ac$&$c$&$a^3c$&$a^2c$&$a^3b$&$a^2b$&$ab$&$b$&$a$&$1$&$a^3$&$a^2$\\
\hline
\end{tabular}
}


Este grupo no es abeliano.

\begin{flalign*}
&\left[1,7,8,0,0\right]
\end{flalign*}

Despu\'es de mostrar el centro del grupo y su derivado, comenzamos por los subgrupos c{\'\i}clicos.


\[
\centro{{\mathsf{MW}}_{10}} = {\left\lbrace{1,a^2,ab,a^3b}\right\rbrace}          \cong          \pdcyclics{2}{2}
\]

Y el subgrupo derivado es


\[
{{\mathsf{MW}}_{10}}^{\prime} = {\left\lbrace{1,a^2}\right\rbrace}          \cong          \cyclicgr{2}
\]


Estamos en condiciones de clasificar los subgrupos c{\'\i}clicos y en relaci\'on con el centro y el derivado:



\begin{flalign*}
&\mathcal{L}_1\left({{\mathsf{MW}}_{10}}\right)=\left\lbrace\right.\\
&\quad {\left\lbrace{1}\right\rbrace} ,
{\left\lbrace{1,a^2}\right\rbrace} ,%dc
{\left\lbrace{1,b}\right\rbrace} ,%n
{\left\lbrace{1,a^2b}\right\rbrace} ,%
{\left\lbrace{1,c}\right\rbrace} ,%
{\left\lbrace{1,ac}\right\rbrace} ,\\%
&\quad {\left\lbrace{1,a^2c}\right\rbrace} ,%
{\left\lbrace{1,a^3c}\right\rbrace} ,%
{\left\lbrace{1,a,a^2,a^3}\right\rbrace} ,%
{\left\lbrace{1,a^2,ab,a^3b}\right\rbrace} ,\\%c
&\quad {\left\lbrace{1,a^2,bc,a^2bc}\right\rbrace} ,%n
{\left\lbrace{1,a^2,abc,a^3bc}\right\rbrace}\\%n
&\left.\right\rbrace\\
&\left[(1,1),dc(1,2),n(1,2),(6,2),c(1,4),n(2,4)\right]
\end{flalign*}


Ahora los subgrupos generados por dos elementos, siguiendo el mismo patr\'on anterior,
sacaremos un nuevo invariante (vac{\'\i}o en este caso):


\begin{flalign*}
&\mathcal{L}_{2}\left({{\mathsf{MW}}_{10}}\right)=\left\lbrace\right.\\
&\quad {\left\lbrace{1,a^2,b,a^2b}\right\rbrace} ,%na
{\left\lbrace{1,a^2,c,a^2c}\right\rbrace} ,%na
{\left\lbrace{1,a^2,ac,a^3c}\right\rbrace} ,\\%na
&\quad {\left\lbrace{1,a,a^2,a^3,b,ab,a^2b,a^3b}\right\rbrace} ,%na
{\left\lbrace{1,a,a^2,a^3,c,ac,a^2c,a^3c}\right\rbrace} ,\\%n
&\quad {\left\lbrace{1,a,a^2,a^3,bc,abc,a^2bc,a^3bc}\right\rbrace} ,%n
{\left\lbrace{1,a^2,b,a^2b,c,a^2c,bc,a^2bc}\right\rbrace} ,\\%n
&\quad {\left\lbrace{1,a^2,b,a^2b,ac,a^3c,abc,a^3bc}\right\rbrace} ,%n
{\left\lbrace{1,a^2,ab,a^3b,c,a^2c,abc,a^3bc}\right\rbrace} ,\\%na
&\quad {\left\lbrace{1,a^2,ab,a^3b,ac,a^3c,bc,a^2bc}\right\rbrace} ,\\%na
&\quad {\mathsf{MW}}_{10}\\
&\left.\right\rbrace\\
&\left[c(1,4),na(1,4),a(4,4),na(5,\pdcyclics{4}{2}),(1,16)\right]
\end{flalign*}

Los \'ordenes de los generadores de la tabla son $o(a)=4$, $o(b)=o(c)=2$.
Podemos observar que todos sus subgrupo propios son abelianos.
Pero el hecho que todo halla salido con dos generadores indica que se puede generar
mediante $2$ elementos del grupo solo.
Vemos que $ab=ba$, $ca=abc$, $bc=cb$. De esto, $ab=cac$.
Luego el grupo es describible como ${\mathsf{MW}}_{10}         \cong         \generadopor{a,b,c\;∣\;bx=xb,a^2x=xa^2,ab=cac}$

Serie central ascendente (comienza en el $\trivialgr$,
y en general
 no tiene porqu\'e alcanzar ${\mathsf{MW}}_{10}$):


\begin{flalign*}
&{\trivialgr} \normsubgr
{\left\lbrace{1,a^2,b,a^2b}\right\rbrace} \normsubgr
{{\mathsf{MW}}_{10}}
\end{flalign*}


Serie central descendente (comienza en el ${\mathsf{MW}}_{10}$,
y en general no tiene proqu\'e alcanzar el $\trivialgr$):


\begin{flalign*}
&{{\mathsf{MW}}_{10}}\triangleright
{\left\lbrace{1,b}\right\rbrace} \triangleright
{\trivialgr}
\end{flalign*}


El invariante que anotamos con lo visto hasta aqu{\'\i}:


\begin{flalign*}
&\left[1,7,8,0,0\right]\\
&\left[
\begin{array}{l}
\left[
(1,1),dc(1,2),c(2,2),(4,2),(4,4)
\right]\\
\left[
c(1,4),na(1,4),a(4,4),na(5,\pdcyclics{4}{2}),(1,16)
\right]\\
\left[\right]\\
\left[\right]
\end{array}
\right]\\
&\left[
\begin{array} {c|c|c|c}
\mathtt{nil}_2 & \trivialgr & {\left\lbrace{1,a^2,b,a^2b}\right\rbrace} & {\mathsf{MW}}_{10} \\
{} & {\mathsf{MW}}_{10} & {\left\lbrace{1,b}\right\rbrace} & \trivialgr
\end{array}
\right]\\
&\left[
\begin{array}{l}
\left[\pdcyclics{2}{2}         \cong         \centro{{\mathsf{MW}}_{10}}\supsetneqq{{\mathsf{MW}}_{10}}^{\prime}         \cong         \cyclicgr{2}\right]\\
\left[\right]
\end{array}
\right]
\end{flalign*}

Nos falta ver cu\'ales sean los cocientes por el centro y por el derivado.

\begin{flalign*}
&{{\mathtt{Aut}}_{\mathtt{Gr}}\entreparentesis{{\mathsf{MW}}_{10}}}         \cong         \nicefrac{{\mathsf{MW}}_{10}}{\centro{{\mathsf{MW}}_{10}}}         \cong         {\pdcyclics{2}{2}}\\
&{{\mathsf{MW}}_{10}}^{\mathtt{ab}}=\nicefrac{{{\mathsf{MW}}_{10}}}{{{\mathsf{MW}}_{10}}^{\prime}}         \cong         {\pdcyclics{4}{2}}
\end{flalign*}

El ret{\'\i}culo de subgrupos maximales (no propios), lo mostramos a continuaci\'on:

\begin{flalign*}
&\mathcal{L}_{maximal} = \left\lbrace\right.\\
&\quad {\left\lbrace{1,a,a^2,a^3,b,ab,a^2b,a^3b}\right\rbrace} ,
{\left\lbrace{1,a,a^2,a^3,c,ac,a^2c,a^3c}\right\rbrace} ,\\
&\quad {\left\lbrace{1,a,a^2,a^3,bc,abc,a^2bc,a^3bc}\right\rbrace} ,
{\left\lbrace{1,a^2,b,a^2b,c,a^2c,bc,a^2bc}\right\rbrace} ,\\
&\quad {\left\lbrace{1,a^2,b,a^2b,ac,a^3c,abc,a^3bc}\right\rbrace} ,
{\left\lbrace{1,a^2,ab,a^3b,c,a^2c,abc,a^3bc}\right\rbrace} ,\\
&\quad {\left\lbrace{1,a^2,ab,a^3b,ac,a^3c,bc,a^2bc}\right\rbrace}\\
&\left.\right\rbrace\\
&\mathbf{\Phi }\entreparentesis{{\mathsf{MW}}_{10}}=\set{1,a^2}
\end{flalign*}

En cuanto al ret{\'\i}culo de subgrupos normales y nilpotentes, lo mostramos a continuaci\'on:


\begin{flalign*}
&{\mathcal{L}}_{norm,nil} = \left\lbrace\right.\\
&\quad \set{1},{\left\lbrace{1,a^2}\right\rbrace},
{\left\lbrace{1,a,a^2,a^3}\right\rbrace},
{\left\lbrace{1,a^2,ab,a^3b}\right\rbrace},
{\left\lbrace{1,a^2,bc,a^2bc}\right\rbrace},
{\left\lbrace{1,a^2,abc,a^3bc}\right\rbrace},\\
&\quad {\left\lbrace{1,a^2}\right\rbrace},
{\left\lbrace{1,a,a^2,a^3}\right\rbrace},
{\left\lbrace{1,a^2,ab,a^3b}\right\rbrace},
{\left\lbrace{1,a^2,bc,a^2bc}\right\rbrace},
{\left\lbrace{1,a^2,abc,a^3bc}\right\rbrace},\\
&\quad {\left\lbrace{1,a^2,b,a^2b}\right\rbrace},
{\left\lbrace{1,a^2,c,a^2c}\right\rbrace},
{\left\lbrace{1,a^2,ac,a^3c}\right\rbrace},
{\left\lbrace{1,a,a^2,a^3,b,ab,a^2b,a^3b}\right\rbrace},\\
&\quad {\left\lbrace{1,a,a^2,a^3,c,ac,a^2c,a^3c}\right\rbrace},
{\left\lbrace{1,a,a^2,a^3,bc,abc,a^2bc,a^3bc}\right\rbrace},
{\left\lbrace{1,a^2,b,a^2b,c,a^2c,bc,a^2bc}\right\rbrace},\\
&\quad {\left\lbrace{1,a^2,b,a^2b,ac,a^3c,abc,a^3bc}\right\rbrace},
{\left\lbrace{1,a^2,ab,a^3b,c,a^2c,abc,a^3bc}\right\rbrace},
{\left\lbrace{1,a^2,ab,a^3b,ac,a^3c,bc,a^2bc}\right\rbrace},\\
&\quad {{\mathsf{MW}}_{10}}\\
&\left.\right\rbrace\\
&\mathbf{\mathsf{F}}\entreparentesis{{\mathsf{MW}}_{10}}={\mathsf{MW}}_{10}
\end{flalign*}

As{\'\i} tenemos que $\frattini{{\mathsf{MW}}_{10}}=\set{1,a^2}$
al haber varios minimales que solo comparten la unidad. A su vez el subgrupo de
Fitting es el completo ya que hay varios maximales normales y nilpotentes de orden $8$,
que necesariamente
generar\'an el completo.


\begin{flalign*}
&\left[1,7,8,0,0\right]\\
&\left[
\begin{array}{l}
\left[
(1,1),dc(1,2),c(2,2),(4,2),(4,4)
\right]\\
\left[
c(1,4),na(1,4),a(4,4),na(5,\pdcyclics{4}{2}),(1,16)
\right]\\
\left[\right]\\
\left[\right]
\end{array}
\right]\\
&\left[
\begin{array} {c|c|c|c}
\mathtt{nil}_2 & \trivialgr & {\left\lbrace{1,a^2,b,a^2b}\right\rbrace} & {\mathsf{MW}}_{10} \\
{} & {\mathsf{MW}}_{10} & {\left\lbrace{1,b}\right\rbrace} & \trivialgr
\end{array}
\right]\\
&\left[
\begin{array}{l}
\left[\pdcyclics{2}{2}         \cong         \centro{{\mathsf{MW}}_{10}}\supsetneqq{{\mathsf{MW}}_{10}}^{\prime}         \cong         \cyclicgr{2}\right]\\
\left[\right]
\end{array}
\right]\\
&\left[
\begin{array}{l}
\left[
{{\mathtt{Aut}}_{\mathtt{Gr}}\entreparentesis{{\mathsf{MW}}_{10}}}         \cong         \nicefrac{{\mathsf{MW}}_{10}}{\centro{{\mathsf{MW}}_{10}}}         \cong         {\pdcyclics{2}{2}},
{{\mathsf{MW}}_{10}}^{\mathtt{ab}}=\nicefrac{{{\mathsf{MW}}_{10}}}{{{\mathsf{MW}}_{10}}^{\prime}}         \cong         {\pdcyclics{4}{2}}
\right]\\
\left[\right]
\end{array}
\right]\\
&\left[
\frattini{{\mathsf{MW}}_{10}}
= {{\mathsf{MW}}_{10}}^{\prime},\fitting{{\mathsf{MW}}_{10}}=
{{\mathsf{MW}}_{10}}
\right]
\end{flalign*}



El siguiente grupo ser\'a, ${\mathsf{MW}}_{11}$.



\smallskip\noindent
\resizebox{\linewidth}{!}{
\begin{tabular}{|c|c|c|c|c|c|c|c|c|c|c|c|c|c|c|c|}
\hline
$1$&$a$&$a^2$&$a^3$&$b$&$ab$&$a^2b$&$a^3b$&$c$&$ac$&$a^2c$&$a^3c$&$bc$&$abc$&$a^2bc$&$a^3bc$\\
\hline
$a$&$a^2$&$a^3$&$1$&$ab$&$a^2b$&$a^3b$&$b$&$ac$&$a^2c$&$a^3c$&$c$&$abc$&$a^2bc$&$a^3bc$&$bc$\\
\hline
$a^2$&$a^3$&$1$&$a$&$a^2b$&$a^3b$&$b$&$ab$&$a^2c$&$a^3c$&$c$&$ac$&$a^2bc$&$a^3bc$&$bc$&$abc$\\
\hline
$a^3$&$1$&$a$&$a^2$&$a^3b$&$b$&$ab$&$a^2b$&$a^3c$&$c$&$ac$&$a^2c$&$a^3bc$&$bc$&$abc$&$a^2bc$\\
\hline
$b$&$ab$&$a^2b$&$a^3b$&$1$&$a$&$a^2$&$a^3$&$bc$&$abc$&$a^2bc$&$a^3bc$&$c$&$ac$&$a^2c$&$a^3c$\\
\hline
$ab$&$a^2b$&$a^3b$&$b$&$a$&$a^2$&$a^3$&$1$&$abc$&$a^2bc$&$a^3bc$&$bc$&$ac$&$a^2c$&$a^3c$&$c$\\
\hline
$a^2b$&$a^3b$&$b$&$ab$&$a^2$&$a^3$&$1$&$a$&$a^2bc$&$a^3bc$&$bc$&$abc$&$a^2c$&$a^3c$&$c$&$ac$\\
\hline
$a^3b$&$b$&$ab$&$a^2b$&$a^3$&$1$&$a$&$a^2$&$a^3bc$&$bc$&$abc$&$a^2bc$&$a^3c$&$c$&$ac$&$a^2c$\\
\hline
$c$&$a^3c$&$a^2c$&$ac$&$bc$&$a^3bc$&$a^2bc$&$abc$&$a^2$&$a$&$1$&$a^3$&$a^2b$&$ab$&$b$&$a^3b$\\
\hline
$ac$&$c$&$a^3c$&$a^2c$&$abc$&$bc$&$a^3bc$&$a^2bc$&$a^3$&$a^2$&$a$&$1$&$a^3b$&$a^2b$&$ab$&$b$\\
\hline
$a^2c$&$ac$&$c$&$a^3c$&$a^2bc$&$abc$&$bc$&$a^3bc$&$1$&$a^3$&$a^2$&$a$&$b$&$a^3b$&$a^2b$&$ab$\\
\hline
$a^3c$&$a^2c$&$ac$&$c$&$a^3bc$&$a^2bc$&$abc$&$bc$&$a$&$1$&$a^3$&$a^2$&$ab$&$b$&$a^3b$&$a^2b$\\
\hline
$bc$&$a^3bc$&$a^2bc$&$abc$&$c$&$a^3c$&$a^2c$&$ac$&$a^2b$&$ab$&$b$&$a^3b$&$a^2$&$a$&$1$&$a^3$\\
\hline
$abc$&$bc$&$a^3bc$&$a^2bc$&$ac$&$c$&$a^3c$&$a^2c$&$a^3b$&$a^2b$&$ab$&$b$&$a^3$&$a^2$&$a$&$1$\\
\hline
$a^2bc$&$abc$&$bc$&$a^3bc$&$a^2c$&$ac$&$c$&$a^3c$&$b$&$a^3b$&$a^2b$&$ab$&$1$&$a^3$&$a^2$&$a$\\
\hline
$a^3bc$&$a^2bc$&$abc$&$bc$&$a^3c$&$a^2c$&$ac$&$c$&$ab$&$b$&$a^3b$&$a^2b$&$a$&$1$&$a^3$&$a^2$\\
\hline
\end{tabular}
}


Este grupo no es abeliano.

\begin{flalign*}
&\left[1,3,12,0,0\right]
\end{flalign*}

Despu\'es de mostrar el centro del grupo y su derivado, comenzamos por los subgrupos c{\'\i}clicos.


\[
\centro{{\mathsf{MW}}_{11}} = {\left\lbrace{1,a^2,b,a^2b}\right\rbrace}          \cong          \pdcyclics{2}{2}
\]

Y el subgrupo derivado es


\[
{{\mathsf{MW}}_{11}}^{\prime} = {\left\lbrace{1,a^2}\right\rbrace}          \cong          \cyclicgr{2}
\]


Estamos en condiciones de clasificar los subgrupos c{\'\i}clicos y en relaci\'on con el centro y el derivado:





\begin{flalign*}
&\mathcal{L}_1\left({{\mathsf{MW}}_{11}}\right)=\left\lbrace\right.\\
&\quad {\left\lbrace{1}\right\rbrace} ,
{\left\lbrace{1,a^2}\right\rbrace} ,%dc
{\left\lbrace{1,b}\right\rbrace} ,%c
{\left\lbrace{1,a^2b}\right\rbrace} ,%c
{\left\lbrace{1,a,a^2,a^3}\right\rbrace} ,\\%n
&\quad {\left\lbrace{1,a^2,c,a^2c}\right\rbrace} ,%n
{\left\lbrace{1,a^2,ac,a^3c}\right\rbrace} ,%n
{\left\lbrace{1,a^2,bc,a^2bc}\right\rbrace} ,\\%n
&{\left\lbrace{1,a^2,abc,a^3bc}\right\rbrace}\\%n
&\left.\right\rbrace\\
&\left[(1,1),dc(1,2),c(2,2),n(4,4)\right]
\end{flalign*}


Ahora los subgrupos generados por dos elementos, siguiendo el mismo patr\'on anterior,
sacaremos un nuevo invariante (vac{\'\i}o en este caso):






\begin{flalign*}
&\mathcal{L}_{2}\left({{\mathsf{MW}}_{11}}\right)=\left\lbrace\right.\\
&\quad {\left\lbrace{1,a^2,b,a^2b}\right\rbrace} ,%na
{\left\lbrace{1,a,a^2,a^3,b,ab,a^2b,a^3b}\right\rbrace} ,\\%na
&\quad {\left\lbrace{1,a,a^2,a^3,c,ac,a^2c,a^3c}\right\rbrace} ,%n
{\left\lbrace{1,a,a^2,a^3,bc,abc,a^2bc,a^3bc}\right\rbrace} ,\\%n
&\quad {\left\lbrace{1,a^2,b,a^2b,c,a^2c,bc,a^2bc}\right\rbrace} ,%na
{\left\lbrace{1,a^2,b,a^2b,ac,a^3c,abc,a^3bc}\right\rbrace} ,\\%na
&\quad {\left\lbrace{1,a^2,ab,a^3b,c,a^2c,abc,a^3bc}\right\rbrace} ,%n
{\left\lbrace{1,a^2,ab,a^3b,ac,a^3c,bc,a^2bc}\right\rbrace} , \\%n
&\quad {\mathsf{MW}}_{11}\\
&\left.\right\rbrace\\
&\left[na(1,4),na(3,8),n(2,\dihgr{4}),n(2,\quaterniongr{1}),(1,16)\right]
\end{flalign*}

Los \'ordenes de los generadores de la tabla son $o(a)=4$, $o(b)=2$ y $o(c)=4$.
Podemos observar que todos sus subgrupo propios son normales.
Vemos que $ab=ba$, $ca=a^3c$, $bc=cb$, $\forall x\; (xc)^2 = a^2$. De esto, $\forall x\; xcaxc = a$.
Luego el grupo es describible como ${\mathsf{MW}}_{11}         \cong         \generadopor{a,b,c\;∣\;a^4=b^2=c^4=1,ba=ab,bc=cb,xcaxc=a}$

Serie central ascendente (comienza en el $\trivialgr$,
y en general
 no tiene porqu\'e alcanzar ${\mathsf{MW}}_{11}$):


\begin{flalign*}
&{\trivialgr} \normsubgr
{\left\lbrace{1,a^2,b,a^2b}\right\rbrace} \normsubgr
{{\mathsf{MW}}_{11}}
\end{flalign*}


Serie central descendente (comienza en el ${\mathsf{MW}}_{11}$,
y en general no tiene proqu\'e alcanzar el $\trivialgr$):


\begin{flalign*}
&{{\mathsf{MW}}_{11}}\triangleright
{\left\lbrace{1,a^2}\right\rbrace} \triangleright
{\trivialgr}
\end{flalign*}


El invariante que anotamos con lo visto hasta aqu{\'\i}:


\begin{flalign*}
&\left[1,3,12,0,0\right]\\
&\left[
\begin{array}{l}
\left[
(1,1),dc(1,2),c(2,2),n(4,4)
\right]\\
\left[
na(1,4),na(3,8),n(2,\dihgr{4}),n(2,\quaterniongr{1}),(1,16)
\right]\\
\left[\right]\\
\left[\right]
\end{array}
\right]\\
&\left[
\begin{array} {c|c|c|c}
\mathtt{nil}_2 & \trivialgr & {\left\lbrace{1,a^2,b,a^2b}\right\rbrace} & {\mathsf{MW}}_{11} \\
{} & {\mathsf{MW}}_{11} & {\left\lbrace{1,a^2}\right\rbrace} & \trivialgr
\end{array}
\right]\\
&\left[
\begin{array}{l}
\left[\pdcyclics{2}{2}         \cong         \centro{{\mathsf{MW}}_{11}}\supsetneqq{{\mathsf{MW}}_{11}}^{\prime}         \cong         \cyclicgr{2}\right]\\
\left[\right]
\end{array}
\right]
\end{flalign*}

Nos falta ver cu\'ales sean los cocientes por el centro y por el derivado.

\begin{flalign*}
&\graut{{\mathsf{MW}}_{11}}         \cong         \nicefrac{{\mathsf{MW}}_{11}}{\centro{{\mathsf{MW}}_{11}}}         \cong         {\pdcyclics{2}{2}}\\
&{{\mathsf{MW}}_{11}}^{\mathtt{ab}}=\nicefrac{{{\mathsf{MW}}_{11}}}{{{\mathsf{MW}}_{11}}^{\prime}}         \cong         {\pdtrescyclics{2}{2}{2}}
\end{flalign*}

El ret{\'\i}culo de subgrupos maximales (no propios), lo mostramos a continuaci\'on:


\begin{flalign*}
&\mathcal{L}_{maximal} = \left\lbrace\right.\\
&\quad {\left\lbrace{1,a,a^2,a^3,b,ab,a^2b,a^3b}\right\rbrace},
{\left\lbrace{1,a,a^2,a^3,c,ac,a^2c,a^3c}\right\rbrace},\\
&\quad {\left\lbrace{1,a,a^2,a^3,bc,abc,a^2bc,a^3bc}\right\rbrace},
{\left\lbrace{1,a^2,b,a^2b,c,a^2c,bc,a^2bc}\right\rbrace},\\
&\quad {\left\lbrace{1,a^2,b,a^2b,ac,a^3c,abc,a^3bc}\right\rbrace},
{\left\lbrace{1,a^2,ab,a^3b,c,a^2c,abc,a^3bc}\right\rbrace},\\
&\quad {\left\lbrace{1,a^2,ab,a^3b,ac,a^3c,bc,a^2bc}\right\rbrace}\\
&\left.\right\rbrace\\
&\frattini{{\mathsf{MW}}_{11}}=\set{1,a^2}
\end{flalign*}

En cuanto al ret{\'\i}culo de subgrupos normales y nilpotentes, lo mostramos a continuaci\'on:

\begin{flalign*}
&\mathcal{L}_{norm,nil} = \left\lbrace\right.\\
&\quad \set{1},{\left\lbrace{1,a^2}\right\rbrace},
{\left\lbrace{1,b}\right\rbrace},
{\left\lbrace{1,a^2b}\right\rbrace},
{\left\lbrace{1,a,a^2,a^3}\right\rbrace},
{\left\lbrace{1,a^2,ab,a^3b}\right\rbrace},\\
&\quad {\left\lbrace{1,a^2,c,a^2c}\right\rbrace},
{\left\lbrace{1,a^2,ac,a^3c}\right\rbrace},
{\left\lbrace{1,a^2,bc,a^2bc}\right\rbrace},
{\left\lbrace{1,a^2,abc,a^3bc}\right\rbrace},\\
&\quad {\left\lbrace{1,a^2,b,a^2b}\right\rbrace},
{\left\lbrace{1,a,a^2,a^3,b,ab,a^2b,a^3b}\right\rbrace},
{\left\lbrace{1,a,a^2,a^3,c,ac,a^2c,a^3c}\right\rbrace},\\
&\quad {\left\lbrace{1,a,a^2,a^3,bc,abc,a^2bc,a^3bc}\right\rbrace},
{\left\lbrace{1,a^2,b,a^2b,c,a^2c,bc,a^2bc}\right\rbrace},
{\left\lbrace{1,a^2,b,a^2b,ac,a^3c,abc,a^3bc}\right\rbrace},\\
&\quad {\left\lbrace{1,a^2,ab,a^3b,c,a^2c,abc,a^3bc}\right\rbrace},
{\left\lbrace{1,a^2,ab,a^3b,ac,a^3c,bc,a^2bc}\right\rbrace},\\
&\quad {{\mathsf{MW}}_{11}}\\
&\left.\right\rbrace\\
&\fitting{{\mathsf{MW}}_{11}}={\mathsf{MW}}_{11}
\end{flalign*}

As{\'\i} tenemos que $\frattini{{\mathsf{MW}}_{11}}=\set{1,a^2}$
al haber varios minimales que solo comparten la unidad. A su vez el subgrupo de
Fitting es el completo ya que hay varios maximales normales y nilpotentes de orden $8$,
que necesariamente
generar\'an el completo.


\begin{flalign*}
&\left[1,3,12,0,0\right]\\
&\left[
\begin{array}{l}
\left[
(1,1),dc(1,2),c(2,2),n(4,4)
\right]\\
\left[
na(1,4),na(3,8),n(2,\dihgr{4}),n(2,\quaterniongr{1}),(1,16)
\right]\\
\left[\right]\\
\left[\right]
\end{array}
\right]\\
&\left[
\begin{array} {c|c|c|c}
\mathtt{nil}_2 & \trivialgr & {\left\lbrace{1,a^2,b,a^2b}\right\rbrace} & {\mathsf{MW}}_{11} \\
{} & {\mathsf{MW}}_{11} & {\left\lbrace{1,a^2}\right\rbrace} & \trivialgr
\end{array}
\right]\\
&\left[
\begin{array}{l}
\left[\pdcyclics{2}{2}         \cong         \centro{{\mathsf{MW}}_{11}}\supsetneqq{{\mathsf{MW}}_{11}}^{\prime}         \cong         \cyclicgr{2}\right]\\
\left[\right]
\end{array}
\right]\\
&\left[
\begin{array}{l}
\left[
\graut{{\mathsf{MW}}_{11}}         \cong         \nicefrac{{\mathsf{MW}}_{11}}{\centro{{\mathsf{MW}}_{11}}}         \cong         {\pdcyclics{2}{2}},
{{\mathsf{MW}}_{11}}^{\mathtt{ab}}=\nicefrac{{{\mathsf{MW}}_{11}}}{{{\mathsf{MW}}_{11}}^{\prime}}         \cong         {\pdtrescyclics{2}{2}{2}}
\right]\\
\left[\right]
\end{array}
\right]\\
&\left[
\frattini{{\mathsf{MW}}_{11}}={{\mathsf{MW}}_{11}}^{\prime},
\fitting{{\mathsf{MW}}_{11}}={{\mathsf{MW}}_{11}}
\right]
\end{flalign*}


El siguiente grupo ser\'a, ${\mathsf{MW}}_{12}$.


\smallskip\noindent
\resizebox{\linewidth}{!}{
\begin{tabular}{|c|c|c|c|c|c|c|c|c|c|c|c|c|c|c|c|}
\hline
$1$&$a$&$a^2$&$a^3$&$b$&$ab$&$a^2b$&$a^3b$&$c$&$ac$&$a^2c$&$a^3c$&$bc$&$abc$&$a^2bc$&$a^3bc$\\
\hline
$a$&$a^2$&$a^3$&$1$&$ab$&$a^2b$&$a^3b$&$b$&$ac$&$a^2c$&$a^3c$&$c$&$abc$&$a^2bc$&$a^3bc$&$bc$\\
\hline
$a^2$&$a^3$&$1$&$a$&$a^2b$&$a^3b$&$b$&$ab$&$a^2c$&$a^3c$&$c$&$ac$&$a^2bc$&$a^3bc$&$bc$&$abc$\\
\hline
$a^3$&$1$&$a$&$a^2$&$a^3b$&$b$&$ab$&$a^2b$&$a^3c$&$c$&$ac$&$a^2c$&$a^3bc$&$bc$&$abc$&$a^2bc$\\
\hline
$b$&$ab$&$a^2b$&$a^3b$&$1$&$a$&$a^2$&$a^3$&$bc$&$abc$&$a^2bc$&$a^3bc$&$c$&$ac$&$a^2c$&$a^3c$\\
\hline
$ab$&$a^2b$&$a^3b$&$b$&$a$&$a^2$&$a^3$&$1$&$abc$&$a^2bc$&$a^3bc$&$bc$&$ac$&$a^2c$&$a^3c$&$c$\\
\hline
$a^2b$&$a^3b$&$b$&$ab$&$a^2$&$a^3$&$1$&$a$&$a^2bc$&$a^3bc$&$bc$&$abc$&$a^2c$&$a^3c$&$c$&$ac$\\
\hline
$a^3b$&$b$&$ab$&$a^2b$&$a^3$&$1$&$a$&$a^2$&$a^3bc$&$bc$&$abc$&$a^2bc$&$a^3c$&$c$&$ac$&$a^2c$\\
\hline
$c$&$abc$&$a^2c$&$a^3bc$&$bc$&$ac$&$a^2bc$&$a^3c$&$a^2$&$a^3b$&$1$&$ab$&$a^2b$&$a^3$&$b$&$a$\\
\hline
$ac$&$a^2bc$&$a^3c$&$bc$&$abc$&$a^2c$&$a^3bc$&$c$&$a^3$&$b$&$a$&$a^2b$&$a^3b$&$1$&$ab$&$a^2$\\
\hline
$a^2c$&$a^3bc$&$c$&$abc$&$a^2bc$&$a^3c$&$bc$&$ac$&$1$&$ab$&$a^2$&$a^3b$&$b$&$a$&$a^2b$&$a^3$\\
\hline
$a^3c$&$bc$&$ac$&$a^2bc$&$a^3bc$&$c$&$abc$&$a^2c$&$a$&$a^2b$&$a^3$&$b$&$ab$&$a^2$&$a^3b$&$1$\\
\hline
$bc$&$ac$&$a^2bc$&$a^3c$&$c$&$abc$&$a^2c$&$a^3bc$&$a^2b$&$a^3$&$b$&$a$&$a^2$&$a^3b$&$1$&$ab$\\
\hline
$abc$&$a^2c$&$a^3bc$&$c$&$ac$&$a^2bc$&$a^3c$&$bc$&$a^3b$&$1$&$ab$&$a^2$&$a^3$&$b$&$a$&$a^2b$\\
\hline
$a^2bc$&$a^3c$&$bc$&$ac$&$a^2c$&$a^3bc$&$c$&$abc$&$b$&$a$&$a^2b$&$a^3$&$1$&$ab$&$a^2$&$a^3b$\\
\hline
$a^3bc$&$c$&$abc$&$a^2c$&$a^3c$&$bc$&$ac$&$a^2bc$&$ab$&$a^2$&$a^3b$&$1$&$a$&$a^2b$&$a^3$&$b$\\
\hline
\end{tabular}
}


Este grupo no es abeliano.

\begin{flalign*}
&\left[1,3,12,0,0\right]
\end{flalign*}

Despu\'es de mostrar el centro del grupo y su derivado, comenzamos por los subgrupos c{\'\i}clicos.


\[
\centro{{\mathsf{MW}}_{12}} = {\left\lbrace{1,a^2,b,a^2b}\right\rbrace}          \cong          \pdcyclics{2}{2}
\]

Y el subgrupo derivado es


\[
{{\mathsf{MW}}_{12}}^{\prime} = {\left\lbrace{1,b}\right\rbrace}          \cong          \cyclicgr{2}
\]


Estamos en condiciones de clasificar los subgrupos c{\'\i}clicos y en relaci\'on con el centro y el derivado:


\begin{flalign*}
&\mathcal{L}_1\left({{\mathsf{MW}}_{12}}\right)=\left\lbrace\right.\\
&\quad {\left\lbrace{1}\right\rbrace} ,
{\left\lbrace{1,a^2}\right\rbrace} ,%c
{\left\lbrace{1,b}\right\rbrace} ,%dc
{\left\lbrace{1,a^2b}\right\rbrace} ,%c
{\left\lbrace{1,a,a^2,a^3}\right\rbrace} ,\\%
&\quad {\left\lbrace{1,a^2,ab,a^3b}\right\rbrace} ,%
{\left\lbrace{1,a^2,c,a^2c}\right\rbrace} ,%
{\left\lbrace{1,a^2,bc,a^2bc}\right\rbrace} ,\\%
&{\left\lbrace{1,b,ac,abc}\right\rbrace} ,%n
{\left\lbrace{1,b,a^3c,a^3bc}\right\rbrace}\\%n
&\left.\right\rbrace\\
&\left[(1,1),dc(1,2),c(2,2),n(2,4),(4,4)\right]
\end{flalign*}


Ahora los subgrupos generados por dos elementos, siguiendo el mismo patr\'on anterior,
sacaremos un nuevo invariante (vac{\'\i}o en este caso):


\begin{flalign*}
&\mathcal{L}_{2}\left({{\mathsf{MW}}_{12}}\right)=\left\lbrace\right.\\
&\quad {\left\lbrace{1,a^2,b,a^2b}\right\rbrace} ,%na
{\left\lbrace{1,a,a^2,a^3,b,ab,a^2b,a^3b}\right\rbrace} ,\\%na
&\quad {\left\lbrace{1,a^2,b,a^2b,c,a^2c,bc,a^2bc}\right\rbrace} ,%na
{\left\lbrace{1,a^2,b,a^2b,ac,a^3c,abc,a^3bc}\right\rbrace} , \\%na
&\quad {\mathsf{MW}}_{12}\\
&\left.\right\rbrace\\
&\left[na(1,4),na(1,\pdcyclics{4}{2}),na(2,\pdtrescyclics{2}{2}{2}),(1,16)\right]
\end{flalign*}


Los \'ordenes de los generadores de la tabla son $o(a)=4$, $o(b)=2$ y $o(c)=4$.
Podemos observar que todos sus subgrupo propios son normales.
Vemos que $ab=ba$, $ca=abc$, $bc=cb$, $(ac)^2 =(a^3c)^2 = (abc)^2=(a^3bc)^2 = b$,
$(ab)^2 =(a^3b)^2 = c^2=(a^2c)^2 = (bc)^2 = (a^2bc)^2 = a^2$.
Luego el grupo es describible como
${\mathsf{MW}}_{12}         \cong         \generadopor{a,b,c\;∣\;a^4=b^2=c^4=1,ba=ab,(ac)^2=(abc)^2=b,(ab)^2=c^2=(bc)^2=a^2}$


Serie central ascendente (comienza en el $\trivialgr$,
y en general
 no tiene porqu\'e alcanzar ${\mathsf{MW}}_{12}$):


\begin{flalign*}
&{\trivialgr} \normsubgr
{\left\lbrace{1,a^2,b,a^2b}\right\rbrace} \normsubgr
{{\mathsf{MW}}_{12}}
\end{flalign*}


Serie central descendente (comienza en el ${\mathsf{MW}}_{12}$,
y en general no tiene proqu\'e alcanzar el $\trivialgr$):


\begin{flalign*}
&{{\mathsf{MW}}_{12}}\triangleright
{\left\lbrace{1,b}\right\rbrace} \triangleright
{\trivialgr}
\end{flalign*}


El invariante que anotamos con lo visto hasta aqu{\'\i}:


\begin{flalign*}
&\left[1,3,12,0,0\right]\\
&\left[
\begin{array}{l}
\left[
(1,1),dc(1,2),c(2,2),n(2,4),(4,4)
\right]\\
\left[
na(1,4),na(1,\pdcyclics{4}{2}),na(2,\pdtrescyclics{2}{2}{2}),(1,16)
\right]\\
\left[\right]\\
\left[\right]
\end{array}
\right]\\
&\left[
\begin{array} {c|c|c|c}
\mathtt{nil}_2 & \trivialgr & {\left\lbrace{1,a^2,b,a^2b}\right\rbrace} & {\mathsf{MW}}_{12} \\
{} & {\mathsf{MW}}_{12} & {\left\lbrace{1,b}\right\rbrace} & \trivialgr
\end{array}
\right]\\
&\left[
\begin{array}{l}
\left[\pdcyclics{2}{2}         \cong         \centro{{\mathsf{MW}}_{12}}\supsetneqq{{\mathsf{MW}}_{12}}^{\prime}         \cong         \cyclicgr{2}\right]\\
\left[\right]
\end{array}
\right]
\end{flalign*}


Nos falta ver cu\'ales sean los cocientes por el centro y por el derivado.


\begin{flalign*}
&\graut{{\mathsf{MW}}_{12}}         \cong         \nicefrac{{\mathsf{MW}}_{12}}{\centro{{\mathsf{MW}}_{12}}}         \cong         {\pdcyclics{2}{2}}\\
&{{\mathsf{MW}}_{12}}^{\mathtt{ab}}=\nicefrac{{{\mathsf{MW}}_{12}}}{{{\mathsf{MW}}_{12}}^{\prime}}         \cong         {\pdcyclics{4}{2}}
\end{flalign*}


El ret{\'\i}culo de subgrupos maximales (propios) la herramienta desarrollada para el trabajo, al igual que el subgrupo de Frattini:


\begin{flalign*}
&\frattini{{\mathsf{MW}}_{12}}=\set{1,a^2,b,a^2b}
\end{flalign*}

En cuanto al ret{\'\i}culo de subgrupos normales y nilpotentes tambi\'en lo da la herramienta. El subgrupo de Fitting es:

\begin{flalign*}
&\fitting{{\mathsf{MW}}_{12}}={\mathsf{MW}}_{12}
\end{flalign*}

As{\'\i} tenemos que $\frattini{{\mathsf{MW}}_{12}}=\set{1,a^2}$
al haber varios minimales que solo comparten la unidad. A su vez el subgrupo de
Fitting es el completo ya que hay varios maximales normales y nilpotentes de orden $8$,
que necesariamente
generar\'an el completo.


\begin{flalign*}
&\left[1,3,12,0,0\right]\\
&\left[
\begin{array}{l}
\left[
(1,1),dc(1,2),c(2,2),n(2,4),(4,4)
\right]\\
\left[
na(1,4),na(1,\pdcyclics{4}{2}),na(2,\pdtrescyclics{2}{2}{2}),(1,16)
\right]\\
\left[\right]\\
\left[\right]
\end{array}
\right]\\
&\left[
\begin{array} {c|c|c|c}
\mathtt{nil}_2 & \trivialgr & {\left\lbrace{1,a^2,b,a^2b}\right\rbrace} & {\mathsf{MW}}_{12} \\
{} & {\mathsf{MW}}_{12} & {\left\lbrace{1,a^2}\right\rbrace} & \trivialgr
\end{array}
\right]\\
&\left[
\begin{array}{l}
\left[\pdcyclics{2}{2}         \cong         \centro{{\mathsf{MW}}_{12}}\supsetneqq{{\mathsf{MW}}_{12}}^{\prime}         \cong         \cyclicgr{2}\right]\\
\left[\right]
\end{array}
\right]\\
&\left[
\begin{array}{l}
\left[
\graut{{\mathsf{MW}}_{12}}         \cong         \nicefrac{{\mathsf{MW}}_{12}}{\centro{{\mathsf{MW}}_{12}}}         \cong         {\pdcyclics{2}{2}},
{{\mathsf{MW}}_{12}}^{\mathtt{ab}}=\nicefrac{{{\mathsf{MW}}_{12}}}{{{\mathsf{MW}}_{12}}^{\prime}}         \cong         {\pdtrescyclics{2}{2}{2}}
\right]\\
\left[\right]
\end{array}
\right]\\
&\left[
\frattini{{\mathsf{MW}}_{12}}=\centro{{\mathsf{MW}}_{12}},
\fitting{{\mathsf{MW}}_{12}}={{\mathsf{MW}}_{12}}
\right]
\end{flalign*}



El siguiente grupo ser\'a, ${\mathsf{MW}}_{13}$.


\smallskip\noindent
\resizebox{\linewidth}{!}{
\begin{tabular}{|c|c|c|c|c|c|c|c|c|c|c|c|c|c|c|c|}
\hline
$1$&$a$&$a^2$&$a^3$&$b$&$ab$&$a^2b$&$a^3b$&$c$&$ac$&$a^2c$&$a^3c$&$bc$&$abc$&$a^2bc$&$a^3bc$\\
\hline
$a$&$a^2$&$a^3$&$1$&$ab$&$a^2b$&$a^3b$&$b$&$ac$&$a^2c$&$a^3c$&$c$&$abc$&$a^2bc$&$a^3bc$&$bc$\\
\hline
$a^2$&$a^3$&$1$&$a$&$a^2b$&$a^3b$&$b$&$ab$&$a^2c$&$a^3c$&$c$&$ac$&$a^2bc$&$a^3bc$&$bc$&$abc$\\
\hline
$a^3$&$1$&$a$&$a^2$&$a^3b$&$b$&$ab$&$a^2b$&$a^3c$&$c$&$ac$&$a^2c$&$a^3bc$&$bc$&$abc$&$a^2bc$\\
\hline
$b$&$ab$&$a^2b$&$a^3b$&$1$&$a$&$a^2$&$a^3$&$bc$&$abc$&$a^2bc$&$a^3bc$&$c$&$ac$&$a^2c$&$a^3c$\\
\hline
$ab$&$a^2b$&$a^3b$&$b$&$a$&$a^2$&$a^3$&$1$&$abc$&$a^2bc$&$a^3bc$&$bc$&$ac$&$a^2c$&$a^3c$&$c$\\
\hline
$a^2b$&$a^3b$&$b$&$ab$&$a^2$&$a^3$&$1$&$a$&$a^2bc$&$a^3bc$&$bc$&$abc$&$a^2c$&$a^3c$&$c$&$ac$\\
\hline
$a^3b$&$b$&$ab$&$a^2b$&$a^3$&$1$&$a$&$a^2$&$a^3bc$&$bc$&$abc$&$a^2bc$&$a^3c$&$c$&$ac$&$a^2c$\\
\hline
$c$&$ac$&$a^2c$&$a^3c$&$bc$&$abc$&$a^2bc$&$a^3bc$&$b$&$ab$&$a^2b$&$a^3b$&$1$&$a$&$a^2$&$a^3$\\
\hline
$ac$&$a^2c$&$a^3c$&$c$&$abc$&$a^2bc$&$a^3bc$&$bc$&$ab$&$a^2b$&$a^3b$&$b$&$a$&$a^2$&$a^3$&$1$\\
\hline
$a^2c$&$a^3c$&$c$&$ac$&$a^2bc$&$a^3bc$&$bc$&$abc$&$a^2b$&$a^3b$&$b$&$ab$&$a^2$&$a^3$&$1$&$a$\\
\hline
$a^3c$&$c$&$ac$&$a^2c$&$a^3bc$&$bc$&$abc$&$a^2bc$&$a^3b$&$b$&$ab$&$a^2b$&$a^3$&$1$&$a$&$a^2$\\
\hline
$bc$&$abc$&$a^2bc$&$a^3bc$&$c$&$ac$&$a^2c$&$a^3c$&$1$&$a$&$a^2$&$a^3$&$b$&$ab$&$a^2b$&$a^3b$\\
\hline
$abc$&$a^2bc$&$a^3bc$&$bc$&$ac$&$a^2c$&$a^3c$&$c$&$a$&$a^2$&$a^3$&$1$&$ab$&$a^2b$&$a^3b$&$b$\\
\hline
$a^2bc$&$a^3bc$&$bc$&$abc$&$a^2c$&$a^3c$&$c$&$ac$&$a^2$&$a^3$&$1$&$a$&$a^2b$&$a^3b$&$b$&$ab$\\
\hline
$a^3bc$&$bc$&$abc$&$a^2bc$&$a^3c$&$c$&$ac$&$a^2c$&$a^3$&$1$&$a$&$a^2$&$a^3b$&$b$&$ab$&$a^2b$\\
\hline
\end{tabular}
}


Teniendo en cuenta que $b=(bc)^2$, esto es $1=cbc$, $c^3=bc$, $c^2=b$ y que es abeliano podemos transformar la tabla:


\smallskip\noindent
\resizebox{\linewidth}{!}{
\begin{tabular}{|c|c|c|c|c|c|c|c|c|c|c|c|c|c|c|c|}
\hline
$1$&$a$&$a^2$&$a^3$&$c^2$&$ac^2$&$a^2c^2$&$a^3c^2$&$c$&$ac$&$a^2c$&$a^3c$&$c^3$&$ac^3$&$a^2c^3$&$a^3c^3$\\
\hline
$a$&$a^2$&$a^3$&$1$&$ac^2$&$a^2c^2$&$a^3c^2$&$c^2$&$ac$&$a^2c$&$a^3c$&$c$&$ac^3$&$a^2c^3$&$a^3c^3$&$c^3$\\
\hline
$a^2$&$a^3$&$1$&$a$&$a^2c^2$&$a^3c^2$&$c^2$&$ac^2$&$a^2c$&$a^3c$&$c$&$ac$&$a^2c^3$&$a^3c^3$&$c^3$&$ac^3$\\
\hline
$a^3$&$1$&$a$&$a^2$&$a^3c^2$&$c^2$&$ac^2$&$a^2c^2$&$a^3c$&$c$&$ac$&$a^2c$&$a^3c^3$&$c^3$&$ac^3$&$a^2c^3$\\
\hline
$c^2$&$ac^2$&$a^2c^2$&$a^3c^2$&$1$&$a$&$a^2$&$a^3$&$c^3$&$ac^3$&$a^2c^3$&$a^3c^3$&$c$&$ac$&$a^2c$&$a^3c$\\
\hline
$ac^2$&$a^2c^2$&$a^3c^2$&$c^2$&$a$&$a^2$&$a^3$&$1$&$ac^3$&$a^2c^3$&$a^3c^3$&$c^3$&$ac$&$a^2c$&$a^3c$&$c$\\
\hline
$a^2c^2$&$a^3c^2$&$c^2$&$ac^2$&$a^2$&$a^3$&$1$&$a$&$a^2c^3$&$a^3c^3$&$c^3$&$ac^3$&$a^2c$&$a^3c$&$c$&$ac$\\
\hline
$a^3c^2$&$c^2$&$ac^2$&$a^2c^2$&$a^3$&$1$&$a$&$a^2$&$a^3c^3$&$c^3$&$ac^3$&$a^2c^3$&$a^3c$&$c$&$ac$&$a^2c$\\
\hline
$c$&$ac$&$a^2c$&$a^3c$&$c^3$&$ac^3$&$a^2c^3$&$a^3c^3$&$c^2$&$ac^2$&$a^2c^2$&$a^3c^2$&$1$&$a$&$a^2$&$a^3$\\
\hline
$ac$&$a^2c$&$a^3c$&$c$&$ac^3$&$a^2c^3$&$a^3c^3$&$c^3$&$ac^2$&$a^2c^2$&$a^3c^2$&$c^2$&$a$&$a^2$&$a^3$&$1$\\
\hline
$a^2c$&$a^3c$&$c$&$ac$&$a^2c^3$&$a^3c^3$&$c^3$&$ac^3$&$a^2c^2$&$a^3c^2$&$c^2$&$ac^2$&$a^2$&$a^3$&$1$&$a$\\
\hline
$a^3c$&$c$&$ac$&$a^2c$&$a^3c^3$&$c^3$&$ac^3$&$a^2c^3$&$a^3c^2$&$c^2$&$ac^2$&$a^2c^2$&$a^3$&$1$&$a$&$a^2$\\
\hline
$c^3$&$ac^3$&$a^2c^3$&$a^3c^3$&$c$&$ac$&$a^2c$&$a^3c$&$1$&$a$&$a^2$&$a^3$&$c^2$&$ac^2$&$a^2c^2$&$a^3c^2$\\
\hline
$ac^3$&$a^2c^3$&$a^3c^3$&$c^3$&$ac$&$a^2c$&$a^3c$&$c$&$a$&$a^2$&$a^3$&$1$&$ac^2$&$a^2c^2$&$a^3c^2$&$c^2$\\
\hline
$a^2c^3$&$a^3c^3$&$c^3$&$ac^3$&$a^2c$&$a^3c$&$c$&$ac$&$a^2$&$a^3$&$1$&$a$&$a^2c^2$&$a^3c^2$&$c^2$&$ac^2$\\
\hline
$a^3c^3$&$c^3$&$ac^3$&$a^2c^3$&$a^3c$&$c$&$ac$&$a^2c$&$a^3$&$1$&$a$&$a^2$&$a^3c^2$&$c^2$&$ac^2$&$a^2c^2$\\
\hline
\end{tabular}
}



Ya podemos decir algo m\'as sobre este grupo, ${\mathsf{MW}}_{13}=\entreangulos{a∣a^4=1}\entreangulos{c∣c^4=1}        \cong        \pdcyclics{4}{4}$.

\begin{flalign*}
&\left[1,3,12,0,0\right]
\end{flalign*}

Despu\'es de mostrar el centro del grupo y su derivado, comenzamos por los subgrupos c{\'\i}clicos.

\[
\centro{{\mathsf{MW}}_{13}} = {\mathsf{MW}}_{13}         \cong        \pdcyclics{4}{4}
\]

Y el subgrupo derivado es


\[
{{\mathsf{MW}}_{13}}^{\prime} = {\left\lbrace{1}\right\rbrace}          \cong          \cyclicgr{1}
\]


Los grupos c{\'\i}clicos son:
\begin{flalign*}
&\mathcal{L}_1\left({{\mathsf{MW}}_{13}}\right)=\left\lbrace\right.\\
&\quad {\left\lbrace{1}\right\rbrace} ,
{\left\lbrace{1}\right\rbrace} ,
{\left\lbrace{1,a^2}\right\rbrace} ,
{\left\lbrace{1,c^2}\right\rbrace} ,
{\left\lbrace{1,a^2c^2}\right\rbrace} ,
{\left\lbrace{1,a,a^2,a^3}\right\rbrace} ,\\%
&\quad {\left\lbrace{1,a^2,ac^2,a^3c^2}\right\rbrace} ,
{\left\lbrace{1,c^2,c,c^3}\right\rbrace} ,
{\left\lbrace{1,c^2,a^2c,a^2c^3}\right\rbrace} ,\\%
&\quad {\left\lbrace{1,a^2c^2,ac,a^3c^3}\right\rbrace} ,
{\left\lbrace{1,a^2c^2,a^3c,ac^3}\right\rbrace}\\%
&\left.\right\rbrace\\
&\left[(1,1),c(3,2),c(6,4)\right]
\end{flalign*}


Ahora los subgrupos generados por dos elementos, siguiendo el mismo patr\'on anterior,
sacaremos un nuevo invariante (vac{\'\i}o en este caso):


\begin{flalign*}
&\mathcal{L}_{2}\left({{\mathsf{MW}}_{13}}\right)=\left\lbrace\right.\\
&\quad {\left\lbrace{1,a^2,c^2,a^2c^2}\right\rbrace} ,
{\left\lbrace{1,a,a^2,a^3,c^2,ac^2,a^2c^2,a^3c^2}\right\rbrace} ,\\
&\quad {\left\lbrace{1,a^2,c^2,a^2c^2,c,a^2c,c^3,a^2c^3}\right\rbrace} ,
{\left\lbrace{1,a^2,c^2,a^2c^2,ac,a^3c,ac^3,a^3c^3}\right\rbrace} ,\\
&\quad {\mathsf{MW}}_{13}\\
&\left.\right\rbrace\\
&\left[c(1,4),c(3,\pdcyclics{4}{2}),c(1,16)\right]
\end{flalign*}


Luego el grupo es describible como
${\mathsf{MW}}_{13}         \cong         \generadopor{a,c\;∣\;a^4=c^4=1,ac=ca}$

El invariante que anotamos con lo visto hasta aqu{\'\i}:


\begin{flalign*}
&\left[1,3,12,0,0\right]\\
&\left[
\begin{array}{l}
\left[
(1,1),c(3,2),c(6,4)
\right]\\
\left[
c(1,4),c(3,\pdcyclics{4}{2}),c(1,16)
\right]\\
\left[\right]\\
\left[\right]
\end{array}
\right]\\
&\left[
\begin{array} {c|c|c}
\mathtt{nil}_1 & \trivialgr  & {\mathsf{MW}}_{13} \\
{} & {\mathsf{MW}}_{13}  & \trivialgr
\end{array}
\right]\\
&\left[
\begin{array}{l}
\left[\pdcyclics{4}{4}         \cong         \centro{{\mathsf{MW}}_{13}}\supsetneqq{{\mathsf{MW}}_{13}}^{\prime}         \cong         \cyclicgr{1}\right]\\
\left[\right]
\end{array}
\right]
\end{flalign*}


Nos falta ver cu\'ales sean los cocientes por el centro y por el derivado.


\begin{flalign*}
&\graut{{\mathsf{MW}}_{13}}         \cong         \nicefrac{{\mathsf{MW}}_{13}}{\centro{{\mathsf{MW}}_{13}}}         \cong         {\trivialgr}\\
&{{\mathsf{MW}}_{13}}^{\mathtt{ab}}=\nicefrac{{{\mathsf{MW}}_{13}}}{{{\mathsf{MW}}_{13}}^{\prime}}         \cong         {\pdcyclics{4}{4}}
\end{flalign*}


El ret{\'\i}culo de subgrupos maximales (propios) la herramienta desarrollada para el trabajo, al igual que el subgrupo de Frattini:


\begin{flalign*}
&\frattini{{\mathsf{MW}}_{13}}=\set{a^2,c^2,a^2c^2}        \cong        \pdcyclics{2}{2}
\end{flalign*}

En cuanto al ret{\'\i}culo de subgrupos normales y nilpotentes tambi\'en lo da la herramienta. El subgrupo de Fitting es:

\begin{flalign*}
&\fitting{{\mathsf{MW}}_{13}}={\mathsf{MW}}_{13}
\end{flalign*}

As{\'\i} tenemos que $\frattini{{\mathsf{MW}}_{13}}=\set{1,a^2}$
al haber varios minimales que solo comparten la unidad. A su vez el subgrupo de
Fitting es el completo ya que hay varios maximales normales y nilpotentes de orden $8$,
que necesariamente
generar\'an el completo.


\begin{flalign*}
&\left[1,3,12,0,0\right]\\
&\left[
\begin{array}{l}
\left[
(1,1),c(3,2),c(6,4)
\right]\\
\left[
c(1,4),c(3,\pdcyclics{4}{2}),c(1,16)
\right]\\
\left[\right]\\
\left[\right]
\end{array}
\right]\\
&\left[
\begin{array} {c|c|c}
\mathtt{nil}_1 & \trivialgr &  {\mathsf{MW}}_{13} \\
{} & {\mathsf{MW}}_{13} &  \trivialgr
\end{array}
\right]\\
&\left[
\begin{array}{l}
\left[\pdcyclics{4}{4}         \cong         \centro{{\mathsf{MW}}_{13}}\supsetneqq{{\mathsf{MW}}_{13}}^{\prime}         \cong         \cyclicgr{1}\right]\\
\left[\right]
\end{array}
\right]\\
&\left[
\begin{array}{l}
\left[
\graut{{\mathsf{MW}}_{13}}         \cong         \nicefrac{{\mathsf{MW}}_{13}}{\centro{{\mathsf{MW}}_{13}}}         \cong         \trivialgr,
{{\mathsf{MW}}_{13}}^{\mathtt{ab}}=\nicefrac{{{\mathsf{MW}}_{13}}}{{{\mathsf{MW}}_{13}}^{\prime}}         \cong         \pdcyclics{4}{4}
\right]\\
\left[\right]
\end{array}
\right]\\
&\left[
\frattini{{\mathsf{MW}}_{13}}        \cong        \pdcyclics{2}{2},
\fitting{{\mathsf{MW}}_{13}}={{\mathsf{MW}}_{13}}
\right]
\end{flalign*}


El siguiente grupo ser\'a, ${\mathsf{MW}}_{14}$ que realmente no viene como extensi\'on, sino como a\~nadido del
caso que hab{\'\i}a sido dejado de lado, cu\'ando el grupo normal de orden $8$ a extender es $\pdtrescyclics{2}{2}{2}$
que solo se puede extender a $\pdquatrocyclics{2}{2}{2}{2}$. As{\'\i}, de entrada, ${\mathsf{MW}}_{14}=\pdquatrocyclics{2}{2}{2}{2}$.

\smallskip\noindent
\resizebox{\linewidth}{!}{
\begin{tabular}{|c|c|c|c|c|c|c|c|c|c|c|c|c|c|c|c|}
\hline
$1$&$a$&$b$&$ab$&$c$&$ac$&$bc$&$abc$&$d$&$ad$&$bd$&$abd$&$cd$&$acd$&$bcd$&$abcd$\\
\hline
$a$&$1$&$ab$&$b$&$ac$&$c$&$abc$&$bc$&$ad$&$d$&$abd$&$bd$&$acd$&$cd$&$abcd$&$bcd$\\
\hline
$b$&$ab$&$1$&$a$&$bc$&$abc$&$c$&$ac$&$bd$&$abd$&$d$&$ad$&$bcd$&$abcd$&$cd$&$acd$\\
\hline
$ab$&$b$&$a$&$1$&$abc$&$bc$&$ac$&$c$&$abd$&$bd$&$ad$&$d$&$abcd$&$bcd$&$acd$&$cd$\\
\hline
$c$&$ac$&$bc$&$abc$&$1$&$a$&$b$&$ab$&$cd$&$acd$&$bcd$&$abcd$&$d$&$ad$&$bd$&$abd$\\
\hline
$ac$&$c$&$abc$&$bc$&$a$&$1$&$ab$&$b$&$acd$&$cd$&$abcd$&$bcd$&$ad$&$d$&$abd$&$bd$\\
\hline
$bc$&$abc$&$c$&$ac$&$b$&$ab$&$1$&$a$&$bcd$&$abcd$&$cd$&$acd$&$bd$&$abd$&$d$&$ad$\\
\hline
$abc$&$bc$&$ac$&$c$&$ab$&$b$&$a$&$1$&$abcd$&$bcd$&$acd$&$cd$&$abd$&$bd$&$ad$&$d$\\
\hline
$d$&$ad$&$bd$&$abd$&$cd$&$acd$&$bcd$&$abcd$&$1$&$a$&$b$&$ab$&$c$&$ac$&$bc$&$abc$\\
\hline
$ad$&$d$&$abd$&$bd$&$acd$&$cd$&$abcd$&$bcd$&$a$&$1$&$ab$&$b$&$ac$&$c$&$abc$&$bc$\\
\hline
$bd$&$abd$&$d$&$ad$&$bcd$&$abcd$&$cd$&$acd$&$b$&$ab$&$1$&$a$&$bc$&$abc$&$c$&$ac$\\
\hline
$abd$&$bd$&$ad$&$d$&$abcd$&$bcd$&$acd$&$cd$&$ab$&$b$&$a$&$1$&$abc$&$bc$&$ac$&$c$\\
\hline
$cd$&$acd$&$bcd$&$abcd$&$d$&$ad$&$bd$&$abd$&$c$&$ac$&$bc$&$abc$&$1$&$a$&$b$&$ab$\\
\hline
$acd$&$cd$&$abcd$&$bcd$&$ad$&$d$&$abd$&$bd$&$ac$&$c$&$abc$&$bc$&$a$&$1$&$ab$&$b$\\
\hline
$bcd$&$abcd$&$cd$&$acd$&$bd$&$abd$&$d$&$ad$&$bc$&$abc$&$c$&$ac$&$b$&$ab$&$1$&$a$\\
\hline
$abcd$&$bcd$&$acd$&$cd$&$abd$&$bd$&$ad$&$d$&$abc$&$bc$&$ac$&$c$&$ab$&$b$&$a$&$1$\\
\hline
\end{tabular}
}


Teniendo en cuenta que ${\mathsf{MW}}_{14}=\generadopor{a}\generadopor{b}\generadopor{c}\generadopor{d}       \cong       \pdquatrocyclics{2}{2}{2}{2}$
ya que tenemos las siguientes relaciones
\begin{flalign*}
& o(a)=o(b)=o(c)=o(d)=2\\
& ab=ba \ly ac=ca \ly ad=da \ly bc=cb \ly bd=db \ly cd=dc\\
&\generadopor{a}       \cap \generadopor{b}=\generadopor{a}       \cap \generadopor{c}=\generadopor{a}       \cap \generadopor{d}=\set{1}\\
&\generadopor{b}       \cap \generadopor{c}=\generadopor{b}       \cap \generadopor{d}=\generadopor{c}       \cap \generadopor{d}=\set{1}
\end{flalign*}

Ya podemos decir algo m\'as sobre este grupo, ${\mathsf{MW}}_{14}=\entreangulos{a∣a^2=1}\entreangulos{b∣b^2=1}\entreangulos{c∣c^4=1}\entreangulos{d∣d^4=1}        \cong        \pdquatrocyclics{2}{2}{2}{2}$.

\begin{flalign*}
&\left[1,15,0,0,0\right]
\end{flalign*}

Despu\'es de mostrar el centro del grupo y su derivado, comenzamos por los subgrupos c{\'\i}clicos.

\[
\centro{{\mathsf{MW}}_{14}} = {\mathsf{MW}}_{14}         \cong        \pdquatrocyclics{2}{2}{2}{2}
\]

Y el subgrupo derivado es


\[
{{\mathsf{MW}}_{14}}^{\prime} = {\left\lbrace{1}\right\rbrace}          \cong          \cyclicgr{1}
\]


Los grupos c{\'\i}clicos son:
\begin{flalign*}
&\mathcal{L}_1\left({{\mathsf{MW}}_{14}}\right)=\left\lbrace\right.\\
&\quad {\left\lbrace{1}\right\rbrace} ,
{\left\lbrace{1,a}\right\rbrace} ,
{\left\lbrace{1,b}\right\rbrace} ,
{\left\lbrace{1,ab}\right\rbrace} ,
{\left\lbrace{1,c}\right\rbrace} ,\\%
&\quad {\left\lbrace{1,ac}\right\rbrace} ,
{\left\lbrace{1,bc}\right\rbrace} ,
{\left\lbrace{1,abc}\right\rbrace} ,
{\left\lbrace{1,d}\right\rbrace} ,
{\left\lbrace{1,ad}\right\rbrace} ,\\%
&\quad {\left\lbrace{1,bd}\right\rbrace} ,
{\left\lbrace{1,abd}\right\rbrace} ,
{\left\lbrace{1,cd}\right\rbrace} ,
{\left\lbrace{1,acd}\right\rbrace} ,\\%
&\quad {\left\lbrace{1,bcd}\right\rbrace} ,
{\left\lbrace{1,abcd}\right\rbrace}\\
&\left.\right\rbrace\\
&\left[(1,1),c(15,2)\right]
\end{flalign*}


Ahora los subgrupos generados por dos elementos, siguiendo el mismo patr\'on anterior,
sacaremos un nuevo invariante (vac{\'\i}o en este caso):

\begin{flalign*}
&\mathcal{L}_{2}\left({{\mathsf{MW}}_{14}}\right)=\left\lbrace\right.\\
&\quad {\left\lbrace{1,a,b,ab}\right\rbrace} ,
{\left\lbrace{1,a,c,ac}\right\rbrace} ,
{\left\lbrace{1,a,bc,abc}\right\rbrace} ,\\
&\quad {\left\lbrace{1,a,d,ad}\right\rbrace} ,
{\left\lbrace{1,a,bd,abd}\right\rbrace} ,
{\left\lbrace{1,a,cd,acd}\right\rbrace} ,\\
&\quad {\left\lbrace{1,a,bcd,abcd}\right\rbrace} ,
{\left\lbrace{1,b,c,bc}\right\rbrace} ,
{\left\lbrace{1,b,ac,abc}\right\rbrace} ,\\
&\quad {\left\lbrace{1,b,d,bd}\right\rbrace} ,
{\left\lbrace{1,b,ad,abd}\right\rbrace} ,
{\left\lbrace{1,b,cd,bcd}\right\rbrace} ,\\
&\quad {\left\lbrace{1,b,acd,abcd}\right\rbrace} ,
{\left\lbrace{1,ab,c,abc}\right\rbrace} ,
{\left\lbrace{1,ab,ac,bc}\right\rbrace} ,\\
&\quad {\left\lbrace{1,ab,d,abd}\right\rbrace} ,
{\left\lbrace{1,ab,ad,bd}\right\rbrace} ,
{\left\lbrace{1,ab,cd,abcd}\right\rbrace} ,\\
&\quad {\left\lbrace{1,ab,acd,bcd}\right\rbrace} ,
{\left\lbrace{1,c,d,cd}\right\rbrace} ,
{\left\lbrace{1,c,ad,acd}\right\rbrace} ,\\
&\quad {\left\lbrace{1,c,bd,bcd}\right\rbrace} ,
{\left\lbrace{1,c,abd,abcd}\right\rbrace} ,
{\left\lbrace{1,ac,d,acd}\right\rbrace} ,\\
&\quad {\left\lbrace{1,ac,ad,cd}\right\rbrace} ,
{\left\lbrace{1,ac,bd,abcd}\right\rbrace} ,
{\left\lbrace{1,ac,abd,bcd}\right\rbrace} ,\\
&\quad {\left\lbrace{1,bc,d,bcd}\right\rbrace} ,
{\left\lbrace{1,bc,ad,abcd}\right\rbrace} ,
{\left\lbrace{1,bc,bd,cd}\right\rbrace} ,\\
&\quad {\left\lbrace{1,bc,abd,acd}\right\rbrace} ,
{\left\lbrace{1,abc,d,abcd}\right\rbrace} ,
{\left\lbrace{1,abc,ad,bcd}\right\rbrace} ,\\
&\quad {\left\lbrace{1,abc,bd,acd}\right\rbrace} ,
{\left\lbrace{1,abc,abd,cd}\right\rbrace}\\
&\left.\right\rbrace\\
&\left[c(35,4)\right]
\end{flalign*}

Ahora el invariante debido a los subgrupos generados por $3$ elementos:


\begin{flalign*}
{\mathcal{L}}_{3}\entreparentesis{{\mathsf{MW}}_{14}}=\left\lbrace\right.\\
&\quad {\left\lbrace{1,a,b,ab,c,ac,bc,abc}\right\rbrace},
{\left\lbrace{1,a,b,ab,d,ad,bd,abd}\right\rbrace},\\
&\quad {\left\lbrace{1,a,b,ab,cd,acd,bcd,abcd}\right\rbrace},
{\left\lbrace{1,a,c,ac,d,ad,cd,acd}\right\rbrace},\\
&\quad {\left\lbrace{1,a,c,ac,bd,abd,bcd,abcd}\right\rbrace},
{\left\lbrace{1,a,bc,abc,d,ad,bcd,abcd}\right\rbrace},\\
&\quad {\left\lbrace{1,a,bc,abc,bd,abd,cd,acd}\right\rbrace},
{\left\lbrace{1,b,c,bc,d,bd,cd,bcd}\right\rbrace},\\
&\quad {\left\lbrace{1,b,c,bc,ad,abd,acd,abcd}\right\rbrace},
{\left\lbrace{1,b,ac,abc,d,bd,acd,abcd}\right\rbrace},\\
&\quad {\left\lbrace{1,b,ac,abc,ad,abd,cd,bcd}\right\rbrace},
{\left\lbrace{1,ab,c,abc,d,abd,cd,abcd}\right\rbrace},\\
&\quad {\left\lbrace{1,ab,c,abc,ad,bd,acd,bcd}\right\rbrace},
{\left\lbrace{1,ab,ac,bc,d,abd,acd,bcd}\right\rbrace},\\
&\quad {\left\lbrace{1,ab,ac,bc,ad,bd,cd,abcd}\right\rbrace}\\
&\left.\right\rbrace\\
&\left[c(15,\pdtrescyclics{2}{2}{2})\right]\\
&\left[c(1,16)\right]\\
\end{flalign*}

Luego el grupo es describible como
${\mathsf{MW}}_{14}         \cong         \generadopor{a,b,c,d\;∣\;a^2=b^2=c^2=d^2=1,ab=ba,ac=ca,ad=da,bc=cb,bd=db,cd=dc}$

El invariante que anotamos con lo visto hasta aqu{\'\i}:


\begin{flalign*}
&\left[1,15,0,0,0\right]\\
&\left[
\begin{array}{l}
\left[
(1,1),c(15,2)
\right]\\
\left[
c(35,4)
\right]\\
\left[
c(15,\pdtrescyclics{2}{2}{2})
\right]\\
\left[
c(1,16)
\right]
\end{array}
\right]\\
&\left[
\begin{array} {c|c|c}
\mathtt{nil}_1 & \trivialgr  & {\mathsf{MW}}_{14} \\
{} & {\mathsf{MW}}_{14}  & \trivialgr
\end{array}
\right]\\
&\left[
\begin{array}{l}
\left[\pdquatrocyclics{2}{2}{2}{2}         \cong         \centro{{\mathsf{MW}}_{14}}\supsetneqq{{\mathsf{MW}}_{14}}^{\prime}         \cong         \cyclicgr{1}\right]\\
\left[\right]
\end{array}
\right]
\end{flalign*}


Nos falta ver cu\'ales sean los cocientes por el centro y por el derivado.


\begin{flalign*}
&\graut{{\mathsf{MW}}_{14}}         \cong         \nicefrac{{\mathsf{MW}}_{14}}{\centro{{\mathsf{MW}}_{14}}}         \cong         {\trivialgr}\\
&{{\mathsf{MW}}_{14}}^{\mathtt{ab}}=\nicefrac{{{\mathsf{MW}}_{14}}}{{{\mathsf{MW}}_{14}}^{\prime}}         \cong         {\pdquatrocyclics{2}{2}{2}{2}}
\end{flalign*}


El ret{\'\i}culo de subgrupos maximales (propios) la herramienta desarrollada para el trabajo, al igual que el subgrupo de Frattini:


\begin{flalign*}
&\frattini{{\mathsf{MW}}_{14}}=\set{1}        \cong        \cyclicgr{1}
\end{flalign*}

En cuanto al ret{\'\i}culo de subgrupos normales y nilpotentes tambi\'en lo da la herramienta. El subgrupo de Fitting es:

\begin{flalign*}
&\fitting{{\mathsf{MW}}_{14}}={\mathsf{MW}}_{14}
\end{flalign*}

As{\'\i} tenemos que $\frattini{{\mathsf{MW}}_{14}}=\set{1}$
al haber varios minimales que solo comparten la unidad. A su vez el subgrupo de
Fitting es el completo ya que hay varios maximales normales y nilpotentes de orden $8$,
que necesariamente
generar\'an el grupo completo.


\begin{flalign*}
&\left[1,15,0,0,0\right]\\
&\left[
\begin{array}{l}
\left[
(1,1),c(15,2)
\right]\\
\left[
c(35,4)
\right]\\
\left[
c(15,\pdtrescyclics{2}{2}{2})
\right]\\
\left[
c(1,16)
\right]
\end{array}
\right]\\
&\left[
\begin{array} {c|c|c}
\mathtt{nil}_1 & \trivialgr  & {\mathsf{MW}}_{14} \\
{} & {\mathsf{MW}}_{14}  & \trivialgr
\end{array}
\right]\\
&\left[
\begin{array}{l}
\left[\pdquatrocyclics{2}{2}{2}{2}         \cong         \centro{{\mathsf{MW}}_{14}}\supsetneqq{{\mathsf{MW}}_{14}}^{\prime}         \cong         \cyclicgr{1}\right]\\
\left[\right]
\end{array}
\right]\\
&\left[
\begin{array}{l}
\left[
\graut{{\mathsf{MW}}_{14}}         \cong         \nicefrac{{\mathsf{MW}}_{14}}{\centro{{\mathsf{MW}}_{14}}}         \cong         \trivialgr,
{{\mathsf{MW}}_{14}}^{\mathtt{ab}}=\nicefrac{{{\mathsf{MW}}_{14}}}{{{\mathsf{MW}}_{14}}^{\prime}}         \cong         \pdquatrocyclics{2}{2}{2}{2}
\right]\\
\left[\right]
\end{array}
\right]\\
&\left[
\frattini{{\mathsf{MW}}_{14}}        \cong        \cyclicgr{1},
\fitting{{\mathsf{MW}}_{14}}={{\mathsf{MW}}_{14}}
\right]
\end{flalign*}

\chapter{Los grupos de orden 18}

Sea $\left|\grG\right| = 18 = 2\cdot3^2$ y $n_p$ es el n\'umero de $p$-subgrupos de Sylow de $\grG$.
Entonces por los teoremas de Sylow $n_3 \vert 2$ ($n_3        \in \set{1,2}$) y $n_3       \equiv        1 mod 3$. Por lo tanto $n_3 = 1$ y
tenemos que existe un \'unico grupo de orden $9$, que llamaremos $\grK$.
De aqu{'\i} que $\grK        \cong        \cyclicgr{9}$ o $\grK        \cong        \pdcyclics{3}{3}$.
Por otra parte, $n_2 \vert 9$ ($n_2       \in \set{1, 3, 9}$).

Ahora bien, todo elemento de $\grK$ tiene orden $3$ o $9$, y cualquier otro elemento que no est\'e en
$\grK$ ha de tener orden $2$ o $6$ (si no es c\'{\i}clico). Por el teorema de Cauchy, existir\'an
elementos de orden $2$, as\'{\i} que podemos considerar $\grN_1=\set{1,a}$ tal que $a$ sea ese elemento
de orden $2$. Sabemos que $\grK      \cup      \grN_1 = \set{\unity}$ y que $\grK\grN_1 = \grG$. Por lo tanto
existir\'a un homomorfismo $\phi : \grN_1 \To \graut{\grK}$ tal que $\grG = \grK {\rtimes}_{\phi } \grN_1$.

Ahora podemos dividir el problema en casos:

\section{Caso 1}

Asumimos $n_2 = 1$.


Entonces existe un $2$-subgrupo de Sylow $\grN$. Ahora, como es \'unico es caracter{\'\i}stico y por lo tanto normal.
Tambi\'en lo era $\grK$, de d\'onde como $\grK∩\grN=\set{\unity}$, $\grG = \grK\grN       \cong       \grK     \times     \grN$.


Salen dos grupos posibles:


\begin{flalign}
&\grG_1       \cong       \cyclicgr{9}       \times       \cyclicgr{2}       \cong       \cyclicgr{18}\\
&\grG_2       \cong       \cyclicgr{3}       \times       \cyclicgr{3}       \times       \cyclicgr{2}       \cong       \pdcyclics{6}{3}
\end{flalign}


\section{Caso 2}

Asumimos $n_2 = 9$.


Entonces existen $9$ $2$-subgrupos de Sylow $\grN$. Cada $2$-subgrupo de Sylow tiene un elemento de orden $2$ diferente.

Cada elemento que no est\'a en $\grK$ tiene orden $2$.

Sea $x       \in \grG\setminus\grK$ (tiene orden $2$). Podemos partir $\grG$ en clases de equivalencia laterales de
$\grK$ y por lo tanto $\grG=\grK\biguplus\grK x$. $\;\forall x       \in \grG\setminus\grK\;\; x^2       \in \grK$. Si $h       \in \grN$,
$hx       \notin       \grK$, de d\'onde su orden es $2$. $(hx)^2=\unity       \Rightarrow       hx=xh^{-1}       \Rightarrow       xhx=h^{-1}$.

\subsection{Caso 2.1}

Asumimos que $\grN       \cong       \cyclicgr{9}$.


Sea $\grN=\generadopor{y}$. Ahora $\forall x\; x^2=\unity       \Rightarrow        yx=x^{-1}y$
y tiene un solo grupo de orden $9$ que es normal y $9$ grupos no normales de orden
$2$ con la relaci\'on de inversi\'on ya dicha:

\begin{flalign}
&\grG_3       \cong       \dihgr{9}
\end{flalign}

Lo modelamos sobre el producto directo $\cyclicgr{9}      \times      \cyclicgr{2}$ con la siguiente acci\'on:
\begin{flalign*}
&\phi \; :\; \cyclicgr{2} \To \graut{\cyclicgr{9}}\; ::\; z^n \mapsto
\begin{cases}
  \generadopor{x      \mapsto x} &       \Leftarrow       n=0 \\
  \generadopor{x      \mapsto x^8} &       \Leftarrow       n=1
\end{cases}
\end{flalign*}


\subsection{Caso 2.2}

Asumimos que $\grN       \cong       \pdcyclics{3}{3}$.

Sea $\grN=\generadopor{y,z}$. $y^3=z^3=\unity$.

Cualquier elemento de orden $2$, digamos $x$, es tal que $xyx=y^2$ y $xzx=z^2$.

Solo cabe un grupo con estas caracter{\'\i}sticas.

\begin{flalign}
&\grG_4       \cong       \generadopor{x,y,z\;\vert\;\;x^3=y^3=z^2=\unity\;\ly\;xy=yx\;\ly\;xz=zx^2\;\ly\;yz=zy^2}
\end{flalign}

Lo modelamos sobre el producto directo $\cyclicgr{3}      \times      \cyclicgr{3}      \times      \cyclicgr{2}$ con la siguiente acci\'on:
\begin{flalign*}
&\phi \; :\; \cyclicgr{2} \To \graut{\pdcyclics{3}{3}}\; ::\; z^n \mapsto
\begin{cases}
  \generadopor{\begin{pmatrix}
                 x      \mapsto x \\
                 y      \mapsto y
               \end{pmatrix}} &       \Leftarrow       n=0 \\
  \generadopor{\begin{pmatrix}
                 x      \mapsto x^2 \\
                 y      \mapsto y^2
               \end{pmatrix}} &       \Leftarrow       n=1
\end{cases}
\end{flalign*}

\section{Caso 3}

Asumimos $n_2 = 3$.



\subsection{Caso 3.1}

Asumimos $\grK       \cong       \cyclicgr{9}$.

Sean $a,b,c$ tres elementos diferentes de orden $2$ que generan los tres grupos de orden $2$.
$\grN_1 = \generadopor{a}$, $\grN_2 = \generadopor{b}$ y $\grN_3 = \generadopor{c}$.

Sea $\grK=\generadopor{x}$ el subgrupo de orden $9$.

El grupo ser\'a $\grG=\set{\unity,a,x,x^2,x^3,x^4,x^5,x^6,x^7,x^8,ax,ax^2,ax^3,ax^4,ax^5,ax^6,ax^7,ax^8}$,
d\'onde tiene haber otros dos elementos $ax^i$ y $ax^j$ de orden $2$, y solamente esos dos.
Adem\'as $xa       \neq        ax$.

Sea $k       \in \set{2,3,4,5,6,7,8}$, y $xa=x^ka$. Entonces $x^iax^j=ax^{ik+j}$ y $ax^iax^j=x^{ik+j}$.

Buscamos un $h,l       \in \set{2,\ldots,8}$ tal que sean los \'unicos que $a^2=(ax^h)^2=(ax^l)^2=\unity$.
Supongamos que existen esos $k,h,l$. Como $9 ∣(k+1)h$ y $9 ∣(k+1)l$, ha de ser $k+1      \in \set{3,6}$,
esto es: $k      \in \set{2,5}$. Adem\'as $h$ y $l$ han de contribuir a esta ecuaci\'on siendo m\'ultiplos de $3$.
Consideraremos $0<i<j<9$. Solo queda, $(i,j)=(3,9)$.

Luego solo existen dos posibilidades: $k=2;h=3;l=6$ o $k=5;h=3;l=6$.

\subsubsection{Caso 3.1.1}

Asumimos $k=2$.


De esta forma $ax$ no puede ser de orden $9$, ni tampoco $2$. Solo puede ser de orden $3$ o $6$.
\[
  xa=ax^2
\]
\[
  (ax)(ax^7)=\unity       \Rightarrow       {(ax)}^{-1}=ax^7
\]
\[
  (ax^4)(ax)=\unity       \Rightarrow       {(ax)}^{-1}=ax^4
\]
\[
  ax^4=ax^7
\]
\[
  1=x^3 \mbox{  contra la hip\'otesis}
\]



\subsubsection{Caso 3.1.2}

Asumimos $k=5$.


De esta forma $ax$ no puede ser de orden $9$, ni tampoco $2$. Solo puede ser de orden $3$ o $6$.
\[
  xa=ax^5
\]
\[
  (ax)(ax^4)=\unity       \Rightarrow       {(ax)}^{-1}=ax^4
\]
\[
  (ax^4)(ax)=a^2x^{4     \cdot 5}x=x^{21-2     \cdot 9}=x^3
\]
\[
  1=x^3 \mbox{  contra la hip\'otesis}
\]


\subsection{Caso 3.2}

Asumimos $\grK       \cong       \pdcyclics{3}{3}$.

Tambi\'en tenemos que $n_2 = 3$ y que $n_9 = 1$.

Coincide con:

\begin{flalign}
&\grG_5        \cong        \dihgr{3}      \times      \cyclicgr{3}       \cong       \cyclicgr{3}      \times      \dihgr{3}       \cong       \cyclicgr{3}      \times      \left(\cyclicgr{3}\rtimes\cyclicgr{2}\right)
\end{flalign}

En la herramienta podemos modelarlo como $\grG_5       \cong       \entreangulos{a,b,z\vert a^3=b^3=z^2=\unity\;\ly\;ab=ba\;\ly\;az=za\;\ly\;zbz=b^2}$.
Se modela sobre un producto directo de $\cyclicgr{3}      \times      \cyclicgr{3}      \times      \cyclicgr{2}$ d\'onde solo hay que variar el comportamiento del
producto de las dos \'ultimas componentes.

\begin{flalign*}
&\phi \; :\; \cyclicgr{2} \To \graut{\pdcyclics{3}{3}}\; ::\; z^n \mapsto
\begin{cases}
  \generadopor{\begin{pmatrix}
                 x      \mapsto x \\
                 y      \mapsto y
               \end{pmatrix}} &       \Leftarrow       n=0 \\
  \generadopor{\begin{pmatrix}
                 x      \mapsto x \\
                 y      \mapsto y^2
               \end{pmatrix}} &       \Leftarrow       n=1
\end{cases}
\end{flalign*}


\section{El grupo G1}



$\grG_1      \cong      \cyclicgr{18}$


La tabla de Cayley es:


\smallskip\noindent
\resizebox{\linewidth}{!}{
\begin{tabular}{|c|c|c|c|c|c|c|c|c|c|c|c|c|c|c|c|c|c|}
\hline
$1$&$a$&$a^2$&$a^3$&$a^4$&$a^5$&$a^6$&$a^7$&$a^8$&$a^9$&$a^{10}$&$a^{11}$&$a^{12}$&$a^{13}$&$a^{14}$&$a^{15}$&$a^{16}$&$a^{17}$\\
\hline
$a$&$a^2$&$a^3$&$a^4$&$a^5$&$a^6$&$a^7$&$a^8$&$a^9$&$a^{10}$&$a^{11}$&$a^{12}$&$a^{13}$&$a^{14}$&$a^{15}$&$a^{16}$&$a^{17}$&$1$\\
\hline
$a^2$&$a^3$&$a^4$&$a^5$&$a^6$&$a^7$&$a^8$&$a^9$&$a^{10}$&$a^{11}$&$a^{12}$&$a^{13}$&$a^{14}$&$a^{15}$&$a^{16}$&$a^{17}$&$1$&$a$\\
\hline
$a^3$&$a^4$&$a^5$&$a^6$&$a^7$&$a^8$&$a^9$&$a^{10}$&$a^{11}$&$a^{12}$&$a^{13}$&$a^{14}$&$a^{15}$&$a^{16}$&$a^{17}$&$1$&$a$&$a^2$\\
\hline
$a^4$&$a^5$&$a^6$&$a^7$&$a^8$&$a^9$&$a^{10}$&$a^{11}$&$a^{12}$&$a^{13}$&$a^{14}$&$a^{15}$&$a^{16}$&$a^{17}$&$1$&$a$&$a^2$&$a^3$\\
\hline
$a^5$&$a^6$&$a^7$&$a^8$&$a^9$&$a^{10}$&$a^{11}$&$a^{12}$&$a^{13}$&$a^{14}$&$a^{15}$&$a^{16}$&$a^{17}$&$1$&$a$&$a^2$&$a^3$&$a^4$\\
\hline
$a^6$&$a^7$&$a^8$&$a^9$&$a^{10}$&$a^{11}$&$a^{12}$&$a^{13}$&$a^{14}$&$a^{15}$&$a^{16}$&$a^{17}$&$1$&$a$&$a^2$&$a^3$&$a^4$&$a^5$\\
\hline
$a^7$&$a^8$&$a^9$&$a^{10}$&$a^{11}$&$a^{12}$&$a^{13}$&$a^{14}$&$a^{15}$&$a^{16}$&$a^{17}$&$1$&$a$&$a^2$&$a^3$&$a^4$&$a^5$&$a^6$\\
\hline
$a^8$&$a^9$&$a^{10}$&$a^{11}$&$a^{12}$&$a^{13}$&$a^{14}$&$a^{15}$&$a^{16}$&$a^{17}$&$1$&$a$&$a^2$&$a^3$&$a^4$&$a^5$&$a^6$&$a^7$\\
\hline
$a^9$&$a^{10}$&$a^{11}$&$a^{12}$&$a^{13}$&$a^{14}$&$a^{15}$&$a^{16}$&$a^{17}$&$1$&$a$&$a^2$&$a^3$&$a^4$&$a^5$&$a^6$&$a^7$&$a^8$\\
\hline
$a^{10}$&$a^{11}$&$a^{12}$&$a^{13}$&$a^{14}$&$a^{15}$&$a^{16}$&$a^{17}$&$1$&$a$&$a^2$&$a^3$&$a^4$&$a^5$&$a^6$&$a^7$&$a^8$&$a^9$\\
\hline
$a^{11}$&$a^{12}$&$a^{13}$&$a^{14}$&$a^{15}$&$a^{16}$&$a^{17}$&$1$&$a$&$a^2$&$a^3$&$a^4$&$a^5$&$a^6$&$a^7$&$a^8$&$a^9$&$a^{10}$\\
\hline
$a^{12}$&$a^{13}$&$a^{14}$&$a^{15}$&$a^{16}$&$a^{17}$&$1$&$a$&$a^2$&$a^3$&$a^4$&$a^5$&$a^6$&$a^7$&$a^8$&$a^9$&$a^{10}$&$a^{11}$\\
\hline
$a^{13}$&$a^{14}$&$a^{15}$&$a^{16}$&$a^{17}$&$1$&$a$&$a^2$&$a^3$&$a^4$&$a^5$&$a^6$&$a^7$&$a^8$&$a^9$&$a^{10}$&$a^{11}$&$a^{12}$\\
\hline
$a^{14}$&$a^{15}$&$a^{16}$&$a^{17}$&$1$&$a$&$a^2$&$a^3$&$a^4$&$a^5$&$a^6$&$a^7$&$a^8$&$a^9$&$a^{10}$&$a^{11}$&$a^{12}$&$a^{13}$\\
\hline
$a^{15}$&$a^{16}$&$a^{17}$&$1$&$a$&$a^2$&$a^3$&$a^4$&$a^5$&$a^6$&$a^7$&$a^8$&$a^9$&$a^{10}$&$a^{11}$&$a^{12}$&$a^{13}$&$a^{14}$\\
\hline
$a^{16}$&$a^{17}$&$1$&$a$&$a^2$&$a^3$&$a^4$&$a^5$&$a^6$&$a^7$&$a^8$&$a^9$&$a^{10}$&$a^{11}$&$a^{12}$&$a^{13}$&$a^{14}$&$a^{15}$\\
\hline
$a^{17}$&$1$&$a$&$a^2$&$a^3$&$a^4$&$a^5$&$a^6$&$a^7$&$a^8$&$a^9$&$a^{10}$&$a^{11}$&$a^{12}$&$a^{13}$&$a^{14}$&$a^{15}$&$a^{16}$\\
\hline
\end{tabular}
}



Los \'ordenes de los elementos del grupo son:



\smallskip\noindent
\resizebox{\linewidth}{!}{
\begin{tabular}{|c|c|c|c|c|c|c|}
\hline
orden & =1 & =2 & =3 & =6 & =9 & =18\\
\hline
elementos & $\set{1}$ & $\set{a^9}$ & $\set{a^6,a^{12}}$ & $\set{a^3,a^{15}}$ & $\set{a^2,a^4,a^8,a^{10},a^{14},a^{16}}$ & $\set{a,a^5,a^7,a^{11},a^{13},a^{17}}$ \\
\hline
cardinal & 1 & 1 & 2 & 2 & 6 & 6 \\
\hline
\end{tabular}
}



Los subgrupos generados por 1 elemento o ninguno (subgrupo trivial) ordenados por sus \'ordenes coinciden en este caso con
el ret\'{\i}culo completo de los subgrupos, siendo adem\'as todos ellos abelianos, normales (característicos) y nilpotentes:



\begin{tabular}{|c|c|c|c|c|c|c|}
\hline
orden & =1 & =2 & =3 & =6 & =9 & =18\\
\hline
elementos & $\set{1}$ & $\generadopor{a^9}$ & $\generadopor{a^6}$ & $\generadopor{a^3}$ & $\generadopor{a^2}$ & $\generadopor{a}$ \\
\hline
isomorfo a & $\trivialgr$,1 & $\cyclicgr{2}$,1 & $\cyclicgr{3}$,1 & $\cyclicgr{6}$,1 & $\cyclicgr{9}$,1 & $\cyclicgr{18}$,1 \\
\hline
\end{tabular}

A partir de aqu\'{\i}, el \'unico subgrupo que merece inter\'es es el subgrupo de Frattini, que es $\frattini{\grG_1}=\generadopor{a^6}$,
 que hace que $\grcoc{\grG_1}{\frattini{\grG_1}}     \cong     \pdcyclics{3}{3}$. En una tabla resumimos los grupos con nombre propio:


\begin{tabular}{|c|c|c|c|c|}
\hline
Funci\'on $\mathfrak{F}\entrellaves{\grG}$ & Centro $\centro{\grG}$ & Derivado $\grderivado{\grG}{1}$ & Frattini $\frattini{\grG}$& Fitting $\fitting{\grG}$\\
\hline
Subgrupo & $\grG_1$ & $\set{1}$ & $\generadopor{a^6}$ & $\grG_1$ \\
\hline
Isomorfo a $\mathfrak{F}\entrellaves{\grG}     \cong     $ & $     \cong     \cyclicgr{18}$ & $     \cong     \trivialgr$ & $     \cong     \cyclicgr{3}$ & $     \cong     \cyclicgr{18}$ \\
\hline
$\grcoc{\grG}{\mathfrak{F}\entrellaves{\grG}}     \cong     $ & $     \cong     \trivialgr$ & $     \cong     \cyclicgr{18}$ & $     \cong     \pdcyclics{3}{3}$ & $     \cong     \trivialgr$ \\
\hline
\end{tabular}

En cuanto a las series centrales, el caso abeliano es trivial y viene ya dado en la anterior tabla
(no hay una $L$-derivada segunda ni un segundo centro). Solo decir que como abeliano es $\nil{1}$.

\section{El grupo G2}



$\grG_2      \cong      \pdcyclics{6}{3}$


La tabla de Cayley es:


\smallskip\noindent
\resizebox{\linewidth}{!}{
\begin{tabular}{|c|c|c|c|c|c|c|c|c|c|c|c|c|c|c|c|c|c|}
\hline
$1$&$a$&$a^2$&$a^3$&$a^4$&$a^5$&$b$&$ab$&$a^2b$&$a^3b$&$a^4b$&$a^5b$&$b^2$&$ab^2$&$a^2b^2$&$a^3b^2$&$a^4b^2$&$a^5b^2$\\
\hline
$a$&$a^2$&$a^3$&$a^4$&$a^5$&$1$&$ab$&$a^2b$&$a^3b$&$a^4b$&$a^5b$&$b$&$ab^2$&$a^2b^2$&$a^3b^2$&$a^4b^2$&$a^5b^2$&$b^2$\\
\hline
$a^2$&$a^3$&$a^4$&$a^5$&$1$&$a$&$a^2b$&$a^3b$&$a^4b$&$a^5b$&$b$&$ab$&$a^2b^2$&$a^3b^2$&$a^4b^2$&$a^5b^2$&$b^2$&$ab^2$\\
\hline
$a^3$&$a^4$&$a^5$&$1$&$a$&$a^2$&$a^3b$&$a^4b$&$a^5b$&$b$&$ab$&$a^2b$&$a^3b^2$&$a^4b^2$&$a^5b^2$&$b^2$&$ab^2$&$a^2b^2$\\
\hline
$a^4$&$a^5$&$1$&$a$&$a^2$&$a^3$&$a^4b$&$a^5b$&$b$&$ab$&$a^2b$&$a^3b$&$a^4b^2$&$a^5b^2$&$b^2$&$ab^2$&$a^2b^2$&$a^3b^2$\\
\hline
$a^5$&$1$&$a$&$a^2$&$a^3$&$a^4$&$a^5b$&$b$&$ab$&$a^2b$&$a^3b$&$a^4b$&$a^5b^2$&$b^2$&$ab^2$&$a^2b^2$&$a^3b^2$&$a^4b^2$\\
\hline
$b$&$ab$&$a^2b$&$a^3b$&$a^4b$&$a^5b$&$b^2$&$ab^2$&$a^2b^2$&$a^3b^2$&$a^4b^2$&$a^5b^2$&$1$&$a$&$a^2$&$a^3$&$a^4$&$a^5$\\
\hline
$ab$&$a^2b$&$a^3b$&$a^4b$&$a^5b$&$b$&$ab^2$&$a^2b^2$&$a^3b^2$&$a^4b^2$&$a^5b^2$&$b^2$&$a$&$a^2$&$a^3$&$a^4$&$a^5$&$1$\\
\hline
$a^2b$&$a^3b$&$a^4b$&$a^5b$&$b$&$ab$&$a^2b^2$&$a^3b^2$&$a^4b^2$&$a^5b^2$&$b^2$&$ab^2$&$a^2$&$a^3$&$a^4$&$a^5$&$1$&$a$\\
\hline
$a^3b$&$a^4b$&$a^5b$&$b$&$ab$&$a^2b$&$a^3b^2$&$a^4b^2$&$a^5b^2$&$b^2$&$ab^2$&$a^2b^2$&$a^3$&$a^4$&$a^5$&$1$&$a$&$a^2$\\
\hline
$a^4b$&$a^5b$&$b$&$ab$&$a^2b$&$a^3b$&$a^4b^2$&$a^5b^2$&$b^2$&$ab^2$&$a^2b^2$&$a^3b^2$&$a^4$&$a^5$&$1$&$a$&$a^2$&$a^3$\\
\hline
$a^5b$&$b$&$ab$&$a^2b$&$a^3b$&$a^4b$&$a^5b^2$&$b^2$&$ab^2$&$a^2b^2$&$a^3b^2$&$a^4b^2$&$a^5$&$1$&$a$&$a^2$&$a^3$&$a^4$\\
\hline
$b^2$&$ab^2$&$a^2b^2$&$a^3b^2$&$a^4b^2$&$a^5b^2$&$1$&$a$&$a^2$&$a^3$&$a^4$&$a^5$&$b$&$ab$&$a^2b$&$a^3b$&$a^4b$&$a^5b$\\
\hline
$ab^2$&$a^2b^2$&$a^3b^2$&$a^4b^2$&$a^5b^2$&$b^2$&$a$&$a^2$&$a^3$&$a^4$&$a^5$&$1$&$ab$&$a^2b$&$a^3b$&$a^4b$&$a^5b$&$b$\\
\hline
$a^2b^2$&$a^3b^2$&$a^4b^2$&$a^5b^2$&$b^2$&$ab^2$&$a^2$&$a^3$&$a^4$&$a^5$&$1$&$a$&$a^2b$&$a^3b$&$a^4b$&$a^5b$&$b$&$ab$\\
\hline
$a^3b^2$&$a^4b^2$&$a^5b^2$&$b^2$&$ab^2$&$a^2b^2$&$a^3$&$a^4$&$a^5$&$1$&$a$&$a^2$&$a^3b$&$a^4b$&$a^5b$&$b$&$ab$&$a^2b$\\
\hline
$a^4b^2$&$a^5b^2$&$b^2$&$ab^2$&$a^2b^2$&$a^3b^2$&$a^4$&$a^5$&$1$&$a$&$a^2$&$a^3$&$a^4b$&$a^5b$&$b$&$ab$&$a^2b$&$a^3b$\\
\hline
$a^5b^2$&$b^2$&$ab^2$&$a^2b^2$&$a^3b^2$&$a^4b^2$&$a^5$&$1$&$a$&$a^2$&$a^3$&$a^4$&$a^5b$&$b$&$ab$&$a^2b$&$a^3b$&$a^4b$\\
\hline
\end{tabular}
}



Los \'ordenes de los elementos del grupo son:



\smallskip\noindent
\resizebox{\linewidth}{!}{
\begin{tabular}{|c|c|c|c|c|c|c|}
\hline
orden & =1 & =2 & =3 & =6 & =9 & =18\\
\hline
elementos & $\set{1}$ & $\set{a^3}$ & $\set{a^2,a^4,b,a^2b,a^4b,b^2,a^2b^2,a^4b^2}$ & $\set{a,a^5,ab,a^3b,a^5b,ab^2,a^3b^2,a^5b^2}$ & ${}$ & ${}$ \\
\hline
cardinal & 1 & 1 & 8 & 8 & 0 & 0 \\
\hline
\end{tabular}
}



Los subgrupos ordenados por sus \'ordenes, todos ellos abelianos,
normales (no necesariamente caracter\'{\i}sticos) y nilpotentes:


\smallskip\noindent
\resizebox{\linewidth}{!}{
\begin{tabular}{|c|c|c|c|c|c|c|}
\hline
orden & =1 & =2 & =3 & =6 & =9 & =18\\
\hline
elementos & $\set{\set{1}}$ & $\set{\generadopor{a^3}}$ & $\set{\generadopor{a^2},\generadopor{b},\generadopor{a^2b},\generadopor{a^4b}}$ & $\set{\generadopor{a},\generadopor{a^3b},\generadopor{ab},\generadopor{ab^2}}$ & $\set{\generadopor{a^2,b}}$ & $\set{\generadopor{a,b}}$ \\
\hline
isomorfo a, card & $\trivialgr$,1 & $\cyclicgr{2}$,1 & $\cyclicgr{3}$,4 & $\cyclicgr{6}$,4 & \pdcyclics{3}{3},1 & \pdcyclics{6}{3},1 \\
\hline
\end{tabular}
}


A partir de aqu\'{\i}, el \'unico subgrupo que merece inter\'es es el subgrupo intersecci\'on de los maximales c\'{\i}clicos, que es $\bigcap_{x     \in \grG_2\ly\generadopor{x}\;maximal\;como\;ciclico}{\generadopor{x}}=\generadopor{a^3}$, que es caracter\'{\i}stico, al igual que el cociente correspondiente que es el de orden $9$ (tambi\'en caracter\'{\i}stico). En una tabla resumimos los grupos con nombre propio:



\begin{tabular}{|c|c|c|c|c|}
\hline
Funci\'on $\mathfrak{F}\entrellaves{\grG}$ & Centro $\centro{\grG}$ & Derivado $\grderivado{\grG}{1}$ & Frattini $\frattini{\grG}$& Fitting $\fitting{\grG}$\\
\hline
Subgrupo & $\grG_2$ & $\set{1}$ & $\set{1}$ & $\grG_2$ \\
\hline
Isomorfo a $\mathfrak{F}\entrellaves{\grG}     \cong     $ & $     \cong     \pdcyclics{6}{3}$ & $     \cong     \trivialgr$ & $     \cong     \trivialgr$ & $     \cong     \pdcyclics{6}{3}$ \\
\hline
$\grcoc{\grG}{\mathfrak{F}\entrellaves{\grG}}     \cong     $ & $     \cong     \trivialgr$ & $     \cong     \pdcyclics{6}{3}$ & $     \cong     \pdcyclics{6}{3}$ & $     \cong     \trivialgr$ \\
\hline
\end{tabular}



En cuanto a las series centrales, el caso abeliano es trivial y viene ya dado en la anterior tabla
(no hay una $L$-derivada segunda ni un segundo centro). Solo decir que como abeliano es $\nil{1}$.


\section{El grupo G3}





$\grG_3      \cong      \dihgr{9}$


La tabla de Cayley es:


\smallskip\noindent
\resizebox{\linewidth}{!}{
\begin{tabular}{|c|c|c|c|c|c|c|c|c|c|c|c|c|c|c|c|c|c|}
\hline
$1$&$x$&$x^2$&$x^3$&$x^4$&$x^5$&$x^6$&$x^7$&$x^8$&$z$&$xz$&$x^2z$&$x^3z$&$x^4z$&$x^5z$&$x^6z$&$x^7z$&$x^8z$\\
\hline
$x$&$x^2$&$x^3$&$x^4$&$x^5$&$x^6$&$x^7$&$x^8$&$1$&$x^8z$&$z$&$xz$&$x^2z$&$x^3z$&$x^4z$&$x^5z$&$x^6z$&$x^7z$\\
\hline
$x^2$&$x^3$&$x^4$&$x^5$&$x^6$&$x^7$&$x^8$&$1$&$x$&$x^7z$&$x^8z$&$z$&$xz$&$x^2z$&$x^3z$&$x^4z$&$x^5z$&$x^6z$\\
\hline
$x^3$&$x^4$&$x^5$&$x^6$&$x^7$&$x^8$&$1$&$x$&$x^2$&$x^6z$&$x^7z$&$x^8z$&$z$&$xz$&$x^2z$&$x^3z$&$x^4z$&$x^5z$\\
\hline
$x^4$&$x^5$&$x^6$&$x^7$&$x^8$&$1$&$x$&$x^2$&$x^3$&$x^5z$&$x^6z$&$x^7z$&$x^8z$&$z$&$xz$&$x^2z$&$x^3z$&$x^4z$\\
\hline
$x^5$&$x^6$&$x^7$&$x^8$&$1$&$x$&$x^2$&$x^3$&$x^4$&$x^4z$&$x^5z$&$x^6z$&$x^7z$&$x^8z$&$z$&$xz$&$x^2z$&$x^3z$\\
\hline
$x^6$&$x^7$&$x^8$&$1$&$x$&$x^2$&$x^3$&$x^4$&$x^5$&$x^3z$&$x^4z$&$x^5z$&$x^6z$&$x^7z$&$x^8z$&$z$&$xz$&$x^2z$\\
\hline
$x^7$&$x^8$&$1$&$x$&$x^2$&$x^3$&$x^4$&$x^5$&$x^6$&$x^2z$&$x^3z$&$x^4z$&$x^5z$&$x^6z$&$x^7z$&$x^8z$&$z$&$xz$\\
\hline
$x^8$&$1$&$x$&$x^2$&$x^3$&$x^4$&$x^5$&$x^6$&$x^7$&$xz$&$x^2z$&$x^3z$&$x^4z$&$x^5z$&$x^6z$&$x^7z$&$x^8z$&$z$\\
\hline
$z$&$xz$&$x^2z$&$x^3z$&$x^4z$&$x^5z$&$x^6z$&$x^7z$&$x^8z$&$1$&$x$&$x^2$&$x^3$&$x^4$&$x^5$&$x^6$&$x^7$&$x^8$\\
\hline
$xz$&$x^2z$&$x^3z$&$x^4z$&$x^5z$&$x^6z$&$x^7z$&$x^8z$&$z$&$x^8$&$1$&$x$&$x^2$&$x^3$&$x^4$&$x^5$&$x^6$&$x^7$\\
\hline
$x^2z$&$x^3z$&$x^4z$&$x^5z$&$x^6z$&$x^7z$&$x^8z$&$z$&$xz$&$x^7$&$x^8$&$1$&$x$&$x^2$&$x^3$&$x^4$&$x^5$&$x^6$\\
\hline
$x^3z$&$x^4z$&$x^5z$&$x^6z$&$x^7z$&$x^8z$&$z$&$xz$&$x^2z$&$x^6$&$x^7$&$x^8$&$1$&$x$&$x^2$&$x^3$&$x^4$&$x^5$\\
\hline
$x^4z$&$x^5z$&$x^6z$&$x^7z$&$x^8z$&$z$&$xz$&$x^2z$&$x^3z$&$x^5$&$x^6$&$x^7$&$x^8$&$1$&$x$&$x^2$&$x^3$&$x^4$\\
\hline
$x^5z$&$x^6z$&$x^7z$&$x^8z$&$z$&$xz$&$x^2z$&$x^3z$&$x^4z$&$x^4$&$x^5$&$x^6$&$x^7$&$x^8$&$1$&$x$&$x^2$&$x^3$\\
\hline
$x^6z$&$x^7z$&$x^8z$&$z$&$xz$&$x^2z$&$x^3z$&$x^4z$&$x^5z$&$x^3$&$x^4$&$x^5$&$x^6$&$x^7$&$x^8$&$1$&$x$&$x^2$\\
\hline
$x^7z$&$x^8z$&$z$&$xz$&$x^2z$&$x^3z$&$x^4z$&$x^5z$&$x^6z$&$x^2$&$x^3$&$x^4$&$x^5$&$x^6$&$x^7$&$x^8$&$1$&$x$\\
\hline
$x^8z$&$z$&$xz$&$x^2z$&$x^3z$&$x^4z$&$x^5z$&$x^6z$&$x^7z$&$x$&$x^2$&$x^3$&$x^4$&$x^5$&$x^6$&$x^7$&$x^8$&$1$\\
\hline
\end{tabular}

}



Los \'ordenes de los elementos del grupo son:




\smallskip\noindent
\resizebox{\linewidth}{!}{
\begin{tabular}{|c|c|c|c|c|c|c|}
\hline
orden & =1 & =2 & =3 & =6 & =9 & =18\\
\hline
elementos & $\set{1}$ & $\set{z,xz,x^2z,x^3z,x^4z,x^5z,x^6z,x^7z,x^8z}$ & $\set{x^3,x^6}$ & $\set{}$ & $\set{x,x^2,x^4,x^5,x^7,x^8}$ & $\set{}$ \\
\hline
cardinal & 1 & 9 & 2 & 0 & 6 & 0 \\
\hline
\end{tabular}
}



El centro del grupo es $\centro{\grG_3}=\set{\unity}$ y el derivado es $\grderivado{\grG_3}{1}=\generadopor{x}     \cong     \cyclicgr{9}$. Esto nos sirve para poder clasificar los subgrupos que a continuaci\'on vamos a listar.



El ret\'{\i}culo de subgrupos queda as\'{\i} clasificado:



\smallskip\noindent
\resizebox{\linewidth}{!}{
\begin{tabular}{|c|c|c|c|c|c|c|}
\hline
orden & =1 & =2 & =3 & =6 & =9 & =18\\
\hline
elementos
&$\set{\set{1}}$
&$\set{\generadopor{z},\generadopor{xz},\generadopor{x^2z},\generadopor{x^3z},\generadopor{x^4z},\generadopor{x^5z},\generadopor{x^6z},\generadopor{x^7z},\generadopor{x^8z}}$
&$\set{\generadopor{x^3}}$
&$\set{\generadopor{x^3,z},\generadopor{x^3,xz},\generadopor{x^3,x^2z}}$
&$\set{\generadopor{x}}$
&$\set{\generadopor{x,z}}$ \\
\hline
isomorfo a, card & $\trivialgr$,1 & $\cyclicgr{2}$,9 & $\cyclicgr{3}$,1 & $\dihgr{3}$,3 & \cyclicgr{9},1 & \dihgr{9},1 \\
\hline
propiedades & cent & cyc & ch,cyc,der & {} & ch,cyc,der & {} \\
\hline
\end{tabular}
}

Para la fila de propiedades de los subgrupos la leyenda es: ``norm'' significa normal, ``ch'' significa caracter\'{\i}stico, ``der'' significa que est\'a contenido en el derivado, ``cent'' significa que est\'a en el centro, ``ab'' que es abeliano, ``cyc'' que es c\'{\i}clico, ``nil'' que es nilpotente.


Los grupos con nombre propio que quedan son solo los de Frattini y de Fitting.


\begin{tabular}{|c|c|c|c|c|}
\hline
Funci\'on $\mathfrak{F}\entrellaves{\grG}$ & Centro $\centro{\grG}$ & Derivado $\grderivado{\grG}{1}$ & Frattini $\frattini{\grG}$& Fitting $\fitting{\grG}$\\
\hline
Subgrupo & $\set{\unity}$ & $\generadopor{x}$ & $\set{\unity}$ & $\generadopor{x}$ \\
\hline
Isomorfo a $\mathfrak{F}\entrellaves{\grG}     \cong     $ & $     \cong     \trivialgr$ & $     \cong     \dihgr{9}$ & $     \cong     \trivialgr$ & $     \cong     \dihgr{9}$ \\
\hline
$\grcoc{\grG}{\mathfrak{F}\entrellaves{\grG}}     \cong     $ & $     \cong     \dihgr{9}$ & $     \cong     \cyclicgr{2}$ & $     \cong     \dihgr{9}$ & $     \cong     \cyclicgr{2}$ \\
\hline
\end{tabular}



En cuanto a las series centrales, hay que tener en cuenta que este grupo es resoluble pero no nilpotente, as\'{\i} que las series centrales se estabilizan antes del pretendido final.

\textbf{Serie Central Ascendente:}
\begin{flalign*}
&\centro{\grG_3}=\set{1}\subgreq\centronum{\grG_3}{2}=\set{1}\subgreq\ldots\subgreq\centronum{\grG_3}{n}=\set{1}\subgreq\ldots\quad n     \in \setN
\end{flalign*}



\textbf{Serie Central Ascendente:}
\begin{flalign*}
&\grG_3\rhd\grLderivado{\grG_3}{1}=\generadopor{x}\unrhd\ldots\unrhd\grLderivado{\grG_3}{n}=\generadopor{x}\unrhd\ldots\quad n     \in \setN
\end{flalign*}



\textbf{Serie Derivada:}
\begin{flalign*}
&\grG_3\rhd\grderivado{\grG_3}{1}=\generadopor{x}\rhd\set{\unity}=\grderivado{\grG_3}{2}     \cong     \trivialgr
\end{flalign*}



\section{El grupo G4}






$\grG_4      \cong      \psd{\left(\pdcyclics{3}{3}\right)}{\cyclicgr{2}}$


La tabla de Cayley es:


\smallskip\noindent
\resizebox{\linewidth}{!}{
\begin{tabular}{|c|c|c|c|c|c|c|c|c|c|c|c|c|c|c|c|c|c|}
\hline
$1$&$a$&$a^2$&$b$&$ab$&$a^2b$&$b^2$&$ab^2$&$a^2b^2$&$c$&$ac$&$a^2c$&$bc$&$abc$&$a^2bc$&$b^2c$&$ab^2c$&$a^2b^2c$\\
\hline
$a$&$a^2$&$1$&$ab$&$a^2b$&$b$&$ab^2$&$a^2b^2$&$b^2$&$a^2c$&$c$&$ac$&$a^2bc$&$bc$&$abc$&$a^2b^2c$&$b^2c$&$ab^2c$\\
\hline
$a^2$&$1$&$a$&$a^2b$&$b$&$ab$&$a^2b^2$&$b^2$&$ab^2$&$ac$&$a^2c$&$c$&$abc$&$a^2bc$&$bc$&$ab^2c$&$a^2b^2c$&$b^2c$\\
\hline
$b$&$ab$&$a^2b$&$b^2$&$ab^2$&$a^2b^2$&$1$&$a$&$a^2$&$b^2c$&$ab^2c$&$a^2b^2c$&$c$&$ac$&$a^2c$&$bc$&$abc$&$a^2bc$\\
\hline
$ab$&$a^2b$&$b$&$ab^2$&$a^2b^2$&$b^2$&$a$&$a^2$&$1$&$a^2b^2c$&$b^2c$&$ab^2c$&$a^2c$&$c$&$ac$&$a^2bc$&$bc$&$abc$\\
\hline
$a^2b$&$b$&$ab$&$a^2b^2$&$b^2$&$ab^2$&$a^2$&$1$&$a$&$ab^2c$&$a^2b^2c$&$b^2c$&$ac$&$a^2c$&$c$&$abc$&$a^2bc$&$bc$\\
\hline
$b^2$&$ab^2$&$a^2b^2$&$1$&$a$&$a^2$&$b$&$ab$&$a^2b$&$bc$&$abc$&$a^2bc$&$b^2c$&$ab^2c$&$a^2b^2c$&$c$&$ac$&$a^2c$\\
\hline
$ab^2$&$a^2b^2$&$b^2$&$a$&$a^2$&$1$&$ab$&$a^2b$&$b$&$a^2bc$&$bc$&$abc$&$a^2b^2c$&$b^2c$&$ab^2c$&$a^2c$&$c$&$ac$\\
\hline
$a^2b^2$&$b^2$&$ab^2$&$a^2$&$1$&$a$&$a^2b$&$b$&$ab$&$abc$&$a^2bc$&$bc$&$ab^2c$&$a^2b^2c$&$b^2c$&$ac$&$a^2c$&$c$\\
\hline
$c$&$ac$&$a^2c$&$bc$&$abc$&$a^2bc$&$b^2c$&$ab^2c$&$a^2b^2c$&$1$&$a$&$a^2$&$b$&$ab$&$a^2b$&$b^2$&$ab^2$&$a^2b^2$\\
\hline
$ac$&$a^2c$&$c$&$abc$&$a^2bc$&$bc$&$ab^2c$&$a^2b^2c$&$b^2c$&$a^2$&$1$&$a$&$a^2b$&$b$&$ab$&$a^2b^2$&$b^2$&$ab^2$\\
\hline
$a^2c$&$c$&$ac$&$a^2bc$&$bc$&$abc$&$a^2b^2c$&$b^2c$&$ab^2c$&$a$&$a^2$&$1$&$ab$&$a^2b$&$b$&$ab^2$&$a^2b^2$&$b^2$\\
\hline
$bc$&$abc$&$a^2bc$&$b^2c$&$ab^2c$&$a^2b^2c$&$c$&$ac$&$a^2c$&$b^2$&$ab^2$&$a^2b^2$&$1$&$a$&$a^2$&$b$&$ab$&$a^2b$\\
\hline
$abc$&$a^2bc$&$bc$&$ab^2c$&$a^2b^2c$&$b^2c$&$ac$&$a^2c$&$c$&$a^2b^2$&$b^2$&$ab^2$&$a^2$&$1$&$a$&$a^2b$&$b$&$ab$\\
\hline
$a^2bc$&$bc$&$abc$&$a^2b^2c$&$b^2c$&$ab^2c$&$a^2c$&$c$&$ac$&$ab^2$&$a^2b^2$&$b^2$&$a$&$a^2$&$1$&$ab$&$a^2b$&$b$\\
\hline
$b^2c$&$ab^2c$&$a^2b^2c$&$c$&$ac$&$a^2c$&$bc$&$abc$&$a^2bc$&$b$&$ab$&$a^2b$&$b^2$&$ab^2$&$a^2b^2$&$1$&$a$&$a^2$\\
\hline
$ab^2c$&$a^2b^2c$&$b^2c$&$ac$&$a^2c$&$c$&$abc$&$a^2bc$&$bc$&$a^2b$&$b$&$ab$&$a^2b^2$&$b^2$&$ab^2$&$a^2$&$1$&$a$\\
\hline
$a^2b^2c$&$b^2c$&$ab^2c$&$a^2c$&$c$&$ac$&$a^2bc$&$bc$&$abc$&$ab$&$a^2b$&$b$&$ab^2$&$a^2b^2$&$b^2$&$a$&$a^2$&$1$\\
\hline
\end{tabular}
}



Los \'ordenes de los elementos del grupo son:



\smallskip\noindent
\resizebox{\linewidth}{!}{
\begin{tabular}{|c|c|c|c|c|c|c|}
\hline
orden & =1 & =2 & =3 & =6 & =9 & =18\\
\hline
elementos & $\set{1}$ & $\set{c,bc,b^2c}$ & $\set{a,a^2,b,ab,a^2b,b^2,ab^2,a^2b^2}$ & $\set{ac,a^2c,abc,a^3bc,ab^2c,a^2b^2c}$ & $\set{}$ & $\set{}$ \\
\hline
cardinal & 1 & 3 & 8 & 6 & 0 & 0 \\
\hline
\end{tabular}
}



El centro del grupo es $\centro{\grG_4}=\set{\unity}    \cong    \trivialgr$ y el derivado es $\grderivado{\grG_4}{1}=\generadopor{a,b}     \cong     \pdcyclics{3}{3}$. Esto nos sirve para poder clasificar los subgrupos que a continuaci\'on vamos a listar.



El ret\'{\i}culo de subgrupos queda as\'{\i} clasificado:



\smallskip\noindent
\resizebox{\linewidth}{!}{
\begin{tabular}{|c|c|c|c|c|c|c|}
\hline
orden & =1 & =2 & =3 & =6 & =9 & =18\\
\hline
elementos
&$\set{\set{1}}$
&$\set{\generadopor{c},\generadopor{ac},\generadopor{a^2c},\generadopor{bc},\generadopor{b^2c},\generadopor{abc},\generadopor{a^2bc},\generadopor{ab^2c},\generadopor{a^2b^2c}}$
&$\set{\generadopor{a},\generadopor{b},\generadopor{ab},\generadopor{a^2b}}$
&$\set{
\generadopor{a,c},
\generadopor{a,bc},
\generadopor{a,b^2c},
\generadopor{b,c},
\generadopor{b,ac},
\generadopor{b,a^2c},
\generadopor{ab,c},
\generadopor{ab,ac},
\generadopor{ab,a^2c},
\generadopor{a^2b,c},
\generadopor{a^2b,ac},
\generadopor{a^2b,a^2c}
}$
&$\set{\generadopor{a,b}}$
&$\set{\generadopor{a,b,c}}$ \\
\hline
isomorfo a, card
&$\trivialgr$,1
&$\cyclicgr{2}$,12
&$\cyclicgr{3}$,4
&$\cyclicgr{6}$,3 $\vert\;\;\dihgr{3}$,1
&$\pdcyclics{3}{3},1$
&$\pd{\cyclicgr{3}}{\dihgr{3}}$,1 \\
\hline
propiedades & {} & (cyc),12 & (cyc,norm,der),3 & 3 & (ab,ch,der),1 & {} \\
\hline
\end{tabular}
}

Para la fila de propiedades de los subgrupos la leyenda es: ``norm'' significa normal, ``ch'' significa caracter\'{\i}stico, ``der'' significa que est\'a contenido en el derivado, ``cent'' significa que est\'a en el centro, ``ab'' que es abeliano, ``cyc'' que es c\'{\i}clico, ``nil'' que es nilpotente.


Los grupos con nombre propio que quedan son solo los de Frattini y de Fitting.


\begin{tabular}{|c|c|c|c|c|}
\hline
Funci\'on $\mathfrak{F}\entrellaves{\grG}$ & Centro $\centro{\grG}$ & Derivado $\grderivado{\grG}{1}$ & Frattini $\frattini{\grG}$& Fitting $\fitting{\grG}$\\
\hline
Subgrupo & $\set{\unity}$ & $\generadopor{a,b}$ & $\set{\unity}$ & $\generadopor{a,b}$ \\
\hline
Isomorfo a $\mathfrak{F}\entrellaves{\grG}     \cong     $ & $     \cong     \trivialgr$ & $     \cong     \pdcyclics{3}{3}$ & $     \cong     \trivialgr$ & $     \cong     \pdcyclics{3}{3}$ \\
\hline
$\grcoc{\grG}{\mathfrak{F}\entrellaves{\grG}}     \cong     $ & $     \cong     \psd{\pdcyclics{3}{3}}{\cyclicgr{2}}$ & $     \cong     \cyclicgr{2}$ & $     \cong     \psd{\pdcyclics{3}{3}}{\cyclicgr{2}}$ & $     \cong     \cyclicgr{2}$ \\
\hline
\end{tabular}





En cuanto a las series centrales, hay que tener en cuenta que este grupo es resoluble pero no nilpotente, as\'{\i} que las series centrales se estabilizan antes del pretendido final.

\textbf{Serie Central Ascendente:}
\begin{flalign*}
&\centro{\grG_4}=\set{1}\subgreq\centronum{\grG_4}{2}=\set{1}\subgreq\ldots\subgreq\centronum{\grG_4}{n}=\set{1}\subgreq\ldots\quad n     \in \setN
\end{flalign*}



\textbf{Serie Central Descendente:}
\begin{flalign*}
&\grG_4\rhd\grLderivado{\grG_4}{1}=\generadopor{a,b}\unrhd\ldots\unrhd\grLderivado{\grG_4}{n}=\generadopor{a,b}\unrhd\ldots\quad n     \in \setN
\end{flalign*}



\textbf{Serie Derivada:}
\begin{flalign*}
&\grG_4\rhd\grderivado{\grG_4}{1}=\generadopor{a,b}\rhd\set{\unity}=\grderivado{\grG_4}{2}     \cong     \trivialgr
\end{flalign*}





\section{El grupo G5}







$\grG_5      \cong      \pd{\cyclicgr{3}}{\dihgr{3}}$


La tabla de Cayley es:


\smallskip\noindent
\resizebox{\linewidth}{!}{
\begin{tabular}{|c|c|c|c|c|c|c|c|c|c|c|c|c|c|c|c|c|c|}
\hline
$1$&$a$&$a^2$&$b$&$ab$&$a^2b$&$b^2$&$ab^2$&$a^2b^2$&$c$&$ac$&$a^2c$&$bc$&$abc$&$a^2bc$&$b^2c$&$ab^2c$&$a^2b^2c$\\
\hline
$a$&$a^2$&$1$&$ab$&$a^2b$&$b$&$ab^2$&$a^2b^2$&$b^2$&$ac$&$a^2c$&$c$&$abc$&$a^2bc$&$bc$&$ab^2c$&$a^2b^2c$&$b^2c$\\
\hline
$a^2$&$1$&$a$&$a^2b$&$b$&$ab$&$a^2b^2$&$b^2$&$ab^2$&$a^2c$&$c$&$ac$&$a^2bc$&$bc$&$abc$&$a^2b^2c$&$b^2c$&$ab^2c$\\
\hline
$b$&$ab$&$a^2b$&$b^2$&$ab^2$&$a^2b^2$&$1$&$a$&$a^2$&$b^2c$&$ab^2c$&$a^2b^2c$&$c$&$ac$&$a^2c$&$bc$&$abc$&$a^2bc$\\
\hline
$ab$&$a^2b$&$b$&$ab^2$&$a^2b^2$&$b^2$&$a$&$a^2$&$1$&$ab^2c$&$a^2b^2c$&$b^2c$&$ac$&$a^2c$&$c$&$abc$&$a^2bc$&$bc$\\
\hline
$a^2b$&$b$&$ab$&$a^2b^2$&$b^2$&$ab^2$&$a^2$&$1$&$a$&$a^2b^2c$&$b^2c$&$ab^2c$&$a^2c$&$c$&$ac$&$a^2bc$&$bc$&$abc$\\
\hline
$b^2$&$ab^2$&$a^2b^2$&$1$&$a$&$a^2$&$b$&$ab$&$a^2b$&$bc$&$abc$&$a^2bc$&$b^2c$&$ab^2c$&$a^2b^2c$&$c$&$ac$&$a^2c$\\
\hline
$ab^2$&$a^2b^2$&$b^2$&$a$&$a^2$&$1$&$ab$&$a^2b$&$b$&$abc$&$a^2bc$&$bc$&$ab^2c$&$a^2b^2c$&$b^2c$&$ac$&$a^2c$&$c$\\
\hline
$a^2b^2$&$b^2$&$ab^2$&$a^2$&$1$&$a$&$a^2b$&$b$&$ab$&$a^2bc$&$bc$&$abc$&$a^2b^2c$&$b^2c$&$ab^2c$&$a^2c$&$c$&$ac$\\
\hline
$c$&$ac$&$a^2c$&$bc$&$abc$&$a^2bc$&$b^2c$&$ab^2c$&$a^2b^2c$&$1$&$a$&$a^2$&$b$&$ab$&$a^2b$&$b^2$&$ab^2$&$a^2b^2$\\
\hline
$ac$&$a^2c$&$c$&$abc$&$a^2bc$&$bc$&$ab^2c$&$a^2b^2c$&$b^2c$&$a$&$a^2$&$1$&$ab$&$a^2b$&$b$&$ab^2$&$a^2b^2$&$b^2$\\
\hline
$a^2c$&$c$&$ac$&$a^2bc$&$bc$&$abc$&$a^2b^2c$&$b^2c$&$ab^2c$&$a^2$&$1$&$a$&$a^2b$&$b$&$ab$&$a^2b^2$&$b^2$&$ab^2$\\
\hline
$bc$&$abc$&$a^2bc$&$b^2c$&$ab^2c$&$a^2b^2c$&$c$&$ac$&$a^2c$&$b^2$&$ab^2$&$a^2b^2$&$1$&$a$&$a^2$&$b$&$ab$&$a^2b$\\
\hline
$abc$&$a^2bc$&$bc$&$ab^2c$&$a^2b^2c$&$b^2c$&$ac$&$a^2c$&$c$&$ab^2$&$a^2b^2$&$b^2$&$a$&$a^2$&$1$&$ab$&$a^2b$&$b$\\
\hline
$a^2bc$&$bc$&$abc$&$a^2b^2c$&$b^2c$&$ab^2c$&$a^2c$&$c$&$ac$&$a^2b^2$&$b^2$&$ab^2$&$a^2$&$1$&$a$&$a^2b$&$b$&$ab$\\
\hline
$b^2c$&$ab^2c$&$a^2b^2c$&$c$&$ac$&$a^2c$&$bc$&$abc$&$a^2bc$&$b$&$ab$&$a^2b$&$b^2$&$ab^2$&$a^2b^2$&$1$&$a$&$a^2$\\
\hline
$ab^2c$&$a^2b^2c$&$b^2c$&$ac$&$a^2c$&$c$&$abc$&$a^2bc$&$bc$&$ab$&$a^2b$&$b$&$ab^2$&$a^2b^2$&$b^2$&$a$&$a^2$&$1$\\
\hline
$a^2b^2c$&$b^2c$&$ab^2c$&$a^2c$&$c$&$ac$&$a^2bc$&$bc$&$abc$&$a^2b$&$b$&$ab$&$a^2b^2$&$b^2$&$ab^2$&$a^2$&$1$&$a$\\
\hline
\end{tabular}
}

Por simple inspecci\'on de la tabla, vemos que $\grG_5 = $ $\left\lbrace\right.\generadopor{a}\generadopor{b,c}$ $\vert$ $b^3=c^2=\unity$ $\ly$ $cbc=b^2$ $    \cong    $ $\pd{\cyclicgr{3}}{\dihgr{3}}\left.\right\rbrace$.



Los \'ordenes de los elementos del grupo son:



\smallskip\noindent
\resizebox{\linewidth}{!}{
\begin{tabular}{|c|c|c|c|c|c|c|}
\hline
orden & =1 & =2 & =3 & =6 & =9 & =18\\
\hline
elementos & $\set{1}$ & $\set{c,bc,b^2c}$ & $\set{a,a^2,b,ab,a^2b,b^2,ab^2,a^2b^2}$ & $\set{ac,a^2c,abc,a^2bc,ab^2c,a^2b^2c}$ & $\set{}$ & $\set{}$ \\
\hline
cardinal & 1 & 2 & 8 & 6 & 0 & 0 \\
\hline
\end{tabular}
}





El centro del grupo es $\centro{\grG_5}=\generadopor{a}    \cong    \cyclicgr{3}$ y el derivado es $\grderivado{\grG_5}{1}=\generadopor{b}     \cong     \cyclicgr{3}$. Esto nos sirve para poder clasificar los subgrupos que a continuaci\'on vamos a listar.



El ret\'{\i}culo de subgrupos queda as\'{\i} clasificado:


\smallskip\noindent
\resizebox{\linewidth}{!}{
\begin{tabular}{|c|c|c|c|c|c|c|}
\hline
orden & =1 & =2 & =3 & =6 & =9 & =18\\
\hline
elementos
&$\set{\set{\unity}}$
&$\set{\generadopor{c},\generadopor{bc},\generadopor{b^2c}}$
&$\set{\generadopor{a}} \vert \set{\generadopor{b}} \vert \set{\generadopor{ab},\generadopor{a^2b}}$
&$\set{\generadopor{c},\generadopor{bc},\generadopor{b^2c}} \vert \set{\generadopor{b,c}}$
&$\set{\generadopor{a,b}}$
&$\set{\generadopor{a,b,c}}$ \\
\hline
isomorfo a, card
&$\trivialgr,1$
&$\cyclicgr{2},3$
&$\cyclicgr{3},1\vert\cyclicgr{3},1\vert\cyclicgr{3},2$
&$\cyclicgr{6},3\vert\dihgr{3},1$
&$\pdcyclics{3}{3},1$
&$\pd{\cyclicgr{3}}{\dihgr{3}},1$ \\
\hline
propiedades
&{}
&(cyc),3
&cyc: (ch,cent),1$\vert$(ch,der),1$\vert$(norm),2
&(cyc),3$\vert$(ch),1
&(ab,ch),1
&{} \\
\hline
\end{tabular}
}

Para la fila de propiedades de los subgrupos la leyenda es: ``norm'' significa normal, ``ch'' significa caracter\'{\i}stico, ``der'' significa que est\'a contenido en el derivado, ``cent'' significa que est\'a en el centro, ``ab'' que es abeliano, ``cyc'' que es c\'{\i}clico, ``nil'' que es nilpotente.


Los grupos con nombre propio que quedan son solo los de Frattini y de Fitting.


\begin{tabular}{|c|c|c|c|c|}
\hline
Funci\'on $\mathfrak{F}\entrellaves{\grG}$
& Centro $\centro{\grG}$
& Derivado $\grderivado{\grG}{1}$
& Frattini $\frattini{\grG}$
& Fitting $\fitting{\grG}$\\
\hline
Subgrupo
& $\generadopor{a}$
& $\generadopor{b}$
& $\set{\unity}$
& $\generadopor{a,b}$ \\
\hline
Isomorfo a $\mathfrak{F}\entrellaves{\grG}     \cong     $
& $     \cong     \cyclicgr{3}$
& $     \cong     \cyclicgr{3}$
& $     \cong     \trivialgr$
& $     \cong     \pdcyclics{3}{3}$ \\
\hline
$\grcoc{\grG}{\mathfrak{F}\entrellaves{\grG}}     \cong     $
& $     \cong     \dihgr{3}$
& $     \cong     \pdcyclics{3}{2}$
& $     \cong     \pd{\cyclicgr{3}}{\dihgr{3}}$
& $     \cong     \cyclicgr{2}$ \\
\hline
\end{tabular}





En cuanto a las series centrales, hay que tener en cuenta que este grupo es resoluble pero no nilpotente, as\'{\i} que las series centrales se estabilizan antes del pretendido final.

\textbf{Serie Central Ascendente:}
\begin{flalign*}
&\set{1}\unlhd\centro{\grG_5}=
\generadopor{a}\unlhd\centronum{\grG_5}{2}=
\generadopor{a}\unlhd\ldots\unrhd\centronum{\grG_5}{n}=
\generadopor{a}\unlhd\ldots
\quad n     \in \setN
\end{flalign*}



\textbf{Serie Central Descendente:}
\begin{flalign*}
&\grG_5\rhd\grLderivado{\grG_5}{1}=
\generadopor{b}\unrhd\ldots\unrhd\grLderivado{\grG_5}{n}=
\generadopor{b}\unrhd\ldots
\quad n     \in \setN
\end{flalign*}



\textbf{Serie Derivada:}
\begin{flalign*}
&\grG_5\rhd\grderivado{\grG_5}{1}=
\generadopor{b}\rhd\set{\unity}=
\grderivado{\grG_5}{2}     \cong     \trivialgr
\end{flalign*}





\chapter{Los grupos de orden 20}

Los razonamientos para obtener los grupos de orden $20$ son muy parecidos a los del cap\'{\i}tulo anterior.

Llamamos $\nu_5 = \left|\psylowlattice{\grG}{5}\right|$ y $\nu_2 = \left|\psylowlattice{\grG}{2}\right|$. Ahora,
por los teoremas de Sylow sabemos que ambos n\'umeros son mayores o iguales a $1$. Adem\'as, $\nu_5 ∣4$ y
$\nu_5       \equiv       1 \mod 5$. As\'{\i}, $\nu_5       \in \set{1,2,4}∩\set{1,6,11,\ldots} = \set{1}$. As\'{\i} que $\nu_5=1$
y el grupo de orden $5$ ser\'a \'unico y por lo tanto caracter\'{\i}stico (y por lo tanto normal). Este grupo solo
puede ser el c\'{\i}clico de orden $5$.

En este cap\'{\i}tulo $\grK\normsubgr\grG$, siendo $\left|\grG\right|=20$, y $\grK$ es el \'unico subgrupo de orden $5$
de $\grG$.

A su vez, $\nu_2      \in \set{1,5}∩\set{1,3,5}=\set{1,5}$.

\section{Caso de 1 solo subgrupo de orden 4}


Lo que asumimos es que $\nu_2 = 1$.


En esta secci\'on $\grN\normsubgr\grG$, es el \'unico subgrupo de orden $4$.


Ahora bien $\grK     \cap \grN=\set{\unity}$, $\grG=\grK\grN$ y por lo tanto
$\grG     \cong     \grK     \times     \grN     \cong     \cyclicgr{5}     \times     \grN$.

Si $\grN     \cong     \cyclicgr{4}$ entonces $\grG     \cong     \cyclicgr{20}$.
Si $\grN     \cong     \pdcyclics{2}{2}$ entonces $\grG     \cong     \pdcyclics{10}{2}$.

\begin{flalign}
\grG_1 &\mathdef \cyclicgr{20}\\
\grG_2 &\mathdef \pdcyclics{10}{2}
\end{flalign}

\section{Caso de exactamente 5 subgrupos de orden 4}


Lo que asumimos es que $\nu_2 = 5$.



Consideramos $\grK={x,x^2,x^3,x^4,\unity}=\generadopor{x}$.



En esta secci\'on $\psylowlattice{\grG}{2}=\set{\grN_1,\grN_2,\grN_3,\grN_4,\grN_5}$.
Estos han de ser conjugados entre s\'{\i}, isomorfos y no normales. Adem\'as,
ning\'un elemento de estos subgrupos que no sea la unidad, conmuta con $x$.


\subsection{Caso de exactamente 5 subgrupos isomorfos a C4}



Sea $\grN_1=\generadopor{a}     \cong      \cyclicgr{4}$.


Como $a^2$ es de orden $2$, habr\'a subgrupo normales de orden $10$.


Podemos desarrollar $\grG=\set{x^i a^j \vert 0\leq i\leq 5 \;\ly\; 0\leq j\leq 4\;\;\ly\;\; xa     \neq      ax\;\ly\;     \exists      mn\;\;ax=x^m a^n}$.



Adem\'as sabemos que deben existir $\set{(i_\kappa,j_\kappa )\vert 1\leq\kappa \leq4\ly (i_\kappa,j_\kappa )     \neq     (0,1)}$
tales que son los \'unicos con la propiedad: $(x^{i_\kappa }a^{j_\kappa })^2=\unity$ (adem\'as de $a$ misma).
Cualquier otro elemento tendr\'a orden $5$.



Planteando estas ecuaciones de divisibilidad nos quedan las siguientes soluciones:



\begin{flalign*}
m=4 &\ly 1\leq i \leq 4\\
n=3 &\ly 1\leq j \leq 4
\end{flalign*}


\begin{flalign*}
m=4 &\ly 1\leq i \leq 4\\
n=1 &\ly j=2
\end{flalign*}


Que nos llevan respectivamente a las ecuaciones $ax=x^4a^3$ y $ax=x^4a$.


La primera nos lleva directamente a la siguiente ecuaci\'on de modelado del grupo $axa=x^4=x^{-1}$ y $xax=a^3=a^{-1}$.
Constituye una acci\'on, $\phi : \cyclicgr{4} \To \graut{\pdcyclics{5}{4}}$ que conforma el siguiente producto semidirecto (que puede reducirse a una acci\'on de $\cyclicgr{2}$ sobre $\cyclicgr{10}$):


\begin{flalign*}
&\grG_3 \mathdef \left\langle x,a\vert o(x)=5,o(a)=4, \star\right\rangle\\
&\quad(x^i,a^j)\star(x^k,a^l)= \left\lbrace
\begin{array}{ccc}
  (x^{i+k},a^{j+l}) &      \Leftarrow      & l\mod 2=0 \\
  (x^{(5-i)+k},a^{(4-j)+l}) &      \Leftarrow      & l\mod 2=1
\end{array}
\right.\\
\end{flalign*}


La anterior ecuaci\'on de divisibilidad nos llevaba a $ax=x^4a$.
Sugiere una acci\'on, $\phi : \cyclicgr{4} \To \graut{\cyclicgr{5}}$
que conforma el siguiente producto semidirecto:


\begin{flalign*}
&\grG_4 \mathdef \left\langle x,a\vert o(x)=5,o(a)=4, \star\right\rangle\\
&\quad(x^i,a^j)\star(x^k,a^l)=
  \begin{cases}
  (x^{i+k},a^{j+l}) &     \Leftarrow      l=0\\
  (x^{2i+k},a^{j+l}) &     \Leftarrow      l=1\\
  (x^{4i+k},a^{j+l}) &     \Leftarrow      l=2\\
  (x^{3i+k},a^{j+l}) &     \Leftarrow      l=3
  \end{cases}
\end{flalign*}

\subsection{Caso de exactamente 5 subgrupos isomorfos a C2xC2}

Asumimos
\[
\grK=\generadopor{x}     \cong      \cyclicgr{5}
\]
\[
\grN_1=\generadopor{a,b}     \cong      \pdcyclics{2}{2}
\]
\[
\grH_1=\generadopor{a}     \cong      \cyclicgr{2}
\]
\[
\grH_2=\generadopor{b}     \cong      \cyclicgr{2}
\]
y por lo tanto habr\'a subgrupos normales de orden $10$. Los de orden $2$ han de ser no normales.

Para esto tenemos el grupo ya estudiado $\dihgr{10}$ y el producto directo $\dihgr{5}     \times     \cyclicgr{2}$.

\begin{flalign}
&\grG_5\mathdef\dihgr{10}\\
&\grG_6\mathdef\pd{\dihgr{5}}{\cyclicgr{2}}
\end{flalign}

\section{Los grupos G1 y G2 de orden 20}

Solo daremos de estos grupos su subgrupo de Frattini, para no engrosar m\'as el trabajo con informaci\'on casi inmediata.

\begin{flalign}
&\grG_1 = \generadopor{a \vert \mathbf{\mathtt{o}}\left(a\right)=20}\\
&\grG_2 = \generadopor{a,b \vert \mathbf{\mathtt{o}}\left(a\right)=10,\mathbf{\mathtt{o}}\left(a\right)=2,ab=ba}
\end{flalign}


\section{El grupo G3 de orden 20}


$\grG_3      \cong      \psdef{\cyclicgr{5}}{\cyclicgr{4}}{\generadopor{a   \mapsto {\generadopor{x   \mapsto x^4}}}}    \cong    \dicgr{5}$


La tabla de Cayley es:


\smallskip\noindent
\resizebox{\linewidth}{!}{
\begin{tabular}{|c|c|c|c|c|c|c|c|c|c|c|c|c|c|c|c|c|c|c|c|}
\hline
$1$&$x$&$x^2$&$x^3$&$x^4$&$a$&$xa$&$x^2a$&$x^3a$&$x^4a$&$a^2$&$xa^2$&$x^2a^2$&$x^3a^2$&$x^4a^2$&$a^3$&$xa^3$&$x^2a^3$&$x^3a^3$&$x^4a^3$\\
\hline
$x$&$x^2$&$x^3$&$x^4$&$1$&$x^4a$&$a$&$xa$&$x^2a$&$x^3a$&$xa^2$&$x^2a^2$&$x^3a^2$&$x^4a^2$&$a^2$&$x^4a^3$&$a^3$&$xa^3$&$x^2a^3$&$x^3a^3$\\
\hline
$x^2$&$x^3$&$x^4$&$1$&$x$&$x^3a$&$x^4a$&$a$&$xa$&$x^2a$&$x^2a^2$&$x^3a^2$&$x^4a^2$&$a^2$&$xa^2$&$x^3a^3$&$x^4a^3$&$a^3$&$xa^3$&$x^2a^3$\\
\hline
$x^3$&$x^4$&$1$&$x$&$x^2$&$x^2a$&$x^3a$&$x^4a$&$a$&$xa$&$x^3a^2$&$x^4a^2$&$a^2$&$xa^2$&$x^2a^2$&$x^2a^3$&$x^3a^3$&$x^4a^3$&$a^3$&$xa^3$\\
\hline
$x^4$&$1$&$x$&$x^2$&$x^3$&$xa$&$x^2a$&$x^3a$&$x^4a$&$a$&$x^4a^2$&$a^2$&$xa^2$&$x^2a^2$&$x^3a^2$&$xa^3$&$x^2a^3$&$x^3a^3$&$x^4a^3$&$a^3$\\
\hline
$a$&$xa$&$x^2a$&$x^3a$&$x^4a$&$a^2$&$xa^2$&$x^2a^2$&$x^3a^2$&$x^4a^2$&$a^3$&$xa^3$&$x^2a^3$&$x^3a^3$&$x^4a^3$&$1$&$x$&$x^2$&$x^3$&$x^4$\\
\hline
$xa$&$x^2a$&$x^3a$&$x^4a$&$a$&$x^4a^2$&$a^2$&$xa^2$&$x^2a^2$&$x^3a^2$&$xa^3$&$x^2a^3$&$x^3a^3$&$x^4a^3$&$a^3$&$x^4$&$1$&$x$&$x^2$&$x^3$\\
\hline
$x^2a$&$x^3a$&$x^4a$&$a$&$xa$&$x^3a^2$&$x^4a^2$&$a^2$&$xa^2$&$x^2a^2$&$x^2a^3$&$x^3a^3$&$x^4a^3$&$a^3$&$xa^3$&$x^3$&$x^4$&$1$&$x$&$x^2$\\
\hline
$x^3a$&$x^4a$&$a$&$xa$&$x^2a$&$x^2a^2$&$x^3a^2$&$x^4a^2$&$a^2$&$xa^2$&$x^3a^3$&$x^4a^3$&$a^3$&$xa^3$&$x^2a^3$&$x^2$&$x^3$&$x^4$&$1$&$x$\\
\hline
$x^4a$&$a$&$xa$&$x^2a$&$x^3a$&$xa^2$&$x^2a^2$&$x^3a^2$&$x^4a^2$&$a^2$&$x^4a^3$&$a^3$&$xa^3$&$x^2a^3$&$x^3a^3$&$x$&$x^2$&$x^3$&$x^4$&$1$\\
\hline
$a^2$&$xa^2$&$x^2a^2$&$x^3a^2$&$x^4a^2$&$a^3$&$xa^3$&$x^2a^3$&$x^3a^3$&$x^4a^3$&$1$&$x$&$x^2$&$x^3$&$x^4$&$a$&$xa$&$x^2a$&$x^3a$&$x^4a$\\
\hline
$xa^2$&$x^2a^2$&$x^3a^2$&$x^4a^2$&$a^2$&$x^4a^3$&$a^3$&$xa^3$&$x^2a^3$&$x^3a^3$&$x$&$x^2$&$x^3$&$x^4$&$1$&$x^4a$&$a$&$xa$&$x^2a$&$x^3a$\\
\hline
$x^2a^2$&$x^3a^2$&$x^4a^2$&$a^2$&$xa^2$&$x^3a^3$&$x^4a^3$&$a^3$&$xa^3$&$x^2a^3$&$x^2$&$x^3$&$x^4$&$1$&$x$&$x^3a$&$x^4a$&$a$&$xa$&$x^2a$\\
\hline
$x^3a^2$&$x^4a^2$&$a^2$&$xa^2$&$x^2a^2$&$x^2a^3$&$x^3a^3$&$x^4a^3$&$a^3$&$xa^3$&$x^3$&$x^4$&$1$&$x$&$x^2$&$x^2a$&$x^3a$&$x^4a$&$a$&$xa$\\
\hline
$x^4a^2$&$a^2$&$xa^2$&$x^2a^2$&$x^3a^2$&$xa^3$&$x^2a^3$&$x^3a^3$&$x^4a^3$&$a^3$&$x^4$&$1$&$x$&$x^2$&$x^3$&$xa$&$x^2a$&$x^3a$&$x^4a$&$a$\\
\hline
$a^3$&$xa^3$&$x^2a^3$&$x^3a^3$&$x^4a^3$&$1$&$x$&$x^2$&$x^3$&$x^4$&$a$&$xa$&$x^2a$&$x^3a$&$x^4a$&$a^2$&$xa^2$&$x^2a^2$&$x^3a^2$&$x^4a^2$\\
\hline
$xa^3$&$x^2a^3$&$x^3a^3$&$x^4a^3$&$a^3$&$x^4$&$1$&$x$&$x^2$&$x^3$&$xa$&$x^2a$&$x^3a$&$x^4a$&$a$&$x^4a^2$&$a^2$&$xa^2$&$x^2a^2$&$x^3a^2$\\
\hline
$x^2a^3$&$x^3a^3$&$x^4a^3$&$a^3$&$xa^3$&$x^3$&$x^4$&$1$&$x$&$x^2$&$x^2a$&$x^3a$&$x^4a$&$a$&$xa$&$x^3a^2$&$x^4a^2$&$a^2$&$xa^2$&$x^2a^2$\\
\hline
$x^3a^3$&$x^4a^3$&$a^3$&$xa^3$&$x^2a^3$&$x^2$&$x^3$&$x^4$&$1$&$x$&$x^3a$&$x^4a$&$a$&$xa$&$x^2a$&$x^2a^2$&$x^3a^2$&$x^4a^2$&$a^2$&$xa^2$\\
\hline
$x^4a^3$&$a^3$&$xa^3$&$x^2a^3$&$x^3a^3$&$x$&$x^2$&$x^3$&$x^4$&$1$&$x^4a$&$a$&$xa$&$x^2a$&$x^3a$&$xa^2$&$x^2a^2$&$x^3a^2$&$x^4a^2$&$a^2$\\
\hline
\end{tabular}
}


Por simple inspecci\'on de la tabla, vemos que no es abeliano, y por construcci\'on:
\[
\grG_3 =\generadopor{x}\generadopor{a} =
\generadopor{x,a \vert x^5=a^4=\unity\ly xa=ax^4\ly xa^2=a^2x} =
\psdef{\cyclicgr{5}}{\cyclicgr{4}}{\generadopor{a   \mapsto {\generadopor{x   \mapsto x^4}}}}
\]



Los \'ordenes de los elementos del grupo son:

\smallskip\noindent
\resizebox{\linewidth}{!}{
\begin{tabular}{|c|c|c|c|c|c|c|}
\hline
orden & =1 & =2 & =4 & =5 & =10 & =20\\
\hline
elementos
& $\set{\unity}$
& $\set{a^2}$
& $\set{a,xa,x^2a,x^3a,x^4a,a^3,xa^3,x^2a^3,x^3a^3,x^4a^3}$
& $\set{x,x^2,x^3,x^4}$
& $\set{xa^2,x^2a^2,x^3a^2,x^4a^2}$
& $\set{}$ \\
\hline
cardinal & 1 & 1 & 10 & 4 & 4 & 0 \\
\hline
\end{tabular}
}





El centro del grupo es $\centro{\grG_3}=\generadopor{a^2}    \cong    \cyclicgr{2}$ y el derivado es $\grderivado{\grG_3}{1}=\generadopor{x}     \cong     \cyclicgr{5}$.
Esto nos sirve para poder clasificar los subgrupos que a continuaci\'on vamos a listar.



El ret\'{\i}culo de subgrupos queda as\'{\i} clasificado:



\smallskip\noindent
\resizebox{\linewidth}{!}{
\begin{tabular}{|c|c|c|c|c|c|c|}
\hline
orden & =1 & =2 & =4 & =5 & =10 & =20\\
\hline
elementos
&$\set{\set{\unity}}$
&$\set{\generadopor{a^2}}$
&$\set{\generadopor{a},\generadopor{xa},\generadopor{x^2a},\generadopor{x^3a},\generadopor{x^4a}}$
&$\set{\generadopor{x}}$
&$\set{\generadopor{xa^2}}$
&$\set{\generadopor{x,a}}$ \\
\hline
isomorfo a, card
&$\trivialgr,1$
&$\cyclicgr{2},1$
&$\cyclicgr{4},5$
&$\cyclicgr{5},1$
&$\cyclicgr{10},1$
&$\psd{\cyclicgr{5}}{\cyclicgr{4}},1$ \\
\hline
propiedades
&{}
&(cyc,ch,cent),1
&(cyc),5
&(cyc,ch,der),1
&(cyc,ch)
&{} \\
\hline
\end{tabular}
}

Para la fila de propiedades de los subgrupos la leyenda es: ``norm'' significa normal, ``ch'' significa caracter\'{\i}stico,
``der'' significa que est\'a contenido en el derivado, ``cent'' significa que est\'a en el centro, ``ab'' que es abeliano,
``cyc'' que es c\'{\i}clico, ``nil'' que es nilpotente.


Los grupos con nombre propio que quedan son solo los de Frattini y de Fitting.


\begin{tabular}{|c|c|c|c|c|}
\hline
Funci\'on $\mathfrak{F}\entrellaves{\grG}$
& Centro $\centro{\grG}$
& Derivado $\grderivado{\grG}{1}$
& Frattini $\frattini{\grG}$
& Fitting $\fitting{\grG}$\\
\hline
Subgrupo
& $\generadopor{a^2}$
& $\generadopor{x}$
& $\set{a^2}$
& $\generadopor{xa^2}$ \\
\hline
Isomorfo a $\mathfrak{F}\entrellaves{\grG}     \cong     $
& $     \cong     \cyclicgr{2}$
& $     \cong     \cyclicgr{5}$
& $     \cong     \cyclicgr{2}$
& $     \cong     \cyclicgr{10}$ \\
\hline
$\grcoc{\grG}{\mathfrak{F}\entrellaves{\grG}}     \cong     $
& $     \cong     \dihgr{5}$
& $     \cong     \cyclicgr{4}$
& $     \cong     \dihgr{5}$
& $     \cong     \cyclicgr{2}$ \\
\hline
\end{tabular}





En cuanto a las series centrales, hay que tener en cuenta que este grupo es resoluble pero no nilpotente, as\'{\i} que las series centrales se estabilizan antes del pretendido final.

\textbf{Serie Central Ascendente:}
\begin{flalign*}
&\set{\unity}\lhd\centro{\grG_3}=
\generadopor{a^2}\unlhd\centronum{\grG_3}{2}=
\generadopor{a^2}\unlhd\ldots\unlhd\centronum{\grG_3}{n}=
\generadopor{a^2}\unlhd\ldots
\quad n     \in \setN
\end{flalign*}



\textbf{Serie Central Descendente:}
\begin{flalign*}
&\grG_3\rhd\grLderivado{\grG_3}{1}=
\generadopor{x}\unrhd\ldots\unrhd\grLderivado{\grG_3}{n}=
\generadopor{x}\unrhd\ldots
\quad n     \in \setN
\end{flalign*}



\textbf{Serie Derivada:}
\begin{flalign*}
&\grG_3\rhd\grderivado{\grG_3}{1}=
\generadopor{x}\rhd\set{\unity}=
\grderivado{\grG_3}{2}     \cong     \trivialgr
\end{flalign*}





\section{El grupo G4 de orden 20}


\[
\grG_4      \cong      \psdef{\cyclicgr{5}}{\cyclicgr{4}}{\generadopor{a^{k}}   \mapsto {\generadopor{x   \mapsto x^{2^k}}}}
\]
\[
  \mathbf{\mathtt{o}}\entreparentesis{x}=5
\]
\[
  \mathbf{\mathtt{o}}\entreparentesis{a}=4
\]

La tabla de Cayley es:


\smallskip\noindent
\resizebox{\linewidth}{!}{
\begin{tabular}{|c|c|c|c|c|c|c|c|c|c|c|c|c|c|c|c|c|c|c|c|}
\hline
$1$&$x$&$x^2$&$x^3$&$x^4$&$a$&$xa$&$x^2a$&$x^3a$&$x^4a$&$a^2$&$xa^2$&$x^2a^2$&$x^3a^2$&$x^4a^2$&$a^3$&$xa^3$&$x^2a^3$&$x^3a^3$&$x^4a^3$\\
\hline
$x$&$x^2$&$x^3$&$x^4$&$1$&$x^2a$&$x^3a$&$x^4a$&$a$&$xa$&$x^4a^2$&$a^2$&$xa^2$&$x^2a^2$&$x^3a^2$&$x^3a^3$&$x^4a^3$&$a^3$&$xa^3$&$x^2a^3$\\
\hline
$x^2$&$x^3$&$x^4$&$1$&$x$&$x^4a$&$a$&$xa$&$x^2a$&$x^3a$&$x^3a^2$&$x^4a^2$&$a^2$&$xa^2$&$x^2a^2$&$xa^3$&$x^2a^3$&$x^3a^3$&$x^4a^3$&$a^3$\\
\hline
$x^3$&$x^4$&$1$&$x$&$x^2$&$xa$&$x^2a$&$x^3a$&$x^4a$&$a$&$x^2a^2$&$x^3a^2$&$x^4a^2$&$a^2$&$xa^2$&$x^4a^3$&$a^3$&$xa^3$&$x^2a^3$&$x^3a^3$\\
\hline
$x^4$&$1$&$x$&$x^2$&$x^3$&$x^3a$&$x^4a$&$a$&$xa$&$x^2a$&$xa^2$&$x^2a^2$&$x^3a^2$&$x^4a^2$&$a^2$&$x^2a^3$&$x^3a^3$&$x^4a^3$&$a^3$&$xa^3$\\
\hline
$a$&$xa$&$x^2a$&$x^3a$&$x^4a$&$a^2$&$xa^2$&$x^2a^2$&$x^3a^2$&$x^4a^2$&$a^3$&$xa^3$&$x^2a^3$&$x^3a^3$&$x^4a^3$&$1$&$x$&$x^2$&$x^3$&$x^4$\\
\hline
$xa$&$x^2a$&$x^3a$&$x^4a$&$a$&$x^2a^2$&$x^3a^2$&$x^4a^2$&$a^2$&$xa^2$&$x^4a^3$&$a^3$&$xa^3$&$x^2a^3$&$x^3a^3$&$x^3$&$x^4$&$1$&$x$&$x^2$\\
\hline
$x^2a$&$x^3a$&$x^4a$&$a$&$xa$&$x^4a^2$&$a^2$&$xa^2$&$x^2a^2$&$x^3a^2$&$x^3a^3$&$x^4a^3$&$a^3$&$xa^3$&$x^2a^3$&$x$&$x^2$&$x^3$&$x^4$&$1$\\
\hline
$x^3a$&$x^4a$&$a$&$xa$&$x^2a$&$xa^2$&$x^2a^2$&$x^3a^2$&$x^4a^2$&$a^2$&$x^2a^3$&$x^3a^3$&$x^4a^3$&$a^3$&$xa^3$&$x^4$&$1$&$x$&$x^2$&$x^3$\\
\hline
$x^4a$&$a$&$xa$&$x^2a$&$x^3a$&$x^3a^2$&$x^4a^2$&$a^2$&$xa^2$&$x^2a^2$&$xa^3$&$x^2a^3$&$x^3a^3$&$x^4a^3$&$a^3$&$x^2$&$x^3$&$x^4$&$1$&$x$\\
\hline
$a^2$&$xa^2$&$x^2a^2$&$x^3a^2$&$x^4a^2$&$a^3$&$xa^3$&$x^2a^3$&$x^3a^3$&$x^4a^3$&$1$&$x$&$x^2$&$x^3$&$x^4$&$a$&$xa$&$x^2a$&$x^3a$&$x^4a$\\
\hline
$xa^2$&$x^2a^2$&$x^3a^2$&$x^4a^2$&$a^2$&$x^2a^3$&$x^3a^3$&$x^4a^3$&$a^3$&$xa^3$&$x^4$&$1$&$x$&$x^2$&$x^3$&$x^3a$&$x^4a$&$a$&$xa$&$x^2a$\\
\hline
$x^2a^2$&$x^3a^2$&$x^4a^2$&$a^2$&$xa^2$&$x^4a^3$&$a^3$&$xa^3$&$x^2a^3$&$x^3a^3$&$x^3$&$x^4$&$1$&$x$&$x^2$&$xa$&$x^2a$&$x^3a$&$x^4a$&$a$\\
\hline
$x^3a^2$&$x^4a^2$&$a^2$&$xa^2$&$x^2a^2$&$xa^3$&$x^2a^3$&$x^3a^3$&$x^4a^3$&$a^3$&$x^2$&$x^3$&$x^4$&$1$&$x$&$x^4a$&$a$&$xa$&$x^2a$&$x^3a$\\
\hline
$x^4a^2$&$a^2$&$xa^2$&$x^2a^2$&$x^3a^2$&$x^3a^3$&$x^4a^3$&$a^3$&$xa^3$&$x^2a^3$&$x$&$x^2$&$x^3$&$x^4$&$1$&$x^2a$&$x^3a$&$x^4a$&$a$&$xa$\\
\hline
$a^3$&$xa^3$&$x^2a^3$&$x^3a^3$&$x^4a^3$&$1$&$x$&$x^2$&$x^3$&$x^4$&$a$&$xa$&$x^2a$&$x^3a$&$x^4a$&$a^2$&$xa^2$&$x^2a^2$&$x^3a^2$&$x^4a^2$\\
\hline
$xa^3$&$x^2a^3$&$x^3a^3$&$x^4a^3$&$a^3$&$x^2$&$x^3$&$x^4$&$1$&$x$&$x^4a$&$a$&$xa$&$x^2a$&$x^3a$&$x^3a^2$&$x^4a^2$&$a^2$&$xa^2$&$x^2a^2$\\
\hline
$x^2a^3$&$x^3a^3$&$x^4a^3$&$a^3$&$xa^3$&$x^4$&$1$&$x$&$x^2$&$x^3$&$x^3a$&$x^4a$&$a$&$xa$&$x^2a$&$xa^2$&$x^2a^2$&$x^3a^2$&$x^4a^2$&$a^2$\\
\hline
$x^3a^3$&$x^4a^3$&$a^3$&$xa^3$&$x^2a^3$&$x$&$x^2$&$x^3$&$x^4$&$1$&$x^2a$&$x^3a$&$x^4a$&$a$&$xa$&$x^4a^2$&$a^2$&$xa^2$&$x^2a^2$&$x^3a^2$\\
\hline
$x^4a^3$&$a^3$&$xa^3$&$x^2a^3$&$x^3a^3$&$x^3$&$x^4$&$1$&$x$&$x^2$&$xa$&$x^2a$&$x^3a$&$x^4a$&$a$&$x^2a^2$&$x^3a^2$&$x^4a^2$&$a^2$&$xa^2$\\
\hline
\end{tabular}
}


Por simple inspecci\'on de la tabla, vemos que no es abeliano, y por construcci\'on:
\[
\grG_4 =\generadopor{x}\generadopor{a} =
\generadopor{x,a \vert x^5=a^4=\unity\ly xa=ax^2\ly x^4a=a^3x} =
\psdef{\cyclicgr{5}}{\cyclicgr{4}}{\generadopor{a^{k}}   \mapsto {\generadopor{x   \mapsto x^{2^k}}}}
\]



Los \'ordenes de los elementos del grupo son:

\smallskip\noindent
\resizebox{\linewidth}{!}{
\begin{tabular}{|c|c|c|c|c|c|c|}
\hline
orden & =1 & =2 & =4 & =5 & =10 & =20\\
\hline
elementos
& $\set{\unity}$
& $\set{a^2,xa^2,x^2a^2,x^3a^2,x^4a^2}$
& $\set{a,xa,x^2a,x^3a,x^4a,a^3,xa^3,x^2a^3,x^3a^3,x^4a^3}$
& $\set{x,x^2,x^3,x^4}$
& $\set{}$
& $\set{}$ \\
\hline
cardinal & 1 & 5 & 10 & 4 & 0 & 0 \\
\hline
\end{tabular}
}





El centro del grupo es $\centro{\grG_4}=\set{\unity}    \cong    \trivialgr$ y el derivado es $\grderivado{\grG_4}{1}=\generadopor{x}     \cong     \cyclicgr{5}$.
Esto nos sirve para poder clasificar los subgrupos que a continuaci\'on vamos a listar.



El ret\'{\i}culo de subgrupos queda as\'{\i} clasificado:



\smallskip\noindent
\resizebox{\linewidth}{!}{
\begin{tabular}{|c|c|c|c|c|c|c|}
\hline
orden & =1 & =2 & =4 & =5 & =10 & =20\\
\hline
elementos
&$\set{\set{\unity}}$
&$\set{\generadopor{a^2},\generadopor{xa^2},\generadopor{x^2a^2},\generadopor{x^3a^2},\generadopor{x^4a^2}}$
&$\set{\generadopor{a},\generadopor{xa},\generadopor{x^2a},\generadopor{x^3a},\generadopor{x^4a}}$
&$\set{\generadopor{x}}$
&$\set{\generadopor{x,a^2}}$
&$\set{\generadopor{x,a}}$ \\
\hline
isomorfo a, card
&$\trivialgr,1$
&$\cyclicgr{2},5$
&$\cyclicgr{4},5$
&$\cyclicgr{5},1$
&$\dihgr{5},1$
&$\psd{\cyclicgr{5}}{\cyclicgr{4}},1$ \\
\hline
propiedades
&(cent)
&(cyc),5
&(cyc),5
&(cyc,ch,der),1
&(ch)
&{} \\
\hline
\end{tabular}
}

Para la fila de propiedades de los subgrupos la leyenda es: ``norm'' significa normal, ``ch'' significa caracter\'{\i}stico,
``der'' significa que est\'a contenido en el derivado, ``cent'' significa que est\'a en el centro, ``ab'' que es abeliano,
``cyc'' que es c\'{\i}clico, ``nil'' que es nilpotente.


Los grupos con nombre propio que quedan son solo los de Frattini y de Fitting.


\begin{tabular}{|c|c|c|c|c|}
\hline
Funci\'on $\mathfrak{F}\entrellaves{\grG}$
& Centro $\centro{\grG}$
& Derivado $\grderivado{\grG}{1}$
& Frattini $\frattini{\grG}$
& Fitting $\fitting{\grG}$\\
\hline
Subgrupo
& $\set{\unity}$
& $\generadopor{x}$
& $\set{\unity}$
& $\generadopor{x,a^2}$ \\
\hline
Isomorfo a $\mathfrak{F}\entrellaves{\grG}     \cong     $
& $     \cong     \trivialgr$
& $     \cong     \cyclicgr{5}$
& $     \cong     \trivialgr$
& $     \cong     \dihgr{5}$ \\
\hline
$\grcoc{\grG}{\mathfrak{F}\entrellaves{\grG}}     \cong     $
& $     \cong     \grG_4$
& $     \cong     \cyclicgr{4}$
& $     \cong     \grG_4$
& $     \cong     \cyclicgr{2}$ \\
\hline
\end{tabular}





En cuanto a las series centrales, hay que tener en cuenta que este grupo es resoluble pero no nilpotente, as\'{\i} que las series centrales se estabilizan antes del pretendido final.

\textbf{Serie Central Ascendente:}
\begin{flalign*}
&\set{\unity}\unlhd\centro{\grG_4}=
\set{\unity}\unlhd\centronum{\grG_4}{2}=
\set{\unity}\unlhd\ldots\unlhd\centronum{\grG_4}{n}=
\set{\unity}\unlhd\ldots
\quad n     \in \setN
\end{flalign*}



\textbf{Serie Central Descendente:}
\begin{flalign*}
&\grG_4\rhd\grLderivado{\grG_4}{1}=
\generadopor{x}\unrhd\ldots\unrhd\grLderivado{\grG_4}{n}=
\generadopor{x}\unrhd\ldots
\quad n     \in \setN
\end{flalign*}



\textbf{Serie Derivada:}
\begin{flalign*}
&\grG_4\rhd\grderivado{\grG_4}{1}=
\generadopor{x}\rhd\set{\unity}=
\grderivado{\grG_4}{2}     \cong     \trivialgr
\end{flalign*}





\section{El grupo G5 de orden 20}


\[
\grG_5      \cong      \dihgr{10}
\]


La tabla de Cayley es:


\smallskip\noindent
\resizebox{\linewidth}{!}{
\begin{tabular}{|c|c|c|c|c|c|c|c|c|c|c|c|c|c|c|c|c|c|c|c|}
\hline
$1$&$x$&$x^2$&$x^3$&$x^4$&$x^5$&$x^6$&$x^7$&$x^8$&$x^9$&$a$&$xa$&$x^2a$&$x^3a$&$x^4a$&$x^5a$&$x^6a$&$x^7a$&$x^8a$&$x^9a$\\
\hline
$x$&$x^2$&$x^3$&$x^4$&$x^5$&$x^6$&$x^7$&$x^8$&$x^9$&$1$&$x^9a$&$a$&$xa$&$x^2a$&$x^3a$&$x^4a$&$x^5a$&$x^6a$&$x^7a$&$x^8a$\\
\hline
$x^2$&$x^3$&$x^4$&$x^5$&$x^6$&$x^7$&$x^8$&$x^9$&$1$&$x$&$x^8a$&$x^9a$&$a$&$xa$&$x^2a$&$x^3a$&$x^4a$&$x^5a$&$x^6a$&$x^7a$\\
\hline
$x^3$&$x^4$&$x^5$&$x^6$&$x^7$&$x^8$&$x^9$&$1$&$x$&$x^2$&$x^7a$&$x^8a$&$x^9a$&$a$&$xa$&$x^2a$&$x^3a$&$x^4a$&$x^5a$&$x^6a$\\
\hline
$x^4$&$x^5$&$x^6$&$x^7$&$x^8$&$x^9$&$1$&$x$&$x^2$&$x^3$&$x^6a$&$x^7a$&$x^8a$&$x^9a$&$a$&$xa$&$x^2a$&$x^3a$&$x^4a$&$x^5a$\\
\hline
$x^5$&$x^6$&$x^7$&$x^8$&$x^9$&$1$&$x$&$x^2$&$x^3$&$x^4$&$x^5a$&$x^6a$&$x^7a$&$x^8a$&$x^9a$&$a$&$xa$&$x^2a$&$x^3a$&$x^4a$\\
\hline
$x^6$&$x^7$&$x^8$&$x^9$&$1$&$x$&$x^2$&$x^3$&$x^4$&$x^5$&$x^4a$&$x^5a$&$x^6a$&$x^7a$&$x^8a$&$x^9a$&$a$&$xa$&$x^2a$&$x^3a$\\
\hline
$x^7$&$x^8$&$x^9$&$1$&$x$&$x^2$&$x^3$&$x^4$&$x^5$&$x^6$&$x^3a$&$x^4a$&$x^5a$&$x^6a$&$x^7a$&$x^8a$&$x^9a$&$a$&$xa$&$x^2a$\\
\hline
$x^8$&$x^9$&$1$&$x$&$x^2$&$x^3$&$x^4$&$x^5$&$x^6$&$x^7$&$x^2a$&$x^3a$&$x^4a$&$x^5a$&$x^6a$&$x^7a$&$x^8a$&$x^9a$&$a$&$xa$\\
\hline
$x^9$&$1$&$x$&$x^2$&$x^3$&$x^4$&$x^5$&$x^6$&$x^7$&$x^8$&$xa$&$x^2a$&$x^3a$&$x^4a$&$x^5a$&$x^6a$&$x^7a$&$x^8a$&$x^9a$&$a$\\
\hline
$a$&$xa$&$x^2a$&$x^3a$&$x^4a$&$x^5a$&$x^6a$&$x^7a$&$x^8a$&$x^9a$&$1$&$x$&$x^2$&$x^3$&$x^4$&$x^5$&$x^6$&$x^7$&$x^8$&$x^9$\\
\hline
$xa$&$x^2a$&$x^3a$&$x^4a$&$x^5a$&$x^6a$&$x^7a$&$x^8a$&$x^9a$&$a$&$x^9$&$1$&$x$&$x^2$&$x^3$&$x^4$&$x^5$&$x^6$&$x^7$&$x^8$\\
\hline
$x^2a$&$x^3a$&$x^4a$&$x^5a$&$x^6a$&$x^7a$&$x^8a$&$x^9a$&$a$&$xa$&$x^8$&$x^9$&$1$&$x$&$x^2$&$x^3$&$x^4$&$x^5$&$x^6$&$x^7$\\
\hline
$x^3a$&$x^4a$&$x^5a$&$x^6a$&$x^7a$&$x^8a$&$x^9a$&$a$&$xa$&$x^2a$&$x^7$&$x^8$&$x^9$&$1$&$x$&$x^2$&$x^3$&$x^4$&$x^5$&$x^6$\\
\hline
$x^4a$&$x^5a$&$x^6a$&$x^7a$&$x^8a$&$x^9a$&$a$&$xa$&$x^2a$&$x^3a$&$x^6$&$x^7$&$x^8$&$x^9$&$1$&$x$&$x^2$&$x^3$&$x^4$&$x^5$\\
\hline
$x^5a$&$x^6a$&$x^7a$&$x^8a$&$x^9a$&$a$&$xa$&$x^2a$&$x^3a$&$x^4a$&$x^5$&$x^6$&$x^7$&$x^8$&$x^9$&$1$&$x$&$x^2$&$x^3$&$x^4$\\
\hline
$x^6a$&$x^7a$&$x^8a$&$x^9a$&$a$&$xa$&$x^2a$&$x^3a$&$x^4a$&$x^5a$&$x^4$&$x^5$&$x^6$&$x^7$&$x^8$&$x^9$&$1$&$x$&$x^2$&$x^3$\\
\hline
$x^7a$&$x^8a$&$x^9a$&$a$&$xa$&$x^2a$&$x^3a$&$x^4a$&$x^5a$&$x^6a$&$x^3$&$x^4$&$x^5$&$x^6$&$x^7$&$x^8$&$x^9$&$1$&$x$&$x^2$\\
\hline
$x^8a$&$x^9a$&$a$&$xa$&$x^2a$&$x^3a$&$x^4a$&$x^5a$&$x^6a$&$x^7a$&$x^2$&$x^3$&$x^4$&$x^5$&$x^6$&$x^7$&$x^8$&$x^9$&$1$&$x$\\
\hline
$x^9a$&$a$&$xa$&$x^2a$&$x^3a$&$x^4a$&$x^5a$&$x^6a$&$x^7a$&$x^8a$&$x$&$x^2$&$x^3$&$x^4$&$x^5$&$x^6$&$x^7$&$x^8$&$x^9$&$1$\\
\hline
\end{tabular}
}


Por simple inspecci\'on de la tabla, vemos que no es abeliano, y por construcci\'on:
\[
\grG_5 =\generadopor{x}\generadopor{a} =
\generadopor{x,a \vert x^5=a^4=\unity\ly axa=x^9} =
\dihgr{10}
\]



Los \'ordenes de los elementos del grupo son:

\smallskip\noindent
\resizebox{\linewidth}{!}{
\begin{tabular}{|c|c|c|c|c|c|c|}
\hline
orden & =1 & =2 & =4 & =5 & =10 & =20\\
\hline
elementos
& $\set{\unity}$
& $\set{x^5,a,xa,x^2a,x^3a,x^4a,x^5a,x^6a,x^7a,x^8a,x^9a}$
& $\set{}$
& $\set{x^2,x^4,x^6,x^8}$
& $\set{x,x^3,x^7,x^9}$
& $\set{}$ \\
\hline
cardinal & 1 & 11 & 0 & 4 & 4 & 0 \\
\hline
\end{tabular}
}





El centro del grupo es $\centro{\grG_5}=\set{x^4}    \cong    \cyclicgr{2}$ y el derivado es $\grderivado{\grG_5}{1}=\generadopor{x^2}    \cong     \cyclicgr{5}$.
Esto nos sirve para poder clasificar los subgrupos que a continuaci\'on vamos a listar.



El ret\'{\i}culo de subgrupos queda as\'{\i} clasificado:



\smallskip\noindent
\resizebox{\linewidth}{!}{
\begin{tabular}{|c|c|c|c|c|c|c|}
\hline
orden & =1 & =2 & =4 & =5 & =10 & =20\\
\hline
elementos
&$\set{\set{\unity}}$
&$\set{\generadopor{x^5}}$
&$\set{
\generadopor{a},\generadopor{xa},\generadopor{x^2a},\generadopor{x^3a},\generadopor{x^4a},
\generadopor{x^5a},\generadopor{x^6a},\generadopor{x^7a},\generadopor{x^8a},\generadopor{x^9a}
}
\vert
\set{
\generadopor{x^5,a},\generadopor{x^5,xa},\generadopor{x^5,x^2a},\generadopor{x^5,x^3a},\generadopor{x^5,x^4a}
}$
&$\set{\generadopor{x^2}}$
&$
\set{
\generadopor{x}
}
\vert
\set{
\generadopor{x^2,a},\generadopor{x^2,xa}
}$
&$\set{\generadopor{x,a}}$ \\
\hline
isomorfo a, card
&$\trivialgr,1$
&$\cyclicgr{2},1\vert\cyclicgr{2},10$
&$\cyclicgr{4},10\vert\pdcyclics{2}{2},5$
&$\cyclicgr{5},1$
&$\cyclicgr{5},1\vert\dihgr{5},2$
&$\dihgr{10},1$\\
\hline
propiedades
&{}
&(cent,ch),1$\vert$(cyc),10
&(cyc),10$\vert${(),5}
&(cyc,der,ch),5
&(cyc,ch),1$\vert$(norm),2
&{} \\
\hline
\end{tabular}
}

Para la fila de propiedades de los subgrupos la leyenda es: ``norm'' significa normal, ``ch'' significa caracter\'{\i}stico,
``der'' significa que est\'a contenido en el derivado, ``cent'' significa que est\'a en el centro, ``ab'' que es abeliano,
``cyc'' que es c\'{\i}clico, ``nil'' que es nilpotente.


Los grupos con nombre propio que quedan son solo los de Frattini y de Fitting.


\begin{tabular}{|c|c|c|c|c|}
\hline
Funci\'on $\mathfrak{F}\entrellaves{\grG}$
& Centro $\centro{\grG}$
& Derivado $\grderivado{\grG}{1}$
& Frattini $\frattini{\grG}$
& Fitting $\fitting{\grG}$\\
\hline
Subgrupo
& $\generadopor{x^5}$
& $\generadopor{x^2}$
& $\set{\unity}$
& $\generadopor{x}$ \\
\hline
Isomorfo a $\mathfrak{F}\entrellaves{\grG}     \cong     $
& $     \cong     \cyclicgr{2}$
& $     \cong     \cyclicgr{5}$
& $     \cong     \trivialgr$
& $     \cong     \dihgr{5}$ \\
\hline
$\grcoc{\grG}{\mathfrak{F}\entrellaves{\grG}}     \cong     $
& $     \cong     \dihgr{5}$
& $     \cong     \pdcyclics{2}{2}$
& $     \cong     \grG_5$
& $     \cong     \cyclicgr{2}$ \\
\hline
\end{tabular}





En cuanto a las series centrales, hay que tener en cuenta que este grupo es resoluble pero no nilpotente, as\'{\i} que las series centrales se estabilizan antes del pretendido final.

\textbf{Serie Central Ascendente:}
\begin{flalign*}
&\set{\unity}\unlhd\centro{\grG_5}=
\generadopor{x^5}{\unity}\unlhd\centronum{\grG_5}{2}=
\generadopor{x^5}\unlhd\ldots\unlhd\centronum{\grG_5}{n}=
\generadopor{x^5}\unlhd\ldots
\quad n     \in \setN
\end{flalign*}



\textbf{Serie Central Descendente:}
\begin{flalign*}
&\grG_5\rhd\grLderivado{\grG_5}{1}=
\generadopor{x^2}\unrhd\ldots\unrhd\grLderivado{\grG_5}{n}=
\generadopor{x^2}\unrhd\ldots
\quad n     \in \setN
\end{flalign*}



\textbf{Serie Derivada:}
\begin{flalign*}
&\grG_5\rhd\grderivado{\grG_5}{1}=
\generadopor{x^2}\rhd\set{\unity}=
\grderivado{\grG_5}{2}     \cong     \trivialgr
\end{flalign*}






\printbibliography



\end{document}
